% =====================================================================
%  LEAN 4 İLE MODERN MANTIK — CİLT I: TEMELLER
%  Ders Kitabı ve Araştırmacı El Kitabı
%
%  Derleme (XeLaTeX):
%     xelatex ana && bibtex ana && xelatex ana && xelatex ana
%  ya da:  latexmk -xelatex ana
% =====================================================================
\documentclass[11pt, a4paper, twoside, openany]{book}

% ---- Kitap bilgileri (buradan değiştirin) ----
\newcommand{\kitapadi}{Lean 4 ile Modern Mantık}
\newcommand{\altbaslik}{Ders Kitabı ve Araştırmacı El Kitabı}
\newcommand{\ciltadi}{Cilt I: Temeller}
\newcommand{\yazar}{İbrahim GÜNEY}
\newcommand{\surum}{Cilt I · Sürüm 1.0 — Ekim 2026}

% =====================================================================
%  kitap-stil.tex — Sayfa düzeni, yazı tipleri, başlıklar, kutular
%  Lean 4 ile Modern Mantık
% =====================================================================

% ---------- Dil ve yazı tipleri (XeLaTeX) ----------
\usepackage{fontspec}
\usepackage{polyglossia}
\setdefaultlanguage{turkish}
\setotherlanguage{english}

\usepackage{amsmath,amsthm}
\usepackage{unicode-math}
\setmainfont{TeX Gyre Pagella}
\setsansfont{TeX Gyre Heros}[Scale=MatchLowercase]
\setmonofont{DejaVu Sans Mono}[Scale=0.84]
\setmathfont{Latin Modern Math}[Scale=MatchLowercase]
\newfontfamily\leanfont{DejaVu Sans Mono}[Scale=0.84]
\newfontfamily\baslikfont{TeX Gyre Heros}[Ligatures=TeX]
\newfontfamily\sembolfont{DejaVu Sans}   % ● ○ ⚠ ✓ ◇ □ gibi semboller için
\newfontfamily\sembolkodfont{DejaVu Sans}[Scale=0.84]  % kod içinde ⋀ ⋁ ⟹ ⊨

% ---------- Sayfa ----------
\usepackage[a4paper, inner=2.6cm, outer=2.4cm, top=2.6cm, bottom=2.8cm,
            headheight=14pt]{geometry}
\usepackage{microtype}
\usepackage[table]{xcolor}
\usepackage{graphicx}
\usepackage[most]{tcolorbox}
\usepackage{tikz}
\usetikzlibrary{arrows.meta,positioning,calc,shapes.geometric,
                fit,backgrounds,decorations.pathreplacing,matrix}
\usepackage{booktabs,array,tabularx,multirow}
\usepackage{enumitem}
\setlist{itemsep=2pt, topsep=4pt}
\usepackage{ebproof}           % doğal çıkarım ağaçları
\usepackage{csquotes}
\usepackage{multicol}
\usepackage[numbers,sort&compress]{natbib}
\usepackage{xurl}

% ---------- Ana renkler ----------
\definecolor{anarenk}{HTML}{3B5BDB}     % lacivert
\definecolor{ikincil}{HTML}{0CA678}     % yeşil
\definecolor{vurgu}{HTML}{D6336C}       % pembe-kırmızı
\definecolor{koyu}{HTML}{1F2A44}
\definecolor{tabloust}{HTML}{3B5BDB}
\definecolor{tablosatir}{HTML}{EDF2FF}

\definecolor{tanimbg}{HTML}{EDF2FF}  \definecolor{tanimframe}{HTML}{3B5BDB}
\definecolor{teorembg}{HTML}{FFF0F6} \definecolor{teoremframe}{HTML}{C2255C}
\definecolor{ornekbg}{HTML}{F3F0FF}  \definecolor{ornekframe}{HTML}{7950F2}
\definecolor{fikirbg}{HTML}{E6FCF5}  \definecolor{fikirframe}{HTML}{0CA678}
\definecolor{dikkatbg}{HTML}{FFF4E6} \definecolor{dikkatframe}{HTML}{E8590C}
\definecolor{arastirmabg}{HTML}{F1F3F5}\definecolor{arastirmaframe}{HTML}{495057}
\definecolor{tarihbg}{HTML}{FFF9DB}  \definecolor{tarihframe}{HTML}{E67700}
\definecolor{alistirmabg}{HTML}{F8F9FA}\definecolor{alistirmaframe}{HTML}{1971C2}

% ---------- Başlıklar ----------
\usepackage{titlesec}
\titleformat{\chapter}[display]
  {\baslikfont\color{koyu}}
  {\begin{tikzpicture}[remember picture, overlay]
     \fill[anarenk] (current page.north west) rectangle
           ([yshift=-3.2cm]current page.north east);
     \node[anchor=south east, text=white, font=\baslikfont\bfseries\fontsize{60}{60}\selectfont]
          at ([xshift=-2.4cm, yshift=-3.0cm]current page.north east) {\thechapter};
     \node[anchor=south west, text=white!85!anarenk, font=\baslikfont\large]
          at ([xshift=2.6cm, yshift=-3.0cm]current page.north west) {\MakeUppercase{\chaptertitlename}};
   \end{tikzpicture}}
  {0pt}
  {\vspace*{1.2cm}\Huge\bfseries\raggedright}
  [\vspace{4pt}{\color{anarenk}\titlerule[1.5pt]}]
\titlespacing*{\chapter}{0pt}{0pt}{24pt}

\titleformat{\section}
  {\baslikfont\Large\bfseries\color{koyu}}
  {\textcolor{anarenk}{\thesection}}{0.8em}{}
\titleformat{\subsection}
  {\baslikfont\large\bfseries\color{koyu}}
  {\textcolor{anarenk}{\thesubsection}}{0.8em}{}
\titleformat{\subsubsection}
  {\baslikfont\normalsize\bfseries\color{koyu}}
  {\textcolor{anarenk}{\thesubsubsection}}{0.8em}{}
\titleformat{\paragraph}[runin]
  {\baslikfont\normalsize\bfseries\color{koyu}}{}{0pt}{}[.]
\setcounter{secnumdepth}{2}
\setcounter{tocdepth}{1}

\titleformat{\part}[display]
  {\centering\baslikfont\color{white}}
  {\begin{tikzpicture}[remember picture, overlay]
     \fill[koyu] (current page.south west) rectangle (current page.north east);
     \fill[anarenk] ([yshift=-9cm]current page.north west) rectangle
                    ([yshift=-9.4cm]current page.north east);
   \end{tikzpicture}\Large KISIM \thepart}
  {12pt}{\Huge\bfseries}

% ---------- Üst bilgi ----------
\usepackage{fancyhdr}
\pagestyle{fancy}
\fancyhf{}
\renewcommand{\chaptermark}[1]{\markboth{#1}{}}
\renewcommand{\sectionmark}[1]{\markright{\thesection\quad #1}}
\fancyhead[LE]{\baslikfont\small\color{black!60}\ifnum\value{chapter}>0 \thechapter\quad\fi\leftmark}
\fancyhead[RO]{\baslikfont\small\color{black!60}\rightmark}
\fancyfoot[LE,RO]{\baslikfont\small\bfseries\color{anarenk}\thepage}
\renewcommand{\headrulewidth}{0.4pt}
\renewcommand{\headrule}{{\color{black!20}\hrule width\headwidth height\headrulewidth}}
\fancypagestyle{plain}{\fancyhf{}\fancyfoot[C]{\baslikfont\small\color{anarenk}\thepage}
  \renewcommand{\headrulewidth}{0pt}}

% ---------- Numaralı kutular (bölüm bazında) ----------
\newcounter{kutusayac}[chapter]
\renewcommand{\thekutusayac}{\thechapter.\arabic{kutusayac}}

\tcbset{
  kitapkutu/.style={
    enhanced, breakable,
    boxrule=0.6pt, arc=3pt, leftrule=4pt,
    fonttitle=\baslikfont\bfseries,
    titlerule=0pt,
    before skip=12pt, after skip=12pt,
    left=8pt, right=8pt,
  }
}

\newtcbtheorem[number within=chapter]{tanim}{Tanım}{kitapkutu,
  colback=tanimbg, colframe=tanimframe, colbacktitle=tanimbg,
  coltitle=tanimframe!80!black, separator sign={.},
  description font=\mdseries}{tnm}
\newtcbtheorem[use counter from=tanim]{teorem}{Teorem}{kitapkutu,
  colback=teorembg, colframe=teoremframe, colbacktitle=teorembg,
  coltitle=teoremframe!80!black, separator sign={.},
  description font=\mdseries}{trm}
\newtcbtheorem[use counter from=tanim]{ornek}{Örnek}{kitapkutu,
  colback=ornekbg, colframe=ornekframe, colbacktitle=ornekbg,
  coltitle=ornekframe!80!black, separator sign={.},
  description font=\mdseries}{orn}

% Numarasız kutular
\newtcolorbox{onemlifikir}[1][Önemli Fikir]{kitapkutu,
  colback=fikirbg, colframe=fikirframe, colbacktitle=fikirbg,
  coltitle=fikirframe!70!black, title={#1}}
\newtcolorbox{gunlukornek}[1][Günlük Hayattan]{kitapkutu,
  colback=ornekbg, colframe=ornekframe, colbacktitle=ornekbg,
  coltitle=ornekframe!70!black, title={#1}}
\newtcolorbox{dikkat}[1][Dikkat: Sık Yapılan Hata]{kitapkutu,
  colback=dikkatbg, colframe=dikkatframe, colbacktitle=dikkatbg,
  coltitle=dikkatframe!80!black, title={{\sembolfont ⚠}\enspace #1}}
\newtcolorbox{tarihce}[1][Tarihçe]{kitapkutu,
  colback=tarihbg, colframe=tarihframe, colbacktitle=tarihbg,
  coltitle=tarihframe!80!black, title={#1}, fontupper=\small}
\newtcolorbox{arastirmanotu}[1][Araştırma Notu]{kitapkutu,
  colback=arastirmabg, colframe=arastirmaframe, colbacktitle=arastirmaframe,
  coltitle=white, title={#1}, fontupper=\small, titlerule=0.5pt,
  attach boxed title to top left={xshift=10pt,yshift=-2.2mm},
  boxed title style={colback=arastirmaframe, arc=2pt, boxrule=0pt}}

% Bölüm başı: hedefler kutusu
\newtcolorbox{bolumhedefleri}{enhanced, breakable,
  colback=white, colframe=anarenk, boxrule=0.8pt, arc=4pt,
  title={\baslikfont\bfseries Bu bölümde öğrenecekleriniz},
  colbacktitle=anarenk, coltitle=white,
  before skip=6pt, after skip=16pt, left=8pt}

% Bölüm sonu özeti
\newtcolorbox{bolumozeti}{enhanced, breakable,
  colback=fikirbg, colframe=ikincil, boxrule=0.8pt, arc=4pt,
  title={\baslikfont\bfseries Bölüm Özeti},
  colbacktitle=ikincil, coltitle=white,
  before skip=12pt, after skip=12pt, left=8pt}

% ---------- Alıştırmalar ----------
\newcounter{alistirma}[chapter]
\renewcommand{\thealistirma}{\thechapter.\arabic{alistirma}}
\newcommand{\zorluk}[1]{%
  {\sembolfont\footnotesize\ifcase#1\or\textcolor{ikincil}{●○○}\or\textcolor{dikkatframe}{●●○}%
  \or\textcolor{vurgu}{●●●}\fi}}
\newenvironment{alistirmalar}{%
  \begin{tcolorbox}[enhanced, breakable, colback=alistirmabg,
     colframe=alistirmaframe, boxrule=0.6pt, arc=3pt, leftrule=4pt,
     title={\baslikfont\bfseries Alıştırmalar}, colbacktitle=alistirmabg,
     coltitle=alistirmaframe!80!black, titlerule=0pt]
  \small
  Zorluk: \zorluk{1}~temel \quad \zorluk{2}~orta \quad \zorluk{3}~ileri.
  Lean alıştırmalarının iskeletleri eşlik eden projedeki
  \lc{Alistirmalar/} klasöründedir; \leantac{sorry} yerine kendi ispatınızı yazın.
  \begin{description}[leftmargin=0pt, itemsep=6pt, font=\normalfont]}
  {\end{description}\end{tcolorbox}}
\newcommand{\alistirma}[1]{\refstepcounter{alistirma}%
  \item[\baslikfont\bfseries\textcolor{alistirmaframe}{\thealistirma}\enspace\zorluk{#1}]}

% ---------- Doğal çıkarım kuralları ----------
\newcommand{\kural}[1]{\textcolor{vurgu}{\textsf{\small #1}}}
\ebproofset{right label template=\kural{\inserttext}, separation=1.4em}

% ---------- Tablolar ----------
\newcommand{\tablobaslik}[1]{\textcolor{white}{\baslikfont\bfseries\small #1}}

% ---------- Bağlantılar ----------
\usepackage[hidelinks]{hyperref}
\hypersetup{colorlinks=true, linkcolor=anarenk, citecolor=ikincil,
            urlcolor=vurgu, pdfauthor={}, pdftitle={Lean 4 ile Modern Mantık}}
\usepackage[english,nameinlink]{cleveref}  % polyglossia-turkish ile uyum için dil açıkça verilir
\crefname{chapter}{Bölüm}{Bölüm}
\crefname{section}{Kesim}{Kesim}
\crefname{subsection}{Kesim}{Kesim}
\crefname{figure}{Şekil}{Şekil}
\crefname{table}{Tablo}{Tablo}
\crefname{tcb@cnt@tanim}{Tanım}{Tanım}
\crefname{tcb@cnt@teorem}{Teorem}{Teorem}
\crefname{tcb@cnt@ornek}{Örnek}{Örnek}
\crefname{alistirma}{Alıştırma}{Alıştırma}
\crefname{part}{Kısım}{Kısım}

% Şekil/tablo başlıkları
\usepackage[font=small, labelfont={bf,color=anarenk}, labelsep=period]{caption}

% Kısaltmalar
\newcommand{\Lean}{Lean~4}
\newcommand{\Mathlib}{\textsf{Mathlib}}
\newcommand{\terim}[2]{\emph{#1} (İng.\ \foreignlanguage{english}{\emph{#2}})}

% ---------- Taslak bölüm planı ----------
\newtcolorbox{taslakplan}{enhanced, breakable,
  colback=white, colframe=black!35, boxrule=0.8pt, arc=4pt,
  borderline={0.4pt}{2pt}{black!20, dashed},
  title={\baslikfont\bfseries Yazım planı \hfill\normalfont\small taslak},
  colbacktitle=black!8, coltitle=black!70,
  before skip=10pt, after skip=10pt, left=8pt}

% Numaralı adım listeleri:  \begin{enumerate}[adimlar] ... \end{enumerate}
\DeclareRobustCommand{\adimno}[1]{\tikz[baseline=(n.base)]\node[circle,
  fill=anarenk, text=white, inner sep=1.4pt, minimum size=1.15em,
  font=\baslikfont\scriptsize\bfseries](n){#1};}
\SetEnumitemKey{adimlar}{label=\adimno{\arabic*}, leftmargin=2em, itemsep=5pt}

% Boş kalan sol sayfalarda üst/alt bilgi gösterme
\makeatletter
\renewcommand{\cleardoublepage}{\clearpage\if@twoside\ifodd\c@page\else
  \hbox{}\thispagestyle{empty}\newpage\if@twocolumn\hbox{}\newpage\fi\fi\fi}
\makeatother
\colorlet{gri}{black!30}

% Nesne dilinin bağlaçları (Bölüm 10–13): metin içinde
\newcommand{\ove}{\mathbin{\textstyle\bigwedge}}
\newcommand{\oveya}{\mathbin{\textstyle\bigvee}}
\newcommand{\oise}{\Longrightarrow}

% ---- Cilt II'ye göndermeler ----
% Kullanım: \ciltiki{4}'te  →  "Cilt II, Bölüm 4'te"
\newcommand{\ciltiki}[1]{Cilt~II, Bölüm~#1}

% ---- Proje kartları (Bölüm 17) ----
\newtcolorbox{projekarti}[2]{kitapkutu,
  colback=white, colframe=ornekframe, colbacktitle=ornekframe,
  coltitle=white, title={#1\enspace\textbar\enspace #2}}
\newcommand{\projebilgi}[3]{%
  \noindent{\baslikfont\small
    \textcolor{ornekframe}{\bfseries Ön koşul:}~#1\hfill
    \textcolor{ornekframe}{\bfseries Zorluk:}~#2\hfill
    \textcolor{ornekframe}{\bfseries Tahmini emek:}~#3}\par\medskip}

% =====================================================================
%  lean-stil.tex — Lean 4 kod renklendirme ve kod kutuları
%  Lean 4 ile Modern Mantık
%
%  Renklendirme minted + Pygments ile yapılır (Unicode ve Türkçe
%  karakterler sorunsuz). Kitaba özel sözcük çözümleyicisi:
%  stil/leankitap.py
%
%  Derleme -shell-escape gerektirir:
%      latexmk -xelatex -shell-escape ana
% =====================================================================

\usepackage[newfloat]{minted}
\usemintedstyle{bw}
\tcbuselibrary{minted}

% Sözcük çözümleyicisi: Pygments'e dosyadan özel çözümleyici yükletmek için
% "-x" bayrağı gerekir; dil adı olarak dosya yolu ve sınıf adı verilir.
\renewcommand{\MintedPygmentize}{pygmentize -x}

% ---------- Renk paleti ----------
\definecolor{leanbg}{HTML}{F7F9FC}      % kod arka planı
\definecolor{leanframe}{HTML}{3B5BDB}   % kod çerçevesi (lacivert)
\definecolor{leankw}{HTML}{7048E8}      % anahtar kelime (mor)
\definecolor{leantac}{HTML}{D6336C}     % taktik (pembe-kırmızı)
\definecolor{leantype}{HTML}{0C8599}    % Prop, Type (camgöbeği)
\definecolor{leanhyp}{HTML}{E8590C}     % hipotez adları (turuncu)
\definecolor{leancom}{HTML}{5C7F4A}     % yorum (zeytin yeşili)
\definecolor{leansym}{HTML}{1971C2}     % → ⊢ ∧ ∨ (mavi)
\definecolor{leancmd}{HTML}{862E9C}     % #check, #eval
\definecolor{leankirik}{HTML}{ADB5BD}   % satır numarası

\definecolor{hedefbg}{HTML}{FFF9DB}     % hedef durumu kutusu
\definecolor{hedefframe}{HTML}{F59F00}
\definecolor{bittibg}{HTML}{EBFBEE}
\definecolor{bittiframe}{HTML}{2F9E44}
\definecolor{hatabg}{HTML}{FFF5F5}
\definecolor{hataframe}{HTML}{E03131}

% ---------- Token renkleri ----------
% Pygments'in "bw" stili yüklendikten sonra, her kod bloğunun başında
% token makroları kitabın renk diliyle değiştirilir (formatcom).
\makeatletter
\newcommand{\leanrenk}[3]{% #1: token kısaltması, #2: renk, #3: ek biçim
  \@namedef{PYG@tok@#1}{\def\PYG@tc####1{\textcolor{#2}{####1}}#3}}
\newcommand{\leanrenkleri}{%
  \leanrenk{k}{leankw}{\let\PYG@bf=\textbf}%          anahtar kelime
  \leanrenk{kt}{leantype}{\let\PYG@bf=\textbf}%       Prop, Type
  \leanrenk{nb}{leantac}{\let\PYG@bf=\textbf}%        taktik
  \leanrenk{nv}{leanhyp}{\let\PYG@bf=\textbf}%        hipotez adı
  \leanrenk{nd}{leancmd}{\let\PYG@bf=\textbf}%        #check
  \leanrenk{o}{leansym}{\let\PYG@bf=\relax}%          → ∧ ∨ ⊢
  \leanrenk{nt}{leankw}{\let\PYG@bf=\relax}%          ⟨ ⟩
  \leanrenk{c}{leancom}{\let\PYG@it=\textit}%         yorum
  \leanrenk{c1}{leancom}{\let\PYG@it=\textit}%
  \leanrenk{cm}{leancom}{\let\PYG@it=\textit}%
  \leanrenk{s2}{leantype}{\let\PYG@it=\relax}%        dizgi
  \leanrenk{s}{leantype}{\let\PYG@it=\relax}%
  \leanrenk{nx}{leansym}{\def\PYG@ff####1{{\sembolkodfont ####1}}}% ⋀ ⋁ ⟹ ⊨
}
\makeatother

\setminted{
  fontsize=\small,
  breaklines, breakanywhere,
  tabsize=2,
  numbersep=8pt,
  formatcom=\leanrenkleri,
}
\renewcommand{\theFancyVerbLine}{\textcolor{leankirik}{\tiny\arabic{FancyVerbLine}}}

% Lean dosyalarından çıkarılan kesitlerin yeri (ana.tex'e göre)
\newcommand{\kesitdizini}{kesitler}

% Derleme başında kesitleri yeniden çıkar (shell-escape açıkken)
\AtBeginDocument{%
  \immediate\write18{python3 ../araclar/kesitleri_cikar.py
    ../lean/ModernMantik \kesitdizini}}

% ---------- Kod kutuları ----------
\tcbset{
  leankutu/.style={
    enhanced, breakable,
    colback=leanbg, colframe=leanframe,
    boxrule=0.8pt, arc=3pt,
    left=4pt, right=6pt, top=6pt, bottom=4pt,
    fonttitle=\sffamily\bfseries\small,
    coltitle=white,
    attach boxed title to top left={xshift=8pt,yshift=-2.2mm},
    boxed title style={colback=leanframe, arc=2pt, boxrule=0pt},
    before skip=12pt, after skip=10pt,
    listing engine=minted,
    minted language=stil/leankitap.py:LeanKitapLexer,
    minted options={linenos, xleftmargin=1.6em, fontsize=\small,
                    breaklines, breakanywhere, numbersep=8pt,
                    formatcom=\leanrenkleri},
  }
}

% Kitap içine doğrudan yazılan kod (derlenmeyen örnekler için):
%   \begin{leankod}[Başlık] ... \end{leankod}
\newtcblisting{leankod}[1][Lean 4]{leankutu, title={#1}, listing only}

% Eşlik eden Lean projesinden ÇEKİLEN kod (tek doğruluk kaynağı):
%   \leandosya[Başlık]{Bolum02}{hipotetik_kiyas}
% Bolum02.lean içindeki  -- ANCHOR: hipotetik_kiyas  ...
%                        -- ANCHOR_END: hipotetik_kiyas  arası basılır.
\newcommand{\leandosya}[3][Lean 4]{%
  \tcbinputlisting{leankutu,
    title={#1\quad{\normalfont\scriptsize\ttfamily ModernMantik/#2.lean}},
    listing only,
    listing file={\kesitdizini/#2/#3.lean}}%
}

% Hedef durumu (Lean'in Infoview paneli):
%   \begin{hedefdurumu}[Başlık] ... \end{hedefdurumu}
\newtcblisting{hedefdurumu}[1][Hedef durumu]{
  enhanced, listing only,
  listing engine=minted, minted language=stil/leankitap.py:LeanKitapLexer,
  minted options={fontsize=\small, formatcom=\leanrenkleri},
  colback=hedefbg, colframe=hedefframe,
  boxrule=0.6pt, arc=2pt, leftrule=3pt,
  left=6pt, top=2pt, bottom=2pt,
  title={\sffamily\footnotesize\bfseries #1},
  coltitle=black!75, colbacktitle=hedefframe!30,
  toptitle=1pt, bottomtitle=1pt,
  before skip=6pt, after skip=6pt,
}

% Lean hata mesajı kutusu
\newtcblisting{leanhata}[1][Lean hata mesajı]{
  enhanced, listing only,
  listing engine=minted, minted language=text,
  minted options={fontsize=\small},
  colback=hatabg, colframe=hataframe,
  boxrule=0.6pt, arc=2pt, leftrule=3pt,
  left=6pt, top=2pt, bottom=2pt,
  title={\sffamily\footnotesize\bfseries #1},
  coltitle=hataframe!80!black, colbacktitle=hataframe!15,
  toptitle=1pt, bottomtitle=1pt,
  before skip=6pt, after skip=6pt,
}

% "İspat tamamlandı" çubuğu
\newcommand{\ispatbitti}{%
  \begin{tcolorbox}[enhanced, colback=bittibg, colframe=bittiframe,
      boxrule=0.6pt, arc=2pt, leftrule=3pt,
      left=6pt, top=3pt, bottom=3pt, before skip=6pt, after skip=8pt]
    {\leanfont\small No goals}\hfill
    \textcolor{bittiframe}{\sffamily\bfseries İspat tamamlandı {\sembolfont ✓}}
  \end{tcolorbox}}

% ---------- Metin içi kısayollar ----------
\newcommand{\leantac}[1]{{\leanfont\textcolor{leantac}{\bfseries #1}}}
\newcommand{\leankw}[1]{{\leanfont\textcolor{leankw}{\bfseries #1}}}
\newcommand{\leanhyp}[1]{{\leanfont\textcolor{leanhyp}{\bfseries #1}}}
\newcommand{\leantype}[1]{{\leanfont\textcolor{leantype}{\bfseries #1}}}
\newcommand{\lc}[1]{{\leanfont #1}}

% İki hedef durumunu yan yana gösterme (önce ⇒ sonra):
%   \begin{hedefgecisi}
%     \begin{hedefdurumu}[Önce] ... \end{hedefdurumu}
%   \sonra
%     \begin{hedefdurumu}[Sonra] ... \end{hedefdurumu}
%   \end{hedefgecisi}
\newenvironment{hedefgecisi}
  {\par\smallskip\noindent\begin{minipage}[t]{0.455\linewidth}\vspace{0pt}}
  {\end{minipage}\par\smallskip}
\newcommand{\sonra}{\end{minipage}\hfill
  \begin{minipage}[t]{0.05\linewidth}\vspace{0.9cm}\centering
    \textcolor{leantac}{\Large$\Rightarrow$}\end{minipage}\hfill
  \begin{minipage}[t]{0.455\linewidth}\vspace{0pt}}


\begin{document}

\frontmatter
% ---------------- Kapak ----------------
\begin{titlepage}
\thispagestyle{empty}
\begin{tikzpicture}[remember picture, overlay]
  % Zemin
  \fill[koyu] (current page.south west) rectangle (current page.north east);
  % Üst şerit
  \fill[anarenk] ([yshift=-7.2cm]current page.north west) rectangle
                 ([yshift=-7.5cm]current page.north east);

  % Arka planda soluk bir doğal çıkarım ağacı motifi
  \begin{scope}[shift={([xshift=-5.2cm, yshift=7.4cm]current page.south east)},
                text=white!25!koyu, font=\leanfont\large]
    \node at (0,3.2)  {[P]};
    \node at (2.2,3.2){P → Q};
    \draw[white!25!koyu, line width=0.8pt] (-0.6,2.8) -- (2.9,2.8);
    \node at (1.1,2.4){Q};
    \node at (3.6,2.4){Q → R};
    \draw[white!25!koyu, line width=0.8pt] (0.5,2.0) -- (4.3,2.0);
    \node at (2.4,1.6){R};
    \draw[white!25!koyu, line width=0.8pt] (1.4,1.2) -- (3.4,1.2);
    \node at (2.4,0.8){P → R};
  \end{scope}

  % Curry–Howard motifi
  \node[anchor=south west, text=white!45!koyu, font=\baslikfont\large,
        align=left]
    at ([xshift=2.4cm, yshift=3.0cm]current page.south west)
    {önermeler = türler\\[2pt] kanıtlar = programlar};

  % Başlık
  \node[anchor=north west, text=white, font=\baslikfont\bfseries\fontsize{40}{46}\selectfont,
        align=left]
    at ([xshift=2.4cm, yshift=-9.0cm]current page.north west)
    {Lean 4 ile\\Modern Mantık};
  \node[anchor=north west, text=white!75!anarenk, font=\baslikfont\Large, align=left]
    at ([xshift=2.4cm, yshift=-13.4cm]current page.north west)
    {\altbaslik};
  \node[anchor=north west, text=koyu, fill=white!85!anarenk, font=\baslikfont\bfseries\Large,
        rounded corners=3pt, inner sep=6pt]
    at ([xshift=2.4cm, yshift=-14.6cm]current page.north west)
    {\ciltadi};

  % Yazar ve sürüm
  \node[anchor=north west, text=white, font=\baslikfont\LARGE]
    at ([xshift=2.4cm, yshift=-3.6cm]current page.north west) {\yazar};
  \node[anchor=north west, text=white!60!koyu, font=\baslikfont\small]
    at ([xshift=2.4cm, yshift=-4.6cm]current page.north west) {\surum};

  % Küçük ⊢ logosu
  \node[anchor=south east, text=anarenk!70!white, font=\leanfont\bfseries\fontsize{70}{70}\selectfont]
    at ([xshift=-2.2cm, yshift=1.6cm]current page.south east) {⊢};
\end{tikzpicture}
\end{titlepage}
\cleardoublepage

\chapter*{Önsöz}
\addcontentsline{toc}{chapter}{Önsöz}
\markboth{Önsöz}{Önsöz}

Mantık, iki bin yılı aşkın bir süre boyunca kâğıt ve kalemle yapılan bir
disiplindi. Aristoteles'in kıyasları, Frege'nin \emph{Begriffsschrift}'i,
Gentzen'in doğal çıkarım sistemi ve Gödel'in eksiklik teoremleri hep bu
gelenekte doğdu. Yirmi birinci yüzyılda ise mantığın kendisi değişmese de onu
\emph{yapma biçimimiz} değişti: Bugün bir ispatı bilgisayara denetletebiliyor,
mantık sistemlerini birer veri yapısı olarak tanımlayıp onlar hakkındaki
teoremleri makinede doğrulayabiliyoruz.

Bu kitap, bu dönüşümü Türkçe okura sıfırdan anlatmak amacıyla yazıldı. Araç
olarak \Lean{} ispat asistanını kullanıyoruz \citep{demoura2021lean4}. \Lean{}
hem bir programlama dili hem de bir ispat denetleyicisidir. İçine yazdığınız
her ispatı, küçük ve güvenilir bir çekirdek adım adım denetler. Bir ispat
\Lean{}'den geçtiyse, o ispatta mantıksal bir boşluk kalmamış demektir.

\section*{Kitabın iki yüzü}

Kitap iki farklı okur kitlesini birlikte gözetiyor:

\begin{description}[font=\baslikfont\bfseries\color{anarenk}]
  \item[Ders kitabı olarak] Kısım~I ve~II, matematik, bilgisayar mühendisliği
  ve felsefe bölümlerinde bir dönemlik ``Mantık ve İspat'' dersine karşılık
  gelir. Önceden programlama bilgisi gerekmez. Her kavram önce kâğıt üzerinde,
  ardından \Lean{}'de, en sonunda da bir görselle anlatılır.
  \item[El kitabı olarak] Kısım~III ve~IV, metamantık (geçerlilik, tamlık,
  kesme eleme) ve tip kuramı ($\lambda$-hesabı, bağımlı tipler, \Lean{}'in
  aksiyomları) konularında araştırmacıya başvuru kaynağı olmayı hedefler.
  Bu bölümlerdeki \emph{Araştırma Notu} kutuları güncel literatüre ve açık
  problemlere yönlendirir. Kısım~V'teki on iki proje ise okuru kendi
  formalizasyon çalışmasına hazırlar.
\end{description}

\section*{İki cilt}

Elinizdeki kitap, iki ciltlik bir dizinin ilk cildidir. \textbf{Cilt~I:
Temeller}, sıfırdan başlayıp önermeler mantığının metateorisine ve
\Lean{}'in kendi tip kuramına kadar ilerler; kendi başına tamamlanmış bir
bütündür. \textbf{Cilt~II: İleri Konular} ise birinci mertebe mantığın
formalizasyonu, kip mantığı, Gödel'in eksiklik teoremleri, küme kuramı,
otomasyon, metaprogramlama ve yapay zekâ destekli teorem ispatını ele
alacaktır. Cilt~I'de ileriye dönük göndermeler ``Cilt~II, Bölüm~$n$''
biçiminde yazılmıştır; Cilt~II'nin planı, Cilt~I'in sonundaki
\emph{Sonsöz}'de yer alır.

\section*{Kitabın omurgası}

Kitabı bir arada tutan tek bir fikir vardır: \textbf{Curry--Howard eşlemesi}.
Bir önerme bir \emph{tür}, o önermenin ispatı ise o türden bir \emph{program}
olarak görülebilir. ``$P$ ise $Q$'' önermesinin bir ispatı, $P$'nin
kanıtını alıp $Q$'nun kanıtını üreten bir fonksiyondur. Bu bakış açısı ilk
bölümlerde sezgisel olarak ortaya çıkar, Kısım~IV'te ise kesin bir matematiksel
teoreme dönüşür.

İkinci omurga, \textbf{yapıcı ve klasik mantık} arasındaki ayrımdır. \Lean{}
varsayılan olarak yapıcıdır: Üçüncü hâlin olanaksızlığını kullanmak
istediğinizde bunu açıkça söylemeniz gerekir. Bu durum, klasik ders
kitaplarında gözden kaçan bir soruyu okurun önüne koyar: \emph{Bir şeyin
yanlış olmadığını göstermek, doğru olduğunu göstermek midir?}

\section*{Eşlik eden Lean projesi}

Kitaptaki \Lean{} kodlarının tamamı, kitapla birlikte dağıtılan
\lc{ModernMantik} projesinden alınır. Kitabın \LaTeX{} kaynağı kodu bu
dosyalardan doğrudan okur; böylece kitapta basılı olan kod ile derlenen kod
her zaman aynıdır. Proje, otomatik bir sürekli tümleştirme (CI) iş akışıyla
her değişiklikte yeniden derlenir.

\vspace{1em}
\begin{flushright}
\yazar\\
\textit{Ekim 2026}
\end{flushright}

\chapter*{Bu Kitap Nasıl Okunur?}
\addcontentsline{toc}{chapter}{Bu Kitap Nasıl Okunur?}
\markboth{Bu Kitap Nasıl Okunur?}{Bu Kitap Nasıl Okunur?}

\section*{Okuma yolları}

Cilt~I beş kısımdan oluşur. Her okurun bütün kısımları sırayla okuması
gerekmez. \Cref{fig:okuma-yollari} üç tipik okuma yolunu gösteriyor.

\begin{figure}[htbp]
\centering
\begin{tikzpicture}[
    kisim/.style={draw, thick, rounded corners=4pt, minimum width=2.15cm,
                  minimum height=1.3cm, align=center, font=\baslikfont\small},
    yol/.style={-{Stealth[length=2.6mm]}, line width=1.3pt},
    etiket/.style={font=\baslikfont\footnotesize\bfseries, anchor=west},
  ]
  \node[kisim, fill=tanimbg, draw=anarenk]          (k1) at (0,0)   {I\\Temeller};
  \node[kisim, fill=tanimbg, draw=anarenk]          (k2) at (2.6,0) {II\\Yüklemler};
  \node[kisim, fill=teorembg, draw=teoremframe]     (k3) at (5.2,0) {III\\Metamantık};
  \node[kisim, fill=teorembg, draw=teoremframe]     (k4) at (7.8,0) {IV\\Tip Kuramı};
  \node[kisim, fill=ornekbg, draw=ornekframe]         (k5) at (10.4,0){V\\Projeler};
  \node[kisim, fill=arastirmabg, draw=arastirmaframe, dashed] (k6) at (13.0,0){Cilt II\\İleri Konular};

  % Lisans dersi
  \draw[yol, ikincil] ([yshift=0.95cm]k1.west) -- ([yshift=0.95cm]k2.east);
  \node[etiket, text=ikincil] at ([xshift=0.15cm, yshift=0.95cm]k2.east)
       {Lisans dersi (1 dönem)};
  % Yüksek lisans
  \draw[yol, teoremframe] ([yshift=-0.95cm]k1.west) -- ([yshift=-0.95cm]k4.east);
  \node[etiket, text=teoremframe] at ([xshift=0.15cm, yshift=-0.95cm]k4.east)
       {Yüksek lisans (2 dönem) + projeler};
  % Araştırmacı
  \draw[yol, arastirmaframe] ([yshift=-1.55cm]k2.west) -- ([yshift=-1.55cm]k4.east);
  \draw[yol, arastirmaframe, dashed] ([yshift=-1.55cm]k4.east) -- ([yshift=-1.55cm]k6.east);
  \node[etiket, text=arastirmaframe, anchor=north] at ([yshift=-1.65cm]k4.south)
       {Araştırmacı: Bölüm 1--2'yi gözden geçirip Kısım III--IV'e, ardından Cilt II'ye};
\end{tikzpicture}
\caption{Cilt~I'in kısımları ve önerilen okuma yolları. Kesikli kutu, ayrı bir kitap olarak hazırlanan Cilt~II'dir.}
\label{fig:okuma-yollari}
\end{figure}

\section*{Kutular ve renkler}

Kitap boyunca aynı renk dili kullanılır. Bir kutunun rengine bakarak ne tür
bir içerikle karşılaştığınızı hemen anlayabilirsiniz.

\begin{center}
\renewcommand{\arraystretch}{1.5}
\begin{tabular}{>{\centering\arraybackslash}m{0.9cm} >{\raggedright\arraybackslash}m{3.2cm} >{\raggedright\arraybackslash}m{8.6cm}}
\tikz\fill[tanimframe, rounded corners=2pt] (0,0) rectangle (0.6,0.4); & \textbf{Tanım} & Yeni bir kavramın kesin tanımı. \\
\tikz\fill[teoremframe, rounded corners=2pt] (0,0) rectangle (0.6,0.4); & \textbf{Teorem} & Kanıtlanmış bir önerme; çoğunun \Lean{} ispatı da verilir. \\
\tikz\fill[ornekframe, rounded corners=2pt] (0,0) rectangle (0.6,0.4); & \textbf{Örnek} & Kavramın somut bir uygulaması. \\
\tikz\fill[leanframe, rounded corners=2pt] (0,0) rectangle (0.6,0.4); & \textbf{Lean kodu} & Eşlik eden projeden alınan, derlenmiş kod. \\
\tikz\fill[hedefframe, rounded corners=2pt] (0,0) rectangle (0.6,0.4); & \textbf{Hedef durumu} & \Lean{}'in o anda gösterdiği hipotezler ve hedef. \\
\tikz\fill[fikirframe, rounded corners=2pt] (0,0) rectangle (0.6,0.4); & \textbf{Önemli Fikir} & Bölümün özünü oluşturan kavrayış. \\
\tikz\fill[dikkatframe, rounded corners=2pt] (0,0) rectangle (0.6,0.4); & \textbf{Dikkat} & Yeni başlayanların sık yaptığı hatalar. \\
\tikz\fill[tarihframe, rounded corners=2pt] (0,0) rectangle (0.6,0.4); & \textbf{Tarihçe} & Kavramın tarihsel arka planı. \\
\tikz\fill[arastirmaframe, rounded corners=2pt] (0,0) rectangle (0.6,0.4); & \textbf{Araştırma Notu} & Güncel literatür, açık problemler, ileri okuma. \\
\end{tabular}
\end{center}

\section*{Lean kodunun renkleri}

\begin{leankod}[Renk örneği]
theorem ornek (P Q : Prop) (hp : P) (hpq : P → Q) : Q := by
  apply hpq   -- taktikler pembe
  exact hp    -- hipotez adları turuncu
#check ornek  -- komutlar mor
\end{leankod}

\noindent
Anahtar kelimeler (\leankw{theorem}, \leankw{fun}, \leankw{by}) mor,
taktikler (\leantac{apply}, \leantac{exact}) pembe, türler
(\leantype{Prop}) camgöbeği, hipotez adları (\leanhyp{hp}) turuncu,
mantık sembolleri (\lc{\textcolor{leansym}{→ ∧ ∨ ⊢}}) mavi, yorumlar
yeşil italik yazılır.

\section*{Alıştırmaların zorluk dereceleri}

\begin{itemize}
  \item \zorluk{1}\quad Temel: Bölümdeki örneklerin doğrudan uygulaması.
  \item \zorluk{2}\quad Orta: Birkaç fikrin birleştirilmesini gerektirir.
  \item \zorluk{3}\quad İleri: Yaratıcılık ya da sonraki bölümlerden bir fikir gerektirir.
\end{itemize}

\tableofcontents

\mainmatter
% ---------------------------------------------------------------------
\part{Temeller}
% =====================================================================
%  Bölüm 1 — Mantık Nedir, İspat Asistanı Nedir?
% =====================================================================
\chapter{Mantık Nedir, İspat Asistanı Nedir?}\label{bol:giris}

\begin{bolumhedefleri}
\begin{itemize}
  \item Bir çıkarımın \emph{geçerli} olmasının ne demek olduğunu ve
        geçerlilik ile doğruluk arasındaki farkı,
  \item Mantığın Aristoteles'ten yapay zekâya uzanan kısa tarihini,
  \item Bir ispat asistanının nasıl çalıştığını ve neden güvenilir olduğunu,
  \item \Lean{}'i kurmayı, \lc{\textcolor{leancmd}{\bfseries\#check}} ve
        \lc{\textcolor{leancmd}{\bfseries\#eval}} komutlarını kullanmayı,
  \item İlk teoreminizi yazmayı ve \Lean{}'in hata mesajlarını okumayı.
\end{itemize}
\end{bolumhedefleri}

% ---------------------------------------------------------------------
\section{Mantık nedir?}\label{sec:mantik-nedir}

Mantık, \emph{doğru akıl yürütmenin} bilimidir. Daha kesin söylersek, mantık
hangi sonuçların hangi öncüllerden \emph{zorunlu olarak} çıktığını inceler.
Bu tanımdaki kilit kelime ``zorunlu olarak''dır. Mantık, öncüllerin doğru
olup olmadığıyla ilgilenmez; öncüller doğru \emph{olsaydı} sonucun da doğru
olmak zorunda olup olmadığıyla ilgilenir.

\begin{tanim}{Çıkarım ve geçerlilik}{cikarim}
Bir \textbf{çıkarım}, bir öncüller kümesi ile bir sonuçtan oluşur.
Öncüllerin hepsinin doğru olduğu her durumda sonuç da doğruysa,
çıkarım \textbf{geçerlidir}. Geçerli olup öncülleri de gerçekten doğru
olan bir çıkarıma \textbf{sağlam} çıkarım denir.
\end{tanim}

\begin{ornek}{Geçerlilik ve doğruluk birbirinden bağımsızdır}{gecerlilik}
\begin{center}
\renewcommand{\arraystretch}{1.3}
\begin{tabular}{p{0.44\linewidth} p{0.44\linewidth}}
\textbf{(a) Geçerli ve sağlam}
& \textbf{(b) Geçerli ama sağlam değil} \\
Bütün insanlar ölümlüdür. & Bütün kuşlar uçar. \\
Sokrates bir insandır. & Penguen bir kuştur. \\
\cmidrule(r){1-1}\cmidrule(l){2-2}
O hâlde Sokrates ölümlüdür. & O hâlde penguen uçar. \\
\end{tabular}
\end{center}
İki çıkarımın \emph{biçimi} aynıdır: ``Bütün $A$'lar $B$'dir; $x$ bir
$A$'dır; o hâlde $x$ bir $B$'dir.'' Bu biçim geçerlidir. (b)'deki sonucun
yanlış olması mantığın değil, ilk öncülün suçudur.
\end{ornek}

Bu örnek, modern mantığın temel fikrini gösteriyor: \textbf{geçerlilik,
içerikle değil biçimle ilgilidir}. Sokrates'in ya da penguenin ne olduğunu
bilmemize gerek yoktur; yalnızca çıkarımın iskeletine bakarız. Bu iskeleti
sembollerle yazdığımızda mantık \emph{formal} (biçimsel) mantık olur:
\[
  \frac{\forall x\,(A(x) \to B(x)) \qquad A(s)}{B(s)}
\]
Kesrin üstündekiler öncüller, altındaki sonuçtur. Kitap boyunca bu gösterimi
çok kullanacağız.

\begin{onemlifikir}
Biçimsel mantıkta bir çıkarımın geçerliliğini denetlemek \emph{mekanik} bir
iştir: Sembollerin anlamını bilmeden, yalnızca kurallara bakarak yapılabilir.
Mekanik olan her iş, bir bilgisayara yaptırılabilir. İspat asistanlarının
varlık nedeni budur.
\end{onemlifikir}

% ---------------------------------------------------------------------
\section{Kısa bir tarihçe}\label{sec:tarihce}

\Cref{fig:zaman-cizelgesi}, bu kitapta karşılaşacağımız fikirlerin
tarihsel sırasını gösteriyor. Dikkat edilirse çizelge üç döneme ayrılıyor:
mantığın \emph{doğuşu}, mantığın \emph{matematikleşmesi} ve mantığın
\emph{makineleşmesi}.

\begin{figure}[htbp]
\centering
\begin{tikzpicture}[
    olay/.style={font=\baslikfont\scriptsize, align=center, text width=2.3cm},
    yil/.style={font=\baslikfont\scriptsize\bfseries},
    nokta/.style={circle, fill=#1, inner sep=2.2pt},
  ]
  % Dönem şeritleri
  \fill[ornekbg]    (0,-0.25) rectangle (2.2,0.25);
  \fill[tanimbg]    (2.2,-0.25) rectangle (8.2,0.25);
  \fill[fikirbg]    (8.2,-0.25) rectangle (14.6,0.25);
  \draw[-{Stealth[length=3mm]}, very thick, koyu] (0,0) -- (14.9,0);
  \node[font=\baslikfont\scriptsize\bfseries, text=ornekframe] at (1.1,-3.0) {Doğuş};
  \node[font=\baslikfont\scriptsize\bfseries, text=anarenk]    at (5.2,-3.0) {Matematikleşme};
  \node[font=\baslikfont\scriptsize\bfseries, text=ikincil]    at (11.4,-3.0) {Makineleşme};

  % Olaylar: (x, yıl, metin, yukarı/aşağı, renk)
  \foreach \x/\y/\t/\d/\c in {
      0.6/MÖ 350/Aristoteles:\\kıyaslar/1/ornekframe,
      2.6/1854/Boole:\\düşünce cebiri/-1/anarenk,
      3.6/1879/Frege:\\Begriffsschrift/1/anarenk,
      4.8/1910/Russell–Whitehead:\\Principia/-1/anarenk,
      5.9/1931/Gödel:\\eksiklik/1/anarenk,
      7.0/1935/Gentzen:\\doğal çıkarım/-1/anarenk,
      8.0/1940/Church:\\tip kuramı/1/anarenk,
      9.1/1967/de Bruijn:\\Automath/-1/ikincil,
      10.0/1969/Howard:\\formüller = tipler/1/ikincil,
      11.0/1988/Coquand–Huet:\\CoC/-1/ikincil,
      12.0/2015/Lean/1/ikincil,
      13.0/2021/Lean 4/-1/ikincil,
      14.2/2024+/YZ destekli\\ispatçılar/1/vurgu}
  {
    \node[nokta=\c] at (\x,0) {};
    \draw[\c, thin] (\x,0) -- (\x,{0.85*\d});
    \node[yil, text=\c] at (\x,{1.05*\d}) {\y};
    \node[olay, text width=1.5cm] at (\x,{1.62*\d+0.38*min(\d,0)}) {\t};
  }
\end{tikzpicture}
\caption{Bu kitaptaki fikirlerin kısa tarihi. Yıllar, ilgili eserin
yayımlandığı ya da sistemin duyurulduğu yıllardır; Howard'ın notları 1969'da
dolaşıma girmiş, 1980'de yayımlanmıştır.}
\label{fig:zaman-cizelgesi}
\end{figure}

\begin{tarihce}[Tarihçe: Frege ve kavram yazısı]
Gottlob Frege 1879'da yayımladığı \emph{Begriffsschrift} (Kavram Yazısı)
adlı küçük kitapta, niceleyicileri (\emph{her}, \emph{bazı}) içeren ilk
tam biçimsel mantık dilini kurdu \citep{frege1879}. Frege'nin amacı,
aritmetiğin tamamen mantıktan türetilebileceğini göstermekti. Bu program
Russell paradoksuyla sarsıldı; ancak Frege'nin kurduğu dil, bugün
kullandığımız yüklemler mantığının doğrudan atasıdır.
\end{tarihce}

% ---------------------------------------------------------------------
\section{Neden bilgisayarla ispat?}\label{sec:neden}

Matematikçiler binlerce yıldır ispatlarını kâğıt üzerinde yazıp birbirlerine
denetletiyorlar. Bu sistem şaşırtıcı derecede iyi çalışır, ama kusursuz
değildir. İspatlar uzadıkça ve karmaşıklaştıkça, hakemlerin her adımı
denetlemesi olanaksızlaşır. Yirminci yüzyılın sonunda bu sorun somut olarak
ortaya çıktı:

\begin{itemize}
  \item \textbf{Dört renk teoremi} (1976), ilk kez büyük ölçüde bilgisayar
  hesabına dayanan bir ispatla kanıtlandı. İspatın tamamı daha sonra Coq
  ispat asistanında baştan sona formalize edildi.
  \item \textbf{Kepler sanısı}nın ispatı o kadar karmaşıktı ki hakemler
  doğruluğundan ``yüzde doksan dokuz emin'' olabildiklerini bildirdi. Bunun
  üzerine başlatılan \emph{Flyspeck} projesi, ispatı HOL Light ve Isabelle
  ispat asistanlarında tamamen doğruladı.
\end{itemize}

Bugün ispat asistanları üç farklı amaca hizmet ediyor:

\begin{center}
\begin{tikzpicture}[
    kutu/.style={draw=#1, fill=#1!8, thick, rounded corners=4pt,
                 text width=4.1cm, minimum height=2.4cm, align=center,
                 font=\small},
  ]
  \node[kutu=anarenk] (a) at (0,0)
    {\baslikfont\bfseries\color{anarenk} Matematiği doğrulamak\\[3pt]
     \normalfont Uzun ve karmaşık ispatların her adımını denetlemek};
  \node[kutu=ikincil] (b) at (4.7,0)
    {\baslikfont\bfseries\color{ikincil} Yazılımı doğrulamak\\[3pt]
     \normalfont Derleyicilerin, işletim sistemlerinin, şifreleme
     kodlarının doğruluğunu ispatlamak};
  \node[kutu=vurgu] (c) at (9.4,0)
    {\baslikfont\bfseries\color{vurgu} Öğrenmek ve öğretmek\\[3pt]
     \normalfont Her adımı anında geri bildirimle öğrenmek};
\end{tikzpicture}
\end{center}

Bu kitap üçüncü amaçla başlıyor, birinciye doğru ilerliyor. İkinci amaç
(yazılım doğrulama) kitabın kapsamı dışında kalsa da, öğreneceğiniz her
teknik orada da doğrudan kullanılabilir.

% ---------------------------------------------------------------------
\section{Bir ispat asistanı nasıl çalışır?}\label{sec:nasil-calisir}

Bir ispat asistanına ``bu teorem doğrudur'' dediğinizde, size
\emph{güvenmez}. Sizden bir ispat ister ve o ispatı adım adım denetler.
\Lean{}'de bu süreç iki katmanda gerçekleşir (\cref{fig:lean-katmanlar}).

\begin{figure}[htbp]
\centering
\begin{tikzpicture}[
    katman/.style={draw=#1, fill=#1!10, thick, rounded corners=4pt,
                   minimum width=3.0cm, minimum height=1.5cm, align=center,
                   font=\baslikfont\small},
    ok/.style={-{Stealth[length=3mm]}, very thick, koyu!70},
    aciklama/.style={font=\scriptsize, text=black!65, align=center,
                     text width=3.0cm},
  ]
  \node[katman=ornekframe] (k) at (0,0) {\bfseries Siz\\ \small taktikler, terimler};
  \node[katman=anarenk]    (e) at (4.9,0) {\bfseries Elaboratör\\ \small büyük, akıllı};
  \node[katman=teoremframe](c) at (9.8,0) {\bfseries Çekirdek\\ \small küçük, güvenilir};
  \node[circle, draw=bittiframe, fill=bittibg, very thick, minimum size=1.1cm,
        font=\sembolfont\Large\bfseries\color{bittiframe}] (t) at (12.6,0) {✓};

  \draw[ok] (k) -- node[above, font=\scriptsize]{kaynak kod} (e);
  \draw[ok] (e) -- node[above, font=\scriptsize]{ispat terimi} (c);
  \draw[ok] (c) -- (t);

  \node[aciklama, below=0.25cm of k] {Okunabilir, kısa ispatlar yazarsınız.};
  \node[aciklama, below=0.25cm of e] {Taktikleri çalıştırır, eksikleri
        tamamlar, ayrıntılı bir ispat terimi üretir.};
  \node[aciklama, below=0.25cm of c] {Yalnızca terimin türünü denetler.
        Hata yaparsa bunu \emph{yakalar}.};

  \draw[decorate, decoration={brace, amplitude=6pt, mirror}, thick, black!40]
    ([yshift=-1.75cm]e.south west) -- ([yshift=-1.75cm]e.south east)
    node[midway, below=8pt, font=\scriptsize, text=black!60, align=center]
    {hata yapabilir\\(ama zarar veremez)};
\end{tikzpicture}
\caption{\Lean{}'in iki katmanlı mimarisi. Güven yalnızca küçük çekirdeğe
dayanır.}
\label{fig:lean-katmanlar}
\end{figure}

Bu tasarımın özü şudur: Elaboratör çok büyük ve karmaşık bir programdır,
dolayısıyla içinde hata olabilir. Ama elaboratörün ürettiği her ispat
terimi, ayrıca küçük bir \emph{çekirdek} tarafından denetlenir. Çekirdek
yanlış bir ispatı kabul etmedikçe, elaboratördeki bir hata yanlış bir
teoremin ``ispatlanmasına'' yol açamaz.

\begin{tanim}{de Bruijn ölçütü}{debruijn}
Bir ispat asistanının ürettiği ispatlar, \textbf{küçük ve bağımsız bir
denetleyici} tarafından yeniden doğrulanabiliyorsa, o sistem \textbf{de
Bruijn ölçütünü} sağlar. Böylece sisteme duyulan güven, yalnızca bu küçük
denetleyiciye duyulan güvene indirgenir.
\end{tanim}

\begin{tarihce}[Tarihçe: Automath]
Hollandalı matematikçi Nicolaas Govert de Bruijn, 1967'de başlattığı
\emph{Automath} projesiyle matematiksel metinleri bilgisayarda denetlemenin
ilk ciddi örneğini verdi \citep{debruijn1970}. Automath ekibi 1970'lerde
Landau'nun \emph{Grundlagen der Analysis} kitabının tamamını formalize etti.
Automath'ta ortaya çıkan ``önermeleri tür olarak görme'' fikri, Howard'ın
çalışmalarıyla birlikte bu kitabın omurgasını oluşturan Curry--Howard
eşlemesinin kaynaklarından biridir.
\end{tarihce}

% ---------------------------------------------------------------------
\section{Lean 4'ü kurmak}\label{sec:kurulum}

\Lean{}'i denemek için hiçbir şey kurmanız gerekmez: Lean'in resmî sitesindeki
\emph{çevrimiçi editör}, tarayıcıda çalışır. Ancak kitaptaki projeyle
çalışmak için yerel bir kurulum öneriyoruz.

\begin{enumerate}[adimlar]
  \item \textbf{Visual Studio Code}'u kurun.
  \item VS Code'un eklenti mağazasından \textbf{Lean 4} eklentisini kurun.
  \item Eklentinin karşılama sayfasındaki adımları izleyin. Eklenti, Lean'in
        sürüm yöneticisi \lc{elan}'ı ve derleme aracı \lc{lake}'i sizin
        için kurar.
  \item Kitabın eşlik eden projesini indirip klasörünü VS Code'da açın ve
        bir terminalde şu komutu çalıştırın:
\begin{center}
\lc{lake build}
\end{center}
  \item \lc{ModernMantik/Bolum01.lean} dosyasını açın. Sağ tarafta
        \emph{Lean Infoview} paneli görünmelidir.
\end{enumerate}

\begin{dikkat}[Dikkat: Kurulum adımları değişebilir]
\Lean{} ve araçları hızla gelişiyor. Yukarıdaki adımlar bu kitabın yazıldığı
dönemde geçerlidir. Bir adım beklediğiniz gibi çalışmazsa, Lean'in resmî
sitesindeki güncel kurulum sayfasına başvurun. Eşlik eden projenin
\lc{lean-toolchain} dosyası, kitabın test edildiği Lean sürümünü belirtir.
\end{dikkat}

% ---------------------------------------------------------------------
\section{İlk adımlar: \#check ve \#eval}\label{sec:check-eval}

\Lean{} hem bir programlama dili hem de bir ispat dilidir. Bu ikili yapıyı
en iyi iki komutla görebiliriz: \lc{\textcolor{leancmd}{\bfseries\#eval}}
bir ifadeyi \emph{hesaplar}, \lc{\textcolor{leancmd}{\bfseries\#check}} ise
bir ifadenin \emph{türünü} söyler.

\leandosya[İlk komutlar]{Bolum01}{check_eval}

Son satır önemli bir şeyi gösteriyor: \lc{2 + 2 = 5} \emph{yanlış} bir
önermedir, ama \Lean{} onu itirazsız kabul eder ve türünün \leantype{Prop}
olduğunu söyler. Çünkü \leantype{Prop}, \emph{önermelerin} türüdür, yani
doğru ya da yanlış olabilen ifadelerin. Bir ifadenin önerme olması, doğru
olduğu anlamına gelmez.

\begin{onemlifikir}[Önemli Fikir: Bir önerme ile onun kanıtı farklı şeylerdir]
\lc{2 + 2 = 4} ve \lc{2 + 2 = 5} ifadelerinin ikisi de birer önermedir;
ikisinin de türü \leantype{Prop}'tur. Aralarındaki fark şudur: Birincisinin
bir \emph{kanıtını} inşa edebiliriz, ikincisinin edemeyiz. \Lean{}'de bir
önermenin doğru olduğunu söylemek, o önermenin bir kanıtını sunmak demektir.
\end{onemlifikir}

\Cref{tab:tur-benzetmesi}, bu fikrin \Lean{}'de nasıl yazıldığını
gösteriyor. Sol sütundaki tanıdık ifadelerle sağ sütundakiler arasında
yapısal bir benzerlik var.

\begin{table}[htbp]
\centering
\renewcommand{\arraystretch}{1.4}
\rowcolors{2}{tablosatir}{white}
\begin{tabular}{>{\raggedright\arraybackslash}p{4.2cm}
                >{\raggedright\arraybackslash}p{4.2cm}}
\rowcolor{tabloust}
\tablobaslik{Veri dünyası} & \tablobaslik{Mantık dünyası} \\
\lc{Nat} bir türdür                & \lc{P} bir önermedir (\lc{P : Prop}) \\
\lc{3 : Nat}                      & \lc{hp : P} \\
``3, bir doğal sayıdır''           & ``\lc{hp}, $P$'nin bir kanıtıdır'' \\
\lc{Nat} boş değildir              & $P$ doğrudur \\
\end{tabular}
\caption{Türler ve önermeler arasındaki benzerlik. Bu tablo, Curry--Howard
eşlemesinin ilk görünümüdür.}
\label{tab:tur-benzetmesi}
\end{table}

Benzetmenin son satırı en derin olanıdır. Bir türün bir elemanını
gösterebiliyorsak o tür ``dolu''dur; bir önermenin bir kanıtını
gösterebiliyorsak o önerme doğrudur. Bu fikri \cref{bol:onermeler1}'de
adım adım geliştirecek, \cref{bol:lambda}'te ise kesin bir teoreme
dönüştüreceğiz.

% ---------------------------------------------------------------------
\section{İlk teorem}\label{sec:ilk-teorem}

Şimdi ilk teoremimizi yazalım. Teoremimiz şunu söylüyor: \emph{Elimizde
$P$'nin bir kanıtı varsa, $P$ doğrudur.} Bu, olabilecek en basit teoremdir;
ama \Lean{}'in bir teoremi nasıl okuduğunu bütün parçalarıyla gösterir.

\leandosya[İlk teorem]{Bolum01}{ilk_teorem}

\Cref{fig:teorem-anatomisi}, bu tek satırın anatomisini gösteriyor.

\begin{figure}[htbp]
\centering
\begin{tikzpicture}[
    parca/.style={font=\leanfont\normalsize, inner sep=2pt},
    etiket/.style={font=\baslikfont\scriptsize, align=center, text width=2.6cm},
  ]
  \node[parca, text=leankw]   (a) at (0,0)    {\bfseries theorem};
  \node[parca]                (b) at (2.05,0) {ilk\_teorem};
  \node[parca]                (c) at (4.35,0) {(P : \textcolor{leantype}{\bfseries Prop})};
  \node[parca]                (d) at (6.85,0) {(\textcolor{leanhyp}{\bfseries hp} : P)};
  \node[parca]                (e) at (8.45,0) {: P};
  \node[parca]                (f) at (9.75,0) {:= \textcolor{leanhyp}{\bfseries hp}};

  \foreach \n/\r in {a/leankw, b/koyu, c/leantype, d/leanhyp, e/vurgu, f/ikincil}
    \draw[\r, thick, rounded corners=2pt] (\n.south west) rectangle (\n.north east);

  \node[etiket, text=leankw]   (ea) at (0,-1.35)   {anahtar kelime};
  \node[etiket, text=koyu]     (eb) at (2.05,1.35) {teoremin adı};
  \node[etiket, text=leantype] (ec) at (4.35,-1.35){``$P$ bir önerme olsun''};
  \node[etiket, text=leanhyp]  (ed) at (6.85,1.35) {varsayım: ``\lc{hp}, $P$'nin kanıtı''};
  \node[etiket, text=vurgu]    (ee) at (8.45,-1.35){iddia: ``$P$ doğrudur''};
  \node[etiket, text=ikincil]  (ef) at (10.2,1.35) {ispat: ``kanıt şudur: \lc{hp}''};

  \foreach \n/\m/\r in {a/ea/leankw, b/eb/koyu, c/ec/leantype,
                        d/ed/leanhyp, e/ee/vurgu, f/ef/ikincil}
    \draw[\r, -{Stealth[length=2mm]}] (\m) -- (\n);
\end{tikzpicture}
\caption{Bir \Lean{} teoreminin parçaları.}
\label{fig:teorem-anatomisi}
\end{figure}

Okunuşu şöyledir: ``$P$ bir önerme olsun ve elimizde $P$'nin \lc{hp} adlı
bir kanıtı bulunsun. O zaman $P$ doğrudur; kanıtı da \lc{hp}'nin
kendisidir.'' İki noktanın ardından gelen \lc{:=} işareti ``ispatı şudur''
diye okunur.

\subsection{Taktik modu}

Yukarıdaki ispat bir \emph{terimdir}: kanıtı doğrudan yazdık. Daha uzun
ispatlarda kanıtı tek seferde yazmak zordur. Bunun yerine \leankw{by}
anahtar kelimesiyle \emph{taktik moduna} geçer, ispatı adım adım inşa
ederiz:

\leandosya[Aynı teorem, taktik moduyla]{Bolum01}{ilk_teorem_taktik}

Taktik modunda \Lean{}, imlecin bulunduğu her noktada \emph{hedef durumunu}
Infoview panelinde gösterir (\cref{fig:infoview}). \leankw{by} kelimesinden
hemen sonra durum şöyledir:

\begin{hedefdurumu}[\texttt{by} sonrası]
P : Prop
hp : P
⊢ P
\end{hedefdurumu}

\noindent \lc{⊢} işaretinin üstündekiler elimizdekilerdir (\emph{bağlam}),
altındaki ise kanıtlamamız gereken şeydir (\emph{hedef}). \leantac{exact hp}
taktiği ``hedefin kanıtı tam olarak \lc{hp}'dir'' der ve ispatı bitirir:

\ispatbitti

\begin{figure}[htbp]
\centering
\begin{tikzpicture}[font=\leanfont\scriptsize]
  % Editör penceresi
  \draw[black!30, fill=white, rounded corners=4pt] (0,0) rectangle (7.6,4.2);
  \fill[black!8, rounded corners=4pt] (0,3.75) rectangle (7.6,4.2);
  \fill[black!8] (0,3.75) rectangle (7.6,3.95);
  \foreach \x/\c in {0.25/hataframe, 0.5/hedefframe, 0.75/bittiframe}
    \fill[fill=\c!80] (\x,3.975) circle (0.07);
  \node[anchor=west, font=\baslikfont\scriptsize, text=black!60] at (1.0,3.975) {Bolum01.lean};
  % Kod satırları
  \foreach \i/\t in {
     1/{\textcolor{leankw}{\bfseries theorem} ilk\_teorem' (P : \textcolor{leantype}{\bfseries Prop})},
     2/{\hspace{3.2em}(\textcolor{leanhyp}{\bfseries hp} : P) : P := \textcolor{leankw}{\bfseries by}},
     3/{\hspace{1em}\textcolor{leantac}{\bfseries exact} \textcolor{leanhyp}{\bfseries hp}}}
  {
    \node[anchor=west, text=black!35] at (0.15,{3.5-0.42*\i}) {\i};
    \node[anchor=west] at (0.55,{3.5-0.42*\i}) {\t};
  }
  % İmleç
  \fill[anarenk] (1.42,1.98) rectangle (1.46,2.28);
  \node[anchor=west, font=\baslikfont\tiny, text=anarenk] at (1.65,1.72) {imleç burada};

  % Infoview
  \draw[black!30, fill=hedefbg!60, rounded corners=4pt] (8.0,0) rectangle (13.4,4.2);
  \fill[hedefframe!30, rounded corners=4pt] (8.0,3.75) rectangle (13.4,4.2);
  \fill[hedefframe!30] (8.0,3.75) rectangle (13.4,3.95);
  \node[anchor=west, font=\baslikfont\scriptsize\bfseries, text=black!70] at (8.15,3.975) {Lean Infoview};
  \node[anchor=west, font=\baslikfont\scriptsize\bfseries] at (8.2,3.3) {{\sembolfont ▾} Tactic state};
  \node[anchor=west, text=black!60] at (8.3,2.85) {1 goal};
  \node[anchor=west] at (8.3,2.4) {P : \textcolor{leantype}{Prop}};
  \node[anchor=west] at (8.3,2.0) {\textcolor{leanhyp}{hp} : P};
  \node[anchor=west] at (8.3,1.55) {\textcolor{leansym}{⊢} P};
  \draw[hedefframe, thick, rounded corners=2pt] (8.2,1.3) rectangle (9.4,1.8);
  \node[anchor=west, font=\baslikfont\tiny, text=hedefframe!80!black] at (9.5,1.55) {hedef};
  \draw[black!40, decorate, decoration={brace, amplitude=4pt}]
       (10.3,2.6) -- (10.3,1.85) node[midway, right=5pt, font=\baslikfont\tiny, text=black!60] {bağlam};

  \draw[-{Stealth[length=2.5mm]}, thick, anarenk, dashed] (3.1,2.13) to[bend left=18] (8.15,2.6);
\end{tikzpicture}
\caption{VS Code'da editör ve Lean Infoview paneli. İmleç \leantac{exact}
satırına gelmeden önceki hedef durumu sağda görünür.}
\label{fig:infoview}
\end{figure}

\subsection{Hesaplamayla ispat}

Bazı önermeler doğrudan hesaplanarak doğrulanabilir. \lc{rfl} (yansıma,
İng.\ \emph{reflexivity}) ``iki taraf hesaplandığında aynıdır'' demektir:

\leandosya[Hesaplamayla ispat]{Bolum01}{rfl_ornek}

Aynı şeyi \lc{2 + 2 = 5} için denerseniz \Lean{} itiraz eder; çünkü iki taraf
hesaplandığında aynı değildir. Bir önermenin doğru olmasının ``\Lean{}'in
kanıtı kabul etmesi'' anlamına geldiğini burada açıkça görüyoruz.

% ---------------------------------------------------------------------
\section{Hata mesajlarını okumak ve \texttt{sorry}}\label{sec:hatalar}

İspat asistanıyla çalışmanın büyük kısmı, hata mesajlarını okumaktan oluşur.
Bunda utanılacak bir şey yoktur: Her hata mesajı, \Lean{}'in sizinle
konuştuğu bir andır. Bilerek yanlış bir ispat yazalım:

\begin{leankod}[Yanlış bir ispat (derlenmez)]
theorem yanlis (P Q : Prop) (hp : P) : Q := hp
\end{leankod}

\Lean{} \lc{hp}'nin altını kırmızıyla çizer ve aşağıdakine benzer bir mesaj
gösterir. (Mesajın tam ifadesi \Lean{} sürümüne göre biraz değişebilir.)

\begin{leanhata}[Lean hata mesajı]
type mismatch
  hp
has type
  P : Prop
but is expected to have type
  Q : Prop
\end{leanhata}

Okunuşu: ``\lc{hp} ifadesinin türü $P$, ama burada $Q$ türünden bir şey
beklenmektedir.'' Mantık diliyle: ``Sen bana $P$'nin kanıtını verdin, ama
ben $Q$'nun kanıtını istedim.'' \Lean{}, $P$'den $Q$'nun çıkacağına dair
hiçbir gerekçe göremediği için ispatı reddeder. Haklıdır da: $P$ doğru,
$Q$ yanlış olabilir.

\paragraph{sorry: ``Buraya sonra döneceğim''}
Bir ispatın bir kısmını henüz yazamıyorsanız, o kısmın yerine \leantac{sorry}
yazabilirsiniz. \Lean{} dosyanın geri kalanını denetlemeye devam eder,
ama \leantac{sorry} içeren her teoremi bir uyarıyla işaretler:

\begin{leankod}[Yarım kalmış bir ispat]
theorem henuz_bitmedi (P Q : Prop) (hp : P) : P ∧ Q := by
  sorry
\end{leankod}

\begin{dikkat}
\leantac{sorry}, bir ispatı \emph{kapatmaz}; yalnızca ertelemeyi sağlar.
\leantac{sorry} içeren bir teorem kanıtlanmış sayılmaz. Bir teoremin
gerçekten kanıtlanıp kanıtlanmadığını \lc{\textcolor{leancmd}{\bfseries\#print axioms}}
komutuyla denetleyebilirsiniz: \leantac{sorry} kullanılmışsa listede
\lc{sorryAx} görünür. Bu komuta \cref{bol:aksiyomlar}'da yeniden döneceğiz.
\end{dikkat}

Kitabın alıştırma dosyalarının hepsi bu yöntemle hazırlanmıştır: Teoremin
ifadesi verilir, ispatın yerinde \leantac{sorry} durur. Sizin göreviniz
\leantac{sorry}'yi gerçek bir ispatla değiştirmektir.

% ---------------------------------------------------------------------
\begin{arastirmanotu}
\textbf{Lean'in kısa geçmişi.} Lean, Leonardo de Moura tarafından Microsoft
Research'te başlatıldı; ilk sistem tanımı 2015'te yayımlandı
\citep{demoura2015lean}. \Lean{}, öncüllerinden farklı olarak büyük ölçüde
\emph{kendi dilinde} yazılmıştır; böylece kullanıcılar dili kendi taktikleri
ve notasyonlarıyla genişletebilir \citep{demoura2021lean4}. Topluluk
tarafından geliştirilen matematik kütüphanesi \Mathlib{}, bugün en büyük
birleşik formal matematik kütüphanelerinden biridir \citep{mathlib2020}.

\textbf{İleri okuma.} \Lean{}'e İngilizce giriş için resmî kaynaklar
\citet{tpil4} ve \citet{logicandproof}'tur. Bilgisayar bilimi
perspektifinden bir giriş için \citet{hitchhiker2026}'ya, programlama
yönü için \citet{fpil}'e bakılabilir. \Lean{}'in tip kuramının kesin
tanımı \citet{carneiro2019}'dadır.
\end{arastirmanotu}

\begin{bolumozeti}
\begin{itemize}
  \item Mantık, çıkarımların \emph{biçimine} bakarak geçerliliklerini
        inceler. Geçerlilik ile doğruluk farklı kavramlardır.
  \item İspat asistanları ispatları mekanik olarak denetler. \Lean{}'de
        güven, küçük bir çekirdeğe dayanır (de Bruijn ölçütü).
  \item \lc{\textcolor{leancmd}{\bfseries\#check}} türü,
        \lc{\textcolor{leancmd}{\bfseries\#eval}} değeri gösterir.
        \leantype{Prop}, önermelerin türüdür.
  \item \lc{hp : P}, ``\lc{hp}, $P$'nin bir kanıtıdır'' diye okunur.
        Bir önermeyi ispatlamak, onun türünden bir terim inşa etmektir.
  \item Taktik modunda (\leankw{by}) \Lean{}, her adımda bağlamı ve hedefi
        gösterir. \leantac{exact} hedefi kapatır, \leantac{sorry} erteler.
\end{itemize}
\end{bolumozeti}

\begin{alistirmalar}
\alistirma{1} Aşağıdaki çıkarımların her biri için geçerli olup
olmadığını ve sağlam olup olmadığını belirleyin:
(a) ``Bütün kediler memelidir. Tekir bir kedidir. O hâlde Tekir bir
memelidir.''
(b) ``Bütün balıklar yüzer. Yunus yüzer. O hâlde yunus bir balıktır.''

\alistirma{1} Aşağıdaki komutların çıktısını önce tahmin edin, sonra
\Lean{}'de deneyin:
\lc{\#check 3}, \lc{\#eval 3 * 7}, \lc{\#check (3 < 7)},
\lc{\#check "merhaba"}, \lc{\#eval "mer" ++ "haba"}.

\alistirma{1} \lc{ilk\_teorem}'in ispatındaki \lc{hp}'yi silip yerine
\leantac{sorry} yazın. Infoview'da ne değişiyor? Sonra \lc{\#print axioms}
komutunu çalıştırın.

\alistirma{2} Aşağıdaki teoremi ispatlayın. Elinizdeki varsayımlardan
hangisi gereksiz?
\begin{center}
\lc{theorem ikinci (P Q : Prop) (hp : P) (hq : Q) : Q := sorry}
\end{center}

\alistirma{2} \cref{sec:hatalar}'teki \lc{yanlis} teoremini \Lean{}'de
yazın. Hata mesajını kendi cümlelerinizle açıklayın. Sonra varsayımlara
tek bir varsayım ekleyerek teoremi ispatlanabilir hâle getirin.

\alistirma{3} (Tartışma) de Bruijn ölçütü neden önemlidir? Bir ispat
asistanının elaboratöründe bir hata bulunduğunu düşünün. Bu hata hangi
koşullarda yanlış bir teoremin kabul edilmesine yol açabilir, hangi
koşullarda açamaz?
\end{alistirmalar}
   % Mantık nedir, ispat asistanı nedir?
% =====================================================================
%  Bölüm 2 — Önermeler Mantığı I: İmplikasyon, Ve, Veya
% =====================================================================
\chapter{Önermeler Mantığı I: İmplikasyon, Ve, Veya}\label{bol:onermeler1}

\begin{bolumhedefleri}
\begin{itemize}
  \item Önermeler mantığının bağlaçlarını ve \Lean{}'deki yazılışlarını,
  \item Doğal çıkarımın \emph{giriş} ve \emph{eleme} kuralları fikrini,
  \item $\to$, $\wedge$, $\vee$ için çıkarım kurallarını hem ağaç hem
        \Lean{} biçiminde,
  \item \leantac{intro}, \leantac{apply}, \leantac{exact},
        \leantac{constructor}, \leantac{obtain}, \leantac{cases},
        \leantac{rintro} taktiklerini,
  \item Her ispatı hem taktik hem terim modunda yazmayı ve
        Curry--Howard eşlemesinin ilk tablosunu.
\end{itemize}
\end{bolumhedefleri}

% ---------------------------------------------------------------------
\section{Önermeler ve bağlaçlar}\label{sec:baglaclar}

\emph{Önermeler mantığı}, önermelerin iç yapısına bakmaz. ``Yağmur
yağıyor'' gibi bir cümleyi bölünmez bir bütün olarak alır ve ona $P$ gibi
bir harf verir. Bu tür önermelere \emph{atomik} önerme denir. Atomik
önermeler \emph{bağlaçlarla} birleştirilerek \emph{bileşik} önermeler
oluşturulur. \Cref{tab:baglaclar} bu kitapta kullanacağımız bağlaçları
gösteriyor.

\begin{table}[htbp]
\centering
\renewcommand{\arraystretch}{1.45}
\rowcolors{2}{tablosatir}{white}
\begin{tabular}{>{\centering\arraybackslash}p{1.6cm}
                >{\raggedright\arraybackslash}p{3.2cm}
                >{\raggedright\arraybackslash}p{3.0cm}
                >{\centering\arraybackslash}p{1.7cm}
                >{\centering\arraybackslash}p{1.6cm}}
\rowcolor{tabloust}
\tablobaslik{Sembol} & \tablobaslik{Okunuşu} & \tablobaslik{Lean'de} &
\tablobaslik{Yazılış} & \tablobaslik{Bölüm} \\
$P \to Q$          & $P$ ise $Q$          & \lc{P → Q}  & \lc{\textbackslash to}  & \ref{bol:onermeler1} \\
$P \wedge Q$       & $P$ ve $Q$           & \lc{P ∧ Q}  & \lc{\textbackslash and} & \ref{bol:onermeler1} \\
$P \vee Q$         & $P$ veya $Q$         & \lc{P ∨ Q}  & \lc{\textbackslash or}  & \ref{bol:onermeler1} \\
$\neg P$           & $P$ değil            & \lc{¬P}     & \lc{\textbackslash not} & \ref{bol:degilleme} \\
$P \leftrightarrow Q$ & $P$ ancak ve ancak $Q$ & \lc{P ↔ Q} & \lc{\textbackslash iff} & \ref{bol:degilleme} \\
$\bot$             & çelişki, yanlış      & \lc{False}  & ---                     & \ref{bol:degilleme} \\
$\top$             & doğru                & \lc{True}   & ---                     & \ref{bol:degilleme} \\
\end{tabular}
\caption{Önermeler mantığının bağlaçları. Bu bölüm ilk üç satırı ele alıyor.}
\label{tab:baglaclar}
\end{table}

\begin{dikkat}[Dikkat: $\to$ sağa doğru birleşir]
\lc{P → Q → R} ifadesi \lc{P → (Q → R)} olarak okunur, \lc{(P → Q) → R}
olarak değil. Yani ``$P$ ise, ($Q$ ise $R$)''. Bu kural, birden fazla
varsayımı olan önermeleri parantezsiz yazmamızı sağlar: \lc{P → Q → R},
``$P$ ve $Q$ verilirse $R$'' demenin \Lean{}'deki doğal yoludur.
\end{dikkat}

% ---------------------------------------------------------------------
\section{Doğal çıkarım: giriş ve eleme}\label{sec:dogal-cikarim-fikri}

Gerhard Gentzen 1935'te, matematikçilerin gerçekte nasıl akıl
yürüttüğünü yansıtan bir ispat sistemi önerdi ve ona \emph{doğal çıkarım}
adını verdi \citep{gentzen1935}. Sistemin zarif fikri şudur: Her bağlaç
için iki tür kural vardır.

\begin{center}
\begin{tikzpicture}[
    kutu/.style={draw=#1, fill=#1!8, very thick, rounded corners=5pt,
                 text width=5.6cm, minimum height=2.5cm, align=center},
  ]
  \node[kutu=ikincil] (g) at (0,0) {%
    {\baslikfont\bfseries\large\color{ikincil} Giriş kuralı (G)}\\[4pt]
    Bir bağlacı içeren bir önermeyi\\ nasıl \textbf{kanıtlarım}?\\[3pt]
    {\small\color{black!60}``$P \wedge Q$'yu göstermek için\\ $P$'yi ve $Q$'yu göster.''}};
  \node[kutu=vurgu] (e) at (6.6,0) {%
    {\baslikfont\bfseries\large\color{vurgu} Eleme kuralı (E)}\\[4pt]
    Bir bağlacı içeren bir kanıtı\\ nasıl \textbf{kullanırım}?\\[3pt]
    {\small\color{black!60}``$P \wedge Q$ elimdeyse\\ $P$'yi de $Q$'yu da kullanabilirim.''}};
  \draw[{Stealth[length=3mm]}-{Stealth[length=3mm]}, very thick, black!40]
    (g.east) -- node[above, font=\baslikfont\scriptsize, text=black!60]{uyum}
    node[below, font=\baslikfont\scriptsize, text=black!60]{(harmony)} (e.west);
\end{tikzpicture}
\end{center}

Giriş ve eleme kuralları birbirini tam olarak dengeler: Eleme kuralları,
giriş kuralının koyduğundan ne fazlasını ne eksiğini geri alır. Bu ilkeye
\emph{uyum} (İng.\ \emph{harmony}) denir ve \cref{bol:dogal-cikarim}'te
ispat normalleştirmesinin temelini oluşturacaktır.

\Lean{}'de her kural bir \emph{taktiğe} karşılık gelir. Kabaca söylersek
giriş kuralları \emph{hedef} üzerinde, eleme kuralları ise
\emph{hipotezler} üzerinde çalışır:

\begin{onemlifikir}
\begin{tabular}{@{}l@{\quad}l}
\textcolor{ikincil}{\bfseries Giriş kuralları} & hedefi parçalar: ``Bunu göstermek için şunu göstereyim.'' \\[2pt]
\textcolor{vurgu}{\bfseries Eleme kuralları}  & hipotezi kullanır: ``Elimde bu var, ondan şu çıkar.''
\end{tabular}
\end{onemlifikir}

\paragraph{Çıkarım ağaçları nasıl okunur?}
Bir kural, çizginin üstündeki öncüllerden altındaki sonuca geçişi gösterir.
Çizginin sağındaki etiket kuralın adıdır. Köşeli parantez içindeki
$[P]^1$ biçimindeki varsayımlar \emph{geçici} varsayımlardır: İspatın bir
noktasında açılır ve üst simgede aynı numarayı taşıyan kuralda
\emph{kapatılırlar}. Kapatılmış bir varsayım artık sonucun dayandığı
varsayımlar arasında değildir.

% ---------------------------------------------------------------------
\section{İmplikasyon (\texorpdfstring{$\to$}{→})}\label{sec:implikasyon}

\subsection{Kurallar}

\begin{center}
\begin{tabular}{c@{\hspace{4em}}c}
\begin{prooftree}
  \hypo{[P]^{1}}
  \ellipsis{}{Q}
  \infer1[$\to$G$^{1}$]{P \to Q}
\end{prooftree}
&
\begin{prooftree}
  \hypo{P \to Q}
  \hypo{P}
  \infer2[$\to$E]{Q}
\end{prooftree}
\\[10pt]
\small\textcolor{ikincil}{\bfseries giriş:} \leantac{intro} &
\small\textcolor{vurgu}{\bfseries eleme (modus ponens):} \leantac{apply}, uygulama
\end{tabular}
\end{center}

\textbf{$\to$ girişi.} $P \to Q$'yu kanıtlamak için $P$'yi geçici olarak
varsayar, bu varsayımla $Q$'ya ulaşırız. Sonuca vardığımızda $P$
varsayımını kapatırız: Artık $P$'yi varsaymıyoruz; ``$P$ ise $Q$''
diyoruz. \Lean{}'de bu adım \leantac{intro} taktiğidir.

\textbf{$\to$ elemesi.} Elimizde $P \to Q$ ve $P$ varsa, $Q$'yu elde
ederiz. Bu, Orta Çağ'dan beri \emph{modus ponens} adıyla bilinen kuraldır.
\Lean{}'de iki biçimi vardır: Hedef $Q$ iken \leantac{apply hpq} hedefi
$P$'ye çevirir (geriye doğru); \lc{hpq hp} yazmak ise doğrudan $Q$'nun
kanıtını üretir (ileriye doğru).

\subsection{Özdeşlik: \texorpdfstring{$P \to P$}{P → P}}

En basit implikasyonla başlayalım.

\leandosya[Özdeşlik]{Bolum02}{ozdeslik}

\noindent \leantac{intro hp} adımından önce ve sonra hedef durumu şöyledir:

\begin{hedefgecisi}
\begin{hedefdurumu}[Önce]
P : Prop
⊢ P → P
\end{hedefdurumu}
\sonra
\begin{hedefdurumu}[\texttt{intro hp} sonrası]
P : Prop
hp : P
⊢ P
\end{hedefdurumu}
\end{hedefgecisi}

\leantac{intro}, implikasyonun sol tarafını hedeften alıp \emph{bağlama}
taşıdı. Bu tam olarak $\to$G kuralının yaptığı şeydir: $P$'yi geçici bir
varsayım olarak açtık. Aynı teoremin terim modundaki ispatı ise bir
fonksiyondur:

\leandosya[Özdeşlik, terim modu]{Bolum02}{ozdeslik_terim}

\noindent \lc{fun hp => hp}, ``bir \lc{hp} al, onu olduğu gibi geri ver''
demektir: özdeşlik fonksiyonu. $P$'nin kanıtını alıp $P$'nin kanıtını
döndüren bir fonksiyon, tam olarak $P \to P$'nin kanıtıdır.

\subsection{Modus ponens}

\leandosya[Modus ponens, taktik modu]{Bolum02}{modus_ponens}
\leandosya[Modus ponens, terim modu]{Bolum02}{modus_ponens_terim}

Terim modundaki ispat, Curry--Howard eşlemesinin belki de en çarpıcı
örneğidir: \textbf{implikasyonu kullanmak, bir fonksiyonu
çağırmaktır}. \lc{hpq}, $P$'nin kanıtını alıp $Q$'nun kanıtını üreten bir
fonksiyondur; \lc{hpq hp} ise bu fonksiyonu \lc{hp} argümanıyla çağırır.

\subsection{Gereksiz bir varsayım: \texorpdfstring{$P \to Q \to P$}{P → Q → P}}

\leandosya[Sabit fonksiyon]{Bolum02}{sabit}

Bu teorem, ``$P$ doğruysa, başka ne varsayarsak varsayalım $P$
doğrudur'' der. \leantac{intro hp \_} iki varsayımı birden içeri alır;
alt çizgi (\lc{\_}) ikinci varsayıma ad vermediğimizi, çünkü onu
kullanmayacağımızı belirtir. Terim olarak bu teorem, ikinci argümanını
görmezden gelen \emph{sabit fonksiyondur}: \lc{fun hp \_ => hp}.

\begin{tarihce}[Tarihçe: K ve S kombinatörleri]
$P \to Q \to P$ ve Alıştırma~\ref{alis:s-kombinator}'deki
$(P \to Q \to R) \to (P \to Q) \to P \to R$ önermeleri, Moses
Schönfinkel ve Haskell Curry'nin \emph{kombinatör mantığındaki} $K$ ve $S$
kombinatörlerinin türleridir \citep{curry1958}. Bu iki önerme, modus
ponens ile birlikte, implikasyon mantığının tamamını türetmeye yeter.
Curry, bu gözlemi 1930'larda yaptığında henüz ``Curry--Howard''
eşlemesinin ilk tohumunu atmış oluyordu.
\end{tarihce}

\subsection{Hipotetik kıyas}\label{sec:hipotetik-kiyas}

Şimdi birden fazla adım gerektiren ilk ispatımıza geçelim. Kural şudur:
\begin{center}
``$P$ ise $Q$'' doğruysa ve ``$Q$ ise $R$'' doğruysa,
o hâlde ``$P$ ise $R$'' de doğrudur.
\end{center}

\begin{gunlukornek}
\begin{tabular}{@{}r@{\;:\;}l}
\textcolor{ornekframe}{$P$} & Yağmur yağıyor. \\
\textcolor{ornekframe}{$Q$} & Yer ıslak. \\
\textcolor{ornekframe}{$R$} & Yer kaygan.
\end{tabular}

\smallskip
``Yağmur yağarsa yer ıslanır'' ve ``yer ıslanırsa kayar'' ise,
``yağmur yağarsa yer kayar'' sonucuna ulaşırız. Hipotetik kıyas,
adımları bir zincirin halkaları gibi birbirine bağlamaktır.
\end{gunlukornek}

\noindent Doğal çıkarım ağacı, iki $\to$E ve bir $\to$G adımından oluşur:
\begin{center}
\begin{prooftree}
  \hypo{[P]^{1}}
  \hypo{P \to Q}
  \infer2[$\to$E]{Q}
  \hypo{Q \to R}
  \infer[rule margin=0.8em]2[$\to$E]{R}
  \infer1[$\to$G$^{1}$]{P \to R}
\end{prooftree}
\end{center}

\noindent \Lean{} ispatı da aynı yapıyı izler:

\leandosya[Hipotetik kıyas]{Bolum02}{hipotetik_kiyas}

\paragraph{Teoremin başlığını okumak}
\Cref{tab:hk-baslik}, ilk satırı parça parça açıklıyor.

\begin{table}[htbp]
\centering
\renewcommand{\arraystretch}{1.35}
\rowcolors{2}{tablosatir}{white}
\begin{tabular}{>{\raggedright\arraybackslash}p{0.33\linewidth}
                >{\raggedright\arraybackslash}p{0.57\linewidth}}
\rowcolor{tabloust}
\tablobaslik{Parça} & \tablobaslik{Anlamı} \\
\leankw{theorem} \lc{hipotetik\_kiyas}
  & ``hipotetik\_kiyas'' adında bir teorem tanımlıyorum. \\
\lc{(P Q R :} \leantype{Prop}\lc{)}
  & $P$, $Q$, $R$ birer \emph{önerme}dir. \\
\lc{(}\leanhyp{h1}\lc{ : P \textcolor{leansym}{→} Q)}
  & Elimde \leanhyp{h1} adında bir kanıt var: ``$P$ ise $Q$''. \\
\lc{(}\leanhyp{h2}\lc{ : Q \textcolor{leansym}{→} R)}
  & Elimde \leanhyp{h2} adında bir kanıt var: ``$Q$ ise $R$''. \\
\lc{: P \textcolor{leansym}{→} R}
  & Kanıtlamam gereken şey, yani \emph{hedef}: ``$P$ ise $R$''. \\
\lc{:=} \leankw{by}
  & İspatı \emph{taktiklerle}, adım adım yapacağım. \\
\end{tabular}
\caption{Hipotetik kıyas teoreminin başlığı.}
\label{tab:hk-baslik}
\end{table}

\paragraph{Kilit fikir: $\to$ bir fonksiyondur}
\lc{h1 : P → Q} bir makinedir: İçine $P$'nin kanıtı girer, dışarı $Q$'nun
kanıtı çıkar. \lc{h2} de aynı biçimde $Q$'yu $R$'ye çevirir.
Hipotetik kıyas, bu iki makineyi uç uca eklemektir:

\begin{center}
\begin{tikzpicture}[
    makine/.style={draw=leanhyp, fill=leanhyp!12, thick, rounded corners=3pt,
                   minimum width=1.6cm, minimum height=0.9cm,
                   font=\leanfont\bfseries\color{leanhyp}},
    kanit/.style={draw=leanframe, fill=leanbg, thick, circle,
                  minimum size=0.9cm, font=\bfseries\color{leanframe}},
    ok/.style={-{Stealth[length=3mm]}, very thick, leansym},
  ]
  \node[kanit] (p) {$P$};
  \node[makine, right=1cm of p] (h1) {h1};
  \node[kanit, right=1cm of h1] (q) {$Q$};
  \node[makine, right=1cm of q] (h2) {h2};
  \node[kanit, right=1cm of h2] (r) {$R$};
  \draw[ok] (p) -- (h1);
  \draw[ok] (h1) -- (q);
  \draw[ok] (q) -- (h2);
  \draw[ok] (h2) -- (r);
  \node[below=2pt of p, font=\leanfont\footnotesize, text=black!60] {p : P};
\end{tikzpicture}
\end{center}

\paragraph{Adım adım ispat}

\noindent\textbf{Başlangıç.} Henüz hiçbir taktik uygulanmadı:
\begin{hedefdurumu}[Başlangıç durumu]
P Q R : Prop
h1 : P → Q
h2 : Q → R
⊢ P → R
\end{hedefdurumu}

\noindent\textbf{Adım 1: \leantac{intro p}.}
``$P \to R$'' kanıtlamak, ``$P$'nin doğru olduğunu varsayayım, $R$'yi
göstereyim'' demektir. \leantac{intro p}, $P$ varsayımını \leanhyp{p}
adıyla bağlama alır. Ağaçtaki $[P]^1$ varsayımı budur.
\begin{hedefdurumu}[\texttt{intro p} sonrası]
P Q R : Prop
h1 : P → Q
h2 : Q → R
p : P
⊢ R
\end{hedefdurumu}

\noindent\textbf{Adım 2: \leantac{apply h2}.}
$R$'yi nasıl elde ederiz? \lc{h2 : Q → R} der ki: ``Bana $Q$ ver, sana $R$
vereyim.'' \leantac{apply h2} şu anlama gelir: ``$R$'yi \leanhyp{h2} ile
elde edeceğim; o hâlde $Q$'yu göstermem yeter.'' Lean burada
\emph{geriye doğru} düşünür:
\begin{hedefdurumu}[\texttt{apply h2} sonrası]
⊢ Q
\end{hedefdurumu}

\noindent\textbf{Adım 3: \leantac{apply h1}.}
Aynı mantıkla, \lc{h1 : P → Q} olduğundan $Q$ için $P$'yi göstermek yeter:
\begin{hedefdurumu}[\texttt{apply h1} sonrası]
⊢ P
\end{hedefdurumu}

\noindent\textbf{Adım 4: \leantac{exact p}.}
Hedef $P$ ve elimizde zaten \lc{p : P} var:
\ispatbitti

\noindent Kolaylık olsun diye Adım~2 ve 3'te yalnızca hedefi gösterdik;
bağlam (\lc{P Q R}, \lc{h1}, \lc{h2}, \lc{p}) değişmeden kalır.

\paragraph{Akışın özeti}
İspat \emph{sondan başa} ilerler: hedeften başlar, elimizdeki hipotezlere
doğru geri gelir. Anlattığı mantık ise \emph{baştan sona} okunur.
\Cref{fig:hk-akis} her iki yönü birlikte gösteriyor.

\begin{figure}[htbp]
\centering
\begin{tikzpicture}[
    hedef/.style={draw=hedefframe, fill=hedefbg, thick, rounded corners=3pt,
                  minimum width=1.55cm, minimum height=0.85cm,
                  font=\leanfont\small},
    bitti/.style={hedef, draw=bittiframe, fill=bittibg},
    kanit/.style={hedef, fill=leanbg, draw=leanframe},
    geri/.style={-{Stealth[length=2.6mm]}, thick, leantac},
    ileri/.style={-{Stealth[length=2.6mm]}, thick, leansym},
    etiket/.style={font=\leanfont\bfseries\scriptsize, yshift=2pt},
  ]
  \node[hedef] (g1) {⊢ P → R};
  \node[hedef, right=1.75cm of g1] (g2) {⊢ R};
  \node[hedef, right=1.75cm of g2] (g3) {⊢ Q};
  \node[hedef, right=1.75cm of g3] (g4) {⊢ P};
  \node[bitti, right=1.75cm of g4] (g5) {bitti};

  \draw[geri] (g1) -- node[above, etiket, text=leantac]{intro p}  (g2);
  \draw[geri] (g2) -- node[above, etiket, text=leantac]{apply h2} (g3);
  \draw[geri] (g3) -- node[above, etiket, text=leantac]{apply h1} (g4);
  \draw[geri] (g4) -- node[above, etiket, text=leantac]{exact p}  (g5);
  \node[above=0.55cm of g3, font=\baslikfont\small\bfseries, text=leantac]
       {Taktiklerin yönü: hedeften geriye};

  \node[bitti, below=1.6cm of g2] (k1) {p : P};
  \node[kanit, below=1.6cm of g3] (k2) {Q};
  \node[kanit, below=1.6cm of g4] (k3) {R};
  \draw[ileri] (k1) -- node[above, etiket, text=leanhyp]{h1} (k2);
  \draw[ileri] (k2) -- node[above, etiket, text=leanhyp]{h2} (k3);
  \node[below=0.55cm of k2, font=\baslikfont\small\bfseries, text=leansym]
       {Mantığın yönü: varsayımdan sonuca};
\end{tikzpicture}
\caption{Hipotetik kıyas ispatında taktiklerin ve mantığın yönü.}
\label{fig:hk-akis}
\end{figure}

\paragraph{Tek satırlık ispat}
Taktik kullanmadan, ispatı doğrudan bir fonksiyon olarak da yazabiliriz:

\leandosya[Hipotetik kıyas, terim modu]{Bolum02}{hipotetik_kiyas_terim}

Okunuşu: ``Bir \leanhyp{p} al. Onu \leanhyp{h1}'e ver, $Q$'nun kanıtı
çıksın. Onu \leanhyp{h2}'ye ver, $R$'nin kanıtı çıksın.'' Bu, matematikteki
\textbf{fonksiyon bileşkesidir}: $h_2 \circ h_1 : P \to R$. Hipotetik kıyas,
mantıkta tam olarak fonksiyonları ard arda uygulamaya karşılık gelir.

% ---------------------------------------------------------------------
\section{Ve (\texorpdfstring{$\wedge$}{∧})}\label{sec:ve}

\subsection{Kurallar}

\begin{center}
\begin{tabular}{c@{\hspace{3em}}c@{\hspace{2.2em}}c}
\begin{prooftree}
  \hypo{P}
  \hypo{Q}
  \infer2[$\wedge$G]{P \wedge Q}
\end{prooftree}
&
\begin{prooftree}
  \hypo{P \wedge Q}
  \infer1[$\wedge$E$_1$]{P}
\end{prooftree}
&
\begin{prooftree}
  \hypo{P \wedge Q}
  \infer1[$\wedge$E$_2$]{Q}
\end{prooftree}
\\[10pt]
\small\textcolor{ikincil}{\bfseries giriş:} \leantac{constructor}, \lc{⟨hp, hq⟩} &
\multicolumn{2}{c}{\small\textcolor{vurgu}{\bfseries eleme:} \lc{h.left}, \lc{h.right}, \leantac{obtain}}
\end{tabular}
\end{center}

``$P$ ve $Q$'' kanıtlamak için ikisini de ayrı ayrı kanıtlarız
($\wedge$G). Elimizde ``$P$ ve $Q$'' varsa, ikisini de ayrı ayrı
kullanabiliriz ($\wedge$E$_1$, $\wedge$E$_2$).

\subsection{Ve'ye giriş}

\leandosya[Ve'ye giriş, taktik modu]{Bolum02}{ve_giris}

\leantac{constructor} taktiği hedefi iki alt hedefe böler. Bu noktada
\Lean{} iki ayrı hedef gösterir:

\begin{hedefdurumu}[\texttt{constructor} sonrası: iki hedef]
case left
P Q : Prop
hp : P
hq : Q
⊢ P

case right
P Q : Prop
hp : P
hq : Q
⊢ Q
\end{hedefdurumu}

Her alt hedefi bir \lc{·} (orta nokta) ile başlayan girintili bir blokta
kapatırız. Terim modunda ise ``ve''nin kanıtı basitçe bir \emph{çifttir}:

\leandosya[Ve'ye giriş, terim modu]{Bolum02}{ve_giris_terim}

\begin{dikkat}[Dikkat: Orta nokta ve girinti]
\lc{·} karakteri sıradan bir nokta (\lc{.}) değildir; VS Code'da
\lc{\textbackslash.} yazıp boşluk tuşuna basarak elde edilir. Bloklardaki
girinti de önemlidir: \Lean{}, girintiye bakarak hangi taktiğin hangi alt
hedefe ait olduğunu anlar. Girintiyi bozarsanız ``\emph{unsolved goals}''
hatası alırsınız.
\end{dikkat}

\subsection{Ve'nin elemesi ve değişme özelliği}

\leandosya[Ve'nin elemesi]{Bolum02}{ve_eleme}

\lc{h.left} ve \lc{h.right} (ya da kısaca \lc{h.1}, \lc{h.2}), çiftin sol
ve sağ bileşenlerini alır. Birden çok bileşeni aynı anda almak için
\leantac{obtain} kullanılır:

\leandosya[Ve'nin değişme özelliği]{Bolum02}{ve_degisme}

\begin{hedefgecisi}
\begin{hedefdurumu}[\texttt{intro h} sonrası]
P Q : Prop
h : P ∧ Q
⊢ Q ∧ P
\end{hedefdurumu}
\sonra
\begin{hedefdurumu}[\texttt{obtain} sonrası]
P Q : Prop
hp : P
hq : Q
⊢ Q ∧ P
\end{hedefdurumu}
\end{hedefgecisi}

\leantac{obtain ⟨hp, hq⟩ := h}, \lc{h} kanıtını iki parçasına ayırır ve
\lc{h}'yi bağlamdan kaldırır. Terim modunda aynı ayırma işlemi
\emph{örüntü eşleme} ile yapılır:

\leandosya[Ve'nin değişme özelliği, terim modu]{Bolum02}{ve_degisme_terim}

\noindent Okunuşu: ``Bir çift al, bileşenlerine \lc{hp} ve \lc{hq} de;
sonra bileşenlerin yerini değiştirerek yeni bir çift döndür.''

% ---------------------------------------------------------------------
\section{Veya (\texorpdfstring{$\vee$}{∨})}\label{sec:veya}

\subsection{Kurallar}

\begin{center}
\begin{tabular}{c@{\hspace{2em}}c@{\hspace{3em}}c}
\begin{prooftree}
  \hypo{P}
  \infer1[$\vee$G$_1$]{P \vee Q}
\end{prooftree}
&
\begin{prooftree}
  \hypo{Q}
  \infer1[$\vee$G$_2$]{P \vee Q}
\end{prooftree}
&
\begin{prooftree}
  \hypo{P \vee Q}
  \hypo{[P]^{1}}
  \ellipsis{}{R}
  \hypo{[Q]^{2}}
  \ellipsis{}{R}
  \infer3[$\vee$E$^{1,2}$]{R}
\end{prooftree}
\\[10pt]
\multicolumn{2}{c}{\small\textcolor{ikincil}{\bfseries giriş:} \lc{Or.inl}, \lc{Or.inr}} &
\small\textcolor{vurgu}{\bfseries eleme:} \leantac{cases}, \leantac{rcases}, \lc{Or.elim}
\end{tabular}
\end{center}

Veya'nın kuralları, ve'ninkilerin bir tür \emph{aynadaki görüntüsüdür}.
``$P$ veya $Q$'' kanıtlamak için ikisinden \emph{birini} kanıtlamak yeter.
Buna karşılık ``$P$ veya $Q$''yu kullanmak daha zahmetlidir: Hangisinin
doğru olduğunu bilmediğimiz için \emph{her iki durumda da} hedefe
ulaşabildiğimizi göstermek zorundayız. Bu \emph{durum ayırma}
yöntemidir (\cref{fig:durum-ayirma}).

\begin{figure}[htbp]
\centering
\begin{tikzpicture}[
    durum/.style={draw=#1, fill=#1!8, thick, rounded corners=4pt,
                  minimum width=2.8cm, minimum height=0.9cm, align=center,
                  font=\small},
    ok/.style={-{Stealth[length=2.6mm]}, thick, #1},
  ]
  \node[durum=leanframe] (h) at (0,0) {\lc{h : P ∨ Q}};
  \node[durum=ikincil]   (a) at (-3.2,-1.7) {1. durum\\ \lc{hp : P}};
  \node[durum=vurgu]     (b) at (3.2,-1.7)  {2. durum\\ \lc{hq : Q}};
  \node[durum=ikincil]   (ar) at (-3.2,-3.4) {$\vdots$\\ \lc{⊢ R} kapatılır};
  \node[durum=vurgu]     (br) at (3.2,-3.4)  {$\vdots$\\ \lc{⊢ R} kapatılır};
  \node[durum=bittiframe] (r) at (0,-5.0) {\lc{R} kanıtlandı};

  \draw[ok=ikincil] (h) -- node[left=8pt, font=\leanfont\scriptsize\bfseries, text=ikincil]{inl} (a);
  \draw[ok=vurgu]   (h) -- node[right=8pt, font=\leanfont\scriptsize\bfseries, text=vurgu]{inr} (b);
  \draw[ok=ikincil] (a) -- (ar);
  \draw[ok=vurgu]   (b) -- (br);
  \draw[ok=bittiframe] (ar) -- (r);
  \draw[ok=bittiframe] (br) -- (r);
\end{tikzpicture}
\caption{Veya'nın elemesi: Her iki durumda da aynı hedefe ulaşılmalıdır.}
\label{fig:durum-ayirma}
\end{figure}

\subsection{Veya'ya giriş}

\leandosya[Veya'ya giriş]{Bolum02}{veya_giris}

\lc{Or.inl} (``sola ekle'', İng.\ \emph{in left}) sol taraftaki önermenin
kanıtından bir veya kanıtı üretir; \lc{Or.inr} sağ taraf için aynısını
yapar. \Lean{}'in güncel sürümlerinde \leantac{left} ve \leantac{right}
taktikleri de aynı işi görür; kullandığınız sürümde tanınmazlarsa
\leantac{apply Or.inl} biçimini kullanın.

\subsection{Veya'nın elemesi ve değişme özelliği}

\leandosya[Veya'nın değişme özelliği]{Bolum02}{veya_degisme}

\leantac{cases h with} satırı \lc{h : P ∨ Q} hipotezini iki duruma ayırır.
Her durum \lc{|} ile başlar ve durumun adını (\lc{inl} ya da \lc{inr})
ve o durumda elde ettiğimiz kanıtın adını verir. Aynı ispat
\leantac{rintro} ile daha kısa yazılabilir; burada \lc{(hp | hq)}
örüntüsü ``ya \lc{hp} ya \lc{hq}'' diye okunur:

\leandosya[rintro ile]{Bolum02}{veya_degisme_rcases}

\noindent Terim modunda durum ayırma \lc{Or.elim} ile yapılır; iki durumun
her biri için bir fonksiyon veririz:

\leandosya[Veya'nın değişme özelliği, terim modu]{Bolum02}{veya_degisme_terim}

\subsection{Dağılma: kuralları birleştirmek}

Şimdi üç bağlacın kurallarını tek bir ispatta kullanalım. Ve'nin veya
üzerine dağılması şöyle ifade edilir:
\[
  P \wedge (Q \vee R) \;\to\; (P \wedge Q) \vee (P \wedge R)
\]

\begin{figure}[htbp]
\centering
\small
\begin{prooftree}
  \hypo{[P \wedge (Q \vee R)]^{3}}
  \infer1[$\wedge$E$_2$]{Q \vee R}
  %
  \hypo{[P \wedge (Q \vee R)]^{3}}
  \infer1[$\wedge$E$_1$]{P}
  \hypo{[Q]^{1}}
  \infer2[$\wedge$G]{P \wedge Q}
  \infer1[$\vee$G$_1$]{(P \wedge Q) \vee (P \wedge R)}
  %
  \hypo{[P \wedge (Q \vee R)]^{3}}
  \infer1[$\wedge$E$_1$]{P}
  \hypo{[R]^{2}}
  \infer2[$\wedge$G]{P \wedge R}
  \infer1[$\vee$G$_2$]{(P \wedge Q) \vee (P \wedge R)}
  %
  \infer3[$\vee$E$^{1,2}$]{(P \wedge Q) \vee (P \wedge R)}
  \infer1[$\to$G$^{3}$]{P \wedge (Q \vee R) \to (P \wedge Q) \vee (P \wedge R)}
\end{prooftree}
\caption{Dağılma yasasının doğal çıkarım ağacı. Üç kural ailesi de
kullanılıyor: $\to$G, $\wedge$G/E ve $\vee$G/E.}
\label{fig:dagilma-agaci}
\end{figure}

\leandosya[Dağılma]{Bolum02}{dagilma}

\Cref{fig:dagilma-agaci} ile \Lean{} ispatını satır satır karşılaştırın:
\leantac{intro h} en alttaki $\to$G$^3$ adımıdır; \leantac{obtain}
$\wedge$E$_1$ ve $\wedge$E$_2$'yi tek seferde uygular; \leantac{cases hqr}
$\vee$E$^{1,2}$'dir; her durumdaki \lc{Or.inl ⟨hp, hq⟩} ise önce
$\wedge$G, sonra $\vee$G uygular. Kâğıt üzerindeki ağaç ile \Lean{} ispatı
aynı mantıksal içeriğin iki farklı yazılışıdır.

% ---------------------------------------------------------------------
\section{Curry--Howard: ilk tablo}\label{sec:curry-howard-1}

Bu bölüm boyunca her ispatı iki biçimde yazdık. Terim modundaki
ispatlara geri dönüp bakarsak, her bağlacın bir \emph{programlama
kavramına} karşılık geldiğini görürüz (\cref{tab:curry-howard-1}).

\begin{table}[htbp]
\centering
\renewcommand{\arraystretch}{1.5}
\rowcolors{2}{tablosatir}{white}
\begin{tabular}{>{\raggedright\arraybackslash}p{2.7cm}
                >{\raggedright\arraybackslash}p{3.5cm}
                >{\raggedright\arraybackslash}p{5.6cm}}
\rowcolor{tabloust}
\tablobaslik{Mantık} & \tablobaslik{Programlama} & \tablobaslik{Lean'deki kanıt} \\
$P \to Q$          & fonksiyon tipi        & \lc{fun hp => ...} \\
$\to$ eleme        & fonksiyon çağırma     & \lc{hpq hp} \\
$P \wedge Q$       & çift (çarpım tipi)    & \lc{⟨hp, hq⟩} \\
$\wedge$ eleme     & bileşen alma          & \lc{h.1}, \lc{h.2} \\
$P \vee Q$         & etiketli seçim (toplam tipi) & \lc{Or.inl hp}, \lc{Or.inr hq} \\
$\vee$ eleme       & durum analizi         & \lc{Or.elim}, \leantac{cases} \\
hipotetik kıyas    & fonksiyon bileşkesi   & \lc{fun p => h2 (h1 p)} \\
\end{tabular}
\caption{Curry--Howard eşlemesinin ilk tablosu. \Cref{bol:degilleme}'te
\lc{True} ve \lc{False}, \cref{bol:yuklemler}'da $\forall$ ve $\exists$
satırları eklenecek.}
\label{tab:curry-howard-1}
\end{table}

\leandosya[Curry ve uncurry]{Bolum02}{curry}

Bu iki teorem, programlamada \emph{currying} adıyla bilinen dönüşümün
mantıktaki karşılığıdır: İki argüman alan bir fonksiyon (\lc{P ∧ Q → R}),
argümanlarını tek tek alan bir fonksiyona (\lc{P → Q → R}) dönüştürülebilir
ve tersi de doğrudur. \lc{P → Q → R} yazımını neden bu kadar sık
kullandığımızın açıklaması da budur.

\begin{onemlifikir}[Önemli Fikir: Curry--Howard eşlemesi]
\begin{center}
\textcolor{leantype}{\bfseries önermeler = türler}, \qquad
\textcolor{leanhyp}{\bfseries kanıtlar = programlar}.
\end{center}
Bu bölümdeki her teorem bir tür, her ispat o türden bir programdı.
Taktik modu bu programı adım adım inşa etmemizi sağlayan bir araçtır;
\Lean{} her taktik ispatını arka planda bir terime çevirir ve çekirdeğe
bu terimi denetletir. Eşlemenin kesin ifadesini \cref{bol:lambda}'te
vereceğiz \citep{howard1980, sorensen2006}.
\end{onemlifikir}

% ---------------------------------------------------------------------
\section{Kanıt ve doğruluk tablosu}\label{sec:dogruluk-tablosu}

Önermeler mantığını belki okulda \emph{doğruluk tablolarıyla}
öğrendiniz. \Cref{fig:dogruluk-tablosu}'deki tablo, $P \to Q$'nun hangi
durumlarda doğru olduğunu gösteriyor.

\begin{figure}[htbp]
\centering
\begin{tikzpicture}
  \matrix[matrix of nodes, nodes={minimum width=1.5cm, minimum height=0.75cm,
          anchor=center, font=\small}, column sep=0pt, row sep=0pt,
          row 1/.style={nodes={fill=tabloust, text=white, font=\baslikfont\bfseries\small}}] (t) {
    $P$ & $Q$ & $P \to Q$ \\
    D & D & |[fill=bittibg, text=bittiframe]| \textbf{D} \\
    D & Y & |[fill=hatabg,  text=hataframe]|  \textbf{Y} \\
    Y & D & |[fill=bittibg, text=bittiframe]| \textbf{D} \\
    Y & Y & |[fill=bittibg, text=bittiframe]| \textbf{D} \\
  };
  \draw[black!25] (t.north west) rectangle (t.south east);
  \draw[black!25] (t-1-3.north west) -- (t-5-3.south west);
  \node[right=0.4cm of t-3-3, font=\small, text=hataframe, align=left]
       {Tek yanlış durum:\\ öncül doğru, sonuç yanlış.};
\end{tikzpicture}
\caption{İmplikasyonun doğruluk tablosu (D: doğru, Y: yanlış).}
\label{fig:dogruluk-tablosu}
\end{figure}

Doğruluk tablosu \emph{anlambilimsel} bir yaklaşımdır: Bir önermenin
doğruluğunu, atomlarına verilebilecek bütün doğruluk değerlerini tek tek
deneyerek belirler. Bu bölümde yaptığımız ise \emph{sözdizimsel} bir
yaklaşımdır: Bir önermeyi, çıkarım kurallarını adım adım uygulayarak
\emph{inşa ederiz}.

İki yaklaşım arasındaki ilişki mantığın en önemli sorularından biridir.
Kanıtlayabildiğimiz her şey doğruluk tablosuna göre doğru mudur
(\emph{geçerlilik})? Doğruluk tablosuna göre doğru olan her şeyi
kanıtlayabilir miyiz (\emph{tamlık})? İlk sorunun cevabı evettir
(\cref{bol:gecerlilik}). İkinci sorunun cevabı ise şaşırtıcıdır:
\emph{Bu bölümün kurallarıyla, hayır.} Alıştırma~\ref{alis:peirce}'deki
Peirce yasası doğruluk tablosuna göre her zaman doğrudur, ama yalnızca bu
bölümün kurallarıyla kanıtlanamaz. Bu gizemi \cref{bol:yapici-klasik}'te
çözeceğiz.

% ---------------------------------------------------------------------
\section{Taktik özeti}\label{sec:taktik-ozeti-2}

\begin{center}
\renewcommand{\arraystretch}{1.45}
\rowcolors{2}{tablosatir}{white}
\begin{tabularx}{\linewidth}{>{\raggedright\arraybackslash}p{3.0cm}
                             >{\raggedright\arraybackslash}p{1.6cm}
                             >{\raggedright\arraybackslash}X
                             >{\raggedright\arraybackslash}p{2.9cm}}
\rowcolor{tabloust}
\tablobaslik{Taktik} & \tablobaslik{Kural} & \tablobaslik{Ne yapar?} & \tablobaslik{Terim karşılığı} \\
\leantac{intro h}                   & $\to$G      & Hedef \lc{A → B} ise \lc{h : A} varsayar, hedef \lc{B} olur. & \lc{fun h => ...} \\
\leantac{apply h}                   & $\to$E      & \lc{h : A → B} ile hedef \lc{B}'yi \lc{A}'ya çevirir. & \lc{h ...} \\
\leantac{exact t}                   & ---         & Hedefi \lc{t} ile kapatır. & \lc{t} \\
\leantac{constructor}               & $\wedge$G   & \lc{⊢ A ∧ B} hedefini iki hedefe böler. & \lc{⟨a, b⟩} \\
\leantac{obtain ⟨a, b⟩ := h}        & $\wedge$E   & \lc{h : A ∧ B}'yi iki hipoteze ayırır. & \lc{h.1}, \lc{h.2} \\
\lc{exact Or.inl a}                 & $\vee$G$_1$ & Veya'nın sol tarafını kanıtlar. & \lc{Or.inl a} \\
\leantac{cases h with ...}          & $\vee$E     & \lc{h : A ∨ B} için iki durum açar. & \lc{Or.elim h f g} \\
\leantac{rintro (a | b)}            & $\to$G, $\vee$E & Varsayımı alır ve hemen durumlara ayırır. & --- \\
\end{tabularx}
\end{center}

% ---------------------------------------------------------------------
\begin{arastirmanotu}
\textbf{Doğal çıkarımın kaynakları.} Gentzen'in 1935 tarihli makalesi
hem doğal çıkarımı hem de \cref{bol:ardisik}'te göreceğimiz ardışık hesabı
tanıtır \citep{gentzen1935}. Doğal çıkarım ispatlarının
\emph{normalleştirilebileceği}, yani giriş--eleme çiftlerinin silinerek
daha dolaysız ispatlara indirgenebileceği \citet{prawitz1965} tarafından
gösterildi. Bu sonuç, Curry--Howard eşlemesi altında programların
\emph{çalıştırılmasına} karşılık gelir \citep{howard1980}.

\textbf{İleri okuma.} Ders kitabı düzeyinde doğal çıkarım için
\citet{vandalen2013}; ispat kuramının ayrıntılı ve modern bir sunumu
için \citet{troelstra2000}; Curry--Howard eşlemesinin kapsamlı bir
incelemesi için \citet{sorensen2006}. \Lean{} odaklı bir anlatım için
\citet{logicandproof} ve \citet{tpil4}.

\textbf{Açık bir soru.} Taktik ispatları okunabilir midir? Uzun taktik
ispatları, ara hedefler görünmediğinde okunması zor metinlere dönüşebilir.
\emph{İspat mühendisliği} alanında okunabilir ve bakımı kolay ispat
yazmaya dair yöntemler hâlâ etkin biçimde tartışılmaktadır;
\ciltiki{9}'da bu konuya \Mathlib{} deneyimi üzerinden döneceğiz.
\end{arastirmanotu}

\begin{bolumozeti}
\begin{itemize}
  \item Doğal çıkarımda her bağlacın bir \textbf{giriş} (kanıtlama) ve bir
        \textbf{eleme} (kullanma) kuralı vardır.
  \item $\to$G = \leantac{intro}; $\to$E = \leantac{apply} ya da fonksiyon
        uygulaması. İmplikasyonun kanıtı bir fonksiyondur.
  \item $\wedge$G = \leantac{constructor} ya da \lc{⟨hp, hq⟩};
        $\wedge$E = \lc{h.1}, \lc{h.2}, \leantac{obtain}. Ve'nin kanıtı bir çifttir.
  \item $\vee$G = \lc{Or.inl}, \lc{Or.inr}; $\vee$E = \leantac{cases},
        \leantac{rintro}. Veya'yı kullanmak durum ayırmayı gerektirir.
  \item Her taktik ispatı arka planda bir terimdir: \textbf{önermeler
        türler, kanıtlar programlardır}.
\end{itemize}
\end{bolumozeti}

\begin{alistirmalar}
\alistirma{1} \lc{P ∧ Q → P} önermesini taktik modunda ve terim modunda
ispatlayın. (\lc{Alistirmalar/Bolum02.lean}, \lc{a2\_1})

\alistirma{1} \lc{P → P ∨ Q} önermesini ispatlayın. Hangi giriş kuralını
kullandınız? (\lc{a2\_2})

\alistirma{1} \lc{(P → Q) → (P → R) → P → Q ∧ R} önermesini ispatlayın ve
doğal çıkarım ağacını çizin. (\lc{a2\_3})

\alistirma{2} Ve'nin birleşme özelliğini ispatlayın:
\lc{P ∧ (Q ∧ R) → (P ∧ Q) ∧ R}. İpucu: \leantac{obtain ⟨hp, hq, hr⟩ := h}
iç içe çiftleri tek seferde açar. (\lc{a2\_4})

\alistirma{2} Veya'nın birleşme özelliğini ispatlayın:
\lc{P ∨ (Q ∨ R) → (P ∨ Q) ∨ R}. Kaç durum ayırmanız gerekti? (\lc{a2\_5})

\alistirma{2} \emph{Durum ayırma ilkesini} ispatlayın:
\lc{(P → R) → (Q → R) → P ∨ Q → R}. Bu önermenin $\vee$E kuralının
kendisi olduğunu gösterin. (\lc{a2\_6})

\alistirma{2} Dağılmanın tersini ispatlayın:
\lc{(P ∧ Q) ∨ (P ∧ R) → P ∧ (Q ∨ R)}. (\lc{a2\_7})

\alistirma{2}\label{alis:s-kombinator} $S$ kombinatörünü terim modunda
ispatlayın: \lc{(P → Q → R) → (P → Q) → P → R}. İspatınızın bir
fonksiyon olarak ne yaptığını tek cümleyle açıklayın. (\lc{a2\_8})

\alistirma{2} \lc{P ∨ (Q ∧ R) → (P ∨ Q) ∧ (P ∨ R)} önermesini
ispatlayın. (\lc{a2\_9})

\alistirma{3} Bir önceki alıştırmanın tersini ispatlayın:
\lc{(P ∨ Q) ∧ (P ∨ R) → P ∨ (Q ∧ R)}. İç içe kaç durum ayırması gerekti?
Daha azıyla yapılabilir mi? (\lc{a2\_10})

\alistirma{3} (Kâğıt-kalem) \Cref{fig:dagilma-agaci}'deki ağaçta hangi
varsayımın hangi kuralda kapandığını gösterin. Aynı ağaçta $[P \wedge
(Q \vee R)]^3$ varsayımının üç kez göründüğüne dikkat edin. \Lean{}
ispatında bu varsayıma karşılık gelen tek bir \lc{h} vardır. Bu fark
neden ortaya çıkıyor?

\alistirma{3}\label{alis:peirce} \textbf{Peirce yasası:}
\lc{((P → Q) → P) → P}. Önce doğruluk tablosuyla her zaman doğru
olduğunu gösterin. Sonra yalnızca bu bölümün araçlarıyla \Lean{}'de
ispatlamayı deneyin ve nerede takıldığınızı not edin. Neden takıldığınızı
\cref{bol:yapici-klasik}'te göreceğiz.
\end{alistirmalar}
   % Önermeler mantığı I: →, ∧, ∨
% =====================================================================
%  Bölüm 3 — Önermeler Mantığı II: Değilleme, Çelişki ve Denklik
% =====================================================================
\chapter{Önermeler Mantığı II: Değilleme, Çelişki ve Denklik}\label{bol:degilleme}

\begin{bolumhedefleri}
\begin{itemize}
  \item \lc{True} ve \lc{False} önermelerini ve Curry--Howard'daki
        karşılıklarını,
  \item Değillemenin \Lean{}'de neden \lc{¬P := P → False} olarak
        tanımlandığını,
  \item Çelişkiden her şeyin çıktığını (\emph{ex falso quodlibet}) ve
        \leantac{exfalso}, \lc{absurd}, \leantac{contradiction} araçlarını,
  \item Denkliği (\lc{↔}), \lc{Iff.mp}, \lc{Iff.mpr} ve \leantac{rw} ile
        yeniden yazmayı,
  \item De Morgan yasalarının hangi yarısının yapıcı olarak ispatlanabildiğini
        ve neden ötekinin ispatlanamadığını.
\end{itemize}
\end{bolumhedefleri}

% ---------------------------------------------------------------------
\section{İki uç: True ve False}\label{sec:true-false}

Önermeler arasında iki özel önerme vardır. \lc{True} her zaman doğru olan,
\lc{False} ise hiçbir zaman doğru olmayan önermedir. Mantık kitaplarında
bunlar genellikle $\top$ (\emph{tepe}) ve $\bot$ (\emph{dip}) ile gösterilir.

\begin{center}
\begin{tabular}{c@{\hspace{5em}}c}
\begin{prooftree}
  \infer0[$\top$G]{\top}
\end{prooftree}
&
\begin{prooftree}
  \hypo{\bot}
  \infer1[$\bot$E]{P}
\end{prooftree}
\\[10pt]
\small\textcolor{ikincil}{\bfseries giriş:} \leantac{trivial}, \lc{True.intro} &
\small\textcolor{vurgu}{\bfseries eleme (ex falso):} \lc{False.elim}, \leantac{exfalso}
\end{tabular}
\end{center}

\noindent İki kural birbirinin aynadaki görüntüsü gibidir:

\begin{itemize}
  \item $\top$'un \textbf{giriş} kuralının hiç öncülü yoktur: \lc{True}'yu
        her zaman, hiçbir şey bilmeden kanıtlayabiliriz. Buna karşılık
        \textbf{eleme} kuralı yoktur: \lc{True}'yu bilmek bize hiçbir bilgi
        vermez.
  \item $\bot$'un \textbf{giriş} kuralı yoktur: \lc{False}'u kanıtlamanın
        doğrudan bir yolu yoktur. Buna karşılık \textbf{eleme} kuralı
        çok güçlüdür: \lc{False}'un bir kanıtı elimize geçerse, \emph{her
        önermeyi} kanıtlayabiliriz.
\end{itemize}

\leandosya[True'nun kanıtı]{Bolum03}{dogru}

\subsection{Ex falso quodlibet}

Latince ``yanlıştan ne istersen'' anlamına gelen bu ilke, $\bot$E
kuralının adıdır.

\leandosya[Ex falso]{Bolum03}{ex_falso}

\noindent Daha ilginç bir ispat, \leantac{cases} taktiğini kullanır:

\leandosya[Ex falso, durum ayırarak]{Bolum03}{ex_falso_cases}

Neden işe yarıyor? \Cref{sec:veya}'te \leantac{cases}'in bir veya
kanıtını \emph{iki} duruma ayırdığını görmüştük, çünkü \lc{P ∨ Q}'nun
kanıtı iki biçimde kurulabilir (\lc{Or.inl} ve \lc{Or.inr}). \lc{False}'un
kanıtı ise \emph{hiçbir} biçimde kurulamaz. Dolayısıyla \lc{h : False}
üzerinde durum ayırınca \emph{sıfır} durum çıkar ve kapatılacak hedef
kalmaz.

\begin{figure}[htbp]
\centering
\begin{tikzpicture}[
    tip/.style={draw=#1, fill=#1!8, thick, rounded corners=4pt,
                minimum width=2.6cm, minimum height=1.0cm,
                font=\leanfont\small, align=center},
    kur/.style={draw=#1, fill=white, thick, rounded corners=3pt,
                font=\leanfont\scriptsize, inner sep=3pt},
    ok/.style={-{Stealth[length=2.3mm]}, thick, #1},
  ]
  % Or
  \node[tip=leanframe] (or) at (0,0) {P ∨ Q};
  \node[kur=ikincil] (o1) at (-1.0,-1.5) {Or.inl};
  \node[kur=ikincil] (o2) at (1.0,-1.5) {Or.inr};
  \draw[ok=ikincil] (or) -- (o1); \draw[ok=ikincil] (or) -- (o2);
  \node[font=\baslikfont\scriptsize, text=black!60] at (0,-2.3) {2 kurucu → 2 durum};
  % True
  \node[tip=bittiframe] (tr) at (5,0) {True};
  \node[kur=bittiframe] (t1) at (5,-1.5) {True.intro};
  \draw[ok=bittiframe] (tr) -- (t1);
  \node[font=\baslikfont\scriptsize, text=black!60] at (5,-2.3) {1 kurucu → 1 durum};
  % False
  \node[tip=hataframe] (fa) at (10,0) {False};
  \node[text=hataframe] at (10,-1.5) {\Large$\emptyset$};
  \node[font=\baslikfont\scriptsize, text=black!60] at (10,-2.3) {0 kurucu → 0 durum};
\end{tikzpicture}
\caption{Bir önermenin \emph{kurucuları}, kanıtını inşa etmenin yollarıdır.
\leantac{cases} her kurucu için bir durum açar. \lc{False}'un kurucusu
yoktur.}
\label{fig:kurucular}
\end{figure}

\begin{onemlifikir}[Önemli Fikir: True ve False birer tiptir]
Curry--Howard eşlemesinde \lc{True}, tek elemanlı tipe (programlamada
\emph{unit}), \lc{False} ise hiç elemanı olmayan \emph{boş tipe} karşılık
gelir. Boş tipten başka bir tipe giden fonksiyonu yazmak kolaydır: Hiçbir
girdi gelmeyeceği için hiçbir çıktı üretmemiz de gerekmez. Ex falso tam
olarak budur. Bu tipler \cref{bol:tumevarimsal-tipler}'da \leankw{inductive}
ile kendimiz tanımlayacağız.
\end{onemlifikir}

% ---------------------------------------------------------------------
\section{Değilleme: \texorpdfstring{$\neg P$}{¬P} ``\texorpdfstring{$P$}{P} ise çelişki'' demektir}\label{sec:degilleme}

\Lean{}'de değilleme yeni bir temel kavram değildir; implikasyon ve
\lc{False} ile \emph{tanımlanır}:

\leandosya[Değillemenin tanımı]{Bolum03}{not_tanim}

\begin{tanim}{Değilleme}{degilleme}
Bir $P$ önermesinin \textbf{değillemesi},
\[
  \neg P \;:=\; P \to \bot
\]
önermesidir. Yani ``$P$ değil'' demek, ``$P$ doğru olsaydı bir çelişkiye
varırdık'' demektir.
\end{tanim}

Bu tanımın önemli bir sonucu var: Değilleme için ayrı kurallar
öğrenmemize gerek yok, çünkü $\neg$ aslında $\to$'dur. $\neg$G kuralı
$\to$G'nin, $\neg$E kuralı da $\to$E'nin özel hâlidir:

\begin{center}
\begin{tabular}{c@{\hspace{4em}}c}
\begin{prooftree}
  \hypo{[P]^{1}}
  \ellipsis{}{\bot}
  \infer1[$\neg$G$^{1}$]{\neg P}
\end{prooftree}
&
\begin{prooftree}
  \hypo{\neg P}
  \hypo{P}
  \infer2[$\neg$E]{\bot}
\end{prooftree}
\\[10pt]
\small\textcolor{ikincil}{\bfseries giriş:} \leantac{intro hp}, hedef \lc{False} olur &
\small\textcolor{vurgu}{\bfseries eleme:} \lc{hnp hp : False}
\end{tabular}
\end{center}

\subsection{Çelişmezlik ilkesi}

Aristoteles'ten beri mantığın temel yasalarından biri sayılan
\emph{çelişmezlik ilkesi}, bir önermenin hem doğru hem yanlış
olamayacağını söyler: $\neg(P \wedge \neg P)$.

\leandosya[Çelişmezlik ilkesi]{Bolum03}{degil_giris}

\begin{hedefgecisi}
\begin{hedefdurumu}[\texttt{intro h} sonrası]
P : Prop
h : P ∧ ¬P
⊢ False
\end{hedefdurumu}
\sonra
\begin{hedefdurumu}[\texttt{obtain} sonrası]
P : Prop
hp : P
hnp : ¬P
⊢ False
\end{hedefdurumu}
\end{hedefgecisi}

\leantac{intro h} adımına dikkat edin: Hedef \lc{¬(P ∧ ¬P)} idi ve
görünürde bir ok yoktu. Ama \lc{¬X}, \lc{X → False}'un kısaltması
olduğu için \leantac{intro} onu açabildi ve hedef \lc{False} oldu.
Son adımda \lc{hnp : P → False} fonksiyonunu \lc{hp} argümanıyla
çağırarak \lc{False}'u elde ettik.

\leandosya[Çelişmezlik, terim modu]{Bolum03}{degil_giris_terim}

\subsection{Çift değillemenin girişi}

\leandosya[\texorpdfstring{$P \to \neg\neg P$}{P → ¬¬P}]{Bolum03}{cift_degil_giris}

\noindent Bu ispatı açık açık yazarsak:
\[
  P \to \neg\neg P
  \;=\; P \to (\neg P \to \bot)
  \;=\; P \to ((P \to \bot) \to \bot).
\]
\leantac{intro hp hnp} iki oku birden açar. Bu teoremin \emph{tersi},
yani $\neg\neg P \to P$, bu bölümde ispatlayamayacağımız ilk önemli
önermedir.

\begin{dikkat}[Dikkat: $\neg\neg P$ ile $P$ aynı şey değildir]
Okulda öğrendiğimiz mantıkta ``değilin değili'' kendisidir. \Lean{}'in
varsayılan mantığında ise $P$'den $\neg\neg P$ çıkar, ama tersi
çıkmaz. $\neg\neg P$ yalnızca ``$P$'nin yanlış olduğunu gösteremezsin''
der; $P$'nin bir kanıtını \emph{vermez}. Bu farkın ne kadar derin
olduğunu \cref{bol:yapici-klasik}'te göreceğiz.
\end{dikkat}

% ---------------------------------------------------------------------
\section{Çelişkiyle çalışmanın üç yolu}\label{sec:celiski}

Bağlamda hem \lc{hp : P} hem \lc{hnp : ¬P} varsa, hedef ne olursa olsun
ispat bitmiştir. \Lean{} bunu söylemek için üç farklı araç sunar:

\leandosya[Aynı ispat, üç araçla]{Bolum03}{celiski_uc_yol}

\begin{center}
\renewcommand{\arraystretch}{1.45}
\rowcolors{2}{tablosatir}{white}
\begin{tabularx}{\linewidth}{>{\raggedright\arraybackslash}p{3.0cm} X}
\rowcolor{tabloust}
\tablobaslik{Araç} & \tablobaslik{Ne zaman kullanılır?} \\
\leantac{exfalso}       & Hedefi \lc{False}'a çevirmek istediğinizde. Ardından çelişkiyi kendiniz kurarsınız. \\
\lc{absurd hp hnp}      & Çelişkili iki terimi doğrudan elinizde tuttuğunuzda (terim modu). \\
\leantac{contradiction} & Bağlamda açık bir çelişki olduğunda: \lc{P} ve \lc{¬P}, \lc{h : False}, ya da \lc{h : a ≠ a} gibi. \\
\end{tabularx}
\end{center}

\leantac{contradiction} en kısa olanıdır, ama yalnızca çelişki bağlamda
\emph{açıkça} duruyorsa çalışır. Çelişkiyi elde etmek için birkaç adım
gerekiyorsa \leantac{exfalso} ile başlamak daha açık bir ispat verir.

% ---------------------------------------------------------------------
\section{Kontrapozisyon ve modus tollens}\label{sec:kontrapozisyon}

``Yağmur yağarsa yer ıslanır'' doğruysa, ``yer ıslak değilse yağmur
yağmamıştır'' da doğrudur. Bu, \emph{kontrapozisyon} ilkesidir:
\[
  (P \to Q) \to (\neg Q \to \neg P).
\]

\leandosya[Kontrapozisyon]{Bolum03}{kontrapozisyon}

\noindent \leantac{intro hpq hnq hp} adımında \emph{üç} varsayım
içeri alınıyor. Üçüncüsü, sonuçtaki \lc{¬P}'nin içinde saklı olan
\lc{P}'dir. Bu adımdan sonra durum şöyledir:

\begin{hedefdurumu}[\texttt{intro hpq hnq hp} sonrası]
P Q : Prop
hpq : P → Q
hnq : ¬Q
hp : P
⊢ False
\end{hedefdurumu}

\noindent Terim modunda kontrapozisyon, yine bir fonksiyon bileşkesidir:
\lc{hpq}'yu uygulayıp sonucu \lc{hnq}'ya veririz.

\leandosya[Kontrapozisyon, terim modu]{Bolum03}{kontrapozisyon_terim}

\begin{figure}[htbp]
\centering
\begin{tikzpicture}[
    kanit/.style={draw=leanframe, fill=leanbg, thick, circle,
                  minimum size=0.95cm, font=\bfseries\color{leanframe}},
    makine/.style={draw=leanhyp, fill=leanhyp!12, thick, rounded corners=3pt,
                   minimum width=1.5cm, minimum height=0.9cm,
                   font=\leanfont\bfseries\color{leanhyp}},
    bos/.style={draw=hataframe, fill=hatabg, thick, circle,
                minimum size=0.95cm, font=\bfseries\color{hataframe}},
    ok/.style={-{Stealth[length=3mm]}, very thick, leansym},
  ]
  \node[kanit] (p) {$P$};
  \node[makine, right=0.9cm of p] (f) {hpq};
  \node[kanit, right=0.9cm of f] (q) {$Q$};
  \node[makine, right=0.9cm of q] (g) {hnq};
  \node[bos, right=0.9cm of g] (b) {$\bot$};
  \draw[ok] (p) -- (f); \draw[ok] (f) -- (q);
  \draw[ok] (q) -- (g); \draw[ok] (g) -- (b);
  \draw[decorate, decoration={brace, amplitude=6pt, mirror}, thick, vurgu]
    ([yshift=-0.75cm]p.west) -- ([yshift=-0.75cm]b.east)
    node[midway, below=9pt, font=\leanfont\small, text=vurgu]
    {fun hp => hnq (hpq hp)  :  P → False  =  ¬P};
\end{tikzpicture}
\caption{Kontrapozisyon: $P$'den $\bot$'a giden yol, $\neg P$'nin kanıtıdır.}
\label{fig:kontrapozisyon}
\end{figure}

Kontrapozisyonun \emph{tersi}, yani $(\neg Q \to \neg P) \to (P \to Q)$,
klasik mantıkta doğrudur ama yapıcı olarak ispatlanamaz. Bunu
\cref{bol:yapici-klasik}'te ispatlayacağız.

\leandosya[Modus tollens]{Bolum03}{modus_tollens}

\noindent Burada daha önce ispatladığımız bir teoremi (\lc{kontrapozisyon})
bir fonksiyon gibi çağırdık. Teoremler, bir kez ispatlandıktan sonra
yeniden kullanılabilen fonksiyonlardır.

% ---------------------------------------------------------------------
\section{Denklik (\texorpdfstring{$\leftrightarrow$}{↔})}\label{sec:denklik}

``$P$ ancak ve ancak $Q$'' önermesi, $P \to Q$ ile $Q \to P$'nin
birlikte doğru olduğunu söyler. \Lean{}'de \lc{P ↔ Q}, bu iki implikasyonun
bir çiftidir; dolayısıyla kuralları \lc{∧}'nınkilere çok benzer.

\begin{center}
\begin{tabular}{c@{\hspace{3em}}c@{\hspace{2em}}c}
\begin{prooftree}
  \hypo{[P]^1}
  \ellipsis{}{Q}
  \hypo{[Q]^2}
  \ellipsis{}{P}
  \infer2[$\leftrightarrow$G$^{1,2}$]{P \leftrightarrow Q}
\end{prooftree}
&
\begin{prooftree}
  \hypo{P \leftrightarrow Q}
  \hypo{P}
  \infer2[$\leftrightarrow$E$_1$]{Q}
\end{prooftree}
&
\begin{prooftree}
  \hypo{P \leftrightarrow Q}
  \hypo{Q}
  \infer2[$\leftrightarrow$E$_2$]{P}
\end{prooftree}
\\[10pt]
\small\textcolor{ikincil}{\bfseries giriş:} \leantac{constructor}, \lc{⟨f, g⟩} &
\multicolumn{2}{c}{\small\textcolor{vurgu}{\bfseries eleme:} \lc{h.mp}, \lc{h.mpr}, \leantac{rw [h]}}
\end{tabular}
\end{center}

\leandosya[Denkliğe giriş]{Bolum03}{iff_giris}
\leandosya[Denkliğe giriş, terim modu]{Bolum03}{iff_giris_terim}
\leandosya[Denkliğin elemesi]{Bolum03}{iff_eleme}

\noindent \lc{mp} adı \emph{modus ponens}'ten gelir; \lc{mpr} ise ``ters
yönde modus ponens'' demektir.

\subsection{Yeniden yazma: \texorpdfstring{\leantac{rw}}{rw}}

Denk önermeler, \emph{her bağlamda} birbirinin yerine kullanılabilir.
\leantac{rw [h]} taktiği, \lc{h : P ↔ Q} denkliğini kullanarak hedefteki
her \lc{P}'yi \lc{Q} ile değiştirir:

\leandosya[Yeniden yazma]{Bolum03}{iff_rw}

\begin{hedefgecisi}
\begin{hedefdurumu}[\texttt{rw [h]} öncesi]
h : P ↔ Q
hqr : Q → R
⊢ P → R
\end{hedefdurumu}
\sonra
\begin{hedefdurumu}[\texttt{rw [h]} sonrası]
h : P ↔ Q
hqr : Q → R
⊢ Q → R
\end{hedefdurumu}
\end{hedefgecisi}

\noindent \leantac{rw} hem eşitliklerle (\lc{a = b}) hem denkliklerle
(\lc{P ↔ Q}) çalışır. Eşitlikle yeniden yazmayı \cref{bol:yuklemler}'da
ayrıntılı olarak göreceğiz. Denklikle yeniden yazma ise \Lean{}'in
\emph{önerme dışsallığı} (\lc{propext}) aksiyomuna dayanır: Denk iki önerme
\emph{eşittir}. Bu aksiyoma \cref{bol:aksiyomlar}'da döneceğiz.

\subsection{Üçlü değilleme}

İki değilleme bir taneye indirgenemez, ama üç değilleme bire indirgenir:

\leandosya[Üçlü değilleme]{Bolum03}{uclu_degilleme}

\noindent Bu ispat, daha önce ispatladığımız \lc{cift\_degilleme\_giris}
teoremini yeniden kullanıyor. Sonuç şaşırtıcıdır: Yapıcı mantıkta
değillemeler dizisi $\neg\neg\neg\neg\neg P$ gibi ne kadar uzarsa uzasın,
yalnızca iki farklı önermeye indirgenir: $\neg P$ ya da $\neg\neg P$
(ve elbette $P$'nin kendisi).

% ---------------------------------------------------------------------
\section{De Morgan yasaları: yapıcılığın sınırı}\label{sec:demorgan}

Augustus De Morgan'ın adını taşıyan yasalar, değillemenin ve ile veya
üzerine nasıl dağıldığını söyler. Klasik mantıkta dört yönün hepsi
doğrudur:
\[
  \neg(P \vee Q) \leftrightarrow \neg P \wedge \neg Q
  \qquad\qquad
  \neg(P \wedge Q) \leftrightarrow \neg P \vee \neg Q
\]

\leandosya[De Morgan: veya'nın değillemesi]{Bolum03}{demorgan_veya}
\leandosya[De Morgan: ve'nin değillemesi, kolay yön]{Bolum03}{demorgan_ve_kolay}

\noindent Dördüncü yönü denediğimizde ise takılırız:

\begin{leankod}[Takıldığımız yer (derlenmez)]
theorem demorgan_ve_zor (P Q : Prop) : ¬(P ∧ Q) → ¬P ∨ ¬Q := by
  intro h
  -- h : ¬(P ∧ Q),  hedef: ⊢ ¬P ∨ ¬Q
  -- Or.inl mi, Or.inr mi? Hangisini seçeceğimizi bilmiyoruz!
  sorry
\end{leankod}

\begin{figure}[htbp]
\centering
\begin{tikzpicture}[
    yon/.style={draw=#1, fill=#1!8, thick, rounded corners=4pt,
                minimum width=3.7cm, minimum height=0.85cm, align=center,
                font=\small},
    ok/.style={-{Stealth[length=2.6mm]}, very thick, #1},
  ]
  \node[yon=anarenk] (a1) at (0,0)   {$\neg(P \vee Q)$};
  \node[yon=anarenk] (a2) at (7,0)   {$\neg P \wedge \neg Q$};
  \node[yon=anarenk] (b1) at (0,-2.1){$\neg(P \wedge Q)$};
  \node[yon=anarenk] (b2) at (7,-2.1){$\neg P \vee \neg Q$};

  \draw[ok=bittiframe] ([yshift=4pt]a1.east) -- node[above, font=\baslikfont\scriptsize\bfseries, text=bittiframe]{yapıcı {\sembolfont ✓}} ([yshift=4pt]a2.west);
  \draw[ok=bittiframe] ([yshift=-4pt]a2.west) -- node[below, font=\baslikfont\scriptsize\bfseries, text=bittiframe]{yapıcı {\sembolfont ✓}} ([yshift=-4pt]a1.east);
  \draw[ok=bittiframe] ([yshift=-4pt]b2.west) -- node[below, font=\baslikfont\scriptsize\bfseries, text=bittiframe]{yapıcı {\sembolfont ✓}} ([yshift=-4pt]b1.east);
  \draw[ok=hataframe, dashed] ([yshift=4pt]b1.east) -- node[above, font=\baslikfont\scriptsize\bfseries, text=hataframe]{yalnızca klasik} ([yshift=4pt]b2.west);
\end{tikzpicture}
\caption{De Morgan yasalarının dört yönü. Kesikli ok, yalnızca klasik
mantıkla ispatlanabilen yöndür.}
\label{fig:demorgan}
\end{figure}

Sorun şudur: Hedef bir \emph{veya} önermesidir ve onu kanıtlamanın tek
yolu (\cref{sec:veya}), iki taraftan \emph{birinin} kanıtını vermektir.
Elimizdeki tek bilgi ise ``$P$ ve $Q$ birlikte doğru olamaz''dır. Bu
bilgi, \emph{hangisinin} yanlış olduğunu söylemez. Gerçek hayattan bir
örnek:

\begin{gunlukornek}
``Bu bilgisayar programı hem hızlı hem doğru olamaz'' dediğimizi
düşünelim. Bundan ``program ya yavaştır ya da yanlıştır'' sonucu çıkar
mı? Klasik mantıkçı ``elbette'' der. Yapıcı mantıkçı ise şunu sorar:
\emph{Hangisi?} Programın yavaş olduğunu mu kanıtlayacaksın, yoksa yanlış
olduğunu mu? İkisinden birini gösteremiyorsan, ``ya o ya bu'' demeye
hakkın yoktur.
\end{gunlukornek}

Bu ``hangisi?'' sorusu, yapıcı mantığın özüdür. Yine de elimiz boş
değildir: Çift değilleme altında De Morgan'ın zor yönü yapıcı olarak
ispatlanabilir (Alıştırma~\ref{alis:demorgan-cift}). Bir sonraki kesim,
bu olgunun çok daha genel bir örneğini gösteriyor.

% ---------------------------------------------------------------------
\section{Üçüncü hâlin olanaksızlığı yanlış olamaz}\label{sec:em-yanlis-olamaz}

Klasik mantığın en ünlü ilkesi, her önermenin ya doğru ya yanlış
olduğunu söyleyen \emph{üçüncü hâlin olanaksızlığıdır} (Latince
\emph{tertium non datur}): $P \vee \neg P$. Bu önermeyi de yapıcı olarak
ispatlayamayız: $P$'yi mi yoksa $\neg P$'yi mi kanıtlayacağımızı bilmeden
\lc{Or.inl} ile \lc{Or.inr} arasında seçim yapamayız. Ama şaşırtıcı bir
şey yapabiliriz: Bu ilkenin \emph{yanlış olamayacağını} ispatlayabiliriz.

\leandosya[\texorpdfstring{$\neg\neg(P \vee \neg P)$}{¬¬(P ∨ ¬P)}]{Bolum03}{em_cift_degil}

Bu ispat, mantıkta bir tür ``hile'' içerir ve adım adım izlenmeye değer:

\begin{enumerate}[adimlar]
  \item \lc{h : ¬(P ∨ ¬P)} varsayıyoruz ve bir çelişki arıyoruz.
  \item \leantac{apply h}: Çelişki elde etmenin bir yolu, \lc{h}'yi
        çürütmek, yani \lc{P ∨ ¬P}'nin bir kanıtını vermek.
  \item \leantac{apply Or.inr}: \emph{Cesur bir seçim yapıyoruz}:
        $\neg P$ tarafını kanıtlayacağız.
  \item \leantac{intro hp}: $\neg P$'yi kanıtlamak için $P$'yi varsayıp
        çelişki arıyoruz.
  \item \leantac{apply h}, \leantac{exact Or.inl hp}: Ama artık elimizde
        $P$ var! Seçimimizi değiştirip \lc{P ∨ ¬P}'nin \emph{sol}
        tarafını kanıtlıyor ve \lc{h} ile çelişki elde ediyoruz.
\end{enumerate}

\begin{gunlukornek}[Bir benzetme: şeytanın teklifi]
Bu ispatın programlama karşılığı, mantıkçılar arasında ``şeytanın
teklifi'' olarak anlatılan bir hikâyeye benzer. Şeytan der ki: ``Ya sana
bir milyar lira vereceğim, ya da bir milyar lira ödersen dilediğin
her şeyi yapacağım. Hangisini yapacağıma ben karar veririm.'' Ve ikinci
seçeneği söyler. Yıllar sonra sen bir milyar lirayı biriktirip getirdiğinde,
şeytan ``Fikrimi değiştirdim: birinci seçenek. İşte paran'' der ve az önce
verdiğin parayı sana geri verir. İspattaki \lc{Or.inr}'den \lc{Or.inl}'ye
geçiş tam olarak budur: Önce $\neg P$ diyor, $P$'nin kanıtı gelince
o kanıtı kullanarak ``aslında $P$'' diyoruz.
\end{gunlukornek}

\leandosya[\texorpdfstring{$\neg\neg(P \vee \neg P)$}{¬¬(P ∨ ¬P)}, terim modu]{Bolum03}{em_cift_degil_terim}

Elimizde artık iki gözlem var: \lc{P ∨ ¬P}'yi yapıcı olarak
ispatlayamıyoruz, ama \lc{¬¬(P ∨ ¬P)}'yi ispatlayabiliyoruz. Arada kalan
tek adım $\neg\neg X \to X$, yani \emph{çift değilleme eleme}dir. Bu adım,
yapıcı mantık ile klasik mantığı ayıran sınırın kendisidir.
\Cref{bol:yapici-klasik} tamamen bu sınıra ayrılmıştır.

% ---------------------------------------------------------------------
\section{Curry--Howard tablosuna eklemeler}\label{sec:curry-howard-2}

\begin{center}
\renewcommand{\arraystretch}{1.5}
\rowcolors{2}{tablosatir}{white}
\begin{tabular}{>{\raggedright\arraybackslash}p{2.7cm}
                >{\raggedright\arraybackslash}p{4.0cm}
                >{\raggedright\arraybackslash}p{5.0cm}}
\rowcolor{tabloust}
\tablobaslik{Mantık} & \tablobaslik{Programlama} & \tablobaslik{Lean'deki kanıt} \\
$\top$              & tek elemanlı tip (\emph{unit}) & \lc{True.intro}, \leantac{trivial} \\
$\bot$              & boş tip                         & --- (kanıt yoktur) \\
$\bot$ eleme        & boş tipten fonksiyon            & \lc{False.elim h}, \leantac{cases h} \\
$\neg P$            & $P$'den boş tipe fonksiyon      & \lc{fun hp => ...} (\lc{False} üretir) \\
$P \leftrightarrow Q$ & iki yönlü dönüşüm çifti       & \lc{⟨f, g⟩} \\
\end{tabular}
\end{center}

% ---------------------------------------------------------------------
\section{Taktik özeti}

\begin{center}
\renewcommand{\arraystretch}{1.45}
\rowcolors{2}{tablosatir}{white}
\begin{tabularx}{\linewidth}{>{\raggedright\arraybackslash}p{3.1cm}
                             >{\raggedright\arraybackslash}p{1.8cm}
                             >{\raggedright\arraybackslash}X}
\rowcolor{tabloust}
\tablobaslik{Taktik / terim} & \tablobaslik{Kural} & \tablobaslik{Ne yapar?} \\
\leantac{trivial}               & $\top$G & \lc{⊢ True} hedefini kapatır. \\
\leantac{exfalso}               & $\bot$E & Hedefi \lc{False}'a çevirir. \\
\lc{False.elim h}               & $\bot$E & \lc{h : False}'tan her şeyi üretir. \\
\lc{absurd hp hnp}              & $\neg$E, $\bot$E & \lc{P} ve \lc{¬P}'den her şeyi üretir. \\
\leantac{contradiction}         & ---     & Bağlamdaki açık çelişkiyi bulur. \\
\leantac{intro hp} (hedef \lc{¬P}) & $\neg$G & \lc{hp : P} varsayar, hedef \lc{False} olur. \\
\leantac{constructor} (hedef \lc{↔}) & $\leftrightarrow$G & İki yönü ayrı hedefler olarak açar. \\
\lc{h.mp}, \lc{h.mpr}           & $\leftrightarrow$E & Denkliği bir yönde kullanır. \\
\leantac{rw [h]}                & ---     & \lc{h : P ↔ Q} ile hedefteki \lc{P}'leri \lc{Q} yapar. \\
\end{tabularx}
\end{center}

% ---------------------------------------------------------------------
\begin{arastirmanotu}
\textbf{Değillemenin tanımı.} $\neg P := P \to \bot$ tanımı, Heyting'in
1930'da sezgici mantığı biçimselleştirirken yaptığı seçime dayanır
\citep{heyting1930}. Bu tanım, değillemeyi bağımsız bir kavram olmaktan
çıkarıp implikasyona indirger ve Curry--Howard eşlemesinin
değillemeye de uzanmasını sağlar \citep{sorensen2006}.

\textbf{İki değilleme sınırı.} Yapıcı önermeler mantığında
$\neg\neg\neg P \leftrightarrow \neg P$ olduğundan, tek bir önerme
değişkeninden yalnızca değilleme ile sonlu sayıda farklı önerme elde
edilir. Tek değişkenli önermelerin tamamı ise sonsuz bir yapı oluşturur;
bu yapı \emph{Rieger--Nishimura kafesi} olarak bilinir.

\textbf{Paraconsistent mantıklar.} Ex falso ilkesi bütün mantıklarda
kabul edilmez. \emph{Paraconsistent} (çelişkiye dayanıklı) mantıklar,
çelişkili bilgiyle çalışan sistemlerde (ör. tutarsız veri tabanları)
her şeyin türetilebilmesini engellemek için bu ilkeyi bırakır. Bu
mantıklar kitabın kapsamı dışındadır.
\end{arastirmanotu}

\begin{bolumozeti}
\begin{itemize}
  \item \lc{True} tek elemanlı, \lc{False} boş tiptir. \lc{False}'tan her
        şey çıkar (\emph{ex falso}).
  \item \lc{¬P}, \lc{P → False}'un kısaltmasıdır. Değillemeyi kanıtlamak
        için \leantac{intro}, kullanmak için fonksiyon uygulaması yeter.
  \item Çelişki üç yoldan kullanılır: \leantac{exfalso}, \lc{absurd},
        \leantac{contradiction}.
  \item \lc{P ↔ Q} iki implikasyonun çiftidir; \lc{h.mp}, \lc{h.mpr} ve
        \leantac{rw} ile kullanılır.
  \item Yapıcı mantıkta $P \to \neg\neg P$ doğrudur, tersi değil. De
        Morgan'ın bir yönü, $P \vee \neg P$ ve $\neg\neg P \to P$
        yapıcı olarak ispatlanamaz; ama $\neg\neg(P \vee \neg P)$
        ispatlanabilir.
\end{itemize}
\end{bolumozeti}

\begin{alistirmalar}
\alistirma{1} \lc{¬False} önermesini ispatlayın. Bu önermenin
\lc{False → False}'tan farkı var mı? (\lc{a3\_1})

\alistirma{1} \lc{P → True} önermesini ispatlayın. (\lc{a3\_2})

\alistirma{1} \lc{P → ¬P → Q} önermesini önce \leantac{exfalso}, sonra
\lc{absurd} ile ispatlayın. (\lc{a3\_3})

\alistirma{1} Denkliğin yansıma özelliğini ispatlayın: \lc{P ↔ P}.
(\lc{a3\_4})

\alistirma{2} Denkliğin simetrisini (\lc{a3\_5}) ve geçişliliğini
(\lc{a3\_6}) ispatlayın.

\alistirma{2} \emph{Disjunktif kıyas}: \lc{P ∨ Q → ¬P → Q}.
``Ya çay ya kahve içeceğim; çay içmeyeceğim; o hâlde kahve içeceğim.''
(\lc{a3\_7})

\alistirma{2} Hiçbir önerme kendi değillemesine denk olamaz:
\lc{¬(P ↔ ¬P)}. İpucu: Önce \leantac{have} ile \lc{¬P}'yi elde edin.
Bu alıştırmanın \emph{yalancı paradoksu} (``Bu cümle yanlıştır'') ile
ilişkisini tartışın. (\lc{a3\_8})

\alistirma{2} \lc{¬(P → Q) → ¬Q} önermesini ispatlayın. (\lc{a3\_9})

\alistirma{2} Çift değillemenin implikasyona dağıldığını ispatlayın:
\lc{¬¬(P → Q) → ¬¬P → ¬¬Q}. (\lc{a3\_10})

\alistirma{3}\label{alis:demorgan-cift} De Morgan'ın zor yönünün çift
değilleme altında yapıcı olduğunu ispatlayın:
\lc{¬(P ∧ Q) → ¬¬(¬P ∨ ¬Q)}. İpucu: \lc{em\_yanlis\_olamaz}
ispatındaki ``fikir değiştirme'' hilesini kullanın. (\lc{a3\_11})

\alistirma{3}\label{alis:peirce-cift} Bölüm~2'deki Peirce yasasının
çift değilleme altında yapıcı olduğunu ispatlayın:
\lc{¬¬(((P → Q) → P) → P)}. (\lc{a3\_12})

\alistirma{3} (Kâğıt-kalem) \lc{em\_yanlis\_olamaz} ispatının doğal
çıkarım ağacını çizin. \lc{h} varsayımı ağaçta kaç kez kullanılıyor?
Bu ispatı, varsayımı yalnızca bir kez kullanan bir ispatla
değiştirmek mümkün müdür? Tartışın.
\end{alistirmalar}
   % Önermeler mantığı II: ¬, False, ↔
% =====================================================================
%  Bölüm 4 — Yapıcı ve Klasik Mantık
% =====================================================================
\chapter{Yapıcı ve Klasik Mantık}\label{bol:yapici-klasik}

\begin{bolumhedefleri}
\begin{itemize}
  \item Brouwer--Heyting--Kolmogorov (BHK) yorumunu ve ``kanıt bir
        inşadır'' fikrini,
  \item \lc{Classical.em} ve \lc{Classical.byContradiction} ile klasik
        ispat yapmayı,
  \item Üçüncü hâlin olanaksızlığı, çift değilleme eleme ve Peirce
        yasasının birbirine denk olduğunu,
  \item \lc{\#print axioms} ile bir ispatın klasik olup olmadığını
        denetlemeyi,
  \item Bir Kripke modeliyle $P \vee \neg P$'nin neden yapıcı olarak
        ispatlanamadığını,
  \item Glivenko teoremini, çift değilleme çevirisini ve karar
        verilebilir önermeleri.
\end{itemize}
\end{bolumhedefleri}

Önceki iki bölümde \Lean{}'in bazı önermeleri ispatlamamıza izin
vermediğini gördük: Peirce yasası (Alıştırma~\ref{alis:peirce}),
De Morgan'ın bir yönü, $\neg\neg P \to P$ ve $P \vee \neg P$. Bu bölüm
bu sınırın ne olduğunu, neden var olduğunu ve nasıl aşılacağını
anlatıyor.

% ---------------------------------------------------------------------
\section{İki kanıt anlayışı}\label{sec:iki-anlayis}

\subsection{Yapıcı olmayan bir ispat}

Şu teoremi ve onun ünlü ispatını ele alalım.

\begin{teorem}{}{irrasyonel-us}
$a^b$ rasyonel olacak şekilde $a$ ve $b$ irrasyonel sayıları vardır.
\end{teorem}

\noindent\textit{İspat.} $x = \sqrt{2}^{\sqrt{2}}$ sayısına bakalım.
Bu sayı ya rasyoneldir ya değildir.
\begin{itemize}
  \item $x$ rasyonelse, $a = b = \sqrt{2}$ alırız. İşimiz biter.
  \item $x$ rasyonel değilse, $a = x$ ve $b = \sqrt{2}$ alırız. O zaman
        $a^b = \bigl(\sqrt{2}^{\sqrt{2}}\bigr)^{\sqrt{2}} = \sqrt{2}^{\,2} = 2$
        olur ve bu rasyoneldir. \hfill\qedsymbol
\end{itemize}

İspat kısa, zarif ve kusursuzdur, ama tuhaf bir eksiği vardır: Sonunda
$a$'nın ve $b$'nin \emph{ne olduğunu} hâlâ bilmiyoruz. İspat, ``ya
$(\sqrt2, \sqrt2)$ ya da $(\sqrt2^{\sqrt2}, \sqrt2)$ çiftinin işe yaradığını''
gösterir, ama hangisinin işe yaradığını söylemez. Bunu söyleyemez, çünkü
ilk adımdaki ``ya rasyoneldir ya değildir'' cümlesi, üçüncü hâlin
olanaksızlığının bir örneğidir ve hangisinin doğru olduğu hakkında hiçbir
bilgi taşımaz.

\begin{tarihce}[Tarihçe: Brouwer ve Hilbert]
Hollandalı matematikçi L.\,E.\,J. Brouwer, 1908'de yayımladığı bir
makalede mantık ilkelerinin, özellikle de üçüncü hâlin olanaksızlığının,
sonsuz kümeler üzerinde güvenilir olmadığını ileri sürdü
\citep{brouwer1908}. Brouwer'e göre matematiksel bir nesnenin var
olduğunu söylemek, onu zihinde \emph{inşa edebilmek} demekti. Bu görüş
\emph{sezgicilik} (İng.\ \emph{intuitionism}) adını aldı.

David Hilbert bu programa sert biçimde karşı çıktı. Ona göre üçüncü hâlin
olanaksızlığını matematikçinin elinden almak, gökbilimcinin elinden
teleskobunu ya da boksörün elinden yumruklarını almak gibiydi. 1920'lerde
yaşanan bu tartışma, matematiğin temelleri üzerine yirminci yüzyılın en
ünlü tartışmalarından biri oldu. Bugün ise iki taraf barışık
yaşıyor: Yapıcı mantık, hesaplanabilirliği gözeten bir \emph{alt sistem}
olarak, klasik mantık ise onun bir \emph{uzantısı} olarak görülüyor.
\end{tarihce}

\subsection{BHK yorumu}

Brouwer'in fikri, Arend Heyting ve Andrey Kolmogorov tarafından kesin bir
biçime kavuşturuldu \citep{heyting1930, kolmogorov1932}. Bugün
\emph{Brouwer--Heyting--Kolmogorov (BHK) yorumu} olarak bilinen bu
yaklaşım, her bağlaç için bir kanıtın \emph{ne olduğunu} söyler
(\cref{tab:bhk}).

\begin{table}[htbp]
\centering
\renewcommand{\arraystretch}{1.5}
\rowcolors{2}{tablosatir}{white}
\begin{tabular}{>{\centering\arraybackslash}p{1.7cm}
                >{\raggedright\arraybackslash}p{6.2cm}
                >{\raggedright\arraybackslash}p{3.6cm}}
\rowcolor{tabloust}
\tablobaslik{Önerme} & \tablobaslik{Bir kanıtı şudur…} & \tablobaslik{Lean'de} \\
$P \wedge Q$      & $P$'nin bir kanıtı \emph{ve} $Q$'nun bir kanıtından oluşan bir çift & \lc{⟨hp, hq⟩} \\
$P \vee Q$        & $P$'nin bir kanıtı \emph{ya da} $Q$'nun bir kanıtı, ve \textbf{hangisi olduğunu gösteren bir etiket} & \lc{Or.inl hp}, \lc{Or.inr hq} \\
$P \to Q$         & $P$'nin her kanıtını $Q$'nun bir kanıtına dönüştüren bir \emph{yöntem} & \lc{fun hp => ...} \\
$\neg P$          & $P$'nin her kanıtını bir çelişkiye dönüştüren bir yöntem & \lc{fun hp => ...} \\
$\bot$            & kanıtı yoktur & --- \\
\end{tabular}
\caption{BHK yorumu. Önceki iki bölümde Lean'de yaptığımız her şey bu
tabloya uyar.}
\label{tab:bhk}
\end{table}

Tablonun kalın yazılmış satırı her şeyi açıklar. BHK yorumunda
$P \vee \neg P$'nin bir kanıtı, ya $P$'nin ya $\neg P$'nin bir kanıtı
olmalıdır, \emph{ve hangisi olduğu belli olmalıdır}. Her önerme için
böyle bir kanıt vermek, matematikteki her sorunu çözebilen genel bir
yöntem bulmak demektir. Böyle bir yöntem yoktur (bu olguyu
\ciltiki{6}'da Turing'in sonuçlarıyla kesinleştireceğiz). Bu
yüzden $P \vee \neg P$, BHK yorumunda genel bir yasa olamaz.

\begin{onemlifikir}[Önemli Fikir: Yapıcı mantık daha az şey söylemez, daha çok şey söyler]
Yapıcı bir ispat, klasik bir ispattan \emph{daha fazla bilgi} taşır.
``$a^b$ rasyonel olan irrasyonel $a, b$ vardır'' cümlesinin yapıcı bir
ispatı, bu sayıları bize \emph{verir}. Klasik ispat ise yalnızca
var olduklarını söyler. Bu nedenle yapıcı mantıkta ispatlanan her şey
klasik mantıkta da doğrudur, ama tersi geçerli değildir.
\end{onemlifikir}

% ---------------------------------------------------------------------
\section{Lean'de klasik mantık}\label{sec:lean-klasik}

\Lean{}'de klasik mantık yasaklanmamıştır; yalnızca \emph{açıkça
istenmesi} gerekir. Kapı, \lc{Classical} ad alanındaki iki teoremdir:

\leandosya[Klasik mantığın iki kapısı]{Bolum04}{em_check}

\subsection{Çift değilleme eleme}

\lc{Classical.em P}, \lc{P ∨ ¬P} türünde bir terimdir. Ona \leantac{cases}
uyguladığımızda, iki durum açılır: $P$ doğru ya da $\neg P$ doğru. Bu
\emph{klasik durum ayırma}, yapıcı mantıkta mümkün olmayan şeylerin
kapısını açar:

\leandosya[Çift değilleme eleme]{Bolum04}{cift_degil_eleme}

\begin{figure}[htbp]
\centering
\begin{tikzpicture}[
    durum/.style={draw=#1, fill=#1!8, thick, rounded corners=4pt,
                  minimum width=3.4cm, minimum height=0.9cm, align=center,
                  font=\small},
    ok/.style={-{Stealth[length=2.6mm]}, thick, #1},
  ]
  \node[durum=teoremframe] (em) at (0,0) {\lc{Classical.em P : P ∨ ¬P}};
  \node[durum=ikincil] (a) at (-3.3,-1.7) {\lc{hp : P}};
  \node[durum=vurgu]   (b) at (3.3,-1.7)  {\lc{hnp : ¬P}};
  \node[durum=ikincil] (a2) at (-3.3,-3.3) {\lc{exact hp}};
  \node[durum=vurgu]   (b2) at (3.3,-3.3)  {\lc{hnnp hnp : False}\\[-1pt]{\scriptsize\lc{absurd hnp hnnp}}};
  \node[durum=bittiframe] (r) at (0,-4.9) {\lc{P} kanıtlandı};
  \draw[ok=ikincil] (em) -- node[left=6pt, font=\leanfont\scriptsize\bfseries, text=ikincil]{inl} (a);
  \draw[ok=vurgu]   (em) -- node[right=6pt, font=\leanfont\scriptsize\bfseries, text=vurgu]{inr} (b);
  \draw[ok=ikincil] (a) -- (a2);  \draw[ok=vurgu] (b) -- (b2);
  \draw[ok=bittiframe] (a2) -- (r); \draw[ok=bittiframe] (b2) -- (r);
  \node[font=\baslikfont\scriptsize, text=teoremframe, align=center] at (0,0.95)
       {klasik mantık burada devreye girer};
\end{tikzpicture}
\caption{\lc{Classical.em} ile durum ayırma. İkinci durumda elimizdeki
\lc{¬¬P} ile \lc{¬P} çelişir.}
\label{fig:klasik-durum}
\end{figure}

\subsection{Çelişki ile ispat}

Klasik matematiğin en sık kullanılan yöntemlerinden biri \emph{çelişki
ile ispattır} (Latince \emph{reductio ad absurdum}): ``$Q$'yu göstermek
için $Q$'nun yanlış olduğunu varsayalım ve bir çelişkiye ulaşalım.''

\leandosya[Çelişki ile ispat]{Bolum04}{celiski_ile_ispat}

\begin{dikkat}[Dikkat: İki tür ``çelişki ile ispat'']
Matematikçiler ``çelişki ile ispat'' derken iki farklı şeyi kastedebilir:
\begin{enumerate}[label=(\alph*)]
  \item \textbf{$\neg P$'yi göstermek için} $P$'yi varsayıp çelişkiye
        ulaşmak. Bu, değillemenin \emph{tanımıdır} ($\neg$G kuralı) ve
        tamamen yapıcıdır. $\sqrt{2}$'nin rasyonel \emph{olmadığının}
        ispatı bu türdendir.
  \item \textbf{$P$'yi göstermek için} $\neg P$'yi varsayıp çelişkiye
        ulaşmak. Bu, aslında $\neg\neg P$'yi gösterip ardından çift
        değillemeyi elemektir ve \emph{klasik} bir adımdır.
\end{enumerate}
\Lean{}'de (a) \leantac{intro} ile, (b) ise
\lc{Classical.byContradiction} ile yapılır. İkisini ayırt etmek, bir
ispatın gerçekten klasik mantığa ihtiyaç duyup duymadığını anlamanın
ilk adımıdır.
\end{dikkat}

\subsection{Eksik kalan ispatlar}

Önceki bölümlerde takıldığımız iki önermeyi artık ispatlayabiliriz:

\leandosya[De Morgan, zor yön]{Bolum04}{demorgan_ve_klasik}
\leandosya[Peirce yasası]{Bolum04}{peirce}

\noindent İkisinde de yöntem aynıdır: \lc{Classical.em P} ile durum ayır,
iki durumda da hedefe ulaş. Bu iki ispat, Bölüm~2'de açık bıraktığımız
Peirce yasası sorusunun cevabıdır: Peirce yasası \emph{klasik} bir
yasadır.

% ---------------------------------------------------------------------
\section{Üç ilke, tek sınır}\label{sec:uc-ilke}

Şimdiye kadar klasik mantığa özgü üç ilkeyle karşılaştık. Bunları
\emph{bütün önermeler için} söylenen birer önerme olarak
tanımlayalım:

\leandosya[Üç klasik ilke]{Bolum04}{ilkeler}

\begin{teorem}{Klasik ilkelerin denkliği}{uc-ilke}
Yapıcı mantıkta aşağıdaki üç ilke birbirine denktir:
\begin{center}
\begin{tabular}{rl}
\textbf{EM} (üçüncü hâlin olanaksızlığı): & $\forall P,\; P \vee \neg P$ \\
\textbf{DNE} (çift değilleme eleme):       & $\forall P,\; \neg\neg P \to P$ \\
\textbf{Peirce yasası}:                    & $\forall P\, Q,\; ((P \to Q) \to P) \to P$
\end{tabular}
\end{center}
Yani yapıcı mantığa bunlardan herhangi birini eklemek, ötekileri de
eklemek demektir. Ortaya çıkan sistem klasik mantıktır.
\end{teorem}

\noindent İspatı bir döngü olarak kuracağız (\cref{fig:uc-ilke}):
EM $\Rightarrow$ DNE $\Rightarrow$ Peirce $\Rightarrow$ EM.

\begin{figure}[htbp]
\centering
\begin{tikzpicture}[
    ilke/.style={draw=teoremframe, fill=teorembg, very thick, rounded corners=5pt,
                 minimum width=2.6cm, minimum height=1.1cm, align=center,
                 font=\baslikfont\bfseries},
    ok/.style={-{Stealth[length=3mm]}, very thick, anarenk},
    etiket/.style={font=\leanfont\scriptsize, text=anarenk, fill=white, inner sep=2pt},
  ]
  \node[ilke] (em)  at (90:2.3)  {EM\\[-1pt]\normalfont\scriptsize $P \vee \neg P$};
  \node[ilke] (dne) at (-30:2.3) {DNE\\[-1pt]\normalfont\scriptsize $\neg\neg P \to P$};
  \node[ilke] (pe)  at (210:2.3) {Peirce\\[-1pt]\normalfont\scriptsize $((P{\to}Q){\to}P){\to}P$};
  \draw[ok] (em) to[bend left=22] node[etiket, right=2pt]{em\_dne} (dne);
  \draw[ok] (dne) to[bend left=22] node[etiket, below=2pt]{dne\_peirce} (pe);
  \draw[ok] (pe) to[bend left=22] node[etiket, left=2pt]{peirce\_em} (em);
  \node[font=\baslikfont\small\bfseries, text=koyu, align=center] at (0,-0.1)
       {hepsi\\yapıcı};
\end{tikzpicture}
\caption{Üç klasik ilkenin denkliği. Her ok, \emph{klasik mantık
kullanmadan} ispatlanmış bir \Lean{} teoremidir.}
\label{fig:uc-ilke}
\end{figure}

\leandosya[EM $\Rightarrow$ DNE]{Bolum04}{em_dne}
\leandosya[DNE $\Rightarrow$ Peirce]{Bolum04}{dne_peirce}
\leandosya[Peirce $\Rightarrow$ EM]{Bolum04}{peirce_em}

\noindent Son ispattaki hileye dikkat edin: Peirce yasasındaki $Q$
herhangi bir önerme olabileceği için \lc{False}'u seçtik. O zaman
\lc{P ∨ ¬P → False}, yani \lc{¬(P ∨ ¬P)} elimize geçer ve
\cref{sec:em-yanlis-olamaz}'deki ``fikir değiştirme'' hilesi işi
bitirir. Döngü tamamlanınca denklikler hemen çıkar:

\leandosya[Denklikler]{Bolum04}{uc_ilke_denk}

\begin{dikkat}[Dikkat: ``Her $P$ için'' ifadesinin önemi]
Teorem, ilkelerin \emph{bütün önermeler} için söylenmiş hâlleri
hakkındadır. Tek bir $P$ için \lc{P ∨ ¬P}'den \lc{¬¬P → P} çıkar, ama
\lc{¬¬P → P}'den \lc{P ∨ ¬P} çıkmaz. \lc{peirce\_em} ispatında Peirce
yasasını $P$ için değil, \lc{P ∨ ¬P} önermesi için kullandığımıza dikkat
edin. Alıştırma~\ref{alis:dne-em} bu konuyu inceliyor.
\end{dikkat}

% ---------------------------------------------------------------------
\section{Bir ispat klasik mi? \texorpdfstring{\#print axioms}{\#print axioms}}\label{sec:print-axioms}

Uzun bir ispatın bir yerinde klasik mantık kullanılıp kullanılmadığını
gözle anlamak zordur. Bir teorem, klasik mantık kullanan başka bir
teoremi çağırıyor olabilir. \Lean{} bu soruyu kesin olarak cevaplar:

\leandosya[Aksiyom denetimi]{Bolum04}{print_axioms}

\lc{em\_dne} teoremi, EM'yi bir \emph{varsayım} olarak alır; hiçbir
aksiyom kullanmaz. \lc{cift\_degilleme\_eleme} ise \lc{Classical.em}'i
çağırdığı için üç aksiyoma dayanır. Bu üç aksiyom, \Lean{}'in bütün
klasik matematiğinin temelidir:

\begin{center}
\renewcommand{\arraystretch}{1.45}
\rowcolors{2}{tablosatir}{white}
\begin{tabularx}{\linewidth}{>{\raggedright\arraybackslash}p{3.3cm} X >{\centering\arraybackslash}p{1.4cm}}
\rowcolor{tabloust}
\tablobaslik{Aksiyom} & \tablobaslik{Kabaca ne söyler?} & \tablobaslik{Bölüm} \\
\lc{propext}          & Denk önermeler eşittir: \lc{(P ↔ Q) → P = Q}. & \ref{bol:aksiyomlar} \\
\lc{Quot.sound}       & Bölüm tiplerinde denk elemanlar eşittir; fonksiyon dışsallığı bundan türetilir. & \ref{bol:aksiyomlar} \\
\lc{Classical.choice} & Boş olmayan her tipten bir eleman seçilebilir (seçim aksiyomu). & \ref{bol:aksiyomlar} \\
\end{tabularx}
\end{center}

\lc{Classical.em}, \Lean{}'de bir aksiyom değil, \emph{teoremdir}: Bu üç
aksiyomdan \emph{Diaconescu teoremi} ile türetilir \citep{diaconescu1975}.
Yani \Lean{}'de klasik mantığın asıl kaynağı seçim aksiyomudur. Bu
şaşırtıcı bağlantıyı \cref{bol:aksiyomlar}'da ayrıntılı olarak
inceleyeceğiz.

% ---------------------------------------------------------------------
\section{Neden ispatlanamaz? Kripke modelleri}\label{sec:kripke-onizleme}

Şimdiye kadar $P \vee \neg P$'yi yapıcı olarak ``ispatlayamadığımızı''
söyledik. Ama bir şeyi ispatlayamamak, onun ispatlanamaz olduğunu
göstermez; belki yeterince akıllıca bir ispat bulamamışızdır. Bir
önermenin bir sistemde ispatlanamadığını göstermek için o sistemde
\emph{geçerli olup da} önermenin \emph{yanlış olduğu} bir model
bulmamız gerekir. Yapıcı mantık için bu modeller, Saul Kripke'nin
1965'te tanıttığı modellerdir \citep{kripke1965}.

\begin{tanim}{Kripke modeli (yapıcı mantık için)}{kripke}
Bir \textbf{Kripke modeli}, şunlardan oluşur:
\begin{itemize}
  \item Bir \emph{bilgi durumları} kümesi $W$ ve bu durumlar
        arasında bir sıralama $\leq$. ($w \leq v$: ``$v$, $w$'den daha
        fazla bilgi içeren bir gelecek durumdur.'')
  \item Her durumda hangi atomik önermelerin \emph{bilindiğini}
        gösteren bir ilişki. Bir kez bilinen bir atom, sonraki bütün
        durumlarda da bilinir (\emph{kalıcılık}).
\end{itemize}
Bileşik önermelerin bir $w$ durumunda bilinmesi ($w \Vdash \varphi$)
şöyle tanımlanır:
\begin{itemize}
  \item $w \Vdash P \wedge Q$ \;ancak ve ancak\; $w \Vdash P$ ve $w \Vdash Q$,
  \item $w \Vdash P \vee Q$ \;ancak ve ancak\; $w \Vdash P$ ya da $w \Vdash Q$,
  \item $w \Vdash P \to Q$ \;ancak ve ancak\; $w$'den sonraki \emph{her}
        $v$ durumunda, $v \Vdash P$ ise $v \Vdash Q$,
  \item $w \Vdash \neg P$ \;ancak ve ancak\; $w$'den sonraki \emph{hiçbir}
        durumda $P$ bilinmiyorsa.
\end{itemize}
\end{tanim}

Bu modellerin arkasındaki sezgi, zaman içinde bilgi biriktiren bir
matematikçidir. $\neg P$'yi bilmek, ``$P$'yi \emph{hiçbir zaman}
öğrenmeyeceğim'' demektir; şu an $P$'yi bilmemek bunun için yetmez.

\begin{figure}[htbp]
\centering
\begin{tikzpicture}[
    dunya/.style={circle, draw=anarenk, very thick, fill=tanimbg,
                  minimum size=1.5cm, font=\baslikfont\bfseries},
    ok/.style={-{Stealth[length=3mm]}, very thick, anarenk},
    bilgi/.style={font=\small, align=left, anchor=west},
  ]
  \node[dunya] (w0) at (0,0) {$w_0$};
  \node[dunya, fill=bittibg, draw=bittiframe] (w1) at (0,3.2) {$w_1$};
  \draw[ok] (w0) -- node[left=4pt, font=\small, text=anarenk]{$\leq$} (w1);

  \node[bilgi] at (1.2,3.2) {%
    $P$ \textcolor{bittiframe}{\textbf{bilinir}} (atom)};
  \node[bilgi] at (1.2,0.45) {%
    $P$ \textcolor{hataframe}{\textbf{bilinmez}}};
  \node[bilgi] at (1.2,-0.15) {%
    $\neg P$ \textcolor{hataframe}{\textbf{bilinmez}}: gelecekte ($w_1$) $P$ öğreniliyor};
  \node[bilgi] at (1.2,-0.75) {%
    $P \vee \neg P$ \textcolor{hataframe}{\textbf{bilinmez}}: iki taraf da bilinmiyor};
  \node[bilgi] at (1.2,-1.35) {%
    $\neg\neg P$ \textcolor{bittiframe}{\textbf{bilinir}}: hiçbir gelecek durumda $\neg P$ yok};
  \draw[black!25, rounded corners=4pt] (-1.0,-1.8) rectangle (9.6,4.15);
\end{tikzpicture}
\caption{İki durumlu bir Kripke modeli. $w_0$ durumunda $\neg\neg P$
bilinir ama $P$ bilinmez; $P \vee \neg P$ de bilinmez.}
\label{fig:kripke-em}
\end{figure}

\Cref{fig:kripke-em}'deki modelde $w_0$ durumu, ``$P$'nin doğru olup
olmadığını henüz bilmiyoruz, ama ileride öğreneceğiz'' durumudur. Bu
durumda:

\begin{itemize}
  \item $P$ bilinmez (henüz öğrenilmedi).
  \item $\neg P$ de bilinmez, çünkü gelecekteki $w_1$ durumunda $P$
        öğreniliyor. ``$P$'yi hiçbir zaman öğrenmeyeceğim'' diyemeyiz.
  \item Dolayısıyla $P \vee \neg P$ bilinmez.
  \item Buna karşılık $\neg\neg P$ bilinir: $w_0$'dan sonraki hiçbir
        durumda $\neg P$ bilinmiyor. O hâlde bu modelde
        $\neg\neg P \to P$ de bilinmez.
\end{itemize}

\begin{teorem}{Kripke modellerine göre geçerlilik}{kripke-gecerlilik}
Yapıcı önermeler mantığında ispatlanabilen her önerme, her Kripke
modelinin her durumunda bilinir \citep{kripke1965}.
\end{teorem}

Bu teoremi kabul edersek sonuç hemen çıkar: $P \vee \neg P$, bir Kripke
modelinin bir durumunda bilinmediğine göre, yapıcı olarak
ispatlanamaz. Teoremin kendisini \ciltiki{2}'de, Kripke modellerini
\Lean{}'de tanımladıktan sonra ispatlayacağız.

% ---------------------------------------------------------------------
\section{Klasik mantık yapıcı mantığın içinde: Glivenko ve Gödel}\label{sec:glivenko}

Klasik mantık yapıcı mantıktan güçlüdür. Ama şaşırtıcı biçimde, klasik
mantık yapıcı mantığın \emph{içine} gömülebilir. Bunun ilk örneğini
\cref{sec:em-yanlis-olamaz}'te gördük: $P \vee \neg P$ yapıcı
değildir, ama $\neg\neg(P \vee \neg P)$ yapıcıdır. Aynı olgu Peirce
yasası için de geçerlidir:

\leandosya[Peirce yasası, çift değilleme altında]{Bolum04}{glivenko_peirce}

\noindent Bu bir rastlantı değildir.

\begin{teorem}{Glivenko, 1929}{glivenko}
Bir $\varphi$ önermesi \emph{klasik önermeler mantığında}
ispatlanabiliyorsa, $\neg\neg\varphi$ önermesi \emph{yapıcı önermeler
mantığında} ispatlanabilir \citep{glivenko1929}.
\end{teorem}

Glivenko teoremi, klasik önermeler mantığının her totolojisinin yapıcı
mantıkta ``yanlış olamaz'' biçiminde zaten bulunduğunu söyler. Yüklemler
mantığında ise $\neg\neg$'yi yalnızca başa koymak yetmez; değillemeleri
önermenin içine yaymak gerekir. Bunu yapan dönüşüm, Gödel ve Gentzen'in
birbirinden bağımsız olarak 1933'te buldukları \emph{çift değilleme
çevirisidir} \citep{godel1933, gentzen1936}.

\begin{tanim}{Gödel--Gentzen çevirisi (önermeler için)}{godel-gentzen}
Her $\varphi$ önermesine bir $\varphi^N$ önermesi şöyle karşılık getirilir:
\begin{center}
\renewcommand{\arraystretch}{1.3}
\begin{tabular}{r@{\;\;:=\;\;}l@{\hspace{3em}}r@{\;\;:=\;\;}l}
$P^N$                 & $\neg\neg P$ \quad ($P$ atomik) &
$(\varphi \wedge \psi)^N$ & $\varphi^N \wedge \psi^N$ \\
$\bot^N$              & $\bot$ &
$(\varphi \to \psi)^N$    & $\varphi^N \to \psi^N$ \\
$(\neg\varphi)^N$     & $\neg\varphi^N$ &
$(\varphi \vee \psi)^N$   & $\neg(\neg\varphi^N \wedge \neg\psi^N)$ \\
\end{tabular}
\end{center}
O zaman $\varphi$ klasik olarak ispatlanabilir ancak ve ancak $\varphi^N$
yapıcı olarak ispatlanabilir.
\end{tanim}

\begin{figure}[htbp]
\centering
\begin{tikzpicture}[
    bolge/.style={draw=#1, fill=#1!7, very thick, rounded corners=8pt},
  ]
  \draw[bolge=anarenk] (0,0) rectangle (11,4.4);
  \node[anchor=north west, font=\baslikfont\bfseries, text=anarenk] at (0.3,4.25) {Klasik mantık};
  \draw[bolge=ikincil] (0.6,0.5) rectangle (5.6,3.3);
  \node[anchor=north west, font=\baslikfont\bfseries\small, text=ikincil] at (0.8,3.15) {Yapıcı mantık};
  \node[draw=anarenk, fill=white, rounded corners=3pt, font=\small, inner sep=5pt] (phi) at (8.6,2.0) {$\varphi$};
  \node[draw=ikincil, fill=white, rounded corners=3pt, font=\small, inner sep=5pt] (phin) at (3.1,1.9) {$\varphi^N$};
  \draw[-{Stealth[length=3mm]}, very thick, vurgu] (phi) to[bend right=20] (phin);
  \node[font=\baslikfont\scriptsize\bfseries, text=vurgu, fill=anarenk!7, inner sep=2pt] at (7.6,3.25) {çift değilleme çevirisi};
  \node[font=\baslikfont\scriptsize, text=black!60, align=center, below=4pt of phi] {klasik olarak\\ispatlanabilir};
  \node[font=\baslikfont\scriptsize, text=black!60, align=center, below=4pt of phin] {yapıcı olarak\\ispatlanabilir};
\end{tikzpicture}
\caption{Çift değilleme çevirisi, klasik mantığı yapıcı mantığın içine gömer.}
\label{fig:cevirisi}
\end{figure}

\begin{onemlifikir}[Önemli Fikir: Klasik mantık tutarlıysa, yapıcı mantık da tutarlıdır — ve tersi]
Çift değilleme çevirisi $\bot$'u $\bot$'a götürür. Dolayısıyla klasik
mantıkta bir çelişki türetilebilseydi, yapıcı mantıkta da türetilebilirdi.
Klasik mantığı kullanmak, yapıcı mantığa göre \emph{yeni bir tutarsızlık
riski} taşımaz. Brouwer ile Hilbert arasındaki tartışmanın bugünkü
barışçıl sonucunun teknik nedeni budur.
\end{onemlifikir}

% ---------------------------------------------------------------------
\section{Karar verilebilir önermeler}\label{sec:decidable}

Bazı önermeler için üçüncü hâlin olanaksızlığı yapıcı olarak doğrudur:
\emph{doğru mu yanlış mı olduğunu hesaplayabildiğimiz} önermeler için.
``$3 < 5$'' önermesinin doğru olup olmadığını bir bilgisayar programı
bize söyleyebilir. Bu tür önermelere \emph{karar verilebilir} denir.

\Lean{}'de bir önermenin karar verilebilir olması, \lc{Decidable P}
türünde bir \emph{örneğin} (İng.\ \emph{instance}) bulunmasıdır. Bu örnek,
$P$'nin doğru mu yanlış mı olduğunu hesaplayan bir algoritmadır.

\leandosya[Karar verilebilirlik]{Bolum04}{decidable}

\begin{table}[htbp]
\centering
\renewcommand{\arraystretch}{1.45}
\rowcolors{2}{tablosatir}{white}
\begin{tabular}{>{\raggedright\arraybackslash}p{3.1cm}
                >{\raggedright\arraybackslash}p{4.3cm}
                >{\raggedright\arraybackslash}p{4.3cm}}
\rowcolor{tabloust}
\tablobaslik{} & \tablobaslik{\lc{Bool}} & \tablobaslik{\lc{Prop}} \\
Elemanları         & yalnızca \lc{true} ve \lc{false}  & sonsuz sayıda önerme \\
Ne anlatır?        & bir \emph{hesaplamanın} sonucunu  & bir \emph{iddiayı} \\
``Ya ya da'' ilkesi & her zaman geçerli (iki eleman var) & yalnızca klasik mantıkla ya da \lc{Decidable} ile \\
Köprü              & \multicolumn{2}{c}{\lc{decide : (P : Prop) → [Decidable P] → Bool}} \\
\end{tabular}
\caption{\lc{Bool} ile \lc{Prop} arasındaki fark.}
\label{tab:bool-prop}
\end{table}

\leandosya[Karar verilebilir önermeler için DNE]{Bolum04}{decidable_dne}

\noindent Bu teoremin \lc{cift\_degilleme\_eleme}'den farkı, köşeli
parantez içindeki \lc{[Decidable P]} varsayımıdır. Bu varsayım
sayesinde \lc{Classical.em} yerine \lc{Decidable.em} kullanılır. Doğal
sayıların eşitliği, sıralaması, sonlu kümeler üzerindeki önermeler gibi
pek çok somut matematiksel önerme karar verilebilirdir ve onlar için
klasik akıl yürütme yapıcı olarak da geçerlidir.

% ---------------------------------------------------------------------
\section{Pratikte: Hangi mantıkla çalışmalı?}\label{sec:pratikte}

\Lean{} kullanıcısı olarak bu tartışmadan çıkarmanız gereken pratik
sonuçlar şunlardır:

\begin{itemize}
  \item \textbf{\Mathlib{} klasiktir.} \Lean{}'in matematik kütüphanesi
        klasik mantığı serbestçe kullanır. Gerçek matematik yaparken
        \lc{Classical.em}'den kaçınmanız gerekmez.
  \item \textbf{Yapıcı ispat daha fazla şey verir.} Bir ispat klasik
        mantık kullanmıyorsa, \Lean{} onu bir program olarak
        \emph{çalıştırabilir}. Bir varlık ispatından, var olan nesneyi
        hesaplayan bir program elde edebilirsiniz. Klasik ispatlar ise
        \leankw{noncomputable} olarak işaretlenmek zorunda kalabilir.
  \item \textbf{Önce yapıcı deneyin.} Bu kitabın Kısım~I ve II'sindeki
        alıştırmaların çoğu yapıcıdır. Klasik mantığa ihtiyaç
        duyduğunuzda, bunun nedenini anlamaya çalışın: Çoğu zaman
        ispatın içinde bir ``hangisi?'' sorusu saklıdır.
  \item \textbf{\lc{\#print axioms} dürüst bir tanıktır.} Bir teoremin
        hangi ilkelere dayandığından emin değilseniz, sorun.
\end{itemize}

% ---------------------------------------------------------------------
\section{Taktik ve terim özeti}

\begin{center}
\renewcommand{\arraystretch}{1.45}
\rowcolors{2}{tablosatir}{white}
\begin{tabularx}{\linewidth}{>{\raggedright\arraybackslash}p{4.6cm} X}
\rowcolor{tabloust}
\tablobaslik{Araç} & \tablobaslik{Ne yapar?} \\
\lc{Classical.em P}                 & \lc{P ∨ ¬P} (klasik). \leantac{cases} ile birlikte klasik durum ayırma. \\
\lc{Classical.byContradiction}      & \lc{(¬P → False) → P}: çelişki ile ispat (klasik). \\
\lc{Decidable.em P}                 & \lc{[Decidable P]} varken \lc{P ∨ ¬P} (yapıcı). \\
\leantac{decide}                    & Karar verilebilir bir önermeyi hesaplayarak ispatlar. \\
\lc{decide P}                       & Karar verilebilir \lc{P}'yi \lc{Bool} değerine çevirir. \\
\lc{\textcolor{leancmd}{\bfseries\#print axioms} t} & \lc{t}'nin dayandığı aksiyomları listeler. \\
\end{tabularx}
\end{center}

% ---------------------------------------------------------------------
\begin{arastirmanotu}
\textbf{Peirce yasası ve programlama.} Timothy Griffin 1990'da, Peirce
yasasının Curry--Howard eşlemesi altında programlama dillerindeki
\emph{call/cc} (``mevcut devamı yakala'') kontrol operatörünün tipine
karşılık geldiğini gösterdi \citep{griffin1990}. Böylece klasik mantık
da bir hesaplama anlamı kazandı: Klasik ispatlar, ``geri dönüp fikir
değiştirebilen'' programlardır. \Cref{sec:em-yanlis-olamaz}'teki
şeytan hikâyesi, Philip Wadler'ın bu fikri anlatmak için kullandığı
benzetmeden uyarlanmıştır \citep{wadler2003}.

\textbf{Yapıcı matematik.} Errett Bishop'un 1967 tarihli kitabı, gerçel
analizin büyük bölümünün yapıcı olarak geliştirilebileceğini gösterdi
\citep{bishop1967}. Yapıcı matematiğin felsefi ve teknik temelleri için
\citet{troelstra1988} ve \citet{dummett2000} temel kaynaklardır.

\textbf{Lean'de yapıcılık.} \Lean{}'in \lc{Prop} evreni
\emph{kanıt-ilgisizdir} (\cref{bol:bagimli}): Bir önermenin iki
kanıtı her zaman eşit sayılır. Bu tasarım, klasik aksiyomların hesaplanan
programlara ``bulaşmasını'' engeller: \lc{Prop} içinde klasik akıl
yürütme, \lc{Type} içindeki hesaplamaları bozmaz. Bu dengenin kesin
ifadesi hâlâ etkin bir araştırma konusudur \citep{carneiro2019}.
\end{arastirmanotu}

\begin{bolumozeti}
\begin{itemize}
  \item BHK yorumunda bir kanıt bir \emph{inşadır}. $P \vee Q$'nun kanıtı
        hangi tarafın doğru olduğunu söylemelidir. Bu yüzden
        $P \vee \neg P$ yapıcı olarak genel bir yasa değildir.
  \item \Lean{}'de klasik mantık \lc{Classical.em} ve
        \lc{Classical.byContradiction} ile açıkça çağrılır.
  \item EM, DNE ve Peirce yasası yapıcı mantıkta birbirine denktir.
        Herhangi birini eklemek klasik mantığı verir.
  \item \lc{\#print axioms}, bir ispatın klasik olup olmadığını gösterir.
        \lc{Classical.em}, seçim aksiyomundan türetilir.
  \item Kripke modelleri, bir önermenin yapıcı olarak ispatlanamadığını
        göstermenin aracıdır.
  \item Glivenko teoremi ve çift değilleme çevirisi, klasik mantığı
        yapıcı mantığın içine gömer.
  \item Karar verilebilir önermeler için üçüncü hâl yapıcı olarak
        geçerlidir.
\end{itemize}
\end{bolumozeti}

\begin{alistirmalar}
\alistirma{1} (Klasik) \lc{(P → Q) → ¬P ∨ Q}. Bu, implikasyonun
``maddi'' tanımıdır. (\lc{a4\_1})

\alistirma{1} (Yapıcı) Bir önceki alıştırmanın tersi:
\lc{¬P ∨ Q → (P → Q)}. Neden bu yön yapıcıdır? (\lc{a4\_2})

\alistirma{2} (Klasik) \lc{¬(P → Q) → P ∧ ¬Q}. İki bileşenden hangisi
klasik mantık gerektiriyor? \lc{\#print axioms} ile kendi
ispatınızı denetleyin. (\lc{a4\_3})

\alistirma{2} (Klasik) Kontrapozisyonun tersini terim modunda
yazın: \lc{(¬Q → ¬P) → (P → Q)}. (\lc{a4\_4})

\alistirma{2} (Klasik) \emph{Dummett ilkesi}: \lc{(P → Q) ∨ (Q → P)}.
İki rastgele önermeden biri her zaman ötekini gerektirir! Bu ilkenin
sezgisel olarak neden tuhaf olduğunu tartışın. (\lc{a4\_5})

\alistirma{2} (Yapıcı) \lc{[Decidable P]} varsayımıyla De Morgan'ın zor
yönünü ispatlayın. \lc{Decidable Q} varsayımı da gerekli midir?
(\lc{a4\_6})

\alistirma{3} (Yapıcı) Glivenko teoremini Dummett ilkesi için
doğrulayın: \lc{¬¬((P → Q) ∨ (Q → P))}. (\lc{a4\_7})

\alistirma{3}\label{alis:dne-em} (Yapıcı) Her önerme için çift
değilleme elemeyi biliyorsak üçüncü hâli de biliriz:
\lc{(∀ X : Prop, ¬¬X → X) → P ∨ ¬P}. Hangi \lc{X} için kullanmanız
gerekti? (\lc{a4\_8})

\alistirma{3} (Kâğıt-kalem) \Cref{fig:kripke-em}'deki modelde Peirce
yasasının ($Q$ yerine $\bot$ alarak) $w_0$ durumunda bilinmediğini
gösterin.

\alistirma{3} (Kâğıt-kalem) Üç durumlu bir Kripke modeli kurarak
Dummett ilkesinin yapıcı olarak ispatlanamadığını gösterin. İpucu:
$w_0$'dan iki farklı yöne dallanan $w_1$ ve $w_2$ durumları düşünün.

\alistirma{3} (Tartışma) Teorem~\ref{trm:irrasyonel-us}'in yapıcı bir ispatı
var mıdır? İpucu: $a = \sqrt{2}$, $b = 2\log_2 3$ seçimini inceleyin.
Bu seçimle ispatın hangi adımları yapıcıdır?
\end{alistirmalar}
   % Yapıcı ve klasik mantık
% =====================================================================
%  Bölüm 5 — Doğal Çıkarım Ağaçları ve Lean Taktikleri
% =====================================================================
\chapter{Doğal Çıkarım Ağaçları ve Lean Taktikleri}\label{bol:dogal-cikarim}

\begin{bolumhedefleri}
\begin{itemize}
  \item Önermeler mantığının doğal çıkarım kurallarının tamamını tek bir
        tabloda görmeyi,
  \item \Lean{}'in hedef ekranının bir \emph{ardışık} ($\Gamma \vdash \varphi$)
        olduğunu ve her taktiğin bir ardışığı başkalarına dönüştürdüğünü,
  \item Geriye doğru (\leantac{apply}, \leantac{suffices}) ve ileriye doğru
        (\leantac{have}) akıl yürütmeyi ve ikisini birlikte kullanmayı,
  \item Bir ispata nereden başlanacağını söyleyen bir strateji haritasını,
  \item İspat normalleştirmeyi: dolambaçlı ispatların nasıl
        sadeleştirildiğini ve bunun neden ``hesaplama'' olduğunu,
  \item Taktik birleştiricilerini: \lc{<;>}, \leantac{all\_goals},
        \leantac{first}, \leantac{refine}.
\end{itemize}
\end{bolumhedefleri}

Bu bölüm Kısım~I'i kapatıyor. Yeni bir bağlaç tanıtmayacağız; bunun
yerine önceki dört bölümde öğrendiklerimizi bir araya getirip
\emph{düzenleyeceğiz}. Amaç, önünüze yeni bir önerme geldiğinde
``nereden başlamalıyım?'' sorusunu sistemli biçimde cevaplayabilmenizdir.

% ---------------------------------------------------------------------
\section{Kural kataloğu}\label{sec:kural-katalogu}

\Cref{fig:kural-katalogu}, Bölüm~2--4'te tanıştığımız bütün çıkarım
kurallarını bir arada gösteriyor. Tabloda iki şeye dikkat edin.
Birincisi, her bağlacın bir \textcolor{ikincil}{\bfseries giriş} ve bir
\textcolor{vurgu}{\bfseries eleme} sütunu vardır. İkincisi, yalnızca son
satırdaki kural (\emph{klasik saçmaya indirgeme}) yapıcı değildir.

\begin{figure}[p]
\centering
\small
\renewcommand{\arraystretch}{1}
\setlength{\tabcolsep}{4pt}
\begin{tabular}{>{\centering\arraybackslash}m{1.2cm}
                >{\centering\arraybackslash}m{5.3cm}
                >{\centering\arraybackslash}m{6.6cm}}
\rowcolor{tabloust}
\tablobaslik{Bağlaç} & \tablobaslik{Giriş (kanıtlamak)} & \tablobaslik{Eleme (kullanmak)} \\[2pt]
% →
$\to$ &
\vspace{2pt}
\begin{prooftree}
  \hypo{[P]^{n}} \ellipsis{}{Q} \infer1[$\to$G$^n$]{P \to Q}
\end{prooftree}\par\smallskip \leantac{intro} &
\vspace{2pt}
\begin{prooftree}
  \hypo{P \to Q} \hypo{P} \infer2[$\to$E]{Q}
\end{prooftree}\par\smallskip \leantac{apply}, \lc{h hp} \\[4pt]
\hline
% ∧
$\wedge$ &
\vspace{2pt}
\begin{prooftree}
  \hypo{P} \hypo{Q} \infer2[$\wedge$G]{P \wedge Q}
\end{prooftree}\par\smallskip \leantac{constructor}, \lc{⟨hp, hq⟩} &
\vspace{2pt}
\begin{prooftree}
  \hypo{P \wedge Q} \infer1[$\wedge$E$_1$]{P}
\end{prooftree}\quad
\begin{prooftree}
  \hypo{P \wedge Q} \infer1[$\wedge$E$_2$]{Q}
\end{prooftree}\par\smallskip \leantac{obtain}, \lc{h.1}, \lc{h.2} \\[4pt]
\hline
% ∨
$\vee$ &
\vspace{2pt}
\begin{prooftree}
  \hypo{P} \infer1[$\vee$G$_1$]{P \vee Q}
\end{prooftree}\quad
\begin{prooftree}
  \hypo{Q} \infer1[$\vee$G$_2$]{P \vee Q}
\end{prooftree}\par\smallskip \lc{Or.inl}, \lc{Or.inr} &
\vspace{2pt}
\begin{prooftree}
  \hypo{P \vee Q}
  \hypo{[P]^{m}} \ellipsis{}{R}
  \hypo{[Q]^{n}} \ellipsis{}{R}
  \infer3[$\vee$E$^{m,n}$]{R}
\end{prooftree}\par\smallskip \leantac{cases}, \leantac{rcases} \\[4pt]
\hline
% ¬
$\neg$ &
\vspace{2pt}
\begin{prooftree}
  \hypo{[P]^{n}} \ellipsis{}{\bot} \infer1[$\neg$G$^n$]{\neg P}
\end{prooftree}\par\smallskip \leantac{intro} &
\vspace{2pt}
\begin{prooftree}
  \hypo{\neg P} \hypo{P} \infer2[$\neg$E]{\bot}
\end{prooftree}\par\smallskip \lc{hnp hp} \\[4pt]
\hline
% ↔
$\leftrightarrow$ &
\vspace{2pt}
\begin{prooftree}
  \hypo{[P]^{m}} \ellipsis{}{Q}
  \hypo{[Q]^{n}} \ellipsis{}{P}
  \infer2[$\leftrightarrow$G$^{m,n}$]{P \leftrightarrow Q}
\end{prooftree}\par\smallskip \leantac{constructor} &
\vspace{2pt}
\begin{prooftree}
  \hypo{P \leftrightarrow Q} \hypo{P} \infer2[$\leftrightarrow$E$_1$]{Q}
\end{prooftree}\quad
\begin{prooftree}
  \hypo{P \leftrightarrow Q} \hypo{Q} \infer2[$\leftrightarrow$E$_2$]{P}
\end{prooftree}\par\smallskip \lc{h.mp}, \lc{h.mpr}, \leantac{rw} \\[4pt]
\hline
% ⊤ ⊥
$\top$, $\bot$ &
\vspace{2pt}
\begin{prooftree}
  \infer0[$\top$G]{\top}
\end{prooftree}\par\smallskip \leantac{trivial} &
\vspace{2pt}
\begin{prooftree}
  \hypo{\bot} \infer1[$\bot$E]{P}
\end{prooftree}\par\smallskip \leantac{exfalso}, \lc{False.elim} \\[4pt]
\hline
\rowcolor{teorembg}
\textcolor{teoremframe}{\bfseries klasik} &
\multicolumn{2}{c}{%
\vspace{2pt}
\begin{prooftree}
  \hypo{[\neg P]^{n}} \ellipsis{}{\bot} \infer1[RAA$^n$]{P}
\end{prooftree}\quad\raisebox{1.2em}{\lc{Classical.byContradiction}}} \\[4pt]
\end{tabular}

\vspace{6pt}
\begin{tikzpicture}
  \node[draw=black!20, rounded corners=3pt, fill=white, font=\footnotesize,
        text width=12.6cm, align=left, inner sep=6pt] {%
    \textbf{Varsayım kuralı.} Bir varsayım, kendi başına bir türetimdir:
    \quad$\begin{prooftree}\infer0{P}\end{prooftree}$ \;(varsayım $P$).
    \Lean{}'de \leantac{assumption} ya da \leantac{exact hp}.
    \quad \textbf{Kapatma.} $[P]^n$ biçimindeki varsayımlar, aynı
    numarayı taşıyan kuralda kapatılır.};
\end{tikzpicture}
\caption{Önermeler mantığı için doğal çıkarımın kural kataloğu ve her
kuralın \Lean{} karşılığı. Pembe satır dışındaki bütün kurallar
yapıcıdır; pembe satır klasik mantığa geçiştir.}
\label{fig:kural-katalogu}
\end{figure}

\begin{tanim}{Doğal çıkarım türetimi}{turetim}
Bir \textbf{türetim}, yaprakları varsayımlar, iç düğümleri
\cref{fig:kural-katalogu}'deki kuralların uygulamaları olan bir
ağaçtır. Kökteki önerme türetimin \textbf{sonucudur}. Bir türetimin
\textbf{açık varsayımları}, hiçbir kuralda kapatılmamış yapraklardır.
$\Gamma$ açık varsayımlar kümesi ve $\varphi$ sonuç ise,
\[
  \Gamma \vdash \varphi
\]
yazarız ve ``$\varphi$, $\Gamma$'dan türetilebilir'' diye okuruz.
Açık varsayımı olmayan bir türetimin sonucuna \textbf{teorem} denir:
$\vdash \varphi$.
\end{tanim}

% ---------------------------------------------------------------------
\section{Lean'in hedef ekranı bir ardışıktır}\label{sec:ardisik-gorunum}

Bölüm~2'den beri her \leantac{intro} adımında varsayımların Infoview
panelinde \lc{⊢} işaretinin \emph{üstüne} taşındığını gördük. Bu
gösterim bir rastlantı değildir: \Lean{}'in hedef ekranı, tam olarak
$\Gamma \vdash \varphi$ biçiminde bir \emph{ardışıktır} (İng.\
\emph{sequent}) (\cref{fig:ardisik}).

\begin{figure}[htbp]
\centering
\begin{tikzpicture}[
    kutu/.style={draw=hedefframe, fill=hedefbg, thick, rounded corners=3pt,
                 font=\leanfont\small, align=left, inner sep=7pt},
    ok/.style={-{Stealth[length=2.6mm]}, thick, #1},
  ]
  \node[kutu] (lean) at (0,0) {%
    P Q : \textcolor{leantype}{\bfseries Prop}\\
    \textcolor{leanhyp}{\bfseries hp} : P\\
    \textcolor{leanhyp}{\bfseries hq} : Q\\[3pt]
    \textcolor{leansym}{⊢} P \textcolor{leansym}{∧} Q};
  \node[font=\Large] (es) at (3.2,0) {$\Longleftrightarrow$};
  \node[font=\Large] (sq) at (6.6,0) {$\underbrace{P,\; Q}_{\Gamma} \;\vdash\; \underbrace{P \wedge Q}_{\varphi}$};
  \draw[decorate, decoration={brace, amplitude=5pt}, thick, leanhyp]
    ([xshift=4pt]lean.north east) -- ([xshift=4pt,yshift=-1.55cm]lean.north east)
    node[midway, right=7pt, font=\baslikfont\scriptsize, text=leanhyp]{bağlam $\Gamma$};
  \node[below=0.35cm of lean, font=\baslikfont\scriptsize, text=black!60] {Lean Infoview};
  \node[below=0.55cm of sq, font=\baslikfont\scriptsize, text=black!60] {ardışık (sequent)};
\end{tikzpicture}
\caption{\Lean{}'in hedef durumu ile mantıktaki ardışık gösterimi aynı
şeydir. Hipotez adları (\lc{hp}, \lc{hq}), varsayımların numaralarının
yerini tutar.}
\label{fig:ardisik}
\end{figure}

Bu bakış açısıyla doğal çıkarım kurallarını, varsayımları ağacın
yapraklarında taşımak yerine her düğümde açıkça yazarak ifade
edebiliriz. Örneğin varsayım kuralı ile implikasyonun kuralları şöyle
olur:

\begin{center}
\begin{tabular}{c@{\hspace{2.6em}}c@{\hspace{2.6em}}c}
\begin{prooftree}
  \infer0[Var]{\Gamma, P \vdash P}
\end{prooftree}
&
\begin{prooftree}
  \hypo{\Gamma, P \vdash Q}
  \infer1[$\to$G]{\Gamma \vdash P \to Q}
\end{prooftree}
&
\begin{prooftree}
  \hypo{\Gamma \vdash P \to Q}
  \hypo{\Gamma \vdash P}
  \infer2[$\to$E]{\Gamma \vdash Q}
\end{prooftree}
\\[10pt]
\small \leantac{assumption} & \small \leantac{intro} & \small \leantac{apply}
\end{tabular}
\end{center}

\noindent Bu biçimde okunduğunda her kural bir \emph{taktiğin tarifidir}:

\begin{onemlifikir}[Önemli Fikir: Taktik = kuralı aşağıdan yukarı uygulamak]
Bir çıkarım kuralı, \emph{yukarıdaki} ardışıklardan \emph{aşağıdakine}
geçişi gösterir. Bir taktik ise aynı kuralı \emph{tersten} uygular:
Aşağıdaki ardışığı (mevcut hedefi) alır ve onu yukarıdaki ardışıklara
(yeni hedeflere) dönüştürür. \leantac{intro}, $\Gamma \vdash P \to Q$
hedefini $\Gamma, P \vdash Q$ hedefine çevirir. \leantac{constructor},
$\Gamma \vdash P \wedge Q$ hedefini iki hedefe böler. Kâğıt üzerinde ağacı
yapraklardan köke doğru \emph{okuruz}; \Lean{}'de ise kökten yapraklara
doğru \emph{inşa ederiz}.
\end{onemlifikir}

\leandosya[Varsayım kuralı]{Bolum05}{varsayim}

\noindent Ardışık biçimindeki bu doğal çıkarım sistemi, \cref{bol:ardisik}'te
göreceğimiz Gentzen'in \emph{ardışık hesabından} farklıdır: Orada sağ
tarafta da birden çok önerme bulunabilir ve eleme kuralları yerine
``sol giriş'' kuralları kullanılır.

% ---------------------------------------------------------------------
\section{Kural--taktik eşleme tablosu}\label{sec:kural-taktik}

\Cref{tab:kural-taktik}, her kuralın taktik modunda (geriye doğru) ve
terim modunda (ileriye doğru) nasıl kullanıldığını özetliyor. Bu tablo,
kitabın geri kalanında sık sık başvuracağınız bir kaynak olacak.

\begin{table}[htbp]
\centering
\small
\renewcommand{\arraystretch}{1.45}
\rowcolors{2}{tablosatir}{white}
\begin{tabularx}{\linewidth}{>{\raggedright\arraybackslash}p{1.6cm}
                             >{\raggedright\arraybackslash}p{3.1cm}
                             >{\raggedright\arraybackslash}X
                             >{\centering\arraybackslash}p{1.0cm}}
\rowcolor{tabloust}
\tablobaslik{Kural} & \tablobaslik{Taktik (hedef üzerinde)} & \tablobaslik{Terim / hipotez üzerinde} & \tablobaslik{Bölüm} \\
Var                 & \leantac{assumption}, \leantac{exact h} & \lc{h} & \ref{bol:giris} \\
$\to$G              & \leantac{intro h}             & \lc{fun h => t} & \ref{bol:onermeler1} \\
$\to$E              & \leantac{apply f}             & \lc{f a} & \ref{bol:onermeler1} \\
$\wedge$G           & \leantac{constructor}         & \lc{⟨a, b⟩} & \ref{bol:onermeler1} \\
$\wedge$E           & ---                           & \lc{h.1}, \lc{h.2}, \leantac{obtain ⟨a, b⟩ := h} & \ref{bol:onermeler1} \\
$\vee$G             & \leantac{apply Or.inl}        & \lc{Or.inl a}, \lc{Or.inr b} & \ref{bol:onermeler1} \\
$\vee$E             & ---                           & \leantac{cases h}, \leantac{rcases h with a | b}, \lc{Or.elim} & \ref{bol:onermeler1} \\
$\neg$G             & \leantac{intro h}             & \lc{fun h => ...} & \ref{bol:degilleme} \\
$\neg$E             & \leantac{apply hn}            & \lc{hn h}, \lc{absurd h hn} & \ref{bol:degilleme} \\
$\leftrightarrow$G  & \leantac{constructor}         & \lc{⟨f, g⟩} & \ref{bol:degilleme} \\
$\leftrightarrow$E  & \leantac{rw [h]}              & \lc{h.mp a}, \lc{h.mpr b} & \ref{bol:degilleme} \\
$\top$G             & \leantac{trivial}             & \lc{True.intro} & \ref{bol:degilleme} \\
$\bot$E             & \leantac{exfalso}             & \lc{False.elim h}, \leantac{cases h} & \ref{bol:degilleme} \\
RAA                 & \leantac{apply Classical.}\newline\leantac{\quad byContradiction} & \lc{Classical.byContradiction f} & \ref{bol:yapici-klasik} \\
\end{tabularx}
\caption{Doğal çıkarım kuralları ile \Lean{} taktiklerinin ve terimlerinin
eşlemesi. Eleme kuralları çoğunlukla \emph{hipotezler} üzerinde, giriş
kuralları ise \emph{hedef} üzerinde çalışır.}
\label{tab:kural-taktik}
\end{table}

% ---------------------------------------------------------------------
\section{Geriye ve ileriye akıl yürütme}\label{sec:geri-ileri}

Bir ispatı iki yönden kurabiliriz (\cref{fig:geri-ileri}):

\begin{itemize}
  \item \textbf{Geriye doğru}: Hedeften başlar, ``bunu göstermek için ne
        yeter?'' diye sorarız. \leantac{apply}, \leantac{intro},
        \leantac{constructor}, \leantac{suffices} bu yönde çalışır.
  \item \textbf{İleriye doğru}: Hipotezlerden başlar, ``bunlardan ne
        çıkar?'' diye sorarız. \leantac{have} ve terim modundaki
        uygulamalar bu yönde çalışır.
\end{itemize}

\begin{figure}[htbp]
\centering
\begin{tikzpicture}[
    dugum/.style={draw=#1, fill=#1!8, thick, rounded corners=3pt,
                  minimum width=1.5cm, minimum height=0.8cm, font=\leanfont\small},
    ok/.style={-{Stealth[length=2.6mm]}, very thick, #1},
  ]
  % hipotezler
  \node[dugum=leanhyp] (p) at (0,0) {p : P};
  \node[dugum=leanframe] (q) at (3.6,0) {Q};
  \node[dugum=leanframe] (r) at (7.2,0) {R};
  \node[dugum=bittiframe] (g) at (10.6,0) {⊢ R};
  \draw[ok=leansym] (p) -- node[above, font=\leanfont\scriptsize\bfseries, text=leanhyp]{q := h1 p} (q);
  \draw[ok=leansym] (q) -- node[above, font=\leanfont\scriptsize\bfseries, text=leanhyp]{r := h2 q} (r);
  \draw[ok=leansym, dashed] (r) -- node[above, font=\leanfont\scriptsize\bfseries, text=leanhyp]{exact r} (g);
  \node[font=\baslikfont\small\bfseries, text=leansym, anchor=west] at (-0.8,0.95) {İleriye doğru (have)};
  % geri
  \node[dugum=hedefframe] (g1) at (10.6,-2.3) {⊢ R};
  \node[dugum=hedefframe] (g2) at (7.2,-2.3) {⊢ Q};
  \node[dugum=hedefframe] (g3) at (3.6,-2.3) {⊢ P};
  \node[dugum=bittiframe] (g4) at (0,-2.3) {p : P};
  \draw[ok=leantac] (g1) -- node[below, font=\leanfont\scriptsize\bfseries, text=leantac]{apply h2} (g2);
  \draw[ok=leantac] (g2) -- node[below, font=\leanfont\scriptsize\bfseries, text=leantac]{apply h1} (g3);
  \draw[ok=leantac, dashed] (g3) -- node[below, font=\leanfont\scriptsize\bfseries, text=leantac]{exact p} (g4);
  \node[font=\baslikfont\small\bfseries, text=leantac, anchor=east] at (11.4,-3.25) {Geriye doğru (apply)};
\end{tikzpicture}
\caption{Aynı ispat iki yönden: İleriye doğru ispat hipotezlerden
hedefe, geriye doğru ispat hedeften hipotezlere yürür. İkisi ortada
buluşur.}
\label{fig:geri-ileri}
\end{figure}

Hipotetik kıyasın geriye doğru ispatını Bölüm~2'de görmüştük:

\leandosya[Geriye doğru]{Bolum05}{kiyas_geri}

\noindent Aynı ispatı ileriye doğru da yazabiliriz. \leantac{have}
taktiği, bağlama yeni bir ara sonuç ekler:

\leandosya[İleriye doğru]{Bolum05}{kiyas_ileri}

\begin{hedefgecisi}
\begin{hedefdurumu}[\texttt{have q} öncesi]
h1 : P → Q
h2 : Q → R
p : P
⊢ R
\end{hedefdurumu}
\sonra
\begin{hedefdurumu}[\texttt{have q} sonrası]
h1 : P → Q
h2 : Q → R
p : P
q : Q
⊢ R
\end{hedefdurumu}
\end{hedefgecisi}

\noindent \leantac{have} hedefi değiştirmez; yalnızca bağlamı
büyütür. Matematik kitaplarındaki ``\ldots olduğundan \ldots elde
ederiz'' cümlelerinin karşılığı budur.

\paragraph{suffices: ``Şunu göstermek yeter''}
Geriye doğru akıl yürütürken ara hedefe bir ad vermek isteyebiliriz.
\leantac{suffices} tam bunu yapar: ``$Q$ elimde olsaydı işi şöyle
bitirirdim; o hâlde $Q$'yu göstermek yeter.''

\leandosya[Yeter koşulu]{Bolum05}{kiyas_suffices}

\paragraph{show: hedefi açıkça yazmak}
\leantac{show} taktiği hiçbir şeyi değiştirmez; hedefin ne olduğunu
okura (ve \Lean{}'e) açıkça söyler. Hedef yazdığınız şeyle uyuşmazsa
\Lean{} hata verir. Uzun ispatlarda \leantac{show}, ispatı okunabilir
kılmanın en ucuz yoludur.

\leandosya[Hedefi açıkça yazmak]{Bolum05}{kiyas_show}

\begin{dikkat}[Dikkat: Hangi yön ne zaman?]
Geriye doğru ispatlar kısadır ama okunması zordur: Okur, ara hedefleri
görmek için \Lean{}'i çalıştırmak zorunda kalır. İleriye doğru ispatlar
uzundur ama kâğıt üzerindeki matematiğe benzer. İyi bir ispat genellikle
ikisini karıştırır: Büyük adımları \leantac{have} ile adlandırır, küçük
boşlukları \leantac{apply} ve \leantac{exact} ile kapatır.
\end{dikkat}

% ---------------------------------------------------------------------
\section{Bir ispat stratejisi}\label{sec:strateji}

Önermeler mantığında bir ispata nereden başlayacağınızı çoğu zaman
\emph{hedefin} ve \emph{hipotezlerin} biçimi söyler. \Cref{fig:strateji}
bu bilgiyi bir harita olarak topluyor.

\begin{figure}[p]
\centering
\begin{tikzpicture}[
    bas/.style={draw=koyu, fill=koyu, text=white, rounded corners=4pt,
                font=\baslikfont\bfseries\small, minimum width=5.6cm,
                minimum height=0.8cm},
    sekil/.style={draw=#1, fill=#1!8, thick, rounded corners=3pt,
                  font=\leanfont\small, minimum width=1.7cm, minimum height=0.75cm},
    eylem/.style={draw=black!20, fill=white, rounded corners=3pt,
                  font=\small, text width=4.3cm, minimum height=0.75cm,
                  align=left, inner sep=4pt},
    ok/.style={-{Stealth[length=2.2mm]}, thick, black!45},
  ]
  % ----- Sol sütun: hedef -----
  \node[bas] at (3.1,0) {1. Hedefe bak: $\vdash$ ?};
  \foreach \i/\s/\e in {
      1/{⊢ A → B}/{\leantac{intro h}},
      2/{⊢ ¬A}/{\leantac{intro h}; hedef \lc{False} olur},
      3/{⊢ A ∧ B}/{\leantac{constructor}},
      4/{⊢ A ↔ B}/{\leantac{constructor}: iki yön},
      5/{⊢ A ∨ B}/{Taraf belliyse \lc{Or.inl}/\lc{Or.inr}; değilse önce 2. adım},
      6/{⊢ True}/{\leantac{trivial}},
      7/{⊢ False}/{Çelişki ara: \lc{hn h}, \leantac{contradiction}},
      8/{⊢ atom}/{2. adıma geç}}
  {
    \node[sekil=ikincil] (h\i) at (0.85,{-1.3*\i}) {\s};
    \node[eylem, anchor=west] (e\i) at (1.95,{-1.3*\i}) {\e};
    \draw[ok] (h\i) -- (e\i);
  }
  % ----- Sağ sütun: hipotezler -----
  \node[bas] at (10.1,0) {2. Hipotezlere bak: h : ?};
  \foreach \i/\s/\e in {
      1/{h : A ∧ B}/{\leantac{obtain ⟨a, b⟩ := h}},
      2/{h : A ∨ B}/{\leantac{rcases h with a | b}: iki durum},
      3/{h : A → B}/{Hedef \lc{B}: \leantac{apply h}; \lc{a : A} varsa: \leantac{have b := h a}},
      4/{h : ¬A}/{Hedef \lc{False}: \leantac{apply h}, sonra \lc{A}'yı göster},
      5/{h : A ↔ B}/{\leantac{rw [h]} ya da \lc{h.mp}, \lc{h.mpr}},
      6/{h : False}/{\leantac{exact h.elim}; bitti},
      7/{h : hedef}/{\leantac{exact h}, \leantac{assumption}}}
  {
    \node[sekil=vurgu] (k\i) at (7.85,{-1.3*\i}) {\s};
    \node[eylem, anchor=west] (f\i) at (8.95,{-1.3*\i}) {\e};
    \draw[ok] (k\i) -- (f\i);
  }
  % ----- Alt: takılınca -----
  \node[draw=teoremframe, fill=teorembg, thick, rounded corners=4pt,
        text width=13.6cm, align=left, font=\small, inner sep=6pt]
     at (6.6,-12.45) {%
     {\baslikfont\bfseries\color{teoremframe} 3. Takıldıysan:}
     Kâğıt üzerinde bir doğruluk tablosu ya da küçük bir Kripke modeli
     dene. Önerme klasik olarak doğru ama yapıcı olarak ispatlanamıyorsa
     (\cref{bol:yapici-klasik}), içinde bir ``hangisi?'' sorusu
     saklıdır: \lc{Classical.em} ile durum ayır. Önerme klasik olarak da
     yanlışsa, bir karşı örnek bul ve ispatlamaya çalışmayı bırak!};
\end{tikzpicture}
\caption{Önermeler mantığında ispat stratejisi. Önce hedefin biçimine,
sonra hipotezlerin biçimine bakılır. Her satır bir giriş (sol) ya da
eleme (sağ) kuralına karşılık gelir.}
\label{fig:strateji}
\end{figure}

Bu haritanın arkasında güçlü bir teorik sonuç yatar: Yapıcı önermeler
mantığında, hedefin biçimine göre giriş kurallarını, hipotezlerin
biçimine göre eleme kurallarını uygulayan ve aynı hedefe dönmemeye
dikkat eden sistemli bir arama, ispat varsa onu \emph{her zaman} bulur. Bu sonucun kesin hâli, \cref{bol:ardisik}'te
göreceğimiz kesme eleme teoreminin bir sonucudur. Yapıcı önermeler
mantığının \emph{karar verilebilir} olmasının nedeni de budur:
Herhangi bir önermenin ispatlanabilir olup olmadığını sonlu bir aramayla
anlayabiliriz.

% ---------------------------------------------------------------------
\section{Kapsamlı bir örnek}\label{sec:kapsamli-ornek}

Stratejiyi bir örnek üzerinde adım adım uygulayalım. İspatlamak
istediğimiz önerme şudur:
\[
  (P \to R) \wedge (Q \to R) \;\leftrightarrow\; (P \vee Q \to R).
\]
``$P$ de $R$'yi gerektiriyorsa $Q$ da gerektiriyorsa, $P$ veya $Q$
$R$'yi gerektirir; tersi de doğrudur.''

\begin{enumerate}[adimlar]
  \item \textbf{Hedef \lc{↔}}: \leantac{constructor} ile iki yöne ayır.
  \item \textbf{(→) yönü, hedef \lc{→}}: \leantac{intro} ile
        \lc{(P → R) ∧ (Q → R)}'yi al. Bu bir \lc{∧} hipotezi:
        \leantac{obtain} ile parçala. Hedef yine \lc{→}: \lc{P ∨ Q}'yu
        al. Bu bir \lc{∨} hipotezi: durum ayır. \leantac{rintro}, bu dört
        adımı tek satırda yapar.
  \item \textbf{Her durumda hedef \lc{R}} (atom): Hipotezlere bak.
        \lc{hpr : P → R} ve \lc{hp : P} var: \lc{hpr hp}.
  \item \textbf{(←) yönü}: \lc{h : P ∨ Q → R}'yi al. Hedef \lc{∧}: iki
        bileşeni ayrı ayrı kur. Her bileşen bir implikasyon: \lc{fun hp => \ldots}.
\end{enumerate}

\leandosya[Durum ayırmanın denkliği]{Bolum05}{durum_iff}

\noindent (→) yönünün doğal çıkarım ağacı \cref{fig:durum-agaci}'de
görülüyor. Ağaçtaki her kural, \Lean{} ispatındaki bir parçaya karşılık
geliyor.

\begin{figure}[htbp]
\centering
\small
\begin{prooftree}
  \hypo{[P \vee Q]^{4}}
  %
  \hypo{[(P \to R) \wedge (Q \to R)]^{3}}
  \infer1[$\wedge$E$_1$]{P \to R}
  \hypo{[P]^{1}}
  \infer2[$\to$E]{R}
  %
  \hypo{[(P \to R) \wedge (Q \to R)]^{3}}
  \infer1[$\wedge$E$_2$]{Q \to R}
  \hypo{[Q]^{2}}
  \infer2[$\to$E]{R}
  %
  \infer3[$\vee$E$^{1,2}$]{R}
  \infer1[$\to$G$^{4}$]{P \vee Q \to R}
  \infer1[$\to$G$^{3}$]{(P \to R) \wedge (Q \to R) \to (P \vee Q \to R)}
\end{prooftree}

\vspace{8pt}
\begin{tikzpicture}[font=\leanfont\scriptsize]
  \node[draw=black!20, rounded corners=3pt, fill=white, align=left, inner sep=6pt]{%
    \textcolor{leantac}{\bfseries rintro} ⟨hpr, hqr⟩ (hp | hq)
    \quad$\rightsquigarrow$\quad {\rmfamily\small $\to$G$^3$, $\wedge$E, $\to$G$^4$, $\vee$E$^{1,2}$}\\[2pt]
    \textcolor{leantac}{\bfseries exact} hpr hp \;/\; \textcolor{leantac}{\bfseries exact} hqr hq
    \quad$\rightsquigarrow$\quad {\rmfamily\small iki $\to$E adımı}};
\end{tikzpicture}
\caption{$(P \to R) \wedge (Q \to R) \to (P \vee Q \to R)$'nin doğal
çıkarım ağacı ve \Lean{} ispatıyla karşılaştırması.}
\label{fig:durum-agaci}
\end{figure}

\noindent Terim modunda aynı ispat iki satıra sığar. (→) yönünde
\lc{Or.elim h hpr hqr} ifadesi, $\vee$E kuralını \emph{doğrudan}
uygular: İki durumun her biri için hazır bir fonksiyonumuz zaten var.

\leandosya[Durum ayırmanın denkliği, terim modu]{Bolum05}{durum_iff_terim}

% ---------------------------------------------------------------------
\section{İspat normalleştirme}\label{sec:normallestirme}

Bazı ispatlar gereksiz dolambaçlar içerir. En basit örnek, bir ``ve''
kanıtını kurup hemen ardından sökmektir:

\leandosya[Dolambaçlı ve dolaysız ispat]{Bolum05}{dolambacli}

\noindent \lc{dolambacli} ispatı, \lc{⟨hp, hq⟩} çiftini kurar ($\wedge$G)
ve hemen sol bileşenini alır ($\wedge$E$_1$). Sonuçta elimize başta
zaten sahip olduğumuz \lc{hp} geçer. Ağaç olarak (\cref{fig:normallestirme}):

\begin{figure}[htbp]
\centering
\begin{tikzpicture}
  \node (sol) at (0,0) {%
    \begin{prooftree}
      \hypo{\mathcal{D}_1}\ellipsis{}{P}
      \hypo{\mathcal{D}_2}\ellipsis{}{Q}
      \infer2[\textcolor{ikincil}{$\wedge$G}]{P \wedge Q}
      \infer1[\textcolor{vurgu}{$\wedge$E$_1$}]{P}
    \end{prooftree}};
  \node (sag) at (5.2,0) {%
    \begin{prooftree}
      \hypo{\mathcal{D}_1}\ellipsis{}{P}
    \end{prooftree}};
  \draw[-{Stealth[length=3mm]}, very thick, anarenk] (1.9,0) -- node[above, font=\baslikfont\scriptsize\bfseries, text=anarenk]{indirgeme} (4.4,0);
  \node[font=\baslikfont\scriptsize, text=black!60] at (0,-1.6) {dolambaç: giriş, hemen ardından eleme};
  \node[font=\baslikfont\scriptsize, text=black!60] at (5.2,-1.6) {normal};

  \node (sol2) at (0,-4.1) {%
    \begin{prooftree}
      \hypo{[P]^{1}}\ellipsis{$\mathcal{D}$}{Q}
      \infer1[\textcolor{ikincil}{$\to$G$^1$}]{P \to Q}
      \hypo{\mathcal{E}}\ellipsis{}{P}
      \infer2[\textcolor{vurgu}{$\to$E}]{Q}
    \end{prooftree}};
  \node (sag2) at (5.2,-4.1) {%
    \begin{prooftree}
      \hypo{\mathcal{E}}\ellipsis{}{P}
      \ellipsis{$\mathcal{D}$}{Q}
    \end{prooftree}};
  \draw[-{Stealth[length=3mm]}, very thick, anarenk] (1.9,-4.1) -- node[above, font=\baslikfont\scriptsize\bfseries, text=anarenk]{indirgeme} (4.4,-4.1);
  \node[font=\baslikfont\scriptsize, text=black!60, align=center] at (5.2,-6.0) {$[P]$ varsayımının yerine\\$\mathcal{E}$ türetimi konur};
\end{tikzpicture}
\caption{İki temel indirgeme. Bir giriş kuralının hemen ardından aynı
bağlacın eleme kuralı geliyorsa, ikisi birlikte silinebilir.}
\label{fig:normallestirme}
\end{figure}

\begin{tanim}{Normal türetim}{normal}
Bir türetimde, bir bağlacın giriş kuralıyla elde edilen bir önermeye
hemen ardından aynı bağlacın eleme kuralı uygulanıyorsa, bu önermeye
\textbf{dolambaç} (İng.\ \emph{detour}, \emph{maximal formula}) denir.
Dolambaç içermeyen türetimlere \textbf{normal} türetim denir.
\end{tanim}

İmplikasyondaki dolambaç da aynı fikri izler: Önce bir fonksiyon kurar,
sonra onu hemen bir argümana uygularız.

\leandosya[İmplikasyonda dolambaç]{Bolum05}{dolambacli_imp}

\noindent Bu ispatın indirgenmiş hâli \lc{hpq hp}'dir: Fonksiyonun
gövdesinde \lc{h} gördüğümüz her yere \lc{hp} koyarız. Programlama
dillerinde buna \textbf{$\beta$-indirgeme}, yani \emph{bir fonksiyonu
çağırmak} denir. Aynı dolambaçları veri dünyasında yazdığımızda
\Lean{} onları \emph{hesaplar}:

\leandosya[Normalleştirme = hesaplama]{Bolum05}{hesaplama}

\begin{teorem}{Normalleştirme, Prawitz 1965}{normallestirme}
Yapıcı önermeler mantığında her türetim, aynı varsayımlardan aynı sonucu
veren bir normal türetime indirgenebilir \citep{prawitz1965}.
\end{teorem}

\begin{onemlifikir}[Önemli Fikir: Normalleştirme = program çalıştırma]
Curry--Howard eşlemesi altında Prawitz'in teoremi şunu söyler: Her
ispat (program) dolambaçları silinerek (hesaplanarak) bir sonuca ulaşır;
hesaplama sonsuza dek sürmez. İspat kuramında ``normalleştirme'', bilgisayar
biliminde ``program çalıştırma'' denen şeyin aynısıdır. Bu bağlantıyı
\cref{bol:lambda}'te $\lambda$-hesabı ile kesinleştireceğiz.
\end{onemlifikir}

Normal türetimlerin önemli bir özelliği vardır: \emph{Alt formül özelliği}.
Normal bir türetimde görünen her önerme, ya sonucun ya da açık
varsayımlardan birinin bir parçasıdır. Yani normal bir ispat ``dışarıdan''
hiçbir yeni kavram getirmez. Bundan çarpıcı bir sonuç çıkar:

\begin{teorem}{Tutarlılık}{tutarlilik}
Yapıcı önermeler mantığında $\bot$ önermesi, varsayımsız olarak
türetilemez: $\nvdash \bot$.
\end{teorem}

\noindent\textit{İspat taslağı.} $\bot$'un bir türetimi olsaydı,
Teorem~\ref{trm:normallestirme}'e göre normal bir türetimi de olurdu. Alt
formül özelliğine göre bu türetimdeki her önerme $\bot$'un bir parçası,
yani $\bot$'un kendisi olmalıdır. Ama $\bot$'un giriş kuralı yoktur ve
açık varsayım da yoktur; dolayısıyla ağacın hiçbir yaprağı olamaz.
Bu bir çelişkidir. \hfill\qedsymbol

\noindent Bu, mantığın \emph{kendi tutarlılığını} kendi araçlarıyla
gösteren ilk ve en basit örnektir. Aynı fikri aritmetiğe uygulama
çabası, Gentzen'i 1936'da aritmetiğin tutarlılık ispatına götürdü
\citep{gentzen1936}; bu hikâyeye \ciltiki{4}'te döneceğiz.

% ---------------------------------------------------------------------
\section{Taktik birleştiricileri}\label{sec:birlestiriciler}

Bazı ispatlarda aynı taktiği birçok hedefe uygulamamız gerekir. Üç
önermenin ``ve''sini kurmak için:

\leandosya[Uzun yol]{Bolum05}{birlestirici_uzun}

\noindent Bu ispatı \emph{taktik birleştiricileri} ile kısaltabiliriz:

\leandosya[Kısa yol]{Bolum05}{birlestirici_kisa}

\begin{center}
\renewcommand{\arraystretch}{1.45}
\rowcolors{2}{tablosatir}{white}
\begin{tabularx}{\linewidth}{>{\raggedright\arraybackslash}p{3.6cm} X}
\rowcolor{tabloust}
\tablobaslik{Birleştirici} & \tablobaslik{Ne yapar?} \\
\lc{t₁ <;> t₂}                     & \lc{t₁}'i uygular, \emph{ortaya çıkan her hedefe} \lc{t₂}'yi uygular. \\
\leantac{all\_goals} \lc{t}        & Açık olan bütün hedeflere \lc{t}'yi uygular. \\
\leantac{first} \lc{| t₁ | t₂}     & Önce \lc{t₁}'i dener; başarısız olursa \lc{t₂}'yi. \\
\leantac{try} \lc{t}               & \lc{t}'yi dener; başarısız olursa hiçbir şey yapmaz. \\
\leantac{repeat} \lc{t}            & \lc{t}'yi başarısız olana dek tekrar tekrar uygular. \\
\leantac{refine} \lc{⟨?\_, ?\_⟩}   & Terimin bir kısmını yazar; \lc{?\_} ile işaretli boşlukları yeni hedef yapar. \\
\end{tabularx}
\end{center}

\leantac{refine}, terim modu ile taktik modunu birleştirir: İspatın
iskeletini bir terim olarak yazar, boşlukları taktiklerle doldururuz.
\leantac{first} ile durum ayırmayı birleştiren bir örnek:

\leandosya[Durum ayırma ve first]{Bolum05}{birlestirici_veya}

\noindent İki durumun her birinde \Lean{} önce \lc{Or.inr h}'yi dener;
türler uyuşmazsa \lc{Or.inl h}'ye geçer.

\begin{dikkat}[Dikkat: Kısa her zaman iyi değildir]
Birleştiriciler ispatı kısaltır ama okunmasını zorlaştırabilir. Bir
ispat \lc{<;>} ile bir satıra indiğinde, okur hangi hedefin hangi
taktikle kapandığını göremez. Bir kural olarak, birleştiricileri
\emph{aynı biçimdeki} hedefler için kullanın (\lc{refine ⟨?\_, ?\_, ?\_⟩
<;> assumption} gibi), farklı mantık gerektiren durumlar için değil.
\end{dikkat}

\paragraph{Otomasyona bir bakış}
Önermeler mantığı karar verilebilir olduğuna göre, bu bölümdeki
ispatları otomatik olarak bulan bir taktik yazılabilir. \Mathlib{}'teki
\leantac{tauto} taktiği tam olarak bunu yapar. \Lean{}'in güncel
sürümlerinde \lc{exact?} komutu ise kütüphanede hedefi kapatan bir
teorem arar. Bu araçları \ciltiki{6}'da tanıyacak,
\ciltiki{7}'de ise kendi önermeler mantığı ispatçımızı
yazacağız. O zamana dek ispatları elle yazmanızı öneririz: Otomasyon,
ne yaptığını anlayan birinin elinde güçlü bir araçtır.

% ---------------------------------------------------------------------
\begin{tarihce}[Tarihçe: Hilbert sistemlerinden Gentzen'e]
Frege'den Hilbert'e kadar mantık sistemleri \emph{aksiyomatik} olarak
kurulurdu: Birkaç aksiyom şeması (örneğin
$P \to (Q \to P)$ ve $(P \to Q \to R) \to (P \to Q) \to P \to R$) ve
tek bir kural (modus ponens). Bu \emph{Hilbert sistemleri} teorik olarak
güçlüdür, ama içlerinde ispat yazmak zahmetlidir: $P \to P$ gibi basit
bir önermenin ispatı bile beş adım tutar. Gentzen 1935'te bu sistemlerin
``gerçek matematiksel akıl yürütmeden çok uzak'' olduğunu yazarak doğal
çıkarımı önerdi \citep{gentzen1935}. Bugün ispat asistanlarının
neredeyse tamamı doğal çıkarım ya da onun bir uzantısı üzerine kuruludur.
Hilbert sisteminin aksiyomlarının Bölüm~2'deki $K$ ve $S$
kombinatörleri olduğuna dikkat edin.
\end{tarihce}

\begin{arastirmanotu}
\textbf{Uyum ve ispat-kuramsal anlambilim.} \Cref{sec:dogal-cikarim-fikri}'te
sözünü ettiğimiz giriş ve eleme kuralları arasındaki \emph{uyum}, Michael
Dummett ve Dag Prawitz'in \emph{ispat-kuramsal anlambilim} programının
temel kavramıdır: Bir bağlacın anlamı, doğruluk tablosuyla değil, giriş
kurallarıyla verilir; eleme kuralları bu anlamla uyumlu olmalıdır.
Normalleştirme teoremi, bu uyumun teknik ifadesidir
\citep{prawitz1965, dummett2000}.

\textbf{Taktik dilleri.} \Lean{}'in taktik dili, kendisi \Lean{}'de
yazılmış bir programlama dilidir; kullanıcılar yeni taktikler
tanımlayabilir \citep{demoura2021lean4, ullrich2022macros}. Taktik
ispatlarının okunabilirliği ve bakımı, büyük formalizasyon projelerinde
pratik bir sorun olarak tartışılmaya devam ediyor; \Mathlib{} bu nedenle
belirli bir ispat stili önerir (\ciltiki{9}).

\textbf{İleri okuma.} Normalleştirme ve alt formül özelliğinin tam
ispatları için \citet{troelstra2000}; Curry--Howard açısından
normalleştirme için \citet{girard1989} ve \citet{sorensen2006}.
\end{arastirmanotu}

\begin{bolumozeti}
\begin{itemize}
  \item Her bağlacın giriş ve eleme kuralları vardır; tek klasik kural
        RAA'dır (\cref{fig:kural-katalogu}).
  \item \Lean{}'in hedef ekranı bir ardışıktır: $\Gamma \vdash \varphi$.
        Bir taktik, bir kuralı aşağıdan yukarı uygular.
  \item İspatlar geriye (\leantac{apply}, \leantac{suffices}) ya da
        ileriye (\leantac{have}) doğru kurulabilir; iyi ispatlar ikisini
        karıştırır.
  \item Strateji: Önce hedefin biçimine (giriş kuralları), sonra
        hipotezlerin biçimine (eleme kuralları) bak.
  \item Dolambaç, bir girişin hemen ardından gelen elemedir. Normalleştirme
        dolambaçları siler; Curry--Howard altında bu, program çalıştırmaktır.
        Normal ispatların alt formül özelliği, tutarlılığı verir.
  \item \lc{<;>}, \leantac{all\_goals}, \leantac{first} ve \leantac{refine}
        ispatları kısaltır; ölçülü kullanılmalıdır.
\end{itemize}
\end{bolumozeti}

\begin{alistirmalar}
\alistirma{1} Varsayımların sırasını değiştirin:
\lc{(P → Q → R) → (Q → P → R)}. Terim modunda bu ispat hangi
fonksiyondur? (\lc{a5\_1})

\alistirma{1} \lc{(P ∧ Q → R) → P → Q → R} önermesini önce yalnızca
\leantac{have} ve \leantac{exact} ile (\lc{a5\_2}), sonra yalnızca
\leantac{apply} ve \leantac{exact} ile (\lc{a5\_3}) ispatlayın. Hangisi
size daha okunabilir geldi?

\alistirma{2} İmplikasyonun ve'ye dağıldığını ispatlayın:
\lc{(P → Q ∧ R) ↔ (P → Q) ∧ (P → R)}. (\lc{a5\_4})

\alistirma{2} \lc{(P ∨ Q) ∧ (P → R) ∧ (Q → R) → R}. Önce
\cref{fig:strateji}'yi kullanarak bir plan yazın, sonra ispatlayın.
(\lc{a5\_5})

\alistirma{2} \lc{P ∧ Q ∧ R → R ∧ Q ∧ P} önermesini \leantac{rintro}
ve tek bir \lc{refine ... <;> assumption} satırıyla ispatlayın.
(\lc{a5\_6})

\alistirma{2} Dört halkalı zinciri önce geriye, sonra ileriye doğru
ispatlayın: \lc{(P → Q) → (Q → R) → (R → S) → P → S}. (\lc{a5\_7})

\alistirma{3} \lc{(P ∨ Q → R) → ¬R → ¬P ∧ ¬Q}. Doğal çıkarım ağacını
da çizin. (\lc{a5\_8})

\alistirma{3} \lc{(P ∨ Q) ∧ (R ∨ S) → (P ∧ R) ∨ (P ∧ S) ∨ (Q ∧ R) ∨ (Q ∧ S)}.
İpucu: \leantac{rintro} ile \lc{⟨hp | hq, hr | hs⟩} örüntüsü dört durumu
birden açar. (\lc{a5\_9})

\alistirma{2} (Kâğıt-kalem) Aşağıdaki türetimi normalleştirin. Kaç
indirgeme adımı gerekti?
\[
\begin{prooftree}
  \hypo{[P]^{1}}
  \hypo{[Q]^{2}}
  \infer2[$\wedge$G]{P \wedge Q}
  \infer1[$\wedge$E$_2$]{Q}
  \infer1[$\to$G$^{1}$]{P \to Q}
  \hypo{P}
  \infer2[$\to$E]{Q}
  \infer1[$\to$G$^{2}$]{Q \to Q}
\end{prooftree}
\]

\alistirma{3} (Kâğıt-kalem) Alt formül özelliğini kullanarak $P \vee Q$'nun
varsayımsız olarak türetilemeyeceğini gösterin ($P$, $Q$ farklı
atomlar). İpucu: Normal bir türetimin son kuralı ne olabilir?

\alistirma{3} (Tartışma) \cref{fig:strateji}'deki haritayı kullanarak
Peirce yasasını ispatlamaya çalışın. Haritanın hangi adımında
takılıyorsunuz? Bu takılma, \cref{bol:yapici-klasik}'teki Kripke
modeliyle nasıl ilişkilidir?
\end{alistirmalar}
   % Doğal çıkarım ağaçları ve taktikler

\part{Yüklemler ve Tümevarım}
% =====================================================================
%  Bölüm 6 — Yüklemler Mantığı: ∀, ∃ ve Eşitlik
% =====================================================================
\chapter{Yüklemler Mantığı: \texorpdfstring{$\forall$}{∀}, \texorpdfstring{$\exists$}{∃} ve Eşitlik}\label{bol:yuklemler}

\begin{bolumhedefleri}
\begin{itemize}
  \item Yüklemleri \emph{önerme değerli fonksiyonlar} olarak görmeyi ve
        \Lean{}'de \lc{α → Prop} türüyle yazmayı,
  \item $\forall$ ve $\exists$ için giriş ve eleme kurallarını, özellikle
        \emph{öz değişken koşulunu},
  \item Niceleyicilerin değilleme ve sırayla nasıl etkileştiğini; hangi
        etkileşimlerin yapıcı, hangilerinin klasik olduğunu,
  \item Serbest ve bağlı değişkenleri, yerine koymayı ve değişken
        yakalama tuzağını,
  \item Eşitliği: \lc{rfl}, \leantac{rw}, \lc{▸}, \lc{calc} ve Leibniz ilkesini,
  \item Berber paradoksunu \Lean{}'de ispatlamayı ve boş evren inceliğini,
  \item Curry--Howard tablosunun bağımlı tiplere genişlemesini:
        $\forall = \Pi$, $\exists = \Sigma$.
\end{itemize}
\end{bolumhedefleri}

Önermeler mantığı, ``Sokrates ölümlüdür'' cümlesini bölünmez bir $P$
harfi olarak görür. Oysa Bölüm~1'deki kıyasın geçerliliği tam da bu
cümlenin \emph{içindeki} yapıya dayanıyordu: ``Bütün insanlar
ölümlüdür; Sokrates bir insandır; o hâlde Sokrates ölümlüdür.'' Bu
çıkarımı ifade etmek için önermelerin içine bakmamız, \emph{bireyler} ve
onların \emph{özellikleri} hakkında konuşmamız gerekir. Bu bölümün konusu
budur.

% ---------------------------------------------------------------------
\section{Önermelerden yüklemlere}\label{sec:yuklemler}

\begin{tanim}{Yüklem}{yuklem}
Bir $\alpha$ \textbf{tanım evreni} üzerindeki bir \textbf{yüklem},
$\alpha$'nın her elemanına bir önerme karşılık getiren bir
fonksiyondur:
\[
  P : \alpha \to \mathsf{Prop}.
\]
$P(a)$, ``$a$, $P$ özelliğini taşır'' diye okunur. İki argümanlı bir
yüklem ($R : \alpha \to \beta \to \mathsf{Prop}$) bir \textbf{ilişkidir}.
\end{tanim}

\Lean{}'de yüklemler sıradan fonksiyonlardır. ``Çifttir'' yüklemi, bir
doğal sayı alıp bir önerme döndüren fonksiyondur:

\leandosya[Yüklemlerin türleri]{Bolum06}{yuklem_tip}

\begin{figure}[htbp]
\centering
\begin{tikzpicture}[
    eleman/.style={circle, fill=anarenk, inner sep=1.8pt},
    onerme/.style={draw=#1, fill=#1!8, rounded corners=3pt, font=\small,
                   inner sep=4pt, minimum width=2.2cm},
    ok/.style={-{Stealth[length=2.2mm]}, thick, black!45},
  ]
  \draw[anarenk, thick, fill=tanimbg, rounded corners=10pt] (-0.9,-2.6) rectangle (0.9,1.0);
  \node[font=\baslikfont\bfseries\small, text=anarenk] at (0,1.35) {$\alpha = \mathbb{N}$};
  \foreach \n/\y in {0/0.5, 1/-0.2, 2/-0.9, 3/-1.6} {
    \node[font=\small] (n\n) at (0,\y) {$\n$};
  }
  \node[font=\small] at (0,-2.2) {$\vdots$};
  \node[onerme=bittiframe] (p0) at (4.6,0.5)  {$0 \% 2 = 0$};
  \node[onerme=hataframe]  (p1) at (4.6,-0.2) {$1 \% 2 = 0$};
  \node[onerme=bittiframe] (p2) at (4.6,-0.9) {$2 \% 2 = 0$};
  \node[onerme=hataframe]  (p3) at (4.6,-1.6) {$3 \% 2 = 0$};
  \foreach \n in {0,1,2,3} \draw[ok] (n\n) -- (p\n);
  \draw[ok, dashed] (0.3,-2.2) -- (3.4,-2.2);
  \node[font=\baslikfont\bfseries\small, text=vurgu] at (2.1,1.35) {$P$};
  \node[font=\baslikfont\bfseries\small, text=koyu] at (4.6,1.35) {$\mathsf{Prop}$};
  \node[font=\baslikfont\scriptsize, text=bittiframe, anchor=west] at (6.0,0.15) {doğru (kanıtı var)};
  \node[font=\baslikfont\scriptsize, text=hataframe, anchor=west] at (6.0,-0.55) {yanlış (kanıtı yok)};
\end{tikzpicture}
\caption{``Çifttir'' yüklemi $P(n) := (n \% 2 = 0)$, her doğal sayıya
bir önerme karşılık getirir. Önermelerin bazıları doğru, bazıları
yanlıştır; ama hepsi birer önermedir.}
\label{fig:yuklem}
\end{figure}

Artık Bölüm~1'deki kıyası tam olarak yazabiliriz:

\leandosya[Sokrates kıyası]{Bolum06}{sokrates}

\noindent İspatın tamamı \lc{h1 sokrates h2}'dir. \lc{h1}'i önce
\lc{sokrates}'e, sonra \lc{h2}'ye uyguluyoruz. ``Bütün insanlar
ölümlüdür'' önermesinin kanıtı, \emph{her birey için} ``insansa
ölümlüdür'' kanıtını üreten bir fonksiyondur.

% ---------------------------------------------------------------------
\section{Evrensel niceleyici (\texorpdfstring{$\forall$}{∀})}\label{sec:forall}

\begin{center}
\begin{tabular}{c@{\hspace{5em}}c}
\begin{prooftree}
  \hypo{\mathcal{D}}
  \ellipsis{}{P(a)}
  \infer1[$\forall$G$^{*}$]{\forall x,\; P(x)}
\end{prooftree}
&
\begin{prooftree}
  \hypo{\forall x,\; P(x)}
  \infer1[$\forall$E]{P(t)}
\end{prooftree}
\\[10pt]
\small\textcolor{ikincil}{\bfseries giriş:} \leantac{intro x} &
\small\textcolor{vurgu}{\bfseries eleme:} \lc{h t}, \leantac{specialize h t}
\end{tabular}
\end{center}

\noindent {\small $^{*}$\,\textbf{Öz değişken koşulu:} $a$, $\mathcal{D}$
türetiminin açık varsayımlarında geçmemelidir.}

\textbf{$\forall$ girişi.} ``Her $x$ için $P(x)$'' göstermek için
\emph{keyfi} bir $a$ seçer ve $P(a)$'yı gösteririz. Matematikçilerin
``keyfi bir $x$ alalım'' cümlesi budur. Kilit sözcük \emph{keyfi}dir:
$a$ hakkında hiçbir özel varsayım yapmamış olmalıyız. Aksi hâlde
``$a = 2$ olsun; $a$ çifttir; o hâlde her sayı çifttir'' gibi bir saçmalık
türetirdik. Bu kısıta \textbf{öz değişken koşulu} (İng.\
\emph{eigenvariable condition}) denir.

\textbf{$\forall$ elemesi.} ``Her $x$ için $P(x)$'' biliyorsak, herhangi bir
$t$ terimi için $P(t)$'yi biliriz. Bu, ``genelden özele'' geçiştir.

\leandosya[$\forall$ girişi]{Bolum06}{forall_giris}

\begin{hedefgecisi}
\begin{hedefdurumu}[Önce]
α : Type
⊢ ∀ (x : α), x = x
\end{hedefdurumu}
\sonra
\begin{hedefdurumu}[\texttt{intro x} sonrası]
α : Type
x : α
⊢ x = x
\end{hedefdurumu}
\end{hedefgecisi}

\noindent \leantac{intro x}, bağlama yeni bir \lc{x : α} ekler. \Lean{},
öz değişken koşulunu kendiliğinden sağlar: \lc{x} yeni bir addır ve
bağlamdaki hiçbir hipotez onun hakkında bir şey söylemez.

\begin{onemlifikir}[Önemli Fikir: $\forall$ bir fonksiyondur, $\to$ onun özel hâlidir]
\leantac{intro}'nun hem $\to$ hem $\forall$ için çalışması bir rastlantı
değildir. \Lean{}'de \lc{∀ x : α, P x} türünde bir kanıt, her \lc{a : α}
için \lc{P a} türünde bir kanıt üreten bir fonksiyondur. Tek fark,
\emph{dönüş türünün argümana bağlı olmasıdır}: \lc{h 3 : P 3},
\lc{h 5 : P 5}. Böyle fonksiyonlara \textbf{bağımlı fonksiyon} denir.
\lc{P → Q} ise dönüş türü argümana bağlı olmayan özel bir durumdur.
\end{onemlifikir}

\leandosya[$\forall$ elemesi]{Bolum06}{forall_eleme}

\noindent İki kuralı birlikte kullanan bir örnek:

\leandosya[$\forall$ ve $\wedge$]{Bolum06}{forall_ve}

\noindent Bazen bir genel hipotezi belirli bir bireye ``indirgemek''
isteriz. \leantac{specialize} taktiği, hipotezi o bireye uygulanmış
hâliyle değiştirir:

\leandosya[\texttt{specialize}]{Bolum06}{specialize}

% ---------------------------------------------------------------------
\section{Varoluşsal niceleyici (\texorpdfstring{$\exists$}{∃})}\label{sec:exists}

\begin{center}
\begin{tabular}{c@{\hspace{5em}}c}
\begin{prooftree}
  \hypo{P(t)}
  \infer1[$\exists$G]{\exists x,\; P(x)}
\end{prooftree}
&
\begin{prooftree}
  \hypo{\exists x,\; P(x)}
  \hypo{[P(a)]^{n}}
  \ellipsis{}{R}
  \infer2[$\exists$E$^{n,*}$]{R}
\end{prooftree}
\\[10pt]
\small\textcolor{ikincil}{\bfseries giriş:} \lc{⟨t, h⟩} &
\small\textcolor{vurgu}{\bfseries eleme:} \leantac{obtain ⟨a, ha⟩ := h}
\end{tabular}
\end{center}

\noindent {\small $^{*}$\,\textbf{Öz değişken koşulu:} $a$, $R$'de ve
$P(a)$ dışındaki açık varsayımlarda geçmemelidir.}

\textbf{$\exists$ girişi.} ``$P$'yi sağlayan bir $x$ vardır'' göstermek
için somut bir \textbf{tanık} $t$ verir ve $P(t)$'yi gösteririz.
BHK yorumuyla uyumlu olarak, bir varlık kanıtı \emph{tanığını taşır}.

\textbf{$\exists$ elemesi.} ``$P$'yi sağlayan bir $x$ vardır'' biliyorsak,
ona bir ad veririz: ``Öyle bir $a$ olsun ki $P(a)$.'' Bu $a$ hakkında
$P(a)$'dan başka hiçbir şey bilmeyiz. Öz değişken koşulu, $a$'nın
sonuca ``sızmasını'' engeller: $a$ yalnızca ispatın içinde yaşayan geçici
bir addır.

\leandosya[$\exists$ girişi]{Bolum06}{exists_giris}

\noindent Köşeli ayraçlar, \lc{∧}'deki gibi bir çift kurar: Birinci
bileşen tanık, ikinci bileşen tanığın özelliğinin kanıtıdır. İkinci
ispatta \leantac{refine} önce tanığı verir; geriye kalan hedef \lc{4 > 3}
olur.

\leandosya[$\exists$ elemesi]{Bolum06}{exists_eleme}

\begin{hedefgecisi}
\begin{hedefdurumu}[\texttt{obtain} öncesi]
h1 : ∀ (x : α), P x → Q x
h2 : ∃ x, P x
⊢ ∃ x, Q x
\end{hedefdurumu}
\sonra
\begin{hedefdurumu}[\texttt{obtain} sonrası]
h1 : ∀ (x : α), P x → Q x
a : α
ha : P a
⊢ ∃ x, Q x
\end{hedefdurumu}
\end{hedefgecisi}

\noindent Terim modunda aynı ispat, \leankw{let} ile örüntü eşleme
kullanır:

\leandosya[$\exists$ elemesi, terim modu]{Bolum06}{exists_eleme_terim}

\begin{figure}[htbp]
\centering
\begin{tikzpicture}[
    evren/.style={draw=#1, thick, fill=#1!6, rounded corners=12pt,
                  minimum width=4.2cm, minimum height=2.8cm},
    nokta/.style={circle, inner sep=2.4pt, fill=#1},
  ]
  % ∀
  \node[evren=ikincil] (A) at (0,0) {};
  \foreach \x/\y in {-1.3/0.7, -0.4/0.9, 0.6/0.6, 1.4/0.8, -1.1/-0.3, 0.0/-0.1,
                     1.0/-0.4, -0.5/-0.9, 0.5/-1.0, 1.5/-0.9}
    \node[nokta=bittiframe] at (\x,\y) {};
  \node[font=\baslikfont\bfseries, text=ikincil] at (0,1.75) {$\forall x,\; P(x)$};
  \node[font=\baslikfont\scriptsize, text=black!60, align=center] at (0,-1.85)
       {\emph{her} eleman $P$'yi sağlar\\kanıt: her eleman için bir kanıt (fonksiyon)};
  % ∃
  \node[evren=vurgu] (B) at (6.4,0) {};
  \foreach \x/\y in {-1.3/0.7, -0.4/0.9, 0.6/0.6, 1.4/0.8, -1.1/-0.3,
                     1.0/-0.4, -0.5/-0.9, 0.5/-1.0, 1.5/-0.9}
    \node[nokta=black!25] at ({6.4+\x},\y) {};
  \node[nokta=bittiframe, minimum size=9pt] (t) at (6.4,-0.1) {};
  \node[font=\baslikfont\scriptsize\bfseries, text=bittiframe, right=1pt of t] {tanık $t$};
  \node[font=\baslikfont\bfseries, text=vurgu] at (6.4,1.75) {$\exists x,\; P(x)$};
  \node[font=\baslikfont\scriptsize, text=black!60, align=center] at (6.4,-1.85)
       {\emph{en az bir} eleman $P$'yi sağlar\\kanıt: bir tanık ve onun kanıtı (çift)};
\end{tikzpicture}
\caption{İki niceleyicinin anlamı ve BHK yorumuna göre kanıtları.}
\label{fig:niceleyiciler}
\end{figure}

% ---------------------------------------------------------------------
\section{Niceleyicilerin etkileşimi}\label{sec:niceleyici-etkilesim}

\subsection{Değilleme ve niceleyiciler}

Niceleyiciler değillemeyle De Morgan yasalarına benzer biçimde etkileşir.
Hatta $\forall$, büyük bir ``ve'', $\exists$ ise büyük bir ``veya''
gibidir: Evren $\{a, b\}$ ise $\forall x\,P(x) = P(a) \wedge P(b)$ ve
$\exists x\,P(x) = P(a) \vee P(b)$ olur. Bölüm~3'teki asimetri burada
da karşımıza çıkar:

\leandosya[``Hiç kimse'' = ``herkes değil'']{Bolum06}{degil_exists}
\leandosya[Yapıcı yön]{Bolum06}{degil_forall_kolay}
\leandosya[Klasik yön]{Bolum06}{degil_forall_klasik}

\noindent Son ispatın neden klasik mantık gerektirdiğini BHK yorumuyla
anlayabiliriz. Hedef bir $\exists$ önermesi olduğuna göre bir tanık
vermemiz gerekir. Elimizdeki tek bilgi ise ``herkesin $P$'yi sağladığı
doğru değildir''dir. Bu bilgi, $P$'yi sağlamayan bireyin \emph{kim}
olduğunu söylemez.

\begin{table}[htbp]
\centering
\renewcommand{\arraystretch}{1.5}
\rowcolors{2}{tablosatir}{white}
\begin{tabular}{>{\centering\arraybackslash}p{5.6cm}
                >{\centering\arraybackslash}p{2.0cm}
                >{\centering\arraybackslash}p{2.0cm}
                >{\raggedright\arraybackslash}p{2.6cm}}
\rowcolor{tabloust}
\tablobaslik{Önerme} & \tablobaslik{$\to$} & \tablobaslik{$\leftarrow$} & \tablobaslik{Lean} \\
$\neg\exists x\,P(x) \leftrightarrow \forall x\,\neg P(x)$ & \textcolor{bittiframe}{yapıcı} & \textcolor{bittiframe}{yapıcı} & \lc{degil\_exists} \\
$\neg\forall x\,P(x) \leftrightarrow \exists x\,\neg P(x)$ & \textcolor{hataframe}{klasik} & \textcolor{bittiframe}{yapıcı} & \lc{degil\_forall} \\
$\forall x\,(P \wedge Q) \leftrightarrow \forall x\,P \wedge \forall x\,Q$ & \textcolor{bittiframe}{yapıcı} & \textcolor{bittiframe}{yapıcı} & Alıştırma~6.4 \\
$\exists x\,(P \vee Q) \leftrightarrow \exists x\,P \vee \exists x\,Q$ & \textcolor{bittiframe}{yapıcı} & \textcolor{bittiframe}{yapıcı} & Alıştırma~6.5 \\
$\exists x\,(P \wedge Q) \to \exists x\,P \wedge \exists x\,Q$ & \textcolor{bittiframe}{yapıcı} & \textcolor{hataframe}{yanlış} & Alıştırma~6.6 \\
$\exists x\,\forall y\,R \to \forall y\,\exists x\,R$ & \textcolor{bittiframe}{yapıcı} & \textcolor{hataframe}{yanlış} & \lc{exists\_forall\_sira} \\
\end{tabular}
\caption{Niceleyicilerin etkileşimleri. ``Yanlış'', o yönün genel olarak
geçerli olmadığı, yani bir karşı örneği bulunduğu anlamına gelir.}
\label{tab:niceleyici-etkilesim}
\end{table}

\subsection{Niceleyicilerin sırası}

Niceleyicilerin sırası anlamı kökten değiştirir:
\begin{center}
\begin{tabular}{rl}
$\forall y\,\exists x\; \mathit{Sever}(x, y)$: & ``Herkesi seven \emph{biri} vardır'' \emph{değil}, ``herkesin bir seveni vardır.'' \\
$\exists x\,\forall y\; \mathit{Sever}(x, y)$: & ``Herkesi seven biri vardır.''
\end{tabular}
\end{center}

\noindent İkincisi birincisinden güçlüdür:

\leandosya[$\exists\forall \Rightarrow \forall\exists$]{Bolum06}{sira}

\noindent Ters yön yanlıştır. Doğal sayılarda ``her sayıdan büyük bir sayı
vardır'' doğrudur, ama ``her sayıdan büyük bir sayı'' yoktur
(\cref{fig:sira}):

\leandosya[Karşı örnek]{Bolum06}{sira_karsi}

\begin{figure}[htbp]
\centering
\begin{tikzpicture}[
    sayi/.style={circle, draw=anarenk, fill=tanimbg, thick, minimum size=0.75cm,
                 font=\small},
    ok/.style={-{Stealth[length=2.2mm]}, thick, #1},
  ]
  \node[font=\baslikfont\bfseries\small, text=bittiframe] at (2.5,1.3) {$\forall y\,\exists x,\; x > y$ \quad (doğru)};
  \foreach \i in {0,...,4} {
    \node[sayi] (y\i) at ({1.25*\i},0) {$\i$};
    \node[sayi, draw=bittiframe, fill=bittibg] (x\i) at ({1.25*\i},-1.6) {$\pgfmathparse{int(\i+1)}\pgfmathresult$};
    \draw[ok=bittiframe] (y\i) -- (x\i);
  }
  \node[font=\small] at (6.25,0) {$\cdots$};
  \node[font=\baslikfont\scriptsize, text=black!60] at (2.5,-2.4) {her $y$ için ayrı bir tanık: $x = y+1$};

  \node[font=\baslikfont\bfseries\small, text=hataframe] at (10.3,1.3) {$\exists x\,\forall y,\; x > y$ \quad (yanlış)};
  \node[sayi, draw=hataframe, fill=hatabg] (X) at (10.3,-1.6) {$x$};
  \foreach \i in {0,...,3} {
    \node[sayi] (z\i) at ({8.4+1.25*\i},0) {};
    \draw[ok=hataframe] (X) -- (z\i);
  }
  \node[font=\baslikfont\scriptsize, text=hataframe, align=center] at (10.3,-2.5)
       {tek bir $x$ hepsinden büyük olamaz:\\$y = x$ seçilince $x > x$ gerekir};
\end{tikzpicture}
\caption{Niceleyicilerin sırası. $\forall\exists$'de tanık $y$'ye bağlı
olabilir; $\exists\forall$'de tek bir tanık herkese yetmelidir.}
\label{fig:sira}
\end{figure}

\begin{dikkat}[Dikkat: Tanık neye bağlıdır?]
$\forall y\,\exists x$ önermesinin kanıtı, Curry--Howard altında, her
$y$ için bir $x$ \emph{üreten bir fonksiyondur}: \lc{fun y => ⟨y + 1, \ldots⟩}.
Tanık $y$'ye bağlıdır. $\exists x\,\forall y$ ise tek bir sabit tanık
ister. Analizdeki ``sürekli'' ile ``düzgün sürekli'' arasındaki farkın
özü de budur: $\delta$, $x$ noktasına bağlı olabilir mi?
\end{dikkat}

% ---------------------------------------------------------------------
\section{Serbest ve bağlı değişkenler}\label{sec:serbest-bagli}

$\forall x,\; x + y = y + x$ önermesinde $x$ ve $y$ değişkenlerinin
rolleri farklıdır. $x$ niceleyiciye \emph{bağlıdır}: Onun yerine başka
bir ad koysak ($\forall z,\; z + y = y + z$) önermenin anlamı değişmez.
$y$ ise \emph{serbesttir}: Önermenin doğruluğu $y$'nin ne olduğuna
bağlıdır.

\begin{tanim}{Serbest ve bağlı değişken}{serbest-bagli}
Bir önermede bir $x$ değişkeninin bir geçişi, bir $\forall x$ ya da
$\exists x$ niceleyicisinin kapsamı içindeyse \textbf{bağlı}, değilse
\textbf{serbesttir}. Bağlı değişkenlerin adları değiştirilerek elde
edilen önermelere \textbf{$\alpha$-denk} önermeler denir ve aynı önerme
sayılırlar.
\end{tanim}

\paragraph{Yerine koyma ve değişken yakalama}
$\forall$E kuralı $\forall x\,P(x)$'ten $P(t)$'ye geçer: $x$'in yerine $t$
terimini koyar. Burada ince bir tuzak vardır (\cref{fig:yakalama}).

\begin{figure}[htbp]
\centering
\begin{tikzpicture}[
    kutu/.style={draw=#1, fill=#1!7, thick, rounded corners=4pt,
                 text width=5.6cm, align=center, inner sep=6pt, font=\small},
    ok/.style={-{Stealth[length=2.6mm]}, very thick, #1},
  ]
  \node[kutu=anarenk] (a) at (0,0)
    {$\forall x,\; \exists y,\; y > x$\\[2pt]{\scriptsize ``Her sayıdan büyük bir sayı vardır.'' (doğru)}};
  \node[kutu=hataframe] (b) at (-3.4,-2.6)
    {$\exists y,\; y > \textcolor{hataframe}{\mathbf{y}}$\\[2pt]{\scriptsize $x$ yerine $y$ konuldu; $y$ \emph{yakalandı}. (yanlış!)}};
  \node[kutu=bittiframe] (c) at (3.4,-2.6)
    {$\exists z,\; z > \textcolor{bittiframe}{\mathbf{y}}$\\[2pt]{\scriptsize Önce bağlı $y$, $z$ yapıldı; sonra yerine koyma. (doğru)}};
  \draw[ok=hataframe] (a) -- node[left=6pt, font=\baslikfont\scriptsize, text=hataframe]{saf yerine koyma} (b);
  \draw[ok=bittiframe] (a) -- node[right=6pt, font=\baslikfont\scriptsize, text=bittiframe]{yakalamadan kaçınan} (c);
\end{tikzpicture}
\caption{Değişken yakalama: $\forall$E ile $x$ yerine $y$ terimini koyarken
içerideki $\exists y$ niceleyicisi, dışarıdan gelen $y$'yi ``yakalar''.}
\label{fig:yakalama}
\end{figure}

Doğru yerine koyma, gerektiğinde bağlı değişkenlerin adlarını değiştirir.
Kâğıt üzerinde bu dikkat gerektiren bir iştir ve mantık kitaplarının
tanımlarında sık rastlanan hataların kaynağıdır. \Lean{} ise bağlı
değişkenleri içeride ad yerine \emph{de Bruijn indeksleriyle} (``kaç
niceleyici yukarıdaki değişken'') tutar; böylece yakalama hiç oluşmaz.
Bu temsili \cref{bol:sozdizim}'te kendi formül tipimizi tanımlarken
kullanacağız.

% ---------------------------------------------------------------------
\section{Eşitlik}\label{sec:esitlik}

Eşitlik, yüklemler mantığının özel bir ilişkisidir. Kuralları şöyledir:

\begin{center}
\begin{tabular}{c@{\hspace{5em}}c}
\begin{prooftree}
  \infer0[$=$G]{t = t}
\end{prooftree}
&
\begin{prooftree}
  \hypo{s = t}
  \hypo{P(s)}
  \infer2[$=$E]{P(t)}
\end{prooftree}
\\[10pt]
\small\textcolor{ikincil}{\bfseries giriş (yansıma):} \lc{rfl} &
\small\textcolor{vurgu}{\bfseries eleme (yerine koyma):} \leantac{rw}, \lc{h ▸ hp}
\end{tabular}
\end{center}

\noindent Yalnızca bu iki kuraldan eşitliğin bütün bilinen özellikleri
türer: simetri, geçişlilik, fonksiyonlarla uyum.

\leandosya[Simetri ve geçişlilik]{Bolum06}{esitlik_temel}
\leandosya[Yeniden yazma]{Bolum06}{esitlik_rw}

\noindent \leantac{rw [h]}, \lc{h : a = b} eşitliğini kullanarak hedefteki
her \lc{a}'yı \lc{b} yapar. Yeniden yazdıktan sonra hedef \lc{c + 1 = c + 1}
gibi apaçık bir eşitliğe dönüşürse, \leantac{rw} onu \lc{rfl} ile
kendiliğinden kapatır. Ters yönde yazmak için \lc{rw [← h]} kullanılır.

\leandosya[Yerine koyma]{Bolum06}{esitlik_subst}
\leandosya[Fonksiyonlarla uyum]{Bolum06}{esitlik_congr}

\paragraph{calc: eşitlik zincirleri}
Matematikte uzun hesapları bir eşitlik zinciri olarak yazarız:
$a = b = c = d$. \leankw{calc} bloğu bunun \Lean{}'deki karşılığıdır.
Her satır bir adımı ve gerekçesini verir; alt çizgi bir önceki satırın
sağ tarafını gösterir:

\leandosya[Eşitlik zinciri]{Bolum06}{calc_ornek}

\noindent \leankw{calc}, \cref{bol:tumevarim}'de aritmetik ispatlarda ve
\lc{<}, \lc{≤} gibi ilişkilerin zincirlerinde çok işimize yarayacak.

\paragraph{Leibniz ilkesi}
Gottfried Leibniz'e göre iki şey, bütün özellikleri ortaksa aynıdır.
Bu ilke, eşitliğin $=$E kuralının tersidir ve \Lean{}'de tek satırla
ispatlanır:

\leandosya[Leibniz ilkesi]{Bolum06}{leibniz}

\noindent İspattaki hile, ``$a$'ya eşit olmak'' özelliğini seçmektir:
$a$ bu özelliği \lc{rfl} ile taşır, o hâlde $b$ de taşır, yani $a = b$.

% ---------------------------------------------------------------------
\section{Berber paradoksu}\label{sec:berber}

\begin{gunlukornek}[Bir köyün berberi]
Bir köyde bir berber, ``kendini tıraş etmeyen herkesi tıraş eder,
kendini tıraş edeni ise tıraş etmez.'' Berber kendini tıraş eder mi?
\begin{itemize}
  \item Ederse, kendini tıraş eden biri olur; o hâlde onu tıraş etmez.
  \item Etmezse, kendini tıraş etmeyen biri olur; o hâlde onu tıraş eder.
\end{itemize}
Her iki durumda da çelişki çıkar. Sonuç: Böyle bir berber \emph{yoktur}.
\end{gunlukornek}

\begin{figure}[htbp]
\centering
\begin{tikzpicture}[
    kisi/.style={circle, draw=anarenk, fill=tanimbg, thick, minimum size=0.9cm,
                 font=\small},
    ok/.style={-{Stealth[length=2.4mm]}, thick, #1},
  ]
  \node[kisi, draw=vurgu, fill=teorembg, minimum size=1.1cm] (b) at (0,0) {$b$};
  \node[kisi] (a) at (-2.6,1.2) {$a$};
  \node[kisi] (c) at (2.6,1.2) {$c$};
  \node[kisi] (d) at (-2.6,-1.2) {$d$};
  \node[kisi] (e) at (2.6,-1.2) {$e$};
  % kendini tıraş edenler: a, c (döngü)
  \draw[ok=black!45] (a) to[loop above, looseness=6] (a);
  \draw[ok=black!45] (c) to[loop above, looseness=6] (c);
  % berber d ve e'yi tıraş eder
  \draw[ok=vurgu] (b) -- (d);
  \draw[ok=vurgu] (b) -- (e);
  \draw[ok=vurgu, dashed] (b) to[loop below, looseness=6] node[below, font=\baslikfont\large\bfseries, text=vurgu]{?} (b);
  \node[font=\baslikfont\scriptsize, text=black!60, align=left, anchor=west] at (4.0,0.6)
       {$a$, $c$: kendini tıraş eder\\$\Rightarrow$ berber onları tıraş etmez};
  \node[font=\baslikfont\scriptsize, text=vurgu, align=left, anchor=west] at (4.0,-0.6)
       {$d$, $e$: kendini tıraş etmez\\$\Rightarrow$ berber onları tıraş eder};
\end{tikzpicture}
\caption{Berber $b$'nin kendine giden oku çizilmeli midir? Çizilse de
çizilmese de kural bozulur.}
\label{fig:berber}
\end{figure}

\leandosya[Berber paradoksu]{Bolum06}{berber}

\noindent İspatın çekirdeği Alıştırma~3.8 ile aynıdır: Hiçbir önerme
kendi değillemesine denk olamaz. \lc{hb b} bize
\lc{TirasEder b b ↔ ¬ TirasEder b b} denkliğini verir; geri kalanı
önermeler mantığıdır. Dikkat edin, ispat \emph{yapıcıdır}: Paradoks,
klasik mantığa hiç ihtiyaç duymaz.

\begin{tarihce}[Tarihçe: Russell'ın mektubu]
Berber hikâyesi, Bertrand Russell'ın 1901'de bulduğu ve 1902'de Frege'ye
bir mektupla bildirdiği paradoksun popüler bir anlatımıdır. Russell'ın
asıl sorusu şuydu: ``Kendini eleman olarak içermeyen bütün kümelerin
kümesi'' kendini içerir mi? Frege'nin aritmetiği mantıktan türetme
programı, sınırsız küme oluşturmaya izin veren bir ilkeye dayandığı için
bu paradoksla çöktü \citep{vanheijenoort1967}. Paradoksa verilen
cevaplardan biri Russell'ın \emph{tip kuramıydı}: Bir kümenin elemanları
ile kümenin kendisi farklı ``tiplerde'' olmalıdır. \Lean{}'in tip
sistemi, bu fikrin yüz yıllık evriminin bir ürünüdür.
\end{tarihce}

% ---------------------------------------------------------------------
\section{Boş evren inceliği}\label{sec:bos-evren}

Klasik mantık kitaplarının çoğunda şu kural geçerlidir:
$\forall x\,P(x) \to \exists x\,P(x)$. ``Herkes $P$ ise, $P$ olan biri
vardır.'' Kulağa apaçık geliyor, ama evrende \emph{hiç kimse yoksa}
ne olur?

\leandosya[Boş ve boş olmayan evrenler]{Bolum06}{bos_evren}

\noindent Boş tip \lc{Empty} üzerinde $\forall x\,P(x)$ \emph{boş yere
doğrudur}: Kontrol edilecek hiç eleman yoktur. Ama $\exists x\,P(x)$
yanlıştır: Tanık bulunamaz. Dolayısıyla $\forall \to \exists$ geçişi,
yalnızca evren boş değilse geçerlidir. \Lean{}'de bu varsayım
\lc{[Inhabited α]} ile açıkça yazılır.

\begin{dikkat}[Dikkat: Geleneksel mantık evreni boş olmayan kabul eder]
Birinci mertebe mantığın geleneksel sunumlarında tanım evreninin boş
olmadığı \emph{sessizce} varsayılır. \Lean{}'de ise tipler boş olabilir
ve bu varsayım her seferinde açıkça yapılmalıdır. Bir kitaptaki ispatı
\Lean{}'e aktarırken \lc{Inhabited} ya da \lc{Nonempty} varsayımının
eksik olduğu yerler, çoğu zaman bu inceliğin gözden kaçtığı yerlerdir.
\end{dikkat}

% ---------------------------------------------------------------------
\section{Curry--Howard: bağımlı tipler}\label{sec:curry-howard-3}

Niceleyiciler, Curry--Howard tablosuna iki yeni satır ekler. Bunlar
tablonun en derin satırlarıdır, çünkü \emph{bağımlı tiplere} karşılık
gelirler (\cref{tab:curry-howard-3}).

\begin{table}[htbp]
\centering
\renewcommand{\arraystretch}{1.55}
\rowcolors{2}{tablosatir}{white}
\begin{tabular}{>{\raggedright\arraybackslash}p{2.6cm}
                >{\raggedright\arraybackslash}p{4.6cm}
                >{\raggedright\arraybackslash}p{4.4cm}}
\rowcolor{tabloust}
\tablobaslik{Mantık} & \tablobaslik{Tip kuramı} & \tablobaslik{Lean'deki kanıt} \\
$P \to Q$                & fonksiyon tipi $P \to Q$                      & \lc{fun hp => ...} \\
$\forall x{:}\alpha,\; P(x)$ & \textbf{bağımlı fonksiyon} tipi $\Pi_{x:\alpha} P(x)$ & \lc{fun x => ...} \\
$P \wedge Q$             & çarpım tipi $P \times Q$                      & \lc{⟨hp, hq⟩} \\
$\exists x{:}\alpha,\; P(x)$ & \textbf{bağımlı çift} tipi $\Sigma_{x:\alpha} P(x)$ & \lc{⟨t, ht⟩} \\
$t = s$                  & eşitlik tipi                                   & \lc{rfl} \\
\end{tabular}
\caption{Curry--Howard eşlemesinin niceleyicilere genişlemesi. $\forall$,
$\to$'nun; $\exists$ ise $\wedge$'nin bağımlı biçimidir.}
\label{tab:curry-howard-3}
\end{table}

\begin{figure}[htbp]
\centering
\begin{tikzpicture}[
    lif/.style={draw=#1, fill=#1!8, thick, rounded corners=2pt,
                minimum width=0.8cm, minimum height=1.7cm},
  ]
  % Π
  \node[font=\baslikfont\bfseries\small, text=ikincil] at (2.2,2.75) {$\Pi_{x} P(x)$: her liften bir seçim};
  \foreach \i in {0,...,4} {
    \node[lif=ikincil] (l\i) at ({1.1*\i},0.9) {};
    \node[font=\scriptsize, text=ikincil] at ({1.1*\i},2.0) {$P(\i)$};
    \node[font=\scriptsize] at ({1.1*\i},-0.2) {$\i$};
  }
  \draw[ikincil, very thick] (0,0.6) -- (1.1,1.25) -- (2.2,0.75) -- (3.3,1.3) -- (4.4,0.8);
  \foreach \i/\y in {0/0.6, 1/1.25, 2/0.75, 3/1.3, 4/0.8} \fill[ikincil] ({1.1*\i},\y) circle (2.6pt);
  \draw[black!40] (-0.5,0.0) -- (4.9,0.0);
  \node[font=\baslikfont\scriptsize, text=black!60] at (2.2,-0.65) {fonksiyon: her $x$ için bir kanıt};
  % Σ
  \node[font=\baslikfont\bfseries\small, text=vurgu] at (9.0,2.75) {$\Sigma_{x} P(x)$: bir lif ve ondan bir nokta};
  \foreach \i in {0,...,4} {
    \node[lif=vurgu, opacity=0.55] (m\i) at ({6.8+1.1*\i},0.9) {};
    \node[font=\scriptsize, text=vurgu] at ({6.8+1.1*\i},2.0) {$P(\i)$};
    \node[font=\scriptsize] at ({6.8+1.1*\i},-0.2) {$\i$};
  }
  \node[lif=vurgu, very thick] at (9.0,0.9) {};
  \fill[vurgu] (9.0,1.05) circle (3pt);
  \node[font=\baslikfont\scriptsize\bfseries, text=vurgu, anchor=west] (et) at (12.0,1.55) {$\langle 2, h\rangle$};
  \draw[vurgu, thin, -{Stealth[length=1.6mm]}] (et.west) to[bend right=15] (9.08,1.1);
  \draw[black!40] (6.3,0.0) -- (11.7,0.0);
  \node[font=\baslikfont\scriptsize, text=black!60] at (9.0,-0.65) {çift: bir tanık ve onun lifinde bir kanıt};
\end{tikzpicture}
\caption{Bağımlı tiplerin geometrik resmi. Her $x$ üzerinde bir ``lif''
$P(x)$ durur. $\Pi$ her liften bir eleman seçer; $\Sigma$ ise bir lif
ve o liften bir eleman seçer.}
\label{fig:pi-sigma}
\end{figure}

\begin{onemlifikir}[Önemli Fikir: Bağımlı tipler mantığın tamamını taşır]
Bölüm~2'de \lc{→}, \lc{∧}, \lc{∨} için başlattığımız tablo, bağımlı
tiplerle yüklemler mantığının tamamını kapsar. \Lean{}'in mantığı,
özünde \emph{bağımlı tip kuramıdır}: Yeni bir mantık kuralı
eklememize gerek kalmadan, yalnızca tipler ve fonksiyonlarla bütün
matematiği ifade edebiliriz. Bu kuramın kesin tanımını
\cref{bol:bagimli}'te vereceğiz \citep{martinlof1984, coquand1988}.
\end{onemlifikir}

% ---------------------------------------------------------------------
\section{Taktik özeti}

\begin{center}
\renewcommand{\arraystretch}{1.45}
\rowcolors{2}{tablosatir}{white}
\begin{tabularx}{\linewidth}{>{\raggedright\arraybackslash}p{3.6cm}
                             >{\raggedright\arraybackslash}p{1.5cm}
                             >{\raggedright\arraybackslash}X}
\rowcolor{tabloust}
\tablobaslik{Taktik / terim} & \tablobaslik{Kural} & \tablobaslik{Ne yapar?} \\
\leantac{intro x}                   & $\forall$G & Keyfi bir \lc{x} alır; hedef \lc{P x} olur. \\
\lc{h t}, \leantac{specialize h t}  & $\forall$E & Genel hipotezi \lc{t}'ye uygular. \\
\lc{⟨t, ht⟩}, \leantac{refine ⟨t, ?\_⟩} & $\exists$G & Tanığı \lc{t} olan bir varlık kanıtı kurar. \\
\leantac{obtain ⟨a, ha⟩ := h}       & $\exists$E & Tanığa ad verir: \lc{a} ve \lc{ha : P a}. \\
\lc{rfl}                            & $=$G       & \lc{t = t} hedefini kapatır. \\
\leantac{rw [h]}, \leantac{rw [← h]} & $=$E      & Hedefte \lc{a}'yı \lc{b}'ye (ya da tersine) çevirir. \\
\lc{h ▸ hp}                         & $=$E       & \lc{P a}'yı \lc{P b}'ye taşır. \\
\lc{h.symm}, \lc{h₁.trans h₂}       & ---        & Simetri, geçişlilik. \\
\lc{congrArg f h}                   & ---        & \lc{a = b}'den \lc{f a = f b}. \\
\leankw{calc}                       & ---        & Eşitlik ya da eşitsizlik zinciri. \\
\end{tabularx}
\end{center}

% ---------------------------------------------------------------------
\begin{arastirmanotu}
\textbf{Tanıklar Prop'tan kaçamaz.} \Lean{}'de \lc{∃ x, P x} bir
\lc{Prop}'tur ve \lc{Prop} kanıt-ilgisizdir. Bu yüzden bir varlık
kanıtından tanığı \emph{bir program olarak} çıkaramayız: \leantac{obtain}
yalnızca hedef de bir \lc{Prop} ise çalışır. Tanığı veri olarak
istediğimizde ya \lc{Type} düzeyindeki bağımlı çifti (\lc{Σ x, P x} ya da
\lc{\{x // P x\}}) kullanırız, ya da hesaplanamayan
\lc{Classical.choose}'a başvururuz. Bu tasarım, \cref{bol:yapici-klasik}'teki
``klasik akıl yürütme hesaplamayı bozmaz'' dengesinin bedelidir
\citep{carneiro2019}.

\textbf{Hilbert'in $\varepsilon$ operatörü.} Hilbert ve Ackermann'ın
mantık sisteminde $\varepsilon x.\,P(x)$ terimi, $P$'yi sağlayan
\emph{bir} nesneyi (varsa) gösterir \citep{hilbert1928}.
\lc{Classical.choose}, bu fikrin tip kuramındaki karşılığıdır.

\textbf{İleri okuma.} Birinci mertebe mantığın klasik bir sunumu için
\citet{vandalen2013}; içki içen paradoksu ve benzeri mantık bulmacaları
için \citet{smullyan1978}; tip kuramı açısından niceleyiciler için
\citet{nederpelt2014}.
\end{arastirmanotu}

\begin{bolumozeti}
\begin{itemize}
  \item Bir yüklem \lc{α → Prop} türünde bir fonksiyondur.
  \item $\forall$G: keyfi bir \lc{x} al (\leantac{intro}); $\forall$E:
        genel kanıtı bir bireye uygula. $\forall$'nın kanıtı bağımlı bir
        fonksiyondur.
  \item $\exists$G: bir tanık ve kanıtını ver (\lc{⟨t, ht⟩}); $\exists$E:
        tanığa ad ver (\leantac{obtain}). $\exists$'nin kanıtı bağımlı
        bir çifttir.
  \item $\neg\exists \leftrightarrow \forall\neg$ yapıcıdır;
        $\neg\forall \to \exists\neg$ klasiktir. $\exists\forall \to
        \forall\exists$ geçerlidir, tersi değildir.
  \item Yerine koyarken bağlı değişkenler yakalanmamalıdır.
  \item Eşitlik: \lc{rfl}, \leantac{rw}, \lc{▸}, \lc{calc}. Leibniz ilkesi
        \Lean{}'de bir teoremdir.
  \item Berber paradoksu yapıcı olarak ispatlanır. $\forall \to \exists$
        yalnızca boş olmayan evrenlerde geçerlidir.
\end{itemize}
\end{bolumozeti}

\begin{alistirmalar}
\alistirma{1} \lc{(∀ x, P x ∧ Q x) → ∀ x, P x}. (\lc{a6\_1})

\alistirma{1} \lc{(∀ x, P x) → (∀ x, Q x) → ∀ x, P x ∧ Q x}. (\lc{a6\_2})

\alistirma{1} \lc{(∃ x, P x ∧ Q x) → ∃ x, P x}. (\lc{a6\_3})

\alistirma{2} $\forall$'nın $\wedge$ üzerine dağıldığını ispatlayın:
\lc{(∀ x, P x ∧ Q x) ↔ (∀ x, P x) ∧ (∀ x, Q x)}. (\lc{a6\_4})

\alistirma{2} $\exists$'nin $\vee$ üzerine dağıldığını ispatlayın:
\lc{(∃ x, P x ∨ Q x) ↔ (∃ x, P x) ∨ (∃ x, Q x)}. (\lc{a6\_5})

\alistirma{2} \lc{(∃ x, P x ∧ Q x) → (∃ x, P x) ∧ (∃ x, Q x)}. Tersinin
neden yanlış olduğunu doğal sayılar üzerinde bir karşı örnekle
gösterin. (\lc{a6\_6})

\alistirma{2} \lc{(∀ x, P x → Q x) → (∃ x, P x) → ∃ x, Q x}. (\lc{a6\_7})

\alistirma{2} \lc{a = b} ve \lc{c = b} ise \lc{a = c}. Önce terim modunda
(\lc{symm}, \lc{trans}), sonra \leantac{rw} ile ispatlayın. (\lc{a6\_8})

\alistirma{2} \lc{a = b → f (f a) = f (f b)}. (\lc{a6\_9})

\alistirma{3} (Klasik) \lc{(¬ ∀ x, P x) ↔ ∃ x, ¬ P x}. Hangi yön klasik
mantık gerektiriyor? (\lc{a6\_10})

\alistirma{3} (Klasik) \textbf{İçki içen paradoksu}: \lc{[Inhabited α]}
varsayımıyla \lc{∃ x, P x → ∀ y, P y}. ``Bir barda öyle biri vardır ki,
o içiyorsa herkes içiyordur.'' Bu önerme neden paradoks gibi görünür?
\lc{Inhabited} varsayımı neden gereklidir? (\lc{a6\_11})

\alistirma{3} Berber paradoksunun genel hâli:
\lc{¬ ∃ x, ∀ y, R x y ↔ ¬ R y y}. Bu teoremden berber paradoksunu
tek satırda elde edin. (\lc{a6\_12})

\alistirma{2} (Kâğıt-kalem) Aşağıdaki önermelerde serbest ve bağlı
değişkenleri belirleyin; ardından $x$ yerine $y + 1$ terimini
yakalamadan koyun:
(a) $\forall y,\; x < y$; \quad (b) $(\exists x,\; P(x)) \wedge Q(x)$.

\alistirma{3} (Kâğıt-kalem) $\forall$G kuralındaki öz değişken koşulu
çiğnenirse $P(a) \to \forall x\,P(x)$ ``ispatlanabilir''. Bu yanlış
türetimi yazın ve hangi adımın kuralı çiğnediğini gösterin.
\end{alistirmalar}
   % Yüklemler mantığı: ∀, ∃, =
% =====================================================================
%  Bölüm 7 — Doğal Sayılar ve Tümevarım
% =====================================================================
\chapter{Doğal Sayılar ve Tümevarım}\label{bol:tumevarim}

\begin{bolumhedefleri}
\begin{itemize}
  \item Peano aksiyomlarını ve onların \Lean{}'de neden \emph{teorem}
        olduğunu,
  \item Doğal sayıları kendi tümevarımsal tipimizle sıfırdan tanımlamayı,
  \item Toplama ve çarpmayı \emph{özyinelemeyle} tanımlamayı ve bu
        tanımların nasıl hesaplandığını,
  \item Tümevarım ilkesini, \leantac{induction} taktiğini ve tümevarım
        hipotezini kullanmayı,
  \item Toplamanın değişme ve birleşme özelliklerini ilk ilkelerden
        ispatlamayı,
  \item Gauss toplam formülünü, güçlü tümevarımı ve pul problemini,
  \item \leantac{simp}, \leantac{omega} ve \leantac{decide} otomasyonunu
        ve sınırlarını.
\end{itemize}
\end{bolumhedefleri}

Şimdiye kadar mantığın \emph{genel} yasalarıyla uğraştık: Her $P$, $Q$
ya da her yüklem için geçerli olan yasalar. Bu bölümde ilk kez belirli bir
matematiksel yapıyı, doğal sayıları ele alıyoruz. Doğal sayılar, mantıkla
matematiğin buluştuğu ilk yerdir. Burada öğreneceğimiz iki fikir,
\emph{özyinelemeli tanım} ve \emph{tümevarımla ispat}, kitabın geri
kalanının neredeyse her sayfasında karşımıza çıkacak.

% ---------------------------------------------------------------------
\section{Peano aksiyomları}\label{sec:peano}

Doğal sayılar nedir? Okulda $0, 1, 2, 3, \ldots$ diye sayarız; ama bu
``$\ldots$'' ne anlama gelir? On dokuzuncu yüzyılın sonunda Richard
Dedekind ve Giuseppe Peano bu soruya kesin bir cevap verdiler
\citep{dedekind1888, peano1889}. Doğal sayıların bütün özellikleri,
birkaç basit ilkeden türetilebilir (\cref{tab:peano}).

\begin{table}[htbp]
\centering
\renewcommand{\arraystretch}{1.5}
\rowcolors{2}{tablosatir}{white}
\begin{tabular}{>{\centering\arraybackslash}p{0.6cm}
                >{\raggedright\arraybackslash}p{6.2cm}
                >{\raggedright\arraybackslash}p{4.8cm}}
\rowcolor{tabloust}
\tablobaslik{} & \tablobaslik{Aksiyom} & \tablobaslik{Lean'de} \\
P1 & $0$ bir doğal sayıdır. & kurucu \lc{sifir} \\
P2 & Her $n$ doğal sayısının bir \emph{ardılı} $S(n)$ vardır. & kurucu \lc{ard} \\
P3 & $0$ hiçbir sayının ardılı değildir: $S(n) \neq 0$. & teorem \lc{sifir\_ard\_degil} \\
P4 & Ardıl fonksiyonu birebirdir: $S(n) = S(m) \Rightarrow n = m$. & teorem \lc{ard\_birebir} \\
P5 & \textbf{Tümevarım}: $0$'ı içeren ve ardıl altında kapalı olan her
     özellik, bütün doğal sayılarda geçerlidir. & teorem \lc{tumevarim} \\
\end{tabular}
\caption{Peano aksiyomları ve \Lean{}'deki karşılıkları. İlk ikisi
tanımın kendisidir; geri kalan üçü, tanımdan \emph{türetilen}
teoremlerdir.}
\label{tab:peano}
\end{table}

\begin{tarihce}[Tarihçe: Dedekind ve Peano]
Dedekind 1888'de yayımladığı \emph{Was sind und was sollen die Zahlen?}
(Sayılar nedir ve ne olmalıdır?) adlı kitabında doğal sayıları
kümeler kuramı içinde tanımladı ve özyinelemeli tanımların neden
anlamlı olduğunu ispatladı \citep{dedekind1888}. Peano bir yıl sonra
aynı fikirleri, kendi geliştirdiği sembolik dille ve aksiyomlar hâlinde
sundu \citep{peano1889}. Bugün kullandığımız $\in$ sembolü ve pek çok
mantık gösterimi Peano'dan gelir. Aksiyomlara Peano'nun adı verilir,
ama Peano kendisi Dedekind'in önceliğini kabul etmiştir.
\end{tarihce}

% ---------------------------------------------------------------------
\section{Doğal sayıları sıfırdan tanımlamak}\label{sec:dogal-tanim}

\Lean{}'in kendi doğal sayıları (\lc{Nat}) vardır; ama ne olduklarını
anlamanın en iyi yolu, onları kendimiz tanımlamaktır:

\leandosya[Kendi doğal sayılarımız]{Bolum07}{dogal_tanim}

\noindent \leankw{inductive} anahtar kelimesi yeni bir tip tanımlar ve
bu tipin elemanlarının \emph{nasıl kurulabileceğini} söyler. \lc{Dogal}
tipinin iki kurucusu vardır: \lc{sifir} ve \lc{ard}. Bu tipin her elemanı
bu iki kurucunun sonlu sayıda uygulanmasıyla elde edilir; başka hiçbir
eleman yoktur (\cref{fig:dogal-zincir}).

\begin{figure}[htbp]
\centering
\begin{tikzpicture}[
    sayi/.style={draw=anarenk, fill=tanimbg, thick, rounded corners=4pt,
                 font=\leanfont\scriptsize, minimum height=0.8cm, inner sep=4pt},
    ok/.style={-{Stealth[length=2.4mm]}, thick, leankw},
  ]
  \node[sayi] (s0) at (0,0) {sifir};
  \node[sayi] (s1) at (2.5,0) {ard sifir};
  \node[sayi] (s2) at (5.6,0) {ard (ard sifir)};
  \node[sayi] (s3) at (9.6,0) {ard (ard (ard sifir))};
  \node (s4) at (12.5,0) {$\cdots$};
  \foreach \a/\b in {s0/s1, s1/s2, s2/s3, s3/s4}
    \draw[ok] (\a) -- node[above, font=\leanfont\scriptsize\bfseries, text=leankw]{ard} (\b);
  \foreach \n/\v in {s0/0, s1/1, s2/2, s3/3}
    \node[below=4pt of \n, font=\baslikfont\small\bfseries, text=anarenk] {$\v$};
\end{tikzpicture}
\caption{\lc{Dogal} tipinin elemanları. Her sayı, \lc{sifir}'dan başlayıp
\lc{ard}'ı yeterince kez uygulayarak elde edilir.}
\label{fig:dogal-zincir}
\end{figure}

Tümevarımsal bir tanımın iki önemli özelliği vardır:
\begin{enumerate}[label=(\roman*)]
  \item \textbf{Kurucular birbirinden farklıdır}: \lc{sifir} ile
        \lc{ard n} asla eşit olamaz.
  \item \textbf{Kurucular birebirdir}: \lc{ard n = ard m} ise \lc{n = m}.
\end{enumerate}

\noindent Bu iki özellik \Lean{}'in tip kuramına yerleşiktir. Dolayısıyla
Peano'nun P3 ve P4 aksiyomları, bizim için \emph{ispatlanabilir
teoremlerdir}:

\leandosya[Peano aksiyomları teorem olarak]{Bolum07}{peano_teorem}

\begin{onemlifikir}[Önemli Fikir: Aksiyom değil, tanım]
Klasik matematikte Peano aksiyomları \emph{varsayılır}. Tip kuramında
ise tek yaptığımız \lc{Dogal} tipini tanımlamaktır; aksiyomların hepsi
bu tanımdan teorem olarak çıkar. Yeni bir varsayım eklemediğimiz için
bu yaklaşım, sistemin tutarlılığı açısından daha güvenlidir. Bu yöntemin
genel hâlini, yani her türlü veri yapısını \leankw{inductive} ile
tanımlamayı, \cref{bol:tumevarimsal-tipler}'da göreceğiz.
\end{onemlifikir}

% ---------------------------------------------------------------------
\section{Özyineleme: toplama ve çarpma}\label{sec:ozyineleme}

Doğal sayılar üzerinde fonksiyonları \emph{özyinelemeyle} tanımlarız:
Fonksiyonun \lc{sifir}'daki değerini ve \lc{ard n}'deki değerini
\lc{n}'deki değeri cinsinden veririz.

\leandosya[Toplama]{Bolum07}{topla_tanim}

\noindent Bu tanım, ikinci argüman üzerinde özyinelemelidir. ``$n + 0 = n$''
ve ``$n + (m+1) = (n+m)+1$'' kurallarını söyler. Bir toplamın nasıl
hesaplandığını adım adım izleyelim (\cref{fig:hesaplama-izi}).

\begin{figure}[htbp]
\centering
\begin{tikzpicture}[
    adim/.style={font=\leanfont\small, anchor=west},
    kural/.style={font=\baslikfont\scriptsize, text=leantac, anchor=west},
  ]
  \node[adim] at (0,0)    {topla iki (ard (ard sifir))};
  \node[adim] at (0,-0.75) {= ard (topla iki (ard sifir))};
  \node[adim] at (0,-1.5) {= ard (ard (topla iki sifir))};
  \node[adim] at (0,-2.25) {= ard (ard iki)};
  \node[adim] at (0,-3.0) {= ard (ard (ard (ard sifir)))};
  \node[kural] at (7.4,-0.75) {2. kural};
  \node[kural] at (7.4,-1.5)  {2. kural};
  \node[kural] at (7.4,-2.25) {1. kural};
  \node[kural] at (7.4,-3.0)  {\lc{iki}'nin tanımı};
  \node[font=\baslikfont\bfseries\small, text=anarenk, anchor=west] at (9.4,-1.5) {$2 + 2 = 4$};
  \draw[anarenk, thick, decorate, decoration={brace, amplitude=6pt}]
       (9.2,0.2) -- (9.2,-3.2);
\end{tikzpicture}
\caption{$2 + 2$'nin hesaplanması. Her adımda tanımdaki iki kuraldan biri
uygulanır; ikinci argüman her adımda küçülür ve hesaplama sona erer.}
\label{fig:hesaplama-izi}
\end{figure}

Çarpma da toplamaya dayanarak aynı biçimde tanımlanır:

\leandosya[Çarpma]{Bolum07}{carp_tanim}

\noindent Tanımlarımızı denemek için, kendi sayılarımızı \Lean{}'in
\lc{Nat}'ına çeviren bir fonksiyon yazıp \lc{\textcolor{leancmd}{\bfseries\#eval}}
ile hesaplatabiliriz:

\leandosya[Hesaplama]{Bolum07}{sayiya}

\begin{dikkat}[Dikkat: Özyineleme her zaman sona ermelidir]
\Lean{}, özyinelemeli her tanımın sona erdiğini denetler. Yukarıdaki
tanımlarda her özyinelemeli çağrı \emph{daha küçük} bir argümanla
(\lc{ard m} yerine \lc{m} ile) yapıldığı için bu denetim kendiliğinden
geçer. Buna \emph{yapısal özyineleme} denir. Sona ermeyen bir tanım
(örneğin \lc{f n := f n + 1}) kabul edilseydi, \lc{f 0 = f 0 + 1}
eşitliğinden \lc{False} türetilebilirdi. Sona erme denetimi, mantığın
tutarlılığının bir parçasıdır.
\end{dikkat}

% ---------------------------------------------------------------------
\section{Tümevarım ilkesi}\label{sec:tumevarim-ilkesi}

\begin{center}
\begin{prooftree}
  \hypo{P(0)}
  \hypo{[P(n)]^{1}}
  \ellipsis{}{P(n+1)}
  \infer2[Tümevarım$^{1,*}$]{\forall n,\; P(n)}
\end{prooftree}
\qquad
{\small \leantac{induction n with \ldots}}
\end{center}

\noindent {\small $^{*}$\,$n$, sağdaki türetimin diğer açık
varsayımlarında geçmemelidir.}

Tümevarım ilkesi, dominoların düşmesine benzetilir
(\cref{fig:domino}). Birinci domino düşüyorsa (\emph{taban adımı}) ve
her domino düştüğünde bir sonrakini de düşürüyorsa (\emph{tümevarım
adımı}), bütün dominolar düşer.

\begin{figure}[htbp]
\centering
\begin{tikzpicture}
  \draw[black!30, thick] (-0.9,0) -- (11.4,0);
  % düşmüş dominolar (sağ alt köşeleri etrafında döndürülmüş)
  \foreach \i in {0,1,2,3} {
    \draw[bittiframe, thick, fill=bittibg, rotate around={-58:({1.25*\i+0.2},0)}]
         ({1.25*\i-0.2},0) rectangle ({1.25*\i+0.2},1.5);
    \node[font=\scriptsize, text=bittiframe] at ({1.25*\i},-0.35) {$P(\i)$};
  }
  % ayakta olanlar
  \foreach \i in {4,5,6,7} {
    \draw[anarenk, thick, fill=tanimbg] ({1.25*\i-0.2},0) rectangle ({1.25*\i+0.2},1.5);
    \node[font=\scriptsize, text=anarenk] at ({1.25*\i},-0.35) {$P(\i)$};
  }
  \node at (10.6,0.7) {$\cdots$};
  \draw[-{Stealth[length=2.6mm]}, very thick, vurgu] (-0.8,2.0) -- (-0.15,1.2);
  \node[font=\baslikfont\scriptsize\bfseries, text=vurgu, anchor=south] at (-0.8,2.0) {taban: $P(0)$};
  \draw[-{Stealth[length=2.6mm]}, very thick, leantac] (4.0,2.3) to[bend left=20] (4.75,1.75);
  \node[font=\baslikfont\scriptsize\bfseries, text=leantac, anchor=south] at (4.0,2.3) {adım: $P(n) \to P(n+1)$};
\end{tikzpicture}
\caption{Tümevarım: İlk domino düşer ve her düşen domino bir sonrakini
düşürür; o hâlde hepsi düşer.}
\label{fig:domino}
\end{figure}

\Lean{}'de tümevarım ilkesi de bir \emph{teoremdir}. \leankw{inductive}
ile her tip tanımladığımızda, \Lean{} o tip için bir \emph{özyineleyici}
(İng.\ \emph{recursor}) üretir:

\leandosya[Özyineleyici]{Bolum07}{rec_check}

\noindent Bu tip, tam olarak tümevarım ilkesidir: \lc{motive} ispatlamak
istediğimiz özellik, ikinci argüman taban adımı, üçüncü argüman
tümevarım adımıdır. Aynı özyineleyici hem \emph{fonksiyon tanımlamak}
(\lc{motive} bir \lc{Type} ise) hem de \emph{ispat yapmak} (\lc{motive}
bir \lc{Prop} ise) için kullanılır. Curry--Howard'ın bir başka
yüzü: \textbf{tümevarımla ispat, özyinelemeli bir fonksiyondur}.

\leandosya[Tümevarım ilkesi]{Bolum07}{tumevarim}

\noindent \leantac{induction n with} satırı iki hedef açar. İkinci
hedefte bağlama \emph{tümevarım hipotezi} \lc{ih : P m} eklenir:

\begin{hedefgecisi}
\begin{hedefdurumu}[\texttt{sifir} durumu]
P : Dogal → Prop
h0 : P sifir
hs : ∀ (n : Dogal), P n → P n.ard
⊢ P sifir
\end{hedefdurumu}
\sonra
\begin{hedefdurumu}[\texttt{ard} durumu]
P : Dogal → Prop
h0 : P sifir
hs : ∀ (n : Dogal), P n → P n.ard
m : Dogal
ih : P m
⊢ P m.ard
\end{hedefdurumu}
\end{hedefgecisi}

% ---------------------------------------------------------------------
\section{İlk tümevarım ispatları}\label{sec:ilk-tumevarim}

\subsection{Bir asimetri}

Toplamanın tanımı ``$n + 0 = n$'' der. Bu eşitlik tanım gereği doğrudur;
\lc{rfl} yeter:

\leandosya[$n + 0 = n$]{Bolum07}{topla_sifir}

\noindent Ya ``$0 + n = n$''? Matematik açısından aynı derecede apaçık
görünür, ama \Lean{} için hiç de öyle değildir: Tanım ikinci argümana
bakar ve \lc{n} bilinmeyen bir sayı olduğu için hiçbir kural
uygulanamaz. Bu eşitlik ancak \emph{tümevarımla} ispatlanabilir:

\leandosya[$0 + n = n$]{Bolum07}{sifir_topla}

\noindent \leantac{show} adımı, hedefteki \lc{topla sifir (ard m)}
ifadesini tanımın ikinci kuralıyla bir adım açar. Ardından tümevarım
hipotezi \lc{ih : topla sifir m = m} ile yeniden yazmak yeterlidir.

\begin{onemlifikir}[Önemli Fikir: Tanımsal ve önermesel eşitlik]
\lc{topla n sifir = n} \emph{tanımsal} bir eşitliktir: \Lean{} iki tarafı
hesaplayarak aynı olduklarını görür ve \lc{rfl} kabul edilir.
\lc{topla sifir n = n} ise \emph{önermesel} bir eşitliktir: Doğrudur,
ama bir ispat gerektirir. Bu ayrım tip kuramının en önemli
kavramlarından biridir ve \cref{bol:bagimli}'te ayrıntılı olarak
incelenecek.
\end{onemlifikir}

Aynı ispatı terim modunda, \emph{özyinelemeli bir fonksiyon} olarak da
yazabiliriz. Kendini daha küçük bir argümanla çağıran bu ``ispat'',
tümevarımın kendisidir:

\leandosya[$0 + n = n$, özyinelemeli ispat]{Bolum07}{sifir_topla_terim}

\subsection{Değişme ve birleşme}

Toplamanın değişme özelliğini ($n + m = m + n$) ispatlamak için önce bir
yardımcı önsav gerekir:

\leandosya[Ardıl ile toplama]{Bolum07}{ard_topla}
\leandosya[Değişme özelliği]{Bolum07}{topla_degisme}
\leandosya[Birleşme özelliği]{Bolum07}{topla_birlesme}

\begin{figure}[htbp]
\centering
\begin{tikzpicture}[
    onsav/.style={draw=#1, fill=#1!8, thick, rounded corners=4pt,
                  font=\leanfont\scriptsize, align=center, inner sep=5pt,
                  minimum width=2.8cm},
    ok/.style={-{Stealth[length=2.4mm]}, thick, black!45},
  ]
  \node[onsav=anarenk] (t0) at (0,0)    {topla\_sifir\\{\rmfamily\small $n + 0 = n$}\\{\rmfamily\scriptsize tanımdan}};
  \node[onsav=ikincil] (s0) at (4.4,1.2)  {sifir\_topla\\{\rmfamily\small $0 + n = n$}};
  \node[onsav=ikincil] (at) at (4.4,-1.2) {ard\_topla\\{\rmfamily\small $(n{+}1) + m = (n{+}m){+}1$}};
  \node[onsav=teoremframe] (dg) at (9.0,0) {topla\_degisme\\{\rmfamily\small $n + m = m + n$}};
  \node[onsav=teoremframe] (bl) at (9.0,-2.5) {topla\_birlesme\\{\rmfamily\small $(a{+}b){+}c = a{+}(b{+}c)$}};
  \draw[ok] (t0) -- (s0);
  \draw[ok] (s0) -- (dg);
  \draw[ok] (at) -- (dg);
  \node[font=\baslikfont\scriptsize, text=black!60, align=center] at (0,-2.3) {her ok:\\``şu önsav kullanılır''};
\end{tikzpicture}
\caption{Toplamanın temel özellikleri arasındaki bağımlılıklar. Birleşme
özelliği hiçbir önsava ihtiyaç duymaz; değişme özelliği ise iki önsav
gerektirir.}
\label{fig:onsav-grafi}
\end{figure}

\Cref{fig:onsav-grafi}, ispatlar arasındaki bağımlılıkları gösteriyor.
Bu küçük grafik, büyük formalizasyon projelerinin nasıl planlandığının
bir minyatürüdür: Önce hangi önsavların gerektiği belirlenir, sonra
bunlar tek tek ispatlanır (\cref{bol:proje}).

\begin{dikkat}[Dikkat: Hangi değişken üzerinde tümevarım?]
Toplama ikinci argümanda özyinelemeli olduğu için, tümevarımı da
genellikle ikinci argüman üzerinde yapmak gerekir. \lc{topla\_degisme}
ispatında tümevarımı \lc{n} üzerinde yapmayı deneyin: Tanım
açılamayacağı için takılırsınız. Doğru değişkeni seçmek, tümevarım
ispatlarındaki en önemli karardır ve çoğunlukla \emph{tanımın} hangi
argümana baktığına göre verilir.
\end{dikkat}

% ---------------------------------------------------------------------
\section{Lean'in doğal sayıları}\label{sec:nat}

\Lean{}'in yerleşik \lc{Nat} tipi, bizim \lc{Dogal} tipimizle tam olarak
aynı biçimde tanımlanmıştır:

\leandosya[\texttt{Nat}'ın tanımı]{Bolum07}{nat_print}

\noindent Tek fark, \Lean{}'in \lc{Nat} için \lc{0}, \lc{1}, \lc{2}
gibi rakamları, \lc{n + 1} gösterimini ve büyük sayılarla verimli
hesaplamayı desteklemesidir. \leantac{induction} taktiği \lc{Nat} üzerinde
de aynı biçimde çalışır; durumlar \lc{zero} ve \lc{succ} adını taşır:

\leandosya[$0 + n = n$, \texttt{Nat} ile]{Bolum07}{nat_sifir_arti}

\noindent Bu ispatta \Lean{}'in kütüphanesindeki \lc{Nat.add\_assoc}
(birleşme) teoremini ters yönde kullandık. Bir sonraki örnekte ise
aritmetik işi \leantac{omega}'ya bırakıyoruz:

\leandosya[İki katı]{Bolum07}{iki_kat}

\subsection{Gauss'un toplamı}

Anlatılan hikâyeye göre küçük Carl Friedrich Gauss'a öğretmeni
1'den 100'e kadar olan sayıları topla demiş, Gauss da birkaç saniye
içinde 5050 cevabını vermiş. Gauss'un fikri, toplamı iki kez yazıp
terimleri eşlemekti (\cref{fig:gauss}).

\begin{figure}[htbp]
\centering
\begin{tikzpicture}[
    top/.style={circle, fill=#1, inner sep=0pt, minimum size=0.36cm},
  ]
  \foreach \satir in {1,...,5} {
    \foreach \sutun in {1,...,6} {
      \pgfmathtruncatemacro{\sinir}{\satir}
      \ifnum\sutun>\sinir
        \node[top=vurgu!55] at ({0.48*\sutun},{-0.48*\satir}) {};
      \else
        \node[top=anarenk] at ({0.48*\sutun},{-0.48*\satir}) {};
      \fi
    }
  }
  \draw[decorate, decoration={brace, amplitude=5pt}, thick] (0.25,-0.2) -- (3.12,-0.2)
       node[midway, above=6pt, font=\small] {$n + 1 = 6$};
  \draw[decorate, decoration={brace, amplitude=5pt, mirror}, thick] (0.18,-0.28) -- (0.18,-2.68)
       node[midway, left=6pt, font=\small] {$n = 5$};
  \node[font=\small, align=left, anchor=west] at (3.8,-1.0)
       {\textcolor{anarenk}{$\bullet$} $1 + 2 + \cdots + n$\\[2pt]
        \textcolor{vurgu!80}{$\bullet$} $n + \cdots + 2 + 1$};
  \node[font=\small, anchor=west] at (3.8,-2.2)
       {$2 \cdot (1 + \cdots + n) = n(n+1)$};
\end{tikzpicture}
\caption{Gauss'un fikri: İki üçgen birleşince $n \times (n+1)$'lik bir
dikdörtgen oluşur.}
\label{fig:gauss}
\end{figure}

\leandosya[Toplam fonksiyonu]{Bolum07}{gauss_tanim}

\noindent \lc{\textcolor{leancmd}{\bfseries\#eval}} ve \leantac{decide},
formülü belirli değerler için denetler; ama \emph{her} $n$ için
doğruluğunu göstermez. Bunun için tümevarım gerekir:

\begin{teorem}{Gauss}{gauss}
Her $n$ doğal sayısı için $2 \cdot (1 + 2 + \cdots + n) = n(n+1)$.
\end{teorem}

\leandosya[Gauss toplamı]{Bolum07}{gauss}

\noindent İspatın tümevarım adımı, kâğıt üzerindeki hesabın
satır satır karşılığıdır:
\begin{align*}
2 \cdot \bigl(T(m) + (m+1)\bigr)
  &= 2 \cdot T(m) + 2(m+1)   && \text{dağılma} \\
  &= m(m+1) + 2(m+1)          && \text{tümevarım hipotezi} \\
  &= (m+2)(m+1)               && \text{ters dağılma} \\
  &= (m+1)(m+2).              && \text{değişme}
\end{align*}

% ---------------------------------------------------------------------
\section{Güçlü tümevarım}\label{sec:guclu}

Bazı ispatlarda $P(n+1)$'i göstermek için yalnızca $P(n)$ değil, daha
önceki bütün değerler gerekir. Örneğin bir sayıyı asal çarpanlarına
ayırırken $n$'nin çarpanlarına bakarız ve bunlar $n-1$'den çok daha
küçük olabilir. \emph{Güçlü tümevarım} bu durumlar içindir
(\cref{fig:guclu}).

\begin{figure}[htbp]
\centering
\begin{tikzpicture}[
    kutu/.style={draw=#1, fill=#1!10, thick, rounded corners=2pt,
                 minimum width=0.8cm, minimum height=0.8cm, font=\scriptsize},
  ]
  \node[font=\baslikfont\small\bfseries, text=anarenk, anchor=west] at (-0.5,1.1) {Sıradan tümevarım};
  \foreach \i in {0,...,4} \node[kutu=gri] (a\i) at ({1.0*\i},0) {$\i$};
  \node[kutu=leantac] (a5) at (5,0) {$5$};
  \draw[-{Stealth[length=2.2mm]}, thick, anarenk] (a4) to[bend left=40] (a5);
  \node[font=\scriptsize, text=black!60, anchor=west] at (5.6,0) {$P(4) \Rightarrow P(5)$};

  \node[font=\baslikfont\small\bfseries, text=ikincil, anchor=west] at (-0.5,-1.1) {Güçlü tümevarım};
  \foreach \i in {0,...,4} \node[kutu=ikincil] (b\i) at ({1.0*\i},-2.2) {$\i$};
  \node[kutu=leantac] (b5) at (5,-2.2) {$5$};
  \foreach \i in {0,...,4} \draw[-{Stealth[length=2.2mm]}, thick, ikincil] (b\i) to[bend right=35] (b5);
  \node[font=\scriptsize, text=black!60, anchor=west] at (5.6,-2.2) {$P(0), \ldots, P(4) \Rightarrow P(5)$};
\end{tikzpicture}
\caption{Sıradan tümevarımda yalnızca bir önceki değer, güçlü
tümevarımda ise bütün önceki değerler kullanılabilir.}
\label{fig:guclu}
\end{figure}

\begin{teorem}{Güçlü tümevarım}{guclu}
$P$ doğal sayılar üzerinde bir yüklem olsun. Her $n$ için ``$m < n$ olan
her $m$ için $P(m)$'' varsayımından $P(n)$ çıkıyorsa, her $n$ için $P(n)$
doğrudur.
\end{teorem}

\noindent Güçlü tümevarımı ayrı bir aksiyom olarak eklememize gerek
yoktur; sıradan tümevarımdan türetilebilir. Hile, daha güçlü bir önermeyi
(``$n$'den küçük her $m$ için $P(m)$'') sıradan tümevarımla ispatlamaktır:

\leandosya[Güçlü tümevarım]{Bolum07}{guclu}

\noindent İspatın ortasındaki \leantac{have} adımına dikkat edin:
\emph{İspatlanacak önermeyi güçlendirmek}, tümevarım ispatlarındaki en
güçlü tekniklerden biridir. Doğrudan \lc{P n}'yi ispatlamaya çalışınca
tümevarım hipotezi yetersiz kalır; güçlendirilmiş önermenin tümevarım
hipotezi ise tam ihtiyacımız olanı verir.

\begin{onemlifikir}[Önemli Fikir: İyi sıralama]
Güçlü tümevarım, \emph{iyi sıralama ilkesine} denktir: Doğal sayıların
boş olmayan her alt kümesinin bir en küçük elemanı vardır. ``En küçük
karşı örnek'' yöntemiyle yapılan bütün ispatlar (``varsayalım ki $P$'yi
sağlamayan sayılar var; en küçüğü $n$ olsun\ldots'') aslında güçlü
tümevarımdır. Bu yöntemi \cref{bol:tumevarimsal-tipler}'da, sona erme
ispatlarında yeniden kullanacağız.
\end{onemlifikir}

% ---------------------------------------------------------------------
\section{Otomasyon: simp, omega, decide}\label{sec:otomasyon-nat}

Doğal sayılarla ilgili ispatların büyük kısmı rutin hesaplardan
oluşur. \Lean{}'in bu işler için güçlü araçları vardır:

\leandosya[Otomasyon örnekleri]{Bolum07}{otomasyon}

\begin{table}[htbp]
\centering
\renewcommand{\arraystretch}{1.5}
\rowcolors{2}{tablosatir}{white}
\begin{tabularx}{\linewidth}{>{\raggedright\arraybackslash}p{2.0cm}
                             >{\raggedright\arraybackslash}X
                             >{\raggedright\arraybackslash}p{4.2cm}}
\rowcolor{tabloust}
\tablobaslik{Taktik} & \tablobaslik{Ne yapar?} & \tablobaslik{Ne yapamaz?} \\
\leantac{decide}  & Değişken içermeyen, karar verilebilir önermeleri hesaplayarak çözer. & Değişken içeren önermeler; çok büyük hesaplar. \\
\leantac{simp}    & Bilinen eşitliklerle (\lc{n + 0 = n} gibi) hedefi sadeleştirir. & Her zaman kapatmaz; sonuç tahmin edilemeyebilir. \\
\leantac{omega}   & Tam sayı ve doğal sayılar üzerinde \emph{doğrusal} aritmetiği (toplama, sabitle çarpma, $<$, $\leq$, $=$) tam olarak karar verir. & Değişkenlerin çarpımı ($x \cdot y$) içeren \emph{doğrusal olmayan} adımlar. \\
\end{tabularx}
\caption{Doğal sayılar için otomasyon araçları.}
\label{tab:otomasyon-nat}
\end{table}

\begin{dikkat}[Dikkat: omega doğrusal olmayan adımları atlayamaz]
Gauss ispatının son adımı olan $m(m+1) + 2(m+1) = (m+1)(m+2)$ eşitliği
\emph{doğrusal değildir}: Değişkenler birbiriyle çarpılıyor. \leantac{omega}
bu tür terimleri yalnızca ``atom'' olarak, yani bilinmeyen sayılar gibi
görür; $m \cdot (m+1)$ ile $(m+1) \cdot (m+2)$ arasındaki ilişkiyi
göremez. Bu yüzden Gauss ispatında dağılma ve değişme adımlarını kendimiz
yazdık. \Mathlib{}'teki \leantac{ring} taktiği bu tür polinom
eşitliklerini kendiliğinden çözer; onu \cref{bol:kumeler}'de tanıyacağız.
\end{dikkat}

\begin{tarihce}[Tarihçe: Presburger aritmetiği]
1929'da Mojżesz Presburger, yalnızca toplamayı içeren doğal sayı
aritmetiğinin \emph{karar verilebilir} ve \emph{tam} olduğunu gösterdi
\citep{presburger1929}: Bu dildeki her önermenin doğru mu yanlış mı
olduğu bir algoritmayla belirlenebilir. Çarpma eklendiğinde ise durum
tamamen değişir: Gödel iki yıl sonra, toplama ve çarpmayı birlikte içeren
aritmetiğin eksik olduğunu gösterecekti (\ciltiki{4}). \leantac{omega}
taktiği, William Pugh'un Presburger aritmetiği için geliştirdiği
``omega testi'' algoritmasının bir çeşitlemesine dayanır \citep{pugh1991}.
\end{tarihce}

% ---------------------------------------------------------------------
\section{Curry--Howard: tümevarım = özyineleme}\label{sec:curry-howard-4}

\begin{center}
\renewcommand{\arraystretch}{1.5}
\rowcolors{2}{tablosatir}{white}
\begin{tabular}{>{\raggedright\arraybackslash}p{4.2cm}
                >{\raggedright\arraybackslash}p{4.6cm}
                >{\raggedright\arraybackslash}p{3.0cm}}
\rowcolor{tabloust}
\tablobaslik{Mantık} & \tablobaslik{Programlama} & \tablobaslik{Lean} \\
tümevarımsal tanım            & veri tipi tanımı                & \leankw{inductive} \\
tümevarımla ispat             & özyinelemeli fonksiyon          & \leantac{induction}, örüntü eşleme \\
taban adımı                   & taban durum                     & \lc{| sifir => \ldots} \\
tümevarım hipotezi            & özyinelemeli çağrının sonucu    & \lc{ih}, \lc{f m} \\
tümevarım ilkesi              & özyineleyici                    & \lc{Dogal.rec} \\
\end{tabular}
\end{center}

% ---------------------------------------------------------------------
\section{Taktik özeti}

\begin{center}
\renewcommand{\arraystretch}{1.45}
\rowcolors{2}{tablosatir}{white}
\begin{tabularx}{\linewidth}{>{\raggedright\arraybackslash}p{4.2cm} X}
\rowcolor{tabloust}
\tablobaslik{Taktik} & \tablobaslik{Ne yapar?} \\
\leantac{induction n with | ... }  & \lc{n} üzerinde tümevarım; her kurucu için bir durum açar. \\
\leantac{injection h with h'}      & \lc{h : c a = c b}'den \lc{h' : a = b} çıkarır. \\
\leantac{cases h} (\lc{h : c₁ = c₂}) & Farklı kurucuların eşitliğinden çelişki çıkarır. \\
\leantac{show t}                   & Hedefi tanımsal olarak eşit bir biçime çevirir (tanımı açar). \\
\leantac{change t at h}            & Hipotezi tanımsal olarak eşit bir biçime çevirir. \\
\leantac{omega}                    & Doğrusal tam sayı aritmetiğini çözer. \\
\leantac{simp}                     & Hedefi bilinen eşitliklerle sadeleştirir. \\
\leantac{decide}                   & Değişken içermeyen önermeleri hesaplar. \\
\end{tabularx}
\end{center}

% ---------------------------------------------------------------------
\begin{arastirmanotu}
\textbf{Natural Number Game.} Kevin Buzzard ve arkadaşlarının hazırladığı
çevrimiçi \emph{Natural Number Game}, bu bölümdeki ispatları (ve çok
daha fazlasını) adım adım bir oyun olarak sunar \citep{nng}. Peano
aksiyomlarından başlayıp $\leq$ ilişkisine ve asal sayılara uzanır.
Bu bölümü bitiren okura özellikle öneririz.

\textbf{Tümevarımın gücü.} Peano aritmetiğinin (PA) gücü, tümevarım
aksiyomunun hangi önermeler için kabul edildiğine bağlıdır. Tümevarımı
yalnızca basit (sınırlı niceleyicili) önermelerle kısıtlayan zayıf
aritmetik sistemleri, karmaşıklık kuramıyla ilişkileri nedeniyle hâlâ
etkin bir araştırma alanıdır. Gentzen'in 1936'daki tutarlılık ispatı,
PA'nın tutarlılığının $\varepsilon_0$ ordinaline kadar
\emph{sonsuz ötesi tümevarımla} gösterilebileceğini ortaya koydu
\citep{gentzen1936}; bu konuya \ciltiki{4}'te döneceğiz.

\textbf{İleri okuma.} Tümevarım ve özyineleme için \citet{avigad2022};
\Lean{}'de özyinelemeli tanımlar ve sona erme için \citet{tpil4} ve
\citet{fpil}.
\end{arastirmanotu}

\begin{bolumozeti}
\begin{itemize}
  \item Doğal sayılar, \lc{sifir} ve \lc{ard} kurucularıyla tümevarımsal
        bir tip olarak tanımlanır. Peano aksiyomları bu tanımdan teorem
        olarak çıkar.
  \item Fonksiyonlar özyinelemeyle tanımlanır; \Lean{} her özyinelemenin
        sona erdiğini denetler.
  \item Tümevarım ilkesi, tipin özyineleyicisidir. \leantac{induction}
        taktiği taban ve adım durumlarını açar; adım durumunda
        tümevarım hipotezi \lc{ih} kullanılır.
  \item \lc{n + 0 = n} tanımsal, \lc{0 + n = n} önermeseldir ve tümevarım
        ister. Tümevarım genellikle tanımın özyinelemeli olduğu argüman
        üzerinde yapılır.
  \item Güçlü tümevarım sıradan tümevarımdan türetilir; anahtar teknik
        önermeyi güçlendirmektir.
  \item \leantac{omega} doğrusal aritmetiği, \leantac{decide} somut
        hesapları çözer. Doğrusal olmayan adımlar elle ya da \Mathlib{}'in
        \leantac{ring} taktiğiyle yapılır.
\end{itemize}
\end{bolumozeti}

\begin{alistirmalar}
\alistirma{1} \lc{topla n (ard sifir) = ard n}. Neden tümevarım
gerekmiyor? (\lc{a7\_1})

\alistirma{2} \lc{carp sifir n = sifir}. Hangi değişken üzerinde
tümevarım yapmalısınız? (\lc{a7\_2})

\alistirma{2} Toplamada sadeleştirme: \lc{topla a c = topla b c → a = b}.
İpucu: Tümevarım adımında \leantac{injection} kullanın. (\lc{a7\_3})

\alistirma{1} \lc{Nat} üzerinde, kütüphane teoremi kullanmadan:
\lc{0 * n = 0}. (\lc{a7\_4})

\alistirma{2} \lc{ucKat n = 3 * n}. (\lc{a7\_5})

\alistirma{2} Tek sayıların toplamı bir karedir:
$1 + 3 + \cdots + (2n-1) = n^2$. Önce \cref{fig:gauss}'e benzer bir
resimle neden doğru olduğunu gösterin, sonra \Lean{}'de ispatlayın.
(\lc{a7\_6})

\alistirma{2} \lc{n < 2 \textasciicircum{} n}. İpucu: \lc{Nat.pow\_succ}. (\lc{a7\_7})

\alistirma{3} \textbf{Pul problemi}: 8 kuruş ve üstündeki her tutar
yalnızca 3 ve 5 kuruşluk pullarla ödenebilir. Güçlü tümevarımla
ispatlayın. Hangi üç durum taban adımı olarak gerekli? (\lc{a7\_8})

\alistirma{2} (Kâğıt-kalem) \lc{carp iki uc} ifadesinin hesaplanmasını
\cref{fig:hesaplama-izi}'deki gibi adım adım yazın.

\alistirma{3} (Kâğıt-kalem) \lc{topla\_degisme}'yi \lc{n} üzerinde
tümevarımla ispatlamaya çalışın. Hangi adımda takılıyorsunuz?
Takılmamak için hangi önsavlar gerekir?

\alistirma{3} (Tartışma) Aşağıdaki ``ispatın'' hatası nedir?
``Bütün atlar aynı renktir. Tümevarımla: Bir attan oluşan her grupta
atlar aynı renktir. $n$ atlık gruplar için doğru olsun; $n+1$ atlık bir
grupta ilk $n$ at aynı renktir, son $n$ at da aynı renktir; iki grup
ortak atlar içerdiği için hepsi aynı renktir.''
\end{alistirmalar}
   % Doğal sayılar ve tümevarım
% =====================================================================
%  Bölüm 8 — Kümeler, İlişkiler ve Fonksiyonlar
% =====================================================================
\chapter{Kümeler, İlişkiler ve Fonksiyonlar}\label{bol:kumeler}

\begin{bolumhedefleri}
\begin{itemize}
  \item \Mathlib{}'i projeye eklemeyi, adlandırma kurallarını ve teorem
        arama araçlarını,
  \item \leantac{ring}, \leantac{linarith}, \leantac{norm\_num},
        \leantac{positivity} ve \leantac{tauto} taktiklerini,
  \item Kümeleri \emph{yüklem} olarak görmeyi ve küme cebirinin
        önermeler mantığının bir yansıması olduğunu,
  \item \leantac{ext} taktiğiyle küme eşitliği ispatlamayı,
  \item İlişkileri, denklik ilişkilerini ve kısmi sıralamaları,
  \item Birebir, örten ve eşleme fonksiyonları; bileşke ve tersleri,
  \item Cantor teoremini ve köşegen yöntemini.
\end{itemize}
\end{bolumhedefleri}

Önceki bölümlerde her şeyi sıfırdan kurduk: bağlaçları, niceleyicileri,
doğal sayıları. Bu yaklaşım, mantığın nasıl çalıştığını anlamak için
vazgeçilmezdir; ama gerçek matematik yapmak için yavaştır. Bu bölümde
\Lean{} topluluğunun yıllardır biriktirdiği büyük matematik kütüphanesi
\Mathlib{}'i kullanmaya başlıyoruz \citep{mathlib2020}. Ardından
matematiğin üç temel dilini, yani kümeleri, ilişkileri ve fonksiyonları
mantık açısından ele alacağız.

% ---------------------------------------------------------------------
\section{Mathlib'e giriş}\label{sec:mathlib}

\Mathlib{}, \Lean{}'de yazılmış, gönüllü bir topluluk tarafından
geliştirilen ve sürekli büyüyen bir matematik kütüphanesidir. Cebirden
analize, topolojiden sayılar kuramına uzanan yüz binlerce tanım ve
teorem içerir. Kütüphaneye eklenen her teorem, \Lean{}'in çekirdeği
tarafından denetlenmiştir.

\paragraph{Kurulum}
Eşlik eden projede \Mathlib{} bağımlılığı bu bölümle birlikte açıldı.
\Mathlib{} belirli bir \Lean{} sürümüyle çalıştığı için kurulumda üç adım
gerekir: projenin \lc{lean-toolchain} dosyasını \Mathlib{}'inkiyle
eşlemek, \lc{lake update} ile bağımlılıkları indirmek ve
\lc{lake exe cache get} ile önceden derlenmiş dosyaları almak. Adımların
tam hâli projenin \lc{README.md} dosyasındadır.

\begin{dikkat}[Dikkat: Önbelleği indirmeyi unutmayın]
\Mathlib{}'i sıfırdan derlemek güçlü bir bilgisayarda bile saatler sürer.
\lc{lake exe cache get} komutu, önceden derlenmiş dosyaları birkaç
dakikada indirir. Bir \Lean{} dosyasının başında \lc{import Mathlib}
satırı varsa ve dosya açılırken bilgisayarınız uzun süre ``derliyor''
görünüyorsa, büyük olasılıkla bu adım atlanmıştır.
\end{dikkat}

\leandosya[Mathlib'e ilk bakış]{Bolum08}{mathlib_ilk}

\subsection{Mathlib'in taktikleri}

\Mathlib{}, çekirdek \Lean{}'in taktiklerine güçlü yeni araçlar ekler
(\cref{tab:mathlib-taktik}).

\begin{table}[htbp]
\centering
\renewcommand{\arraystretch}{1.45}
\rowcolors{2}{tablosatir}{white}
\begin{tabularx}{\linewidth}{>{\raggedright\arraybackslash}p{2.4cm}
                             >{\raggedright\arraybackslash}X
                             >{\raggedright\arraybackslash}p{4.2cm}}
\rowcolor{tabloust}
\tablobaslik{Taktik} & \tablobaslik{Ne yapar?} & \tablobaslik{Örnek hedef} \\
\leantac{ring}        & Değişmeli halkalarda polinom eşitliklerini kanıtlar. & $(a+b)^2 = a^2 + 2ab + b^2$ \\
\leantac{linarith}    & Hipotezlerden doğrusal eşitsizlikleri türetir. & $x < y,\; y < 3 \vdash x < 3$ \\
\leantac{norm\_num}   & Somut sayısal ifadeleri hesaplar; asallık gibi özellikleri denetler. & \lc{Nat.Prime 97} \\
\leantac{positivity}  & Bir ifadenin pozitif ya da negatif olmadığını gösterir. & $0 \leq x^2 + 1$ \\
\leantac{tauto}       & Önermeler mantığının totolojilerini (klasik olarak) kanıtlar. & $P \wedge Q \to Q \vee R$ \\
\leantac{ext}         & İki küme ya da fonksiyonun eşitliğini, elemanlarının eşitliğine indirger. & $s \cap t = t \cap s$ \\
\lc{exact?}           & Kütüphanede hedefi doğrudan kapatan bir teorem arar. & \lc{a + b = b + a} \\
\end{tabularx}
\caption{\Mathlib{}'in sık kullanılan taktikleri.}
\label{tab:mathlib-taktik}
\end{table}

Bölüm~7'de Gauss formülünün son adımını, dağılma ve değişme teoremlerini
tek tek uygulayarak elle yapmıştık. \leantac{ring} bu işi tek adımda
yapar:

\leandosya[Gauss, \texttt{ring} ile]{Bolum08}{gauss_ring}

\noindent \Mathlib{}, sonlu toplamlar için matematikteki $\sum$
gösterimini de sunar. Aynı formül, kitaplarda yazıldığı biçimiyle:

\leandosya[Gauss, toplam gösterimiyle]{Bolum08}{gauss_finset}

\subsection{Teorem bulmak}

\Mathlib{}'de aradığınız teoremin büyük olasılıkla zaten var olduğunu
varsayın. Onu bulmanın üç yolu vardır:

\begin{enumerate}[adimlar]
  \item \textbf{Adını tahmin etmek.} \Mathlib{} katı bir adlandırma
        düzeni izler (\cref{tab:adlandirma}). ``Çarpma değişmelidir''
        teoremi \lc{mul\_comm}, ``bir kümenin kesişimine üyelik'' teoremi
        \lc{Set.mem\_inter\_iff}'tir.
  \item \textbf{\lc{exact?} ile aramak.} Hedefin yerine \lc{exact?}
        yazdığınızda \Lean{}, kütüphanede hedefi kapatan bir teorem arar
        ve bulursa önerir.
  \item \textbf{Çevrimiçi araçlar.} \Mathlib{}'in belgelendirme
        sayfaları ve teoremleri tiplerine göre arayan \emph{Loogle}
        arama motoru \citep{loogle}.
\end{enumerate}

\begin{leankod}[exact? ile arama (çıktı bir öneridir)]
example (a b : ℕ) : a + b = b + a := by
  exact?   -- Try this: exact Nat.add_comm a b
\end{leankod}

\begin{table}[htbp]
\centering
\renewcommand{\arraystretch}{1.4}
\rowcolors{2}{tablosatir}{white}
\begin{tabular}{>{\raggedright\arraybackslash}p{3.6cm}
                >{\raggedright\arraybackslash}p{4.6cm}
                >{\raggedright\arraybackslash}p{3.4cm}}
\rowcolor{tabloust}
\tablobaslik{Ad} & \tablobaslik{İfade} & \tablobaslik{Okunuşu} \\
\lc{mul\_comm}          & $a \cdot b = b \cdot a$          & çarpma, değişme \\
\lc{add\_assoc}         & $(a+b)+c = a+(b+c)$              & toplama, birleşme \\
\lc{mul\_add}           & $a(b+c) = ab + ac$               & çarpım, toplam üzerine \\
\lc{add\_zero}          & $a + 0 = a$                      & toplama, sıfır \\
\lc{le\_trans}          & $a \leq b \to b \leq c \to a \leq c$ & $\leq$, geçişme \\
\lc{Set.mem\_inter\_iff} & $x \in s \cap t \leftrightarrow x \in s \wedge x \in t$ & üyelik, kesişim, ancak ve ancak \\
\end{tabular}
\caption{\Mathlib{}'in adlandırma düzeninden örnekler. Ad, ifadenin
sembollerini soldan sağa okur; \lc{\_iff} bir denkliği, \lc{\_comm}
değişmeyi, \lc{\_assoc} birleşmeyi bildirir.}
\label{tab:adlandirma}
\end{table}

% ---------------------------------------------------------------------
\section{Kümeler yüklemdir}\label{sec:kume-yuklem}

\Mathlib{}'de bir küme, bir yüklemden başka bir şey değildir:

\begin{tanim}{Küme}{kume}
$\alpha$ tipinin bir \textbf{alt kümesi}, $\alpha$ üzerinde bir yüklemdir:
\[
  \mathsf{Set}\;\alpha \;:=\; \alpha \to \mathsf{Prop}.
\]
$x \in s$, ``$s(x)$ doğrudur'' demektir. $\{x \mid P(x)\}$ gösterimi,
$P$ yükleminin kendisidir.
\end{tanim}

\leandosya[Kümeler ve üyelik]{Bolum08}{kume_tanim}

\noindent İkinci örnekteki \lc{Iff.rfl}, ``$a \in \{x \mid P(x)\}$'' ile
``$P(a)$''nın \emph{tanım gereği} aynı olduğunu gösterir. \Cref{fig:kume-yuklem},
bu bakış açısını \cref{fig:yuklem}'deki yüklem resmiyle birleştiriyor.

\begin{figure}[htbp]
\centering
\begin{tikzpicture}[
    eleman/.style={circle, inner sep=2.2pt, fill=#1},
  ]
  \draw[anarenk, thick, fill=tanimbg, rounded corners=12pt] (0,0) rectangle (5.2,3.4);
  \node[font=\baslikfont\bfseries\small, text=anarenk, anchor=north west] at (0.15,3.35) {$\alpha = \mathbb{N}$};
  \draw[vurgu, very thick, fill=vurgu!12, rounded corners=10pt] (2.4,0.35) rectangle (4.9,2.45);
  \node[font=\baslikfont\bfseries\small, text=vurgu, anchor=south east] at (4.9,2.45) {$s = \{n \mid n \% 2 = 1\}$};
  \foreach \x/\y/\n in {2.9/1.9/1, 3.9/1.9/3, 2.9/0.9/5, 3.9/0.9/7}
    { \node[eleman=vurgu] at (\x,\y) {}; \node[font=\scriptsize, right=1pt] at (\x,\y) {$\n$}; }
  \foreach \x/\y/\n in {0.6/1.9/0, 1.5/1.9/2, 0.6/0.9/4, 1.5/0.9/6}
    { \node[eleman=black!35] at (\x,\y) {}; \node[font=\scriptsize, right=1pt] at (\x,\y) {$\n$}; }
  % karakteristik fonksiyon
  \node[font=\small, anchor=west, align=left] at (5.8,2.1)
    {$s(1) = (1 \% 2 = 1)$ \;\textcolor{bittiframe}{doğru}\\[2pt]
     $s(2) = (2 \% 2 = 1)$ \;\textcolor{hataframe}{yanlış}};
  \node[font=\baslikfont\scriptsize, text=black!60, anchor=west, align=left] at (5.8,0.8)
    {küme = her elemana bir önerme\\karşılık getiren fonksiyon};
\end{tikzpicture}
\caption{Bir küme, elemanlarını ``doğru'' önermelere götüren bir
yüklemdir.}
\label{fig:kume-yuklem}
\end{figure}

\begin{dikkat}[Dikkat: Tip kuramında kümeler homojendir]
Klasik küme kuramında her şey bir kümedir ve herhangi iki nesne
arasında $\in$ sorusu sorulabilir: ``$3 \in \pi$ mi?'' gibi anlamsız
sorular bile. Tip kuramında ise \lc{Set α} yalnızca \lc{α} tipindeki
elemanlardan oluşur; \lc{x ∈ s} yazabilmek için \lc{x : α} olmalıdır.
Bu kısıt, Russell paradoksunu (\cref{sec:berber}) daha baştan engeller:
``Kendini içermeyen kümelerin kümesi'' tanımlanamaz, çünkü bir küme
kendi tipinin elemanı değildir.
\end{dikkat}

% ---------------------------------------------------------------------
\section{Küme cebiri = önermeler mantığı}\label{sec:kume-cebiri}

Kümeler yüklem olduğuna göre, küme işlemleri de mantık bağlaçlarıdır
(\cref{tab:kume-mantik}).

\begin{table}[htbp]
\centering
\renewcommand{\arraystretch}{1.5}
\rowcolors{2}{tablosatir}{white}
\begin{tabular}{>{\centering\arraybackslash}p{1.6cm}
                >{\raggedright\arraybackslash}p{2.6cm}
                >{\raggedright\arraybackslash}p{4.6cm}
                >{\centering\arraybackslash}p{1.6cm}}
\rowcolor{tabloust}
\tablobaslik{Küme} & \tablobaslik{Adı} & \tablobaslik{Tanımı ($x \in \ldots$ ancak ve ancak)} & \tablobaslik{Mantık} \\
$s \cap t$       & kesişim      & $x \in s \wedge x \in t$ & $\wedge$ \\
$s \cup t$       & birleşim     & $x \in s \vee x \in t$   & $\vee$ \\
$s^{c}$          & tümleyen     & $\neg(x \in s)$          & $\neg$ \\
$s \setminus t$  & fark         & $x \in s \wedge \neg(x \in t)$ & $\wedge\neg$ \\
$s \subseteq t$  & alt küme     & her $x$ için $x \in s \to x \in t$ & $\to$ \\
$\emptyset$      & boş küme     & $\bot$                   & \lc{False} \\
\lc{univ}        & evrensel küme & $\top$                  & \lc{True} \\
\end{tabular}
\caption{Küme işlemleri ile mantık bağlaçları arasındaki sözlük.}
\label{tab:kume-mantik}
\end{table}

\begin{figure}[htbp]
\centering
\begin{tikzpicture}
  \def\A{(0,0) circle (0.75)}
  \def\B{(0.85,0) circle (0.75)}
  \foreach \i/\baslik in {0/{$s \cap t$}, 1/{$s \cup t$}, 2/{$s \setminus t$}, 3/{$s^{c}$}} {
    \begin{scope}[shift={({3.6*\i},0)}]
      \draw[black!20, rounded corners=4pt, fill=white] (-1.2,-1.1) rectangle (2.05,1.1);
      \ifcase\i
        \begin{scope}\clip \A; \fill[anarenk!45] \B; \end{scope}
      \or
        \fill[anarenk!45] \A; \fill[anarenk!45] \B;
      \or
        \begin{scope}\clip \A; \fill[anarenk!45] \A; \end{scope}
        \fill[white] \B;
      \or
        \fill[anarenk!45, rounded corners=4pt] (-1.2,-1.1) rectangle (2.05,1.1);
        \fill[white] \A;
      \fi
      \draw[thick, anarenk] \A;
      \draw[thick, vurgu] \B;
      \node[font=\scriptsize\bfseries, text=anarenk] at (-0.35,0.5) {$s$};
      \node[font=\scriptsize\bfseries, text=vurgu] at (1.2,0.5) {$t$};
      \node[font=\small] at (0.42,-1.4) {\baslik};
    \end{scope}
  }
\end{tikzpicture}
\caption{Dört temel küme işleminin Venn şemaları.}
\label{fig:venn}
\end{figure}

\begin{onemlifikir}[Önemli Fikir: Küme eşitlikleri mantık denklikleridir]
Bu sözlük, küme cebirindeki her yasanın önermeler mantığındaki bir
yasanın kılık değiştirmiş hâli olduğunu söyler. $s \cap (t \cup u) =
(s \cap t) \cup (s \cap u)$ dağılma yasası, Bölüm~2'deki
$P \wedge (Q \vee R) \leftrightarrow (P \wedge Q) \vee (P \wedge R)$
dağılma yasasıdır. George Boole'un 1854'te keşfettiği bu ortak yapıya
\emph{Boole cebiri} denir.
\end{onemlifikir}

\subsection{Alt küme ve küme eşitliği}

$s \subseteq t$, ``her $x$ için $x \in s$ ise $x \in t$'' demektir. Bir
alt küme ispatı, dolayısıyla \leantac{intro} ile başlar:

\leandosya[Alt kümenin geçişliliği]{Bolum08}{altkume}

\noindent İki kümenin eşitliği ise \emph{dışsallık} ilkesine dayanır:
Aynı elemanlara sahip iki küme eşittir. \leantac{ext x} taktiği,
\lc{s = t} hedefini \lc{x ∈ s ↔ x ∈ t} hedefine çevirir:

\leandosya[Kesişimin değişme özelliği]{Bolum08}{kesisim_degisme}

\begin{hedefgecisi}
\begin{hedefdurumu}[\texttt{ext x} öncesi]
s t : Set α
⊢ s ∩ t = t ∩ s
\end{hedefdurumu}
\sonra
\begin{hedefdurumu}[\texttt{ext x} sonrası]
s t : Set α
x : α
⊢ x ∈ s ∩ t ↔ x ∈ t ∩ s
\end{hedefdurumu}
\end{hedefgecisi}

\noindent \lc{x ∈ s ∩ t}, tanım gereği \lc{x ∈ s ∧ x ∈ t} olduğundan
\leantac{rintro ⟨hs, ht⟩} onu doğrudan parçalayabilir. Dağılma yasasının
ispatı, Bölüm~2'deki ispatın neredeyse harfi harfine aynısıdır:

\leandosya[Dağılma]{Bolum08}{dagilma}

\noindent Aynı ispatı \Mathlib{}'in araçlarıyla kısaltabiliriz:
\leantac{simp only} üyelik ifadelerini mantık bağlaçlarına çevirir,
\leantac{tauto} da ortaya çıkan totolojiyi kanıtlar:

\leandosya[Dağılma, otomasyonla]{Bolum08}{dagilma_simp}

\leandosya[De Morgan, kümeler için]{Bolum08}{tumleyen}

\noindent Bu ispat, Bölüm~3'teki \lc{demorgan\_veya} ispatının küme
diliyle yazılmış hâlidir. Kesişimin tümleyeni için De Morgan yasası ise,
Bölüm~3'teki ``zor yön'' nedeniyle klasik mantık gerektirir
(Alıştırma~\ref{alis:kume-demorgan}).

% ---------------------------------------------------------------------
\section{İlişkiler}\label{sec:iliskiler}

\begin{tanim}{İlişki}{iliski}
$\alpha$ üzerinde bir \textbf{ilişki}, iki argümanlı bir yüklemdir:
$R : \alpha \to \alpha \to \mathsf{Prop}$. $R\,a\,b$ ``$a$, $b$ ile
$R$ ilişkisindedir'' diye okunur.
\end{tanim}

\begin{table}[htbp]
\centering
\renewcommand{\arraystretch}{1.5}
\rowcolors{2}{tablosatir}{white}
\begin{tabular}{>{\raggedright\arraybackslash}p{2.6cm}
                >{\raggedright\arraybackslash}p{5.4cm}
                >{\raggedright\arraybackslash}p{3.6cm}}
\rowcolor{tabloust}
\tablobaslik{Özellik} & \tablobaslik{Tanım} & \tablobaslik{Örnek} \\
yansıma      & $\forall a,\; R\,a\,a$                                  & $=$, $\leq$, $\mid$ \\
simetri      & $\forall a\,b,\; R\,a\,b \to R\,b\,a$                   & $=$, ``kardeştir'' \\
ters simetri & $\forall a\,b,\; R\,a\,b \to R\,b\,a \to a = b$         & $\leq$, $\mid$ ($\mathbb{N}$'de) \\
geçişme      & $\forall a\,b\,c,\; R\,a\,b \to R\,b\,c \to R\,a\,c$    & $=$, $<$, $\leq$, $\mid$ \\
\end{tabular}
\caption{İlişkilerin temel özellikleri.}
\label{tab:iliski-ozellik}
\end{table}

Bu özelliklerin iki önemli bileşimi vardır:
\begin{itemize}
  \item \textbf{Denklik ilişkisi}: yansıma + simetri + geçişme.
        ``Aynı kalanı vermek'', ``aynı doğum gününe sahip olmak'' gibi
        ``bir bakıma aynı olmak'' ilişkileri.
  \item \textbf{Kısmi sıralama}: yansıma + ters simetri + geçişme.
        $\leq$, $\subseteq$ ve bölünebilirlik ($\mid$) gibi.
\end{itemize}

\subsection{Denklik ilişkileri}

``$n$'ye bölündüğünde aynı kalanı vermek'' bir denklik ilişkisidir.
\Lean{}'de bir denklik ilişkisi, üç özelliği alan olarak içeren bir
\lc{Equivalence} yapısıdır:

\leandosya[Mod $n$ denkliği]{Bolum08}{mod_denk}

\noindent Burada \leankw{where} ile bir yapının alanlarını tek tek
dolduruyoruz. Her alanın ispatı, eşitliğin ilgili özelliğinden gelir:
\lc{modDenk n a b}, tanım gereği \lc{a \% n = b \% n} eşitliği olduğu için
yansıma \lc{rfl}, simetri \lc{Eq.symm}, geçişme \lc{Eq.trans} ile
gösterilir.

Bir denklik ilişkisi, kümeyi \emph{denklik sınıflarına} ayırır: Birbirine
denk olan elemanlar aynı sınıftadır ve her eleman tam olarak bir sınıfa
aittir (\cref{fig:denklik-siniflari}).

\begin{figure}[htbp]
\centering
\begin{tikzpicture}[
    sinif/.style={draw=#1, fill=#1!10, very thick, rounded corners=10pt,
                  minimum width=3.4cm, minimum height=2.2cm},
    sayi/.style={font=\small},
  ]
  \node[sinif=anarenk]     (k0) at (0,0)   {};
  \node[sinif=ikincil]     (k1) at (4.0,0) {};
  \node[sinif=vurgu]       (k2) at (8.0,0) {};
  \node[font=\baslikfont\bfseries\small, text=anarenk, above=2pt of k0] {kalan 0};
  \node[font=\baslikfont\bfseries\small, text=ikincil, above=2pt of k1] {kalan 1};
  \node[font=\baslikfont\bfseries\small, text=vurgu, above=2pt of k2] {kalan 2};
  \node[sayi, align=center] at (k0) {$0 \quad 3 \quad 6$\\[3pt]$9 \quad 12 \;\cdots$};
  \node[sayi, align=center] at (k1) {$1 \quad 4 \quad 7$\\[3pt]$10 \quad 13 \;\cdots$};
  \node[sayi, align=center] at (k2) {$2 \quad 5 \quad 8$\\[3pt]$11 \quad 14 \;\cdots$};
\end{tikzpicture}
\caption{``3'e bölündüğünde aynı kalanı vermek'' ilişkisi, $\mathbb{N}$'yi
üç denklik sınıfına ayırır.}
\label{fig:denklik-siniflari}
\end{figure}

\subsection{Kısmi sıralamalar}

Bölünebilirlik, doğal sayılar üzerinde bir kısmi sıralamadır:

\leandosya[Bölünebilirlik]{Bolum08}{bolunme}

\noindent ``Kısmi'' sözcüğü, her iki elemanın karşılaştırılamayabileceğini
belirtir: Ne $2 \mid 3$ ne de $3 \mid 2$ doğrudur. Kısmi sıralamalar,
\emph{Hasse şemalarıyla} çizilir: $a \mid b$ ise $b$, $a$'nın üstünde
durur ve arada başka bir eleman yoksa ikisi bir çizgiyle birleştirilir
(\cref{fig:hasse}).

\begin{figure}[htbp]
\centering
\begin{tikzpicture}[
    d/.style={circle, draw=anarenk, fill=tanimbg, thick, minimum size=0.75cm, font=\small},
  ]
  \node[d] (1) at (0,0) {1};
  \node[d] (2) at (-1.6,1.4) {2};
  \node[d] (3) at (1.6,1.4) {3};
  \node[d] (4) at (-2.6,2.8) {4};
  \node[d] (6) at (0,2.8) {6};
  \node[d, draw=vurgu, fill=teorembg] (12) at (-1.3,4.2) {12};
  \foreach \a/\b in {1/2, 1/3, 2/4, 2/6, 3/6, 4/12, 6/12}
    \draw[thick, black!45] (\a) -- (\b);
  \node[font=\baslikfont\scriptsize, text=black!60, align=left, anchor=west] at (2.6,2.2)
    {12'nin bölenleri\\\emph{kısmi} sıralama:\\4 ile 6 karşılaştırılamaz};
\end{tikzpicture}
\caption{12'nin bölenlerinin bölünebilirlik sıralamasındaki Hasse şeması.}
\label{fig:hasse}
\end{figure}

% ---------------------------------------------------------------------
\section{Fonksiyonlar}\label{sec:fonksiyonlar}

\begin{tanim}{Birebir, örten, eşleme}{fonk-tur}
$f : \alpha \to \beta$ bir fonksiyon olsun.
\begin{itemize}
  \item $f$ \textbf{birebirdir} (İng.\ \emph{injective}):
        $\forall a\,b,\; f(a) = f(b) \to a = b$.
  \item $f$ \textbf{örtendir} (İng.\ \emph{surjective}):
        $\forall b,\; \exists a,\; f(a) = b$.
  \item $f$ \textbf{eşlemedir} (İng.\ \emph{bijective}): hem birebir hem örten.
\end{itemize}
\end{tanim}

\begin{figure}[htbp]
\centering
\begin{tikzpicture}[
    e/.style={circle, fill=#1, inner sep=1.8pt},
    ok/.style={-{Stealth[length=2mm]}, thick, black!50},
  ]
  \foreach \k/\baslik/\renk in {0/{birebir, örten değil}/anarenk, 1/{örten, birebir değil}/ikincil, 2/{eşleme}/vurgu} {
    \begin{scope}[shift={({4.4*\k},0)}]
      \draw[\renk, thick, fill=\renk!7, rounded corners=8pt] (-0.5,-1.6) rectangle (0.5,1.3);
      \draw[\renk, thick, fill=\renk!7, rounded corners=8pt] (1.9,-1.6) rectangle (2.9,1.3);
      \node[font=\small] at (0,1.6) {$\alpha$};
      \node[font=\small] at (2.4,1.6) {$\beta$};
      \node[font=\baslikfont\scriptsize\bfseries, text=\renk] at (1.2,-2.0) {\baslik};
    \end{scope}
  }
  % birebir, örten değil: 3 → 4
  \foreach \i/\y in {1/0.8, 2/0.0, 3/-0.8} \node[e=anarenk] (a\i) at (0,\y) {};
  \foreach \i/\y in {1/0.9, 2/0.3, 3/-0.3, 4/-0.9} \node[e=anarenk] (b\i) at (2.4,\y) {};
  \draw[ok] (a1) -- (b1); \draw[ok] (a2) -- (b2); \draw[ok] (a3) -- (b4);
  % örten, birebir değil: 4 → 3
  \foreach \i/\y in {1/0.9, 2/0.3, 3/-0.3, 4/-0.9} \node[e=ikincil] (c\i) at (4.4,\y) {};
  \foreach \i/\y in {1/0.8, 2/0.0, 3/-0.8} \node[e=ikincil] (d\i) at (6.8,\y) {};
  \draw[ok] (c1) -- (d1); \draw[ok] (c2) -- (d1); \draw[ok] (c3) -- (d2); \draw[ok] (c4) -- (d3);
  % eşleme: 3 → 3
  \foreach \i/\y in {1/0.8, 2/0.0, 3/-0.8} {
    \node[e=vurgu] (f\i) at (8.8,\y) {};
    \node[e=vurgu] (g\i) at (11.2,\y) {};
  }
  \draw[ok] (f1) -- (g2); \draw[ok] (f2) -- (g3); \draw[ok] (f3) -- (g1);
\end{tikzpicture}
\caption{Üç fonksiyon türü. Birebir fonksiyonda iki ok aynı yere gitmez;
örten fonksiyonda her hedef bir ok alır; eşlemede ikisi birden.}
\label{fig:fonksiyon-turleri}
\end{figure}

\Mathlib{}'de bu kavramlar \lc{Function.Injective}, \lc{Function.Surjective}
ve \lc{Function.Bijective} adını taşır. Tanımları, yukarıdaki
formüllerin harfi harfine aynısıdır; ispatlar da Bölüm~6'nın niceleyici
kurallarıyla yapılır:

\leandosya[İkiyle çarpma]{Bolum08}{birebir_ornek}

\noindent Birebirlik ispatı \leantac{intro a b h} ile başlar
($\forall$G, $\to$G); örten olmamanın ispatı ise varsayılan örtenliği
1 sayısına uygular ($\forall$E) ve tanığı \leantac{obtain} ile alır
($\exists$E).

\subsection{Bileşke}

\leandosya[Bileşke]{Bolum08}{bileske}

\noindent \lc{g ∘ f}, ``önce \lc{f}, sonra \lc{g}'' fonksiyonudur:
\lc{(g ∘ f) a = g (f a)}. Bu eşitlik tanımsal olduğu için \leantac{show}
ile hedefi açık biçimde yazabildik.

\subsection{Tersler}

\leandosya[Sol ters ve birebirlik]{Bolum08}{sol_ters}

\noindent İspattaki üç yeniden yazma adımı şu hesabı yapar:
$a = g(f(a)) = g(f(b)) = b$. Tersine, sağ tersi olan bir fonksiyon
örtendir (Alıştırma~\ref{alis:sag-ters}). Bu sonuçların \emph{tersleri}
ise ilginç bir biçimde asimetriktir (\cref{tab:tersler}).

\begin{table}[htbp]
\centering
\renewcommand{\arraystretch}{1.5}
\rowcolors{2}{tablosatir}{white}
\begin{tabular}{>{\raggedright\arraybackslash}p{5.2cm}
                >{\raggedright\arraybackslash}p{6.4cm}}
\rowcolor{tabloust}
\tablobaslik{Önerme} & \tablobaslik{Ne gerekir?} \\
sol tersi var $\Rightarrow$ birebir      & yalnızca mantık (yapıcı) \\
sağ tersi var $\Rightarrow$ örten        & yalnızca mantık (yapıcı) \\
birebir $\Rightarrow$ sol tersi var      & $\alpha$ boş olmamalı ve klasik mantık \\
örten $\Rightarrow$ sağ tersi var        & \textbf{seçim aksiyomu}: her $b$ için bir $a$ \emph{seçmek} \\
\end{tabular}
\caption{Tersler ve fonksiyon türleri.}
\label{tab:tersler}
\end{table}

\noindent Son satır özellikle önemlidir. $f$ örtense, her $b$ için
$f(a) = b$ olan \emph{bir} $a$ vardır; ama sağ tersi kurmak için her $b$
için \emph{belirli bir} $a$ seçmemiz gerekir. Sonsuz sayıda seçimi tek
seferde yapmak, seçim aksiyomunun tam olarak izin verdiği şeydir.
\Mathlib{}'de bu ters, \leankw{noncomputable} olarak işaretlenmiş
\lc{Function.surjInv} fonksiyonudur. Seçim aksiyomunu \cref{bol:aksiyomlar}'da
ayrıntılı olarak inceleyeceğiz.

% ---------------------------------------------------------------------
\section{Cantor teoremi}\label{sec:cantor}

\begin{teorem}{Cantor, 1891}{cantor}
Hiçbir $\alpha$ tipi için $\alpha$'dan $\mathsf{Set}\;\alpha$'ya örten bir
fonksiyon yoktur. Yani bir kümenin alt kümeleri, kümenin kendisinden
\emph{daha çoktur}.
\end{teorem}

\noindent Georg Cantor'un ünlü \emph{köşegen yöntemi}, berber
paradoksuyla aynı yapıya sahiptir. $f : \alpha \to \mathsf{Set}\;\alpha$
herhangi bir fonksiyon olsun ve şu kümeyi düşünelim:
\[
  D \;=\; \{x \mid x \notin f(x)\}.
\]
$D$, ``kendi görüntüsüne ait olmayan'' elemanların kümesidir. $f$ örten
olsaydı, $f(a) = D$ olan bir $a$ bulunurdu. O zaman $a \in D$ ancak ve
ancak $a \notin f(a) = D$ olurdu: Berber kendini tıraş eder mi?
(\cref{fig:kosegen})

\begin{figure}[htbp]
\centering
\begin{tikzpicture}[
    hucre/.style={draw=black!20, minimum size=0.75cm, font=\small},
    kosegen/.style={hucre, fill=vurgu!20, draw=vurgu, thick},
  ]
  \foreach \x/\v in {1/a_1, 2/a_2, 3/a_3, 4/a_4}
    \node[font=\small] at ({0.75*\x},0.75) {$\v$};
  \foreach \y/\v in {1/{f(a_1)}, 2/{f(a_2)}, 3/{f(a_3)}, 4/{f(a_4)}}
    \node[font=\small, anchor=east] at (0.25,{-0.75*(\y-1)}) {$\v$};
  \foreach \y/\satir in {1/{1,0,1,1}, 2/{0,0,1,0}, 3/{1,1,1,0}, 4/{0,1,0,0}} {
    \foreach \v [count=\x] in \satir {
      \ifnum\x=\y
        \node[kosegen] at ({0.75*\x},{-0.75*(\y-1)}) {$\v$};
      \else
        \node[hucre] at ({0.75*\x},{-0.75*(\y-1)}) {$\v$};
      \fi
    }
  }
  \node[font=\small, anchor=east] at (0.25,-3.25) {$D$};
  \foreach \v [count=\x] in {0,1,0,1}
    \node[hucre, fill=bittibg, draw=bittiframe, thick] at ({0.75*\x},-3.25) {$\v$};
  \draw[-{Stealth[length=2.4mm]}, thick, vurgu] (3.6,-1.1) to[bend left=30] (3.6,-3.1);
  \node[font=\baslikfont\scriptsize, text=vurgu, align=left, anchor=west] at (3.9,-2.0)
       {köşegeni\\tersine çevir};
  \node[font=\baslikfont\scriptsize, text=black!60, align=left, anchor=west] at (5.6,-0.6)
       {1: $a_j \in f(a_i)$\\0: $a_j \notin f(a_i)$};
  \node[font=\baslikfont\scriptsize, text=bittiframe, align=left, anchor=west] at (5.6,-3.25)
       {$D$, her satırdan en az\\bir yerde (köşegende) farklıdır};
\end{tikzpicture}
\caption{Cantor'un köşegen yöntemi. $D$ kümesi, $i$. satırla $i$.
sütunda farklı olacak biçimde kurulur; bu yüzden hiçbir satıra eşit
değildir.}
\label{fig:kosegen}
\end{figure}

\leandosya[Cantor teoremi]{Bolum08}{cantor}

\begin{tarihce}[Tarihçe: Sonsuzluğun büyüklükleri]
Cantor 1891'de yayımladığı kısa bir makalede köşegen yöntemini tanıttı
ve bir kümenin alt kümelerinin her zaman kümeden ``daha büyük'' olduğunu
gösterdi \citep{cantor1891}. Bu sonuç, farklı büyüklükte sonsuzlukların
var olduğu anlamına geliyordu ve döneminde büyük tartışmalara yol açtı.
Aynı köşegen fikri sonraki yarım yüzyılda Russell paradoksunda (1901),
Gödel'in eksiklik teoreminde (1931) ve Turing'in durma probleminde (1936)
yeniden karşımıza çıkacaktı.
\end{tarihce}

% ---------------------------------------------------------------------
\section{Taktik özeti}

\begin{center}
\renewcommand{\arraystretch}{1.45}
\rowcolors{2}{tablosatir}{white}
\begin{tabularx}{\linewidth}{>{\raggedright\arraybackslash}p{4.4cm} X}
\rowcolor{tabloust}
\tablobaslik{Taktik} & \tablobaslik{Ne yapar?} \\
\leantac{ext x}                    & Küme (ya da fonksiyon) eşitliğini elemanlara indirger. \\
\leantac{simp only [l₁, l₂]}       & Yalnızca verilen eşitliklerle sadeleştirir. \\
\leantac{tauto}                    & Önermeler mantığı totolojilerini kanıtlar (klasik). \\
\leantac{ring}                     & Değişmeli halkalarda polinom eşitlikleri. \\
\leantac{linarith}                 & Doğrusal eşitsizlikler. \\
\leantac{norm\_num}                & Sayısal hesaplar, asallık. \\
\leantac{positivity}               & Pozitiflik ve negatif olmama. \\
\lc{exact?}, \lc{apply?}           & Kütüphanede teorem arar. \\
\end{tabularx}
\end{center}

% ---------------------------------------------------------------------
\begin{arastirmanotu}
\textbf{Köşegen argümanlarının ortak yapısı.} F.\,William Lawvere 1969'da,
Cantor, Russell, Gödel ve Turing'in köşegen argümanlarının hepsinin
kategori kuramındaki tek bir sabit nokta teoreminin örnekleri olduğunu
gösterdi \citep{lawvere1969}. Bu sonucun matematik altyapısı ağır
olmayan, okunabilir bir sunumu için \citet{yanofsky2003}'ya bakılabilir.
Lawvere'in teoremi \Lean{}'de formalize edilmeye elverişli, kısa ve
öğretici bir projedir.

\textbf{Mathlib ve kütüphane tasarımı.} \Mathlib{}'in yapısı, adlandırma
düzeni ve ``tip sınıfı'' hiyerarşisi \citet{mathlib2020}'de anlatılır.
Büyük bir formal kütüphanede kavramların hangi genellikte tanımlanacağı,
hangi teoremlerin \leantac{simp} lemması yapılacağı gibi sorular aktif
tartışma konularıdır; bu konulara \ciltiki{9}'da döneceğiz.

\textbf{İleri okuma.} \Mathlib{}'le matematik yapmaya kapsamlı bir giriş
için \citet{mil}; kümeler ve fonksiyonlar bölümü bu bölümün doğal
devamıdır.
\end{arastirmanotu}

\begin{bolumozeti}
\begin{itemize}
  \item \Mathlib{}, \Lean{}'in büyük matematik kütüphanesidir. Teoremler
        katı bir adlandırma düzeni izler; \lc{exact?} ve Loogle arama
        araçlarıdır.
  \item \leantac{ring}, \leantac{linarith}, \leantac{norm\_num},
        \leantac{positivity}, \leantac{tauto} rutin işleri otomatikleştirir.
  \item Bir küme, bir yüklemdir: \lc{Set α := α → Prop}. Küme işlemleri
        mantık bağlaçlarıdır; küme cebiri önermeler mantığıdır.
  \item Küme eşitliği \leantac{ext} ile elemanlara indirgenir.
  \item Denklik ilişkileri sınıflara ayırır; kısmi sıralamalar Hasse
        şemalarıyla çizilir.
  \item Sol ters birebirliği, sağ ters örtenliği verir. Örten bir
        fonksiyonun sağ tersini kurmak seçim aksiyomu ister.
  \item Cantor teoremi: Hiçbir küme kendi alt kümeleri kadar büyük
        değildir. Köşegen yöntemi, berber paradoksunun genellemesidir.
\end{itemize}
\end{bolumozeti}

\begin{alistirmalar}
\alistirma{1} \lc{s ∩ s = s}. (\lc{a8\_1})

\alistirma{1} \lc{s ⊆ s ∪ t}. (\lc{a8\_2})

\alistirma{2} Birleşimin kesişim üzerine dağılması:
\lc{s ∪ (t ∩ u) = (s ∪ t) ∩ (s ∪ u)}. Önce elle, sonra \leantac{simp only}
ve \leantac{tauto} ile ispatlayın. (\lc{a8\_3})

\alistirma{2}\label{alis:kume-demorgan} (Klasik) Kesişimin tümleyeni:
\lc{(s ∩ t)ᶜ = sᶜ ∪ tᶜ}. Hangi yön klasik mantık gerektirir? Bölüm~3 ile
ilişkisini açıklayın. (\lc{a8\_4})

\alistirma{2} \lc{s ⊆ t → s ∩ u ⊆ t ∩ u}. (\lc{a8\_5})

\alistirma{2} Bölünebilirliğin geçişliliğini, \lc{dvd\_trans} kullanmadan
tanımdan ispatlayın. \lc{a ∣ b}, \lc{∃ c, b = a * c} demektir.
(\lc{a8\_6})

\alistirma{2} \lc{g ∘ f} örtense \lc{g} örtendir. (\lc{a8\_7})

\alistirma{2} \lc{g ∘ f} birebirse \lc{f} birebirdir. (\lc{a8\_8})

\alistirma{3}\label{alis:sag-ters} Sağ tersi olan fonksiyon örtendir:
\lc{(∀ b, f (k b) = b) → Surjective f}. (\lc{a8\_9})

\alistirma{3} Tam sayılar üzerinde ``farkı 3'ün katı olmak''
ilişkisinin bir denklik ilişkisi olduğunu ispatlayın. Kaç denklik sınıfı
vardır? (\lc{a8\_10})

\alistirma{3} Köşegen yöntemi: \lc{ℕ}'den \lc{ℕ → ℕ}'ye örten fonksiyon
yoktur. İpucu: \lc{fun n => F n n + 1} fonksiyonunu düşünün. Bu sonuç,
``bütün programlar sayılabilir ama bütün fonksiyonlar sayılamaz''
gözleminin özüdür. (\lc{a8\_11})

\alistirma{2} (Kâğıt-kalem) $\{1, 2, 3, 4, 6, 8, 12, 24\}$ kümesi için
bölünebilirlik sıralamasının Hasse şemasını çizin.

\alistirma{3} (Tartışma) \Cref{tab:tersler}'ün üçüncü satırında
``$\alpha$ boş olmamalı'' koşulu neden gerekir? $\alpha$ boşsa ve
$\beta$ boş değilse ne olur?
\end{alistirmalar}
   % Kümeler, ilişkiler, fonksiyonlar
% =====================================================================
%  Bölüm 9 — Tümevarımsal Tipler
% =====================================================================
\chapter{Tümevarımsal Tipler}\label{bol:tumevarimsal-tipler}

\begin{bolumhedefleri}
\begin{itemize}
  \item \leankw{inductive} bildiriminin parçalarını: kurucular,
        parametreler ve indeksler,
  \item \lc{And}, \lc{Or}, \lc{True}, \lc{False}, \lc{Exists} ve
        \lc{Eq}'nun birer tümevarımsal tip olduğunu, onları kendimiz
        tanımlayarak,
  \item \leantac{cases} taktiğinin her kurucu için bir durum açtığını ve
        bunun eleme kurallarını nasıl ürettiğini,
  \item Listeler ve ikili ağaçlar üzerinde özyineleme ve yapısal
        tümevarımı,
  \item Tümevarımsal yüklemleri ve \emph{kural tümevarımını},
  \item Pozitiflik koşulunu ve neden gerekli olduğunu.
\end{itemize}
\end{bolumhedefleri}

Bölüm~7'de doğal sayıları \leankw{inductive} ile tanımladık ve Peano
aksiyomlarının bu tanımdan teorem olarak çıktığını gördük. Bu bölümde aynı
fikri sonuna kadar götürüyoruz. Göreceğiz ki \Lean{}'in mantığında
\emph{temel} olan çok az şey vardır: Fonksiyon tipleri ($\forall$, $\to$)
ve tümevarımsal tipler. Bölüm~2'den beri kullandığımız \lc{∧}, \lc{∨},
\lc{False}, \lc{∃} ve hatta eşitlik, hepsi \leankw{inductive} ile
tanımlanmış sıradan tiplerdir.

% ---------------------------------------------------------------------
\section{\texorpdfstring{\texttt{inductive}}{inductive}'in anatomisi}\label{sec:inductive-anatomi}

Bir tümevarımsal tip bildirimi, bir tipin adını ve o tipin elemanlarını
kurmanın \emph{bütün} yollarını, yani \textbf{kurucuları} verir
(\cref{fig:inductive-anatomi}).

\begin{figure}[htbp]
\centering
\begin{tikzpicture}[
    parca/.style={font=\leanfont\small, inner sep=2pt},
    etiket/.style={font=\baslikfont\scriptsize, align=center, text width=2.8cm},
  ]
  \node[parca, anchor=west] (l1) at (0,0)    {\textcolor{leankw}{\bfseries inductive} \textcolor{koyu}{Liste} \textcolor{leanhyp}{(α : Type)} \textcolor{leankw}{\bfseries where}};
  \node[parca, anchor=west] (l2) at (0.4,-0.65) {| \textcolor{vurgu}{bos}  : Liste α};
  \node[parca, anchor=west] (l3) at (0.4,-1.3) {| \textcolor{vurgu}{ekle} : \textcolor{ikincil}{α → Liste α} → Liste α};
  \draw[koyu, thick, rounded corners=2pt] (1.6,0.22) rectangle (2.45,-0.22);
  \draw[leanhyp, thick, rounded corners=2pt] (2.5,0.22) rectangle (3.95,-0.22);
  \draw[vurgu, thick, rounded corners=2pt] (0.6,-0.43) rectangle (1.15,-0.87);
  \draw[vurgu, thick, rounded corners=2pt] (0.6,-1.08) rectangle (1.25,-1.52);
  \draw[ikincil, thick, rounded corners=2pt] (1.55,-1.08) rectangle (3.45,-1.52);
  \node[etiket, text=koyu] (e1) at (2.0,1.25) {tipin adı};
  \node[etiket, text=leanhyp] (e2) at (5.4,1.0) {\textbf{parametre}: bütün kurucularda sabit};
  \node[etiket, text=vurgu] (e3) at (-2.3,-1.0) {\textbf{kurucular}: elemanları kurmanın tek yolları};
  \node[etiket, text=ikincil] (e4) at (6.3,-1.3) {\textbf{alanlar}: kurucunun aldığı argümanlar; \lc{Liste α} geçtiği için \emph{özyinelemeli}};
  \draw[koyu, -{Stealth[length=2mm]}] (e1) -- (2.02,0.25);
  \draw[leanhyp, -{Stealth[length=2mm]}] (e2) -- (3.6,0.25);
  \draw[vurgu, -{Stealth[length=2mm]}] (e3.east) -- (0.55,-0.65);
  \draw[vurgu, -{Stealth[length=2mm]}] (e3.east) -- (0.55,-1.3);
  \draw[ikincil, -{Stealth[length=2mm]}] (e4.west) -- (3.5,-1.3);
\end{tikzpicture}
\caption{Bir tümevarımsal tip bildiriminin parçaları.}
\label{fig:inductive-anatomi}
\end{figure}

\noindent Yalnızca alansız kuruculardan oluşan en basit tipler
\emph{numaralandırma} tipleridir:

\leandosya[Numaralandırma tipi]{Bolum09}{renk}

\noindent \lc{cases r <;> rfl} ispatı, üç durumun her birini ayrı ayrı
hesaplar. Sonlu ve küçük tipler üzerindeki önermeler böyle
``kaba kuvvetle'' ispatlanabilir.

Bir kurucunun alanları olabilir. \lc{Belki} tipi, ``bir değer ya var ya
yok'' durumunu temsil eder; programlamada sıfıra bölme gibi hataları
güvenle ifade etmenin yoludur:

\leandosya[Ya bir değer ya hiç]{Bolum09}{belki}

% ---------------------------------------------------------------------
\section{Mantık bağlaçları tümevarımsal tiplerdir}\label{sec:baglac-inductive}

Şimdi bölümün ana fikrine geliyoruz. Bölüm~2--6'da kullandığımız bütün
bağlaçları, \Lean{}'in kütüphanesine bakmadan kendimiz tanımlayabiliriz:

\leandosya[Kendi bağlaçlarımız]{Bolum09}{baglaclar}

\noindent Her tanım, ilgili bağlacın \textbf{giriş kurallarını} kurucu
olarak listeler:

\begin{table}[htbp]
\centering
\renewcommand{\arraystretch}{1.5}
\rowcolors{2}{tablosatir}{white}
\begin{tabular}{>{\centering\arraybackslash}p{1.6cm}
                >{\raggedright\arraybackslash}p{3.4cm}
                >{\centering\arraybackslash}p{1.6cm}
                >{\raggedright\arraybackslash}p{4.6cm}}
\rowcolor{tabloust}
\tablobaslik{Bağlaç} & \tablobaslik{Kurucular (giriş)} & \tablobaslik{Durum} & \tablobaslik{\texttt{cases} ne verir? (eleme)} \\
$P \wedge Q$       & \lc{intro : P → Q → \_}             & 1 & iki kanıt: \lc{hp}, \lc{hq} \\
$P \vee Q$         & \lc{sol : P → \_}, \lc{sag : Q → \_}  & 2 & ya \lc{hp} ya \lc{hq} \\
$\top$             & \lc{tanik : \_}                      & 1 & hiçbir şey (bilgi yok) \\
$\bot$             & (yok)                               & 0 & hedef hemen kapanır \\
$\exists x,\,P\,x$ & \lc{intro (w) (h : P w)}            & 1 & bir tanık \lc{w} ve \lc{h : P w} \\
$a = b$            & \lc{refl : a = a}                   & 1 & \lc{b} yerine \lc{a} konur \\
\end{tabular}
\caption{Mantık bağlaçları ve kurucuları. Giriş kuralları kuruculardır;
eleme kuralları, kurucular üzerinde \leantac{cases} ile otomatik olarak
elde edilir.}
\label{tab:baglac-kurucu}
\end{table}

\begin{onemlifikir}[Önemli Fikir: Eleme kuralları kendiliğinden gelir]
Doğal çıkarımda her bağlaç için giriş \emph{ve} eleme kurallarını ayrı
ayrı yazmıştık. Tümevarımsal tiplerde ise yalnızca giriş kurallarını
(kurucuları) yazarız; eleme kuralı, \Lean{}'in ürettiği özyineleyiciden
kendiliğinden çıkar. \leantac{cases} ``bu kanıt hangi kurucuyla kurulmuş
olabilir?'' diye sorar ve her olasılık için bir durum açar. Bölüm~5'te
sözünü ettiğimiz giriş--eleme \emph{uyumu} (harmony), tümevarımsal
tiplerde tasarım gereği sağlanır.
\end{onemlifikir}

\leandosya[Kendi bağlaçlarımızla ispatlar]{Bolum09}{baglac_ispat}

\noindent \lc{Yanlis.elim'} ispatına dikkat edin: \lc{Yanlis}'ın hiç
kurucusu olmadığı için \leantac{cases h} \emph{sıfır} durum açar ve
ispat biter. Bölüm~3'teki ``ex falso'' ilkesinin tip kuramındaki açıklaması
budur (\cref{fig:kurucular}). \Lean{}'in kendi tanımları da bizimkilerle
aynıdır:

\leandosya[\texttt{Or}'un tanımı]{Bolum09}{print_or}

\subsection{Eşitlik de bir tümevarımsal tiptir}

En şaşırtıcı örnek eşitliktir. \lc{a = b} önermesi, tek kurucusu
\lc{refl : a = a} olan bir tiptir:

\leandosya[Eşitliği tanımlamak]{Bolum09}{esit}

\noindent Bu tanımda \lc{a} bir \emph{parametre}, ikinci argüman ise bir
\emph{indekstir}. Fark önemlidir:

\begin{tanim}{Parametre ve indeks}{parametre-indeks}
Bir tümevarımsal tip bildiriminde, bütün kurucularda \emph{aynı} kalan
argümanlara \textbf{parametre}, kurucudan kurucuya \emph{değişebilen}
argümanlara \textbf{indeks} denir. Parametreler iki noktadan önce, indeksler
tip imzasında (\lc{: α → Prop}) yazılır. İndeksli tiplere
\textbf{tümevarımsal aileler} denir.
\end{tanim}

\lc{Esit a} bir ailedir: \lc{Esit a b} tipleri, \lc{b}'ye göre değişir.
Ama ailenin tek kurucusu yalnızca \lc{Esit a a} tipinde bir eleman
üretir. Bu yüzden \lc{h : Esit a b} üzerinde \leantac{cases} yaptığımızda
\Lean{}, \lc{b}'nin \lc{a} olması gerektiğini anlar ve her yerde \lc{b}'yi
\lc{a} ile değiştirir. Bölüm~6'daki yerine koyma kuralı ($=$E), bu
eleme ilkesinin kendisidir.

% ---------------------------------------------------------------------
\section{Listeler}\label{sec:listeler}

Liste, programlamanın en temel veri yapısıdır. Bir liste ya boştur ya da
bir baş eleman ile bir kuyruk listeden oluşur:

\leandosya[Listenin tanımı]{Bolum09}{liste_tanim}

\begin{figure}[htbp]
\centering
\begin{tikzpicture}[
    hucre/.style={draw=anarenk, fill=tanimbg, thick, minimum width=0.9cm,
                  minimum height=0.75cm, font=\small},
    isaret/.style={draw=anarenk, fill=white, thick, minimum width=0.5cm,
                   minimum height=0.75cm},
    ok/.style={-{Stealth[length=2.2mm]}, thick, leankw},
  ]
  \foreach \i/\v in {0/1, 1/2, 2/3} {
    \node[hucre] (v\i) at ({2.6*\i},0) {$\v$};
    \node[isaret, right=0pt of v\i] (p\i) {};
    \node[font=\leanfont\scriptsize, text=leankw, above=2pt of v\i] {ekle};
  }
  \node[draw=vurgu, fill=teorembg, thick, font=\leanfont\small, minimum height=0.75cm] (b) at (7.8,0) {bos};
  \draw[ok] (p0.center) -- (v1.west);
  \draw[ok] (p1.center) -- (v2.west);
  \draw[ok] (p2.center) -- (b.west);
  \node[font=\leanfont\scriptsize, text=black!60, below=0.45cm of v1, anchor=north]
       {ekle 1 (ekle 2 (ekle 3 bos))};
\end{tikzpicture}
\caption{\lc{[1, 2, 3]} listesi: her \lc{ekle} hücresi bir değer ve
kuyruğa bir işaret taşır.}
\label{fig:liste}
\end{figure}

\noindent Listeler üzerindeki fonksiyonlar, doğal sayılarda olduğu gibi
özyinelemeyle tanımlanır. Aslında \lc{Dogal} tipi, elemanları önemsiz
olan bir listedir: \lc{sifir} \lc{bos}'a, \lc{ard} da \lc{ekle}'ye karşılık
gelir. Uzunluk fonksiyonu bu benzerliği tam olarak ortaya koyar.

\leandosya[Liste fonksiyonları]{Bolum09}{liste_fonk}

\subsection{Yapısal tümevarım}

Bir liste özelliğini ispatlamak için \emph{yapısal tümevarım}
kullanırız: Özelliği boş liste için gösteririz; ardından \lc{xs} için
doğru olduğunu varsayıp \lc{ekle x xs} için gösteririz. Doğal sayılardaki
asimetri burada da görülür: \lc{birlestir bos ys = ys} tanımdan çıkar,
ama \lc{birlestir xs bos = xs} tümevarım ister.

\leandosya[Boş listeyle birleştirme]{Bolum09}{birlestir_bos}
\leandosya[Uzunluk ve birleştirme]{Bolum09}{uzunluk_birlestir}
\leandosya[Birleştirmenin birleşme özelliği]{Bolum09}{birlestir_birlesme}

\noindent Şimdi listelerin en ünlü teoremlerinden birine geliyoruz:
Bir listeyi iki kez ters çevirmek onu değiştirmez. Önce bir önsav
gerekir: Birleştirilmiş bir listenin tersi, terslerin \emph{ters sırada}
birleştirilmesidir.

\leandosya[Birleştirmenin tersi]{Bolum09}{ters_birlestir}
\leandosya[İki kez ters çevirmek]{Bolum09}{ters_ters}

\noindent \lc{ters\_ters} ispatındaki \leankw{calc} bloğu, kâğıt üzerinde
yazacağımız eşitlik zincirinin aynısıdır. Birinci ve dördüncü adımlar
\lc{rfl}'dir, yani yalnızca tanımları açar. İkinci adım önsavı, üçüncü
adım tümevarım hipotezini kullanır.

\begin{dikkat}[Dikkat: Programın doğruluğu bir teoremdir]
\lc{ters\_ters}, bir \emph{programın} doğruluğu hakkında bir teoremdir.
Yazılım testleri bir fonksiyonu sonlu sayıda girdiyle denetler; bu
teorem ise \emph{bütün} listeler için doğruluğu kanıtlar. İspat
asistanlarının yazılım doğrulamadaki gücü buradan gelir. Ama dikkat:
Teorem yalnızca \emph{söylediği şeyi} kanıtlar. Hiçbir şey yapmayan
\lc{fun xs => xs} fonksiyonu da \lc{ters\_ters}'i sağlar! Doğru
özellikleri seçmek, doğrulamanın en zor kısmıdır.
\end{dikkat}

% ---------------------------------------------------------------------
\section{İkili ağaçlar}\label{sec:agaclar}

Bir kurucunun \emph{birden fazla} özyinelemeli alanı olabilir. İkili
ağaçlarda her düğüm, bir sol ve bir sağ alt ağaç taşır:

\leandosya[İkili ağaç]{Bolum09}{agac}
\leandosya[Boyut ve ayna]{Bolum09}{agac_fonk}

\begin{figure}[htbp]
\centering
\begin{tikzpicture}[
    d/.style={circle, draw=anarenk, fill=tanimbg, thick, minimum size=0.7cm, font=\small},
    ypr/.style={rectangle, fill=black!25, minimum size=0.18cm, inner sep=0pt},
    level distance=1.0cm,
    level 1/.style={sibling distance=2.0cm},
    level 2/.style={sibling distance=1.0cm},
    edge from parent/.style={draw, black!45, thick},
  ]
  \node[d] at (0,0) {4}
    child {node[d] {2}
      child {node[d] {1} child {node[ypr] {}} child {node[ypr] {}}}
      child {node[ypr] {}}}
    child {node[d] {7}
      child {node[ypr] {}}
      child {node[ypr] {}}};
  \node[font=\baslikfont\small\bfseries, text=anarenk] at (0,0.75) {$t$};

  \draw[-{Stealth[length=3mm]}, very thick, vurgu] (2.6,-1.2) -- node[above, font=\leanfont\small, text=vurgu]{ayna} (4.6,-1.2);

  \begin{scope}[xshift=7.2cm]
  \node[d, draw=vurgu, fill=teorembg] at (0,0) {4}
    child {node[d, draw=vurgu, fill=teorembg] {7}
      child {node[ypr] {}}
      child {node[ypr] {}}}
    child {node[d, draw=vurgu, fill=teorembg] {2}
      child {node[ypr] {}}
      child {node[d, draw=vurgu, fill=teorembg] {1} child {node[ypr] {}} child {node[ypr] {}}}};
  \node[font=\baslikfont\small\bfseries, text=vurgu] at (0,0.75) {\lc{ayna} $t$};
  \end{scope}
\end{tikzpicture}
\caption{Bir ağaç ve aynası. Gri kareler \lc{yaprak}'tır. Ayna, her
düğümde sol ve sağ alt ağaçları yer değiştirir.}
\label{fig:agac-ayna}
\end{figure}

\noindent Bir ağaç üzerinde tümevarım yaparken \lc{dugum} durumunda
\emph{iki} tümevarım hipotezi elde ederiz: biri sol, biri sağ alt ağaç
için.

\leandosya[Aynanın aynası]{Bolum09}{ayna_ayna}
\leandosya[Ayna boyutu korur]{Bolum09}{boyut_ayna}

\begin{onemlifikir}[Önemli Fikir: Tümevarım ilkesi tipin şeklini izler]
Her tümevarımsal tipin tümevarım ilkesi, o tipin tanımından mekanik
olarak okunur: Her kurucu için bir durum vardır ve kurucunun her
özyinelemeli alanı için bir tümevarım hipotezi gelir. Doğal sayılarda
bir (\lc{ard}), listelerde bir (\lc{ekle}), ağaçlarda iki (\lc{dugum})
hipotez vardır. Bölüm~7'deki domino resmi, ağaçlarda dallanan bir
domino dizisine dönüşür.
\end{onemlifikir}

% ---------------------------------------------------------------------
\section{Tümevarımsal yüklemler}\label{sec:tumevarimsal-yuklem}

Tümevarımsal tipler yalnızca veri değil, \emph{önermeler} de
tanımlayabilir. ``$n$ çifttir'' önermesini, bir sayının çift olduğunu
kanıtlamanın kurallarını vererek tanımlayalım:

\leandosya[Çift sayılar]{Bolum09}{cift_tanim}

\noindent Bu tanım, iki çıkarım kuralına karşılık gelir. \lc{Cift 4}'ün
bir kanıtı, bu kurallarla kurulmuş bir türetim ağacıdır
(\cref{fig:cift-turetim}).

\begin{figure}[htbp]
\centering
\begin{tabular}{c@{\hspace{4em}}c@{\hspace{4em}}c}
\begin{prooftree}
  \infer0[\lc{sifir}]{\mathsf{Cift}\;0}
\end{prooftree}
&
\begin{prooftree}
  \hypo{\mathsf{Cift}\;n}
  \infer1[\lc{iki\_ekle}]{\mathsf{Cift}\;(n+2)}
\end{prooftree}
&
\begin{prooftree}
  \infer0[\lc{sifir}]{\mathsf{Cift}\;0}
  \infer1[\lc{iki\_ekle}]{\mathsf{Cift}\;2}
  \infer1[\lc{iki\_ekle}]{\mathsf{Cift}\;4}
\end{prooftree}
\\[10pt]
\small kural 1 & \small kural 2 & \small \lc{Cift 4}'ün türetimi
\end{tabular}
\caption{\lc{Cift} yükleminin kuralları ve \lc{Cift 4}'ün türetimi.
\lc{Cift 1}'in hiçbir türetimi yoktur: Ne kural 1'in ne de kural 2'nin
sonucu 1'dir.}
\label{fig:cift-turetim}
\end{figure}

\leandosya[Çiftlik ispatları]{Bolum09}{cift_ornek}

\noindent \lc{bir\_cift\_degil} ispatında \leantac{cases h}, \lc{Cift 1}
kanıtının hangi kuralla kurulmuş olabileceğini sorar. \lc{sifir} kuralı
\lc{Cift 0} üretir; \lc{1 = 0} olmadığından elenir. \lc{iki\_ekle} kuralı
\lc{Cift (n + 2)} üretir; \lc{1 = n + 2} olacak bir doğal sayı
olmadığından o da elenir. Hiç durum kalmaz.

\leandosya[Her $2n$ çifttir]{Bolum09}{cift_iki_kat}

\subsection{Kural tümevarımı}

Tersini göstermek, yani \lc{Cift n} ise \lc{n}'nin bir \lc{2 * k} olduğunu
ispatlamak için, \lc{n} üzerinde değil, \lc{Cift n} \emph{kanıtı} üzerinde
tümevarım yaparız. Buna \textbf{kural tümevarımı} denir: Türetim ağacının
her kuralı için bir durum açılır.

\leandosya[Kural tümevarımı]{Bolum09}{kural_tumevarimi}

\noindent Kural tümevarımı, Kısım~III'ün temel aracı olacak.
\cref{bol:gecerlilik}'de bir ispat sistemini tümevarımsal bir yüklem
olarak tanımlayacak ve geçerlilik teoremini tam olarak bu yöntemle,
türetimin son kuralına göre durum ayırarak ispatlayacağız.

\subsection{Yansımalı-geçişli kapanış}

Son örneğimiz, bilgisayar biliminin her yerinde karşımıza çıkan bir
yapıdır. $R$ bir ``tek adım'' ilişkisi olsun (bir programın bir adım
çalışması, bir oyunda bir hamle gibi). ``$a$'dan $c$'ye sıfır ya da daha
fazla adımda ulaşılabilir'' ilişkisi, $R$'nin \emph{yansımalı-geçişli
kapanışıdır}:

\leandosya[Yansımalı-geçişli kapanış]{Bolum09}{kapanis}

\begin{figure}[htbp]
\centering
\begin{tikzpicture}[
    d/.style={circle, draw=anarenk, fill=tanimbg, thick, minimum size=0.75cm, font=\small},
    ok/.style={-{Stealth[length=2.4mm]}, thick, #1},
  ]
  \foreach \i/\n in {0/a, 1/b, 2/c, 3/d} \node[d] (n\i) at ({2.2*\i},0) {$\n$};
  \draw[ok=leanhyp] (n0) -- node[above, font=\scriptsize]{$R$} (n1);
  \draw[ok=leanhyp] (n1) -- node[above, font=\scriptsize]{$R$} (n2);
  \draw[ok=leanhyp] (n2) -- node[above, font=\scriptsize]{$R$} (n3);
  \draw[ok=ikincil, dashed] (n0) to[bend right=30] node[below, font=\leanfont\scriptsize, text=ikincil]{Kapanis R a d} (n3);
  \draw[ok=vurgu] (n0) to[loop above, looseness=6] node[above, font=\leanfont\scriptsize, text=vurgu]{refl} (n0);
\end{tikzpicture}
\caption{$R$-adımlarından oluşan bir yol. \lc{Kapanis R a d}, $a$'dan $d$'ye
böyle bir yolun var olduğunu söyler; \lc{refl}, sıfır adımlık yoldur.}
\label{fig:kapanis}
\end{figure}

\noindent \lc{Kapanis.gecisme} ispatı, iki yolun uç uca eklenebileceğini
gösterir. Bu ispat da kural tümevarımıdır: İlk yolun uzunluğu üzerinde
tümevarım yaparız. İspatın programlama karşılığı, iki bağlı listeyi
birleştirmektir.

% ---------------------------------------------------------------------
\section{Pozitiflik koşulu}\label{sec:pozitiflik}

\Lean{} her \leankw{inductive} bildirimini kabul etmez. Aşağıdaki tanımı
düşünün:

\begin{leankod}[Reddedilen bir tanım (derlenmez)]
inductive Kotu where
  | yap : (Kotu → False) → Kotu
\end{leankod}

\begin{leanhata}[Lean hata mesajı]
(kernel) arg #1 of 'Kotu.yap' has a non positive occurrence of the datatypes being declared
\end{leanhata}

\noindent Bu tanım kabul edilseydi bir çelişki türetebilirdik. Kabaca:
\lc{f : Kotu → False} fonksiyonunu ``\lc{yap g} verildiğinde
\lc{g (yap g)} döndür'' diye tanımlayalım. O zaman \lc{f (yap f)} hesabı,
\lc{f (yap f)}'ye döner ve sonsuza dek sürer; mantık olarak da
\lc{f (yap f) : False} bir kanıt olur. Bu, Bölüm~6'daki berber
paradoksunun ve Bölüm~8'deki Cantor köşegeninin bir başka yüzüdür.

\begin{tanim}{Kesin pozitiflik}{pozitiflik}
Bir tümevarımsal bildirimde, tanımlanan tip kurucuların argümanlarında
yalnızca \textbf{kesin pozitif} konumlarda geçebilir: Hiçbir okun
\emph{sol} tarafında bulunmamalıdır. \lc{ekle : α → Liste α → Liste α}
geçerlidir; \lc{yap : (Kotu → False) → Kotu} geçersizdir.
\end{tanim}

\noindent Pozitiflik koşulu ve Bölüm~7'deki sona erme koşulu, \Lean{}'in
mantığını tutarlı kılan iki önemli kısıttır. İkisi de aynı tehlikeye karşı
korur: \emph{Kendine gönderme yaparak} bir çelişki üretmek.

% ---------------------------------------------------------------------
\section{Curry--Howard: tümevarımsal tipler}\label{sec:curry-howard-5}

\begin{center}
\renewcommand{\arraystretch}{1.5}
\rowcolors{2}{tablosatir}{white}
\begin{tabular}{>{\raggedright\arraybackslash}p{3.3cm}
                >{\raggedright\arraybackslash}p{4.6cm}
                >{\raggedright\arraybackslash}p{3.8cm}}
\rowcolor{tabloust}
\tablobaslik{Mantık} & \tablobaslik{Programlama} & \tablobaslik{Lean} \\
$P \wedge Q$           & kayıt (çift)                       & \lc{structure}, \lc{And} \\
$P \vee Q$             & etiketli birleşim                   & \lc{Or}, \lc{Sum} \\
$\bot$                 & boş tip                            & \lc{False}, \lc{Empty} \\
$\top$                 & birim tip                          & \lc{True}, \lc{Unit} \\
$\exists x,\, P\,x$    & bağımlı çift                       & \lc{Exists}, \lc{Sigma} \\
$a = b$                & özdeşlik tipi                      & \lc{Eq} \\
tümevarımsal yüklem    & tümevarımsal aile                  & \lc{Cift}, \lc{Kapanis} \\
kural tümevarımı       & türetim üzerinde özyineleme        & \leantac{induction h} \\
\end{tabular}
\end{center}

% ---------------------------------------------------------------------
\section{Taktik özeti}

\begin{center}
\renewcommand{\arraystretch}{1.45}
\rowcolors{2}{tablosatir}{white}
\begin{tabularx}{\linewidth}{>{\raggedright\arraybackslash}p{4.4cm} X}
\rowcolor{tabloust}
\tablobaslik{Taktik} & \tablobaslik{Ne yapar?} \\
\leantac{cases h with | c₁ ... | c₂ ...} & Her kurucu için bir durum; imkânsız durumları eler. \\
\leantac{induction t with ...}          & Yapısal tümevarım; her özyinelemeli alan için bir hipotez. \\
\leantac{induction h with ...}          & Kural tümevarımı (\lc{h} bir tümevarımsal yüklemin kanıtı). \\
\lc{t <;> rfl}                          & Sonlu durumların her birini hesaplayarak kapatır. \\
\leankw{calc}                           & Eşitlik zinciri; adımlar \lc{rfl}, önsav ya da \leantac{rw} olabilir. \\
\end{tabularx}
\end{center}

% ---------------------------------------------------------------------
\begin{arastirmanotu}
\textbf{Tümevarımsal aileler.} İndeksli tümevarımsal tipler, Peter Dybjer
tarafından Martin-Löf tip kuramının bir uzantısı olarak kesin biçimde
tanımlandı \citep{dybjer1994}. Coq ve \Lean{}'in tümevarımsal tipleri,
Christine Paulin-Mohring'in İnşalar Hesabına eklediği tümevarımsal
tanımlara dayanır \citep{paulin1993}. \Lean{}'in çekirdeğindeki
tümevarımsal tiplerin kesin kuralları için \citet{carneiro2019}.

\textbf{Pozitifliğin ötesi.} Kesin pozitiflik güvenlidir ama bazen
gereğinden kısıtlayıcıdır. Karşılıklı (\emph{mutual}) ve iç içe
(\emph{nested}) tümevarımsal tipler, \emph{tümevarımsal-özyinelemeli}
tanımlar ve \emph{eşözyinelemeli} (\emph{coinductive}) tipler, bu
sınırları genişleten araştırma konularıdır. Homotopi tip kuramının
\emph{yüksek tümevarımsal tipleri}, kurucuların yalnızca elemanlar
değil, elemanlar arasında eşitlikler de üretebildiği bir genellemedir
\citep{hott2013}.
\end{arastirmanotu}

\begin{bolumozeti}
\begin{itemize}
  \item Bir \leankw{inductive} bildirimi, bir tipin bütün elemanlarını
        kurmanın yollarını (kurucuları) listeler. Parametreler sabittir,
        indeksler değişebilir.
  \item \lc{And}, \lc{Or}, \lc{True}, \lc{False}, \lc{Exists} ve \lc{Eq}
        tümevarımsal tiplerdir. Kurucular giriş kurallarıdır; eleme
        kuralları \leantac{cases} ile kendiliğinden gelir.
  \item Listeler ve ağaçlar üzerinde özyineleme ve yapısal tümevarım
        yapılır; tümevarım ilkesi tipin şeklini izler.
  \item Tümevarımsal yüklemler, kurallarla tanımlanan önermelerdir.
        Kanıtları türetim ağaçlarıdır; üzerlerinde kural tümevarımı
        yapılır.
  \item Pozitiflik koşulu, kendine gönderme yoluyla çelişki üretilmesini
        engeller.
\end{itemize}
\end{bolumozeti}

\begin{alistirmalar}
\alistirma{1} Kendi \lc{Ve} tipimizin sol elemesini ispatlayın. (\lc{a9\_1})

\alistirma{1} Kendi \lc{Veya} tipimiz için durum ayırma ilkesini
ispatlayın. (\lc{a9\_2})

\alistirma{2} \lc{(∀ x, P x → Q x) → Var P → Var Q}. (\lc{a9\_3})

\alistirma{2} Ters çevirme uzunluğu korur. İpucu: \lc{uzunluk\_birlestir}.
(\lc{a9\_4})

\alistirma{2} Bir ikili ağaçta yaprak sayısı, düğüm sayısından bir
fazladır. (\lc{a9\_5})

\alistirma{2} \lc{¬ Cift 3}. Kaç kez \leantac{cases} gerekti? Neden?
(\lc{a9\_6})

\alistirma{2} \emph{Ters çevirme}: \lc{Cift (n + 2) → Cift n}.
(\lc{a9\_7})

\alistirma{3} İki çift sayının toplamı çifttir. Kural tümevarımını hangi
kanıt üzerinde yapmalısınız? (\lc{a9\_8})

\alistirma{1} Tek bir \lc{R}-adımı kapanışın içindedir. (\lc{a9\_9})

\alistirma{2} (Kâğıt-kalem) Doğal sayılar üzerinde ``$n$ tektir''
yüklemini iki kuralla tanımlayın. \lc{Tek 5}'in türetim ağacını çizin.

\alistirma{3} (Kâğıt-kalem) Bir \emph{palindrom}, tersinden okunuşu
kendisine eşit olan listedir. Palindromları tümevarımsal bir yüklem
olarak tanımlayan üç kural yazın. İpucu: boş liste, tek elemanlı liste
ve ``\lc{x}, ortada bir palindrom, \lc{x}''.

\alistirma{3} (Tartışma) \cref{sec:pozitiflik}'deki \lc{Kotu} tipi
kabul edilseydi \lc{False}'un bir kanıtını nasıl yazardınız? Bu argümanı
Cantor teoremiyle karşılaştırın.
\end{alistirmalar}
   % Tümevarımsal tipler

\part{Mantığın Matematiği}
% =====================================================================
%  Bölüm 10 — Önermeler Mantığının Sözdizimi ve Anlambilimi
% =====================================================================
\chapter{Sözdizimi ve Anlambilim}\label{bol:sozdizim}

\begin{bolumhedefleri}
\begin{itemize}
  \item \emph{Nesne dili} ile \emph{üst dil} ayrımını ve Kısım~III'ün
        neden bu ayrıma dayandığını,
  \item Önermeler mantığının formüllerini \Lean{}'de bir tümevarımsal
        tip olarak tanımlamayı,
  \item Formüller üzerinde yapısal özyinelemeyi: boyut, atomlar,
        yerine koyma,
  \item Değerlemeleri ve Tarski'nin doğruluk tanımını,
  \item Totoloji, sağlanabilirlik ve mantıksal sonuç ($\models$)
        kavramlarını,
  \item Doğruluk tablolarına dayanan bir karar yöntemini ve sınırlarını,
  \item Doğruluk tablosu anlambiliminin \emph{klasik} olduğunu.
\end{itemize}
\end{bolumhedefleri}

% ---------------------------------------------------------------------
\section{Kısım III'e giriş: mantık hakkında mantık}\label{sec:kisim3-giris}

Şimdiye kadar mantığı \emph{kullandık}: \Lean{}'de önermeler yazdık ve
onları ispatladık. Kısım~III'te mantığın kendisini bir matematiksel
\emph{nesne} olarak inceleyeceğiz. Soracağımız sorular, mantık
\emph{içinde} değil, mantık \emph{hakkında} sorulardır:

\begin{itemize}
  \item İspatlayabildiğimiz her şey doğru mudur? (\emph{geçerlilik},
        \cref{bol:gecerlilik})
  \item Doğru olan her şeyi ispatlayabilir miyiz? (\emph{tamlık},
        \cref{bol:tamlik})
  \item Her ispat dolaysız bir ispata dönüştürülebilir mi?
        (\emph{kesme eleme}, \cref{bol:ardisik})
\end{itemize}

\noindent Bu sorulara verilen cevaplara \emph{metamantık} denir.
\Cref{fig:buyuk-resim}, Kısım~III'ün büyük resmini gösteriyor.

\begin{figure}[htbp]
\centering
\begin{tikzpicture}[
    kutu/.style={draw=#1, fill=#1!8, very thick, rounded corners=6pt,
                 text width=4.2cm, align=center, inner sep=8pt},
    ok/.style={-{Stealth[length=3mm]}, very thick, #1},
  ]
  \node[kutu=anarenk] (s) at (0,0)
    {{\baslikfont\bfseries\color{anarenk} SÖZDİZİMİ}\\[3pt]
     $\Gamma \vdash A$\\[2pt]
     {\small ``$A$, $\Gamma$'dan \emph{türetilebilir}''}\\[2pt]
     {\scriptsize çıkarım kuralları, ispat ağaçları}};
  \node[kutu=ikincil] (a) at (10,0)
    {{\baslikfont\bfseries\color{ikincil} ANLAMBİLİM}\\[3pt]
     $\Gamma \models A$\\[2pt]
     {\small ``$A$, $\Gamma$'nın \emph{mantıksal sonucudur}''}\\[2pt]
     {\scriptsize değerlemeler, doğruluk tabloları}};
  \draw[ok=teoremframe] ([yshift=0.55cm]s.east) -- node[above, font=\baslikfont\small\bfseries, text=teoremframe]{Geçerlilik (Bölüm 11)} ([yshift=0.55cm]a.west);
  \draw[ok=vurgu] ([yshift=-0.55cm]a.west) -- node[below, font=\baslikfont\small\bfseries, text=vurgu]{Tamlık (Bölüm 12)} ([yshift=-0.55cm]s.east);
  \node[font=\baslikfont\scriptsize, text=black!60] at (5,1.3) {ispatlanabilen her şey doğrudur};
  \node[font=\baslikfont\scriptsize, text=black!60] at (5,-1.3) {doğru olan her şey ispatlanabilir};
  \node[font=\baslikfont\small\bfseries, text=koyu] at (5,-2.4) {Bu bölüm: iki tarafın da tanımları};
\end{tikzpicture}
\caption{Kısım III'ün büyük resmi. Sözdizimi ile anlambilim, mantığın
iki yüzüdür; geçerlilik ve tamlık teoremleri bu iki yüzü birbirine bağlar.}
\label{fig:buyuk-resim}
\end{figure}

% ---------------------------------------------------------------------
\section{Nesne dili ve üst dil}\label{sec:nesne-ust}

Mantık hakkında konuşurken iki dili birbirinden ayırmak zorundayız:

\begin{tanim}{Nesne dili ve üst dil}{nesne-ust}
\textbf{Nesne dili}, incelediğimiz mantık sistemidir: formülleri,
bağlaçları, çıkarım kuralları. \textbf{Üst dil} ise nesne dili hakkında
konuşmak ve onun özelliklerini ispatlamak için kullandığımız dildir.
\end{tanim}

\noindent Kâğıt üzerindeki mantık kitaplarında iki dil de aynı sembollerle
yazıldığı için karışıklık kolayca ortaya çıkar. ``$A \to B$ ispatlanabilirse
ve $A$ ispatlanabilirse $B$ ispatlanabilir'' cümlesindeki ilk ``$\to$''
nesne dilinin bir bağlacıdır; ``ise''ler ise üst dilin parçasıdır.

\Lean{}'de bu ayrım kendiliğinden ortaya çıkar (\cref{fig:nesne-ust}).
Üst dil \Lean{}'in kendisidir: \lc{Prop}, \lc{∀}, \lc{→}. Nesne dili ise
\Lean{}'de tanımlayacağımız bir \emph{veri tipidir}: \lc{Formul}. Bir
formül, bir liste ya da bir ağaç gibi sıradan bir veridir; onu
hesaplayabilir, dönüştürebilir ve onun hakkında teoremler
ispatlayabiliriz.

\begin{figure}[htbp]
\centering
\begin{tikzpicture}[
    katman/.style={draw=#1, fill=#1!6, very thick, rounded corners=8pt},
  ]
  \draw[katman=anarenk] (0,0) rectangle (11.4,4.2);
  \node[anchor=north west, font=\baslikfont\bfseries, text=anarenk] at (0.2,4.1) {ÜST DİL: Lean};
  \node[anchor=north west, font=\small, text=black!70, align=left] at (0.2,3.55)
       {\lc{Prop}, \lc{∀}, \lc{∃}, \lc{→}, \lc{theorem}\\[1pt] \lc{Totoloji A}, \lc{Sonuc Γ A}};
  \draw[katman=ikincil] (5.2,0.4) rectangle (11.0,3.3);
  \node[anchor=north west, font=\baslikfont\bfseries, text=ikincil] at (5.4,3.2) {NESNE DİLİ: \lc{Formul}};
  \node[anchor=north west, font=\small, align=left] at (5.4,2.6)
       {\lc{atom 0}, \lc{bot}\\[2pt]
        $p \oise q$, \quad $p \ove q$, \quad $p \oveya {\sim}p$\\[2pt]
        {\scriptsize sıradan veri: hesaplanır, dönüştürülür}};
  \draw[-{Stealth[length=3mm]}, very thick, vurgu] (3.4,1.2) -- node[above, font=\baslikfont\scriptsize\bfseries, text=vurgu, align=center]{hakkında\\teorem} (5.1,1.2);
  \node[font=\leanfont\scriptsize, align=left, anchor=west] at (0.2,1.2)
       {theorem peirce\_totoloji :\\\quad Totoloji peirce};
\end{tikzpicture}
\caption{Üst dil (\Lean{}) ile nesne dili (\lc{Formul} tipi). Nesne
dilinin formülleri, üst dilde birer veridir.}
\label{fig:nesne-ust}
\end{figure}

% ---------------------------------------------------------------------
\section{Formüllerin sözdizimi}\label{sec:formul-sozdizim}

Önermeler mantığının formülleri, kâğıt üzerinde genellikle şöyle bir
\emph{dilbilgisiyle} tanımlanır:
\[
  A, B \;::=\; p_n \;\mid\; \bot \;\mid\; A \wedge B \;\mid\; A \vee B \;\mid\; A \to B
\]
``Bir formül ya bir atomdur, ya $\bot$'dur, ya da iki formülün
bağlaçla birleştirilmesidir.'' Bu tanımın \Lean{}'deki karşılığı,
harfi harfine bir tümevarımsal tiptir:

\leandosya[Formüller]{Bolum10}{formul_tanim}

\noindent Atomları doğal sayılarla adlandırıyoruz: $p_0, p_1, p_2, \ldots$
Değillemeyi ayrı bir kurucu olarak eklemiyoruz; Bölüm~3'teki gibi
$\neg A := A \to \bot$ olarak tanımlıyoruz. Okunabilirlik için nesne
dilinin bağlaçlarına \Lean{}'in kendi bağlaçlarından \emph{farklı}
semboller veriyoruz:

\leandosya[Gösterim ve tanımlanmış bağlaçlar]{Bolum10}{gosterim}

\begin{center}
\renewcommand{\arraystretch}{1.4}
\rowcolors{2}{tablosatir}{white}
\begin{tabular}{>{\centering\arraybackslash}p{2.4cm}
                >{\centering\arraybackslash}p{2.6cm}
                >{\raggedright\arraybackslash}p{6.4cm}}
\rowcolor{tabloust}
\tablobaslik{Nesne dili} & \tablobaslik{Üst dil (Lean)} & \tablobaslik{Fark} \\
$A \ove B$     & \lc{P ∧ Q}  & biri bir \lc{Formul} (veri), öteki bir \lc{Prop} \\
$A \oveya B$   & \lc{P ∨ Q}  & aynı \\
$A \oise B$    & \lc{P → Q}  & aynı \\
${\sim}A$      & \lc{¬P}     & aynı \\
\end{tabular}
\end{center}

\begin{dikkat}[Dikkat: Formüller önerme değildir]
$p \ove q$ yazdığımızda bir \emph{önerme} değil, bir \emph{veri} (bir
sözdizim ağacı) kurarız. Bu ağaç kendi başına ne doğrudur ne yanlış;
\lc{Formul} tipinin bir elemanıdır, tıpkı \lc{[1, 2, 3]}'ün bir liste
olması gibi. Bir formülün doğru olup olmadığını sorabilmek için önce ona
bir \emph{anlam} vermemiz gerekir (\cref{sec:anlambilim}).
\end{dikkat}

\subsection{Sözdizim ağaçları}

Bir formül, aslında bir ağaçtır. Parantezler ve öncelik kuralları
yalnızca bu ağacı tek satıra yazmanın yollarıdır
(\cref{fig:sozdizim-agaci}).

\leandosya[Örnek formüller]{Bolum10}{ornek_formul}

\begin{figure}[htbp]
\centering
\begin{tikzpicture}[
    baglac/.style={circle, draw=leansym, fill=leansym!10, thick,
                   minimum size=0.8cm, font=\small},
    atom/.style={rectangle, rounded corners=3pt, draw=leanhyp, fill=leanhyp!10,
                 thick, minimum size=0.6cm, font=\small},
    level distance=1.05cm,
    level 1/.style={sibling distance=3.0cm},
    level 2/.style={sibling distance=1.9cm},
    level 3/.style={sibling distance=1.2cm},
    edge from parent/.style={draw, black!45, thick},
  ]
  \node[baglac] {$\Rightarrow$}
    child {node[baglac] {$\Rightarrow$}
      child {node[baglac] {$\Rightarrow$}
        child {node[atom] {$p$}}
        child {node[atom] {$q$}}}
      child {node[atom] {$p$}}}
    child {node[atom] {$p$}};
  \node[font=\small, anchor=west, align=left] at (2.8,-1.2)
       {$((p \oise q) \oise p) \oise p$\\[3pt]
        {\scriptsize Peirce yasası: 3 bağlaç, 4 atom geçişi}};
\end{tikzpicture}
\caption{Peirce yasasının sözdizim ağacı. Her iç düğüm bir bağlaç,
her yaprak bir atomdur.}
\label{fig:sozdizim-agaci}
\end{figure}

\subsection{Formüller üzerinde özyineleme}

\lc{Formul} bir tümevarımsal tip olduğu için, Bölüm~9'daki ağaçlar
gibi, üzerinde fonksiyonlar özyinelemeyle tanımlanır ve özellikler
yapısal tümevarımla ispatlanır. Her kurucu için bir durum vardır;
ikili bağlaçlarda iki tümevarım hipotezi gelir.

\leandosya[Formüller üzerinde özyineleme]{Bolum10}{formul_ozyineleme}

% ---------------------------------------------------------------------
\section{Anlambilim: doğruluk}\label{sec:anlambilim}

Bir formüle anlam vermek, atomlarının doğru mu yanlış mı olduğunu
söylemekle başlar.

\begin{tanim}{Değerleme ve doğruluk değeri}{degerleme}
Bir \textbf{değerleme}, her atoma bir doğruluk değeri veren bir
fonksiyondur: $v : \mathbb{N} \to \{\mathsf{true}, \mathsf{false}\}$.
Bir $A$ formülünün $v$ altındaki \textbf{doğruluk değeri} $\lBrack A
\rBrack_v$, formülün yapısı üzerinde özyinelemeyle tanımlanır:
\begin{align*}
\lBrack p_n \rBrack_v &= v(n) &
\lBrack A \ove B \rBrack_v &= \lBrack A \rBrack_v \;\&\&\; \lBrack B \rBrack_v \\
\lBrack \bot \rBrack_v &= \mathsf{false} &
\lBrack A \oveya B \rBrack_v &= \lBrack A \rBrack_v \;\|\; \lBrack B \rBrack_v \\
& & \lBrack A \oise B \rBrack_v &= \mathord{!}\lBrack A \rBrack_v \;\|\; \lBrack B \rBrack_v
\end{align*}
\end{tanim}

\leandosya[Tarski'nin doğruluk tanımı]{Bolum10}{deger}

\noindent Bu tanımın ilginç bir yanı vardır: Nesne dilindeki ``ve''nin
anlamını üst dildeki ``ve'' (\lc{\&\&}) ile veriyoruz. Bu bir kısır döngü
değildir; çünkü üst dilin ``ve''sini zaten anlıyoruz ve onu, henüz anlamı
olmayan bir sembole ($\ove$) anlam vermek için kullanıyoruz.

\begin{tarihce}[Tarihçe: Tarski ve doğruluk]
Polonyalı mantıkçı Alfred Tarski 1933'te yayımladığı ünlü çalışmasında,
biçimsel bir dilde ``doğruluk'' kavramının nasıl tanımlanabileceğini
gösterdi \citep{tarski1933}. Tarski'nin anahtar fikri, nesne dili ile
üst dili kesin biçimde ayırmaktı: Bir dilin doğruluk kavramı, o dilin
\emph{kendi içinde} değil, daha güçlü bir üst dilde tanımlanmalıdır. Bu
kısıtlama, yalancı paradoksunu (``Bu cümle yanlıştır'') engeller.
Tarski'nin \emph{doğruluğun tanımlanamazlığı teoremi}, Gödel'in eksiklik
teoremiyle birlikte \ciltiki{4}'ün konusudur.
\end{tarihce}

\leandosya[Değillemenin değeri]{Bolum10}{deger_degil}

\subsection{Doğruluk tabloları}

Bir formülün doğruluk tablosu, atomlarına verilebilecek bütün değerler
için \lc{deger}'in sonucunu listeler. \Cref{fig:peirce-tablo}, Peirce
yasasının tablosunu gösteriyor.

\begin{figure}[htbp]
\centering
\begin{tikzpicture}
  \matrix[matrix of nodes, nodes={minimum width=1.9cm, minimum height=0.68cm,
          anchor=center, font=\small}, column sep=0pt, row sep=0pt,
          row 1/.style={nodes={fill=tabloust, text=white, font=\baslikfont\bfseries\scriptsize}}] (t) {
    $p$ & $q$ & $p \oise q$ & $(p \oise q) \oise p$ & Peirce \\
    1 & 1 & 1 & 1 & |[fill=bittibg, text=bittiframe]| \textbf{1} \\
    1 & 0 & 0 & 1 & |[fill=bittibg, text=bittiframe]| \textbf{1} \\
    0 & 1 & 1 & 0 & |[fill=bittibg, text=bittiframe]| \textbf{1} \\
    0 & 0 & 1 & 0 & |[fill=bittibg, text=bittiframe]| \textbf{1} \\
  };
  \draw[black!25] (t.north west) rectangle (t.south east);
  \draw[black!25] (t-1-3.north west) -- (t-5-3.south west);
\end{tikzpicture}
\caption{Peirce yasası $((p \oise q) \oise p) \oise p$'nin doğruluk
tablosu. Son sütun her satırda 1'dir: Formül bir totolojidir.}
\label{fig:peirce-tablo}
\end{figure}

% ---------------------------------------------------------------------
\section{Totoloji ve mantıksal sonuç}\label{sec:totoloji}

\begin{tanim}{Totoloji, sağlanabilirlik, mantıksal sonuç}{totoloji}
\begin{itemize}
  \item $A$ bir \textbf{totolojidir} ($\models A$): her $v$ için
        $\lBrack A \rBrack_v = \mathsf{true}$.
  \item $A$ \textbf{sağlanabilirdir}: en az bir $v$ için
        $\lBrack A \rBrack_v = \mathsf{true}$.
  \item $A$, $\Gamma$'nın \textbf{mantıksal sonucudur} ($\Gamma \models A$):
        $\Gamma$'daki bütün formülleri doğru yapan her $v$, $A$'yı da
        doğru yapar.
\end{itemize}
\end{tanim}

\leandosya[Tanımlar]{Bolum10}{totoloji_tanim}

\noindent Bu tanımlar \emph{üst dilde} birer önermedir: \lc{Totoloji A}
bir \lc{Prop}'tur ve onu \Lean{}'de ispatlayabiliriz. Bir totolojiyi
ispatlamanın en doğrudan yolu, değerleme üzerinde durum ayırmaktır;
bu, doğruluk tablosunu satır satır kontrol etmenin ta kendisidir:

\leandosya[Totolojiler]{Bolum10}{totoloji_ornek}

\begin{figure}[htbp]
\centering
\begin{tikzpicture}
  \draw[black!30, thick, fill=white, rounded corners=8pt] (0,0) rectangle (8.6,4.0);
  \node[anchor=north west, font=\baslikfont\small\bfseries, text=black!70] at (0.15,3.95) {bütün değerlemeler};
  \draw[ikincil, very thick, fill=ikincil!10] (5.0,2.0) ellipse (2.8cm and 1.55cm);
  \node[font=\baslikfont\small\bfseries, text=ikincil] at (6.9,3.15) {$A$ doğru};
  \draw[anarenk, very thick, fill=anarenk!15] (4.6,1.85) ellipse (1.55cm and 0.95cm);
  \node[font=\baslikfont\small\bfseries, text=anarenk] at (4.6,1.85) {$\Gamma$ doğru};
  \node[font=\small, anchor=west, align=left] at (9.0,2.4)
       {$\Gamma \models A$:\\$\Gamma$'yı doğrulayan\\her değerleme\\$A$'yı da doğrular};
\end{tikzpicture}
\caption{Mantıksal sonuç: $\Gamma$'nın modelleri, $A$'nın modellerinin
içindedir.}
\label{fig:sonuc}
\end{figure}

\noindent Mantıksal sonuç ilişkisinin bir örneği, modus ponens
kuralının anlambilimsel karşılığıdır:

\leandosya[Anlambilimsel modus ponens]{Bolum10}{sonuc_ornek}

\noindent Bir varsayım eklemek bir mantıksal sonucu bozmaz. Bu özelliğe
\emph{tekdüzelik} (İng.\ \emph{monotonicity}) denir:

\leandosya[Tekdüzelik]{Bolum10}{zayiflatma}

% ---------------------------------------------------------------------
\section{Bir karar yöntemi}\label{sec:karar-yontemi}

Bir formülün totoloji olup olmadığını \emph{hesaplayarak} anlayabiliriz:
İçinde geçen atomlara verilebilecek bütün değerleri tek tek deneriz.
$n$ atom için $2^n$ değerleme vardır:

\leandosya[Doğruluk tablosu yöntemi]{Bolum10}{karar}

\begin{figure}[htbp]
\centering
\begin{tikzpicture}[
    d/.style={circle, draw=anarenk, fill=tanimbg, thick, minimum size=0.55cm, font=\scriptsize},
    y/.style={rectangle, rounded corners=2pt, draw=#1, fill=#1!12, thick, font=\scriptsize, minimum width=1.0cm},
    level distance=1.1cm,
    level 1/.style={sibling distance=4.2cm},
    level 2/.style={sibling distance=2.1cm},
    edge from parent/.style={draw, black!45, thick},
  ]
  \node[d] {$p$}
    child {node[d] {$q$}
      child {node[y=bittiframe] {$\mathsf{T}$} edge from parent node[left, font=\scriptsize]{1}}
      child {node[y=bittiframe] {$\mathsf{T}$} edge from parent node[right, font=\scriptsize]{0}}
      edge from parent node[left=2pt, font=\scriptsize]{$p = 1$}}
    child {node[d] {$q$}
      child {node[y=bittiframe] {$\mathsf{T}$} edge from parent node[left, font=\scriptsize]{1}}
      child {node[y=bittiframe] {$\mathsf{T}$} edge from parent node[right, font=\scriptsize]{0}}
      edge from parent node[right=2pt, font=\scriptsize]{$p = 0$}};
  \node[font=\small, anchor=west, align=left] at (4.6,-1.1)
       {$n$ atom $\Rightarrow$ $2^n$ yaprak\\[2pt]
        {\scriptsize 10 atom: 1024 değerleme}\\
        {\scriptsize 50 atom: $\approx 10^{15}$ değerleme}};
\end{tikzpicture}
\caption{Doğruluk tablosu yöntemi bütün değerlemeleri dener. Peirce
yasası için 4 yaprağın hepsi ``doğru''dur. Yöntemin maliyeti atom
sayısıyla üstel olarak büyür.}
\label{fig:karar-agaci}
\end{figure}

\noindent \lc{totolojiMi} fonksiyonu yalnızca \lc{0, 1, …, n} atomlarını
dener; sonsuz sayıdaki öteki atomların değerlerini görmezden gelir. Bunun
neden yeterli olduğunu bir sonraki kesimde ispatlayacağız: Bir formülün
değeri yalnızca içinde geçen atomlara bağlıdır.

\begin{dikkat}[Dikkat: \#eval bir ispat değildir]
\lc{\textcolor{leancmd}{\bfseries\#eval} totolojiMi peirce} \lc{true}
sonucunu verir; ama bu, \lc{Totoloji peirce} önermesinin bir \emph{ispatı}
değildir. \lc{totolojiMi}'nin kendisi bir programdır ve hatalı olabilir.
Programı bir ispata dönüştürmek için onun \emph{doğruluğunu} ispatlamak
gerekir: \lc{totolojiMi A = true → Totoloji A}. Bu ispat, kendi
ispatlarını yapan karar yöntemlerinin (\emph{yansıma} tekniğinin)
temelidir; \ciltiki{6}'da tamamlayacağız.
\end{dikkat}

\begin{arastirmanotu}[Araştırma Notu: SAT ve P = NP]
Bir formülün \emph{sağlanabilir} olup olmadığını belirleme problemi
(SAT), Stephen Cook'un 1971'de gösterdiği gibi \emph{NP-tam} ilk problemdir
\citep{cook1971}. Doğruluk tablosu yöntemi üstel zaman alır; bundan
temelde daha iyi bir algoritmanın var olup olmadığı, bilgisayar biliminin
en ünlü açık problemi olan P = NP sorusuna denktir. Yine de modern
SAT çözücüleri, milyonlarca değişkenli endüstriyel formülleri
pratikte çözebilmektedir; bu konuya \ciltiki{6}'da döneceğiz.
\end{arastirmanotu}

% ---------------------------------------------------------------------
\section{İki önemli önsav}\label{sec:yerellik}

\subsection{Yerellik}

\begin{teorem}{Yerellik}{yerellik}
İki değerleme bir $A$ formülünde geçen bütün atomlarda aynı değeri
veriyorsa, $A$'ya da aynı değeri verir.
\end{teorem}

\leandosya[Yerellik]{Bolum10}{yerellik}

\noindent İspat, formülün yapısı üzerinde tümevarımdır. Her ikili bağlaç
durumunda, alt formüllerin atomları ana formülün atomlarının bir
alt kümesi olduğu için tümevarım hipotezleri uygulanabilir. Üç ikili
bağlacın durumları neredeyse aynıdır; bu tekrar, \ciltiki{7}'de
kendi taktiklerimizi yazmamız için iyi bir neden olacak.

\subsection{Yerine koyma}

Bir totolojide atomların yerine \emph{herhangi} formüller koyarsak yine
bir totoloji elde ederiz. Örneğin $p \oise p$ bir totoloji olduğuna göre
$(q \ove r) \oise (q \ove r)$ de bir totolojidir.

\leandosya[Yerine koyma]{Bolum10}{yerine_koy}

\noindent \lc{deger\_yerineKoy} önsavı, yerine koymanın anlamını tam
olarak söyler: $\sigma$ ile yerine koyduktan sonra $v$ altında
değerlendirmek, ``her atomu $\sigma$'daki formülün $v$ altındaki
değerine götüren'' yeni bir değerleme altında değerlendirmekle aynıdır.

% ---------------------------------------------------------------------
\section{Doğruluk tabloları klasiktir}\label{sec:klasik-anlambilim}

\lc{\textcolor{leancmd}{\bfseries\#eval} totolojiMi ucuncuHal} ve
\lc{totolojiMi peirce} sonuçları \lc{true}'dur. Oysa Bölüm~4'te bu iki
formülün \emph{yapıcı} olarak ispatlanamadığını görmüştük. Bu bir çelişki
değildir; bize anlambilimimiz hakkında önemli bir şey söyler:

\begin{onemlifikir}[Önemli Fikir: Her mantığın kendi anlambilimi vardır]
Doğruluk değerlerinin yalnızca iki eleman (\lc{true}, \lc{false})
içeren bir kümeden geldiği anlambilim, \emph{klasik} mantığın
anlambilimidir. Bu anlambilimde her formül ya doğru ya yanlıştır;
dolayısıyla $p \oveya {\sim}p$ her zaman doğrudur. Yapıcı mantığın
anlambilimi ise Bölüm~4'te önizlediğimiz Kripke modelleridir;
\ciltiki{2}'de onları da \Lean{}'de tanımlayacağız. Geçerlilik ve
tamlık teoremleri her zaman \emph{bir ispat sistemi} ile \emph{bir
anlambilim} arasındaki ilişki hakkındadır.
\end{onemlifikir}

\noindent Bu yüzden Bölüm~11 ve 12'de inceleyeceğimiz ispat sistemi,
klasik doğal çıkarım olacaktır: Bölüm~5'teki kural kataloğu ve RAA kuralı.

% ---------------------------------------------------------------------
\begin{arastirmanotu}
\textbf{Formalizasyonda sözdizimi.} Önermeler mantığında değişken
bağlama olmadığı için sözdizimi basittir. Yüklemler mantığında ise
niceleyicilerin bağladığı değişkenleri temsil etmek zordur; Bölüm~6'daki
değişken yakalama sorunu formalizasyonda ciddi bir mühendislik
problemine dönüşür. Yaygın çözümler de Bruijn indeksleri, ``yerel
adsız'' (\emph{locally nameless}) temsil ve nominal tekniklerdir.
Lean~4'teki Foundation projesi, Gödel teoremlerini formalize ederken
yerel adsız temsili kullanır \citep{saitou2026, foundation}.

\textbf{Derin ve sığ gömme.} Bu bölümde nesne dilini bir veri tipi
olarak tanımladık; buna \emph{derin gömme} (İng.\ \emph{deep embedding})
denir. Önceki bölümlerde ise mantığı doğrudan \Lean{}'in \lc{Prop}'u
içinde kullanıyorduk; buna \emph{sığ gömme} (\emph{shallow embedding})
denir. Derin gömme, mantık \emph{hakkında} teorem ispatlamamızı sağlar;
sığ gömme ise mantık \emph{içinde} ispat yapmayı kolaylaştırır.

\textbf{İleri okuma.} Önermeler mantığının anlambilimi ve metamantığı
için \citet{vandalen2013}; formalizasyon odaklı bir sunum için
\citet{harrison2009}.
\end{arastirmanotu}

\begin{bolumozeti}
\begin{itemize}
  \item Nesne dili incelenen mantıktır; üst dil, onun hakkında konuşulan
        dildir. \Lean{}'de üst dil \Lean{}'in kendisi, nesne dili ise
        \lc{Formul} veri tipidir.
  \item Formüller tümevarımsal bir tiptir; üzerlerinde özyineleme ve
        yapısal tümevarım yapılır.
  \item Bir değerleme atomlara doğruluk değeri verir; \lc{deger},
        Tarski'nin doğruluk tanımıdır.
  \item Totoloji, sağlanabilirlik ve mantıksal sonuç ($\Gamma \models A$)
        üst dilde tanımlanan önermelerdir.
  \item Doğruluk tablosu yöntemi totolojileri hesaplayarak bulur, ama
        maliyeti üsteldir. \lc{\#eval} bir ispat değildir.
  \item Yerellik ve yerine koyma önsavları, yapısal tümevarımla ispatlanır.
  \item İki değerli anlambilim klasiktir; yapıcı mantığın anlambilimi
        Kripke modelleridir.
\end{itemize}
\end{bolumozeti}

\begin{alistirmalar}
\alistirma{1} ${\sim}A$'nın boyutu, $A$'nın boyutundan bir fazladır.
(\lc{a10\_1})

\alistirma{2} $p \ove q \oise q \ove p$ bir totolojidir. (\lc{a10\_2})

\alistirma{2} $p \ove q \models p$. (\lc{a10\_3})

\alistirma{2} $\Gamma \models A$ ise $\Gamma \models A \oveya B$.
(\lc{a10\_4})

\alistirma{2} Bir totoloji her $\Gamma$'nın mantıksal sonucudur.
(\lc{a10\_5})

\alistirma{3} Anlambilimsel modus ponens: $\Gamma \models A \oise B$ ve
$\Gamma \models A$ ise $\Gamma \models B$. (\lc{a10\_6})

\alistirma{3} \textbf{Anlambilimsel çıkarım teoremi}: $A, \Gamma \models B$
ise $\Gamma \models A \oise B$. Bu teorem, Bölüm~5'teki $\to$G kuralının
anlambilimsel karşılığıdır. (\lc{a10\_7})

\alistirma{2} $\bot$ sağlanabilir değildir. (\lc{a10\_8})

\alistirma{1} \lc{\#eval} ile aşağıdaki formüllerin totoloji olup
olmadığını denetleyin: $(p \oise q) \oise ({\sim}q \oise {\sim}p)$,
$({\sim}q \oise {\sim}p) \oise (p \oise q)$, $(p \oise q) \oveya (q \oise p)$.
Hangileri Bölüm~4'e göre yapıcı olarak ispatlanamaz?

\alistirma{2} (Kâğıt-kalem) $\{\oise, \bot\}$ bağlaçlarının tek başına
yeterli olduğunu gösterin: $A \ove B$ ve $A \oveya B$'yi, yalnızca $\oise$
ve $\bot$ kullanarak, bütün değerlemelerde aynı değeri alan formüllerle
ifade edin.

\alistirma{3} (Proje) \lc{totolojiMi A = true → Totoloji A} önermesini
ispatlayın. İpucu: Önce \lc{tumDegerlemeler n}'nin, \lc{0, …, n-1}
atomlarında her olası değer birleşimini içerdiğini gösterin; sonra
yerellik önsavını kullanın.
\end{alistirmalar}
   % Sözdizimi ve anlambilim
% =====================================================================
%  Bölüm 11 — Geçerlilik Teoremi
% =====================================================================
\chapter{Geçerlilik Teoremi}\label{bol:gecerlilik}

\begin{bolumhedefleri}
\begin{itemize}
  \item Bir ispat sistemini \Lean{}'de tümevarımsal bir yüklem olarak
        tanımlamayı,
  \item Nesne dilinde türetimler kurmayı ve onları ağaç olarak okumayı,
  \item Geçerlilik teoremini (``türetilebilen her şey doğrudur'') kural
        tümevarımıyla tam olarak ispatlamayı,
  \item Geçerlilikten tutarlılığı ve türetilemezlik sonuçlarını
        çıkarmayı,
  \item Yapısal kuralları (zayıflatma) ve bir ispat sisteminin
        anlambilimine göre neyi ``yakalayıp'' neyi yakalayamadığını.
\end{itemize}
\end{bolumhedefleri}

Bölüm~10'da mantığın iki yüzünden birini, anlambilimi tanımladık:
$\Gamma \models A$. Bu bölümde öteki yüzü, sözdizimini tanımlıyor ve
ikisi arasındaki ilk köprüyü kuruyoruz:
\[
  \Gamma \vdash A \quad\Longrightarrow\quad \Gamma \models A.
\]
Bu, kitabın ilk tam \emph{metamantık} teoremi olacak: Bir mantık sistemi
\emph{hakkında}, o sistemin dışından ispatlanan bir teorem.

% ---------------------------------------------------------------------
\section{Bir ispat sistemi, veri olarak}\label{sec:ispat-sistemi}

Bölüm~5'te doğal çıkarım kurallarını ardışık biçiminde, yani her
düğümde bağlamı açıkça yazarak ifade etmiştik. O kuralları şimdi
\Lean{}'de bir \emph{tümevarımsal yüklem} olarak yazıyoruz. Her kural bir
kurucu, her türetim ağacı ise bu yüklemin bir kanıtı olacak:

\leandosya[Klasik doğal çıkarım]{Bolum11}{turetilir}

\noindent \lc{Turetilir Γ A}, kâğıt üzerinde $\Gamma \vdash A$ diye
yazdığımız şeydir. Bağlamı (\lc{Γ}) bir liste olarak tutuyoruz;
yeni bir varsayım listenin başına eklenir. \Cref{fig:kural-sistemi},
on bir kuralın tamamını ardışık biçiminde gösteriyor.

\begin{figure}[htbp]
\centering
\setlength{\tabcolsep}{16pt}
\renewcommand{\arraystretch}{3.0}
\begin{tabular}{ccc}
\begin{prooftree}
  \hypo{A \in \Gamma}
  \infer1[Var]{\Gamma \vdash A}
\end{prooftree}
&
\begin{prooftree}
  \hypo{\Gamma \vdash A}
  \hypo{\Gamma \vdash B}
  \infer2[$\wedge$G]{\Gamma \vdash A \ove B}
\end{prooftree}
&
\begin{prooftree}
  \hypo{\Gamma \vdash A \ove B}
  \infer1[$\wedge$E$_{1,2}$]{\Gamma \vdash A \quad (B)}
\end{prooftree}
\\
\begin{prooftree}
  \hypo{\Gamma \vdash A \quad (B)}
  \infer1[$\vee$G$_{1,2}$]{\Gamma \vdash A \oveya B}
\end{prooftree}
&
\multicolumn{2}{c}{%
\begin{prooftree}
  \hypo{\Gamma \vdash A \oveya B}
  \hypo{A, \Gamma \vdash C}
  \hypo{B, \Gamma \vdash C}
  \infer3[$\vee$E]{\Gamma \vdash C}
\end{prooftree}}
\\
\begin{prooftree}
  \hypo{A, \Gamma \vdash B}
  \infer1[$\to$G]{\Gamma \vdash A \oise B}
\end{prooftree}
&
\begin{prooftree}
  \hypo{\Gamma \vdash A \oise B}
  \hypo{\Gamma \vdash A}
  \infer2[$\to$E]{\Gamma \vdash B}
\end{prooftree}
&
\begin{prooftree}
  \hypo{\Gamma \vdash \bot}
  \infer1[$\bot$E]{\Gamma \vdash A}
\end{prooftree}
\\
\multicolumn{3}{c}{%
\colorbox{teorembg}{\begin{prooftree}
  \hypo{{\sim}A, \Gamma \vdash \bot}
  \infer1[RAA]{\Gamma \vdash A}
\end{prooftree}\quad\raisebox{0.4em}{\textcolor{teoremframe}{\baslikfont\scriptsize\bfseries klasik kural}}}}
\\
\end{tabular}
\caption{\lc{Turetilir} yükleminin on bir kuralı. Her kural, \Lean{}
tanımındaki bir kurucuya karşılık gelir.}
\label{fig:kural-sistemi}
\end{figure}

\begin{onemlifikir}[Önemli Fikir: Türetim = kanıt terimi]
\lc{Turetilir Γ A} türünde bir \Lean{} terimi, bir türetim ağacının
ta kendisidir: Her kurucu bir düğüm, kurucunun argümanları o düğümün
çocuklarıdır. Bölüm~9'daki \lc{Cift} yükleminde \lc{Cift 4}'ün kanıtı
nasıl bir türetim ağacıysa, \lc{Turetilir [] peirce}'in kanıtı da bir
doğal çıkarım ağacıdır. Böylece ``türetim ağaçları üzerinde tümevarım'',
\Lean{}'de sıradan bir kural tümevarımına dönüşür.
\end{onemlifikir}

% ---------------------------------------------------------------------
\section{Nesne dilinde türetimler}\label{sec:nesne-turetim}

En basit türetim, özdeşlik yasasıdır: $\vdash A \oise A$.

\leandosya[Özdeşlik]{Bolum11}{ornek_ozdeslik}

\begin{center}
\begin{prooftree}
  \hypo{A \in [A]}
  \infer1[Var]{A \vdash A}
  \infer1[$\to$G]{\vdash A \oise A}
\end{prooftree}
\end{center}

\noindent Bu ispat ile Bölüm~2'deki \lc{theorem ozdeslik (P : Prop) : P → P}
arasındaki farka dikkat edin. Orada \Lean{}'in \emph{kendi} mantığında bir
önerme ispatlıyorduk (sığ gömme). Burada ise nesne dilinin bir formülünün,
nesne dilinin ispat sisteminde \emph{türetilebilir} olduğunu, üst dilde
ispatlıyoruz (derin gömme).

\subsection{Peirce yasasının türetimi}

Bölüm~4'te Peirce yasasının ancak klasik mantıkla ispatlanabildiğini
görmüştük. Nesne dilindeki ispat sistemimiz klasik olduğu için (RAA
kuralı sayesinde), Peirce yasası burada türetilebilir:

\leandosya[Peirce yasasının türetimi]{Bolum11}{ornek_peirce}

\begin{figure}[htbp]
\centering
\footnotesize
\begin{prooftree}
  \infer0[Var]{\Gamma_1 \vdash {\sim}p}
  \infer0[Var]{\Gamma_1 \vdash (p \oise q) \oise p}
  \infer0[Var]{\Gamma_2 \vdash {\sim}p}
  \infer0[Var]{\Gamma_2 \vdash p}
  \infer2[$\to$E]{\Gamma_2 \vdash \bot}
  \infer1[$\bot$E]{\Gamma_2 \vdash q}
  \infer1[$\to$G]{\Gamma_1 \vdash p \oise q}
  \infer2[$\to$E]{\Gamma_1 \vdash p}
  \infer2[$\to$E]{\Gamma_1 \vdash \bot}
  \infer1[\textcolor{teoremframe}{RAA}]{\Gamma_0 \vdash p}
  \infer1[$\to$G]{\vdash ((p \oise q) \oise p) \oise p}
\end{prooftree}

\vspace{6pt}
\normalsize
\begin{tabular}{r@{\;=\;}l}
$\Gamma_0$ & $[(p \oise q) \oise p]$ \\
$\Gamma_1$ & ${\sim}p :: \Gamma_0$ \\
$\Gamma_2$ & $p :: \Gamma_1$
\end{tabular}
\caption{Peirce yasasının türetim ağacı. Ağacın her düğümü, \Lean{}
ispatındaki bir \leantac{apply} ya da \leantac{exact} adımına karşılık
gelir. Klasik kural RAA yalnızca bir kez kullanılır.}
\label{fig:peirce-turetim}
\end{figure}

\noindent \Lean{} ispatındaki her \leantac{apply} adımı ağaçta bir
düğümü, her \lc{by decide} ise bir varsayımın bağlamda bulunduğunu
denetleyen bir hesabı temsil eder. \lc{by decide}'ın burada işe yaraması,
\lc{Formul} tipinin karar verilebilir eşitliğe sahip olmasındandır
(Bölüm~10'daki \leankw{deriving} \lc{DecidableEq}).

% ---------------------------------------------------------------------
\section{Geçerlilik teoremi}\label{sec:gecerlilik-teoremi}

\begin{teorem}{Geçerlilik}{gecerlilik}
Her $\Gamma$ ve $A$ için, $\Gamma \vdash A$ ise $\Gamma \models A$.
Yani türetilebilen her formül, varsayımlarının mantıksal sonucudur.
\end{teorem}

\subsection{İspatın fikri}

$\Gamma \vdash A$'nın bir türetimi verilsin. Türetim üzerinde kural
tümevarımı yapıyoruz. Her kural için şunu göstermeliyiz: Kuralın
\emph{öncülleri} geçerliyse (tümevarım hipotezleri), \emph{sonucu} da
geçerlidir. Başka bir deyişle, \textbf{her kural doğruluğu korur}
(\cref{fig:dogruluk-korunur}).

\begin{figure}[htbp]
\centering
\begin{tikzpicture}[
    dugum/.style={draw=bittiframe, fill=bittibg, thick, rounded corners=3pt,
                  font=\small, inner sep=4pt, minimum height=0.7cm},
    kural/.style={font=\baslikfont\scriptsize\bfseries, text=leantac},
    ok/.style={-{Stealth[length=2.4mm]}, thick, bittiframe},
  ]
  \node[dugum] (y1) at (-2.6,3.0) {$\Gamma \models A$};
  \node[dugum] (y2) at (-0.2,3.0) {$\Gamma \models A \oise B$};
  \node[dugum] (y3) at (2.6,3.0) {$\Gamma \models C$};
  \node[dugum] (o1) at (-1.4,1.5) {$\Gamma \models B$};
  \node[dugum] (k) at (0.6,0.0) {$\Gamma \models B \ove C$};
  \draw[ok] (y1) -- (o1); \draw[ok] (y2) -- (o1);
  \draw[ok] (o1) -- (k); \draw[ok] (y3) -- (k);
  \node[kural] at (-0.2,2.2) {$\to$E};
  \node[kural] at (1.4,0.85) {$\wedge$G};
  \node[font=\baslikfont\scriptsize, text=black!60, align=left, anchor=west] at (3.8,2.3)
       {yapraklar: Var\\(varsayım doğrudur)};
  \node[font=\baslikfont\scriptsize, text=black!60, align=left, anchor=west] at (3.8,0.6)
       {her kural:\\öncüller geçerliyse\\sonuç da geçerlidir};
  \draw[decorate, decoration={brace, amplitude=6pt, mirror}, thick, black!35]
       (-3.6,-0.4) -- (3.4,-0.4) node[midway, below=8pt, font=\baslikfont\small\bfseries, text=bittiframe]{doğruluk ağaç boyunca aşağı akar};
\end{tikzpicture}
\caption{Geçerlilik ispatının fikri. Yapraklardaki varsayımlar doğrudur;
her kural, öncüllerindeki doğruluğu sonucuna taşır. Dolayısıyla kökteki
formül de doğrudur.}
\label{fig:dogruluk-korunur}
\end{figure}

\noindent Her kuralın durumu, Bölüm~10'daki \lc{deger} tanımından
kolayca çıkar (\cref{tab:gecerlilik-durumlar}).

\begin{table}[htbp]
\centering
\small
\renewcommand{\arraystretch}{1.45}
\rowcolors{2}{tablosatir}{white}
\begin{tabularx}{\linewidth}{>{\raggedright\arraybackslash}p{1.6cm} X}
\rowcolor{tabloust}
\tablobaslik{Kural} & \tablobaslik{Neden doğruluğu korur? ($v$, $\Gamma$'yı doğrulayan bir değerleme)} \\
Var       & $A \in \Gamma$ ve $v$, $\Gamma$'nın her formülünü doğrular. \\
$\wedge$G & $A$ ve $B$ doğruysa \lc{true \&\& true = true}. \\
$\wedge$E & \lc{a \&\& b = true} ise \lc{a = true} ve \lc{b = true}. \\
$\vee$G   & $A$ doğruysa \lc{true || b = true}. \\
$\vee$E   & $A \oveya B$ doğruysa $A$ ya da $B$ doğrudur; o durumda genişletilmiş bağlam $v$ altında doğrudur ve tümevarım hipotezi $C$'yi verir. \\
$\to$G    & $A$ yanlışsa $A \oise B$ doğrudur; $A$ doğruysa $A{::}\Gamma$ doğrulanır ve hipotez $B$'yi verir. \\
$\to$E    & \lc{!a || b = true} ve \lc{a = true} ise \lc{b = true}. \\
$\bot$E   & $\bot$ hiçbir değerlemede doğru değildir: öncül \emph{boş yere} geçerlidir. \\
RAA       & $A$ yanlış olsaydı ${\sim}A$ doğru olurdu ve hipotez $\bot$'u doğrulardı. \textbf{Bu adım, $A$'nın değerinin yalnızca iki seçenekten biri olmasına dayanır.} \\
\end{tabularx}
\caption{Geçerlilik ispatının on bir durumu.}
\label{tab:gecerlilik-durumlar}
\end{table}

\subsection{Lean ispatı}

Önce, bağlama doğru bir formül eklemenin bağlamın doğruluğunu
bozmadığını gösteren küçük bir önsava ihtiyacımız var. Bu önsav,
bağlamı genişleten üç kuralda ($\vee$E, $\to$G, RAA) kullanılacak:

\leandosya[Bağlamı genişletmek]{Bolum11}{baglam_ekle}

\noindent Şimdi teoremin kendisi:

\leandosya[Geçerlilik teoremi]{Bolum11}{gecerlilik}

\noindent İspatta dikkat edilmesi gereken birkaç nokta:

\begin{itemize}
  \item \leantac{induction h with} satırındaki \lc{@varsayim Γ A hA}
        gibi örüntüler, kurucunun \emph{örtük} argümanlarına (\lc{Γ},
        \lc{A}, \lc{B}) ad vermemizi sağlar. Bu adlar, \leantac{show}
        ile hedefi açık yazabilmek için gereklidir.
  \item İkili bağlaçların durumlarında \leantac{simp}, Bölüm~10'daki
        \lc{Bool} işlemlerini (\lc{\&\&}, \lc{||}, \lc{!}) mantık
        bağlaçlarına çevirir: \lc{(a \&\& b) = true} hipotezi
        \lc{a = true ∧ b = true} olur.
  \item $\to$G ve RAA durumlarında \lc{deger v A}'nın \lc{true} mu
        \lc{false} mu olduğuna göre durum ayırırız. İki değerli
        anlambilimin gücü buradan gelir.
\end{itemize}

\begin{dikkat}[Dikkat: RAA neden geçerli?]
RAA kuralının geçerlilik ispatı, \lc{cases hA : deger v A} satırına
dayanır: Bir formülün değeri ya \lc{true} ya da \lc{false}'tur. Bu,
\Lean{}'in üst dilinde klasik mantık \emph{kullandığımız} anlamına
gelmez: \lc{Bool} üzerinde durum ayırmak tamamen yapıcıdır, çünkü
\lc{Bool}'un iki kurucusu vardır. Klasik olan, nesne dilinin \emph{anlambilimidir}: Her formüle
iki değerden birini veriyoruz. RAA bu anlambilime göre geçerlidir;
Bölüm~4'teki Kripke anlambilimine göre ise geçerli değildir.
\end{dikkat}

% ---------------------------------------------------------------------
\section{Sonuçlar}\label{sec:gecerlilik-sonuclar}

Geçerlilik teoreminin en önemli kullanımı, bir şeyin \emph{türetilemez}
olduğunu göstermektir. Bir formülü türetmeye çalışıp başarısız olmak hiçbir
şey kanıtlamaz; ama formülü yanlış yapan bir değerleme bulursak, geçerlilik
teoremi formülün hiçbir zaman türetilemeyeceğini garanti eder
(\cref{fig:karsi-model}).

\leandosya[Tutarlılık ve türetilemezlik]{Bolum11}{sonuclar}

\begin{figure}[htbp]
\centering
\begin{tikzpicture}[
    kutu/.style={draw=#1, fill=#1!8, thick, rounded corners=5pt,
                 text width=2.6cm, align=center, inner sep=6pt, font=\small},
    ok/.style={-{Stealth[length=3mm]}, very thick, #1},
  ]
  \node[kutu=hataframe] (k) at (0,0) {$A$'yı yanlış yapan\\bir değerleme $v$};
  \node[kutu=ikincil]  (s) at (5.2,0) {$\not\models A$};
  \node[kutu=anarenk]  (t) at (10.4,0) {$\nvdash A$\\{\scriptsize hiçbir türetim yok}};
  \draw[ok=hataframe] (k) -- node[above, font=\baslikfont\scriptsize]{tanım} (s);
  \draw[ok=teoremframe] (s) -- node[above, font=\baslikfont\scriptsize\bfseries, text=teoremframe]{geçerlilik} node[below, font=\baslikfont\scriptsize]{(kontrapozisyon)} (t);
\end{tikzpicture}
\caption{Türetilemezliği göstermenin yolu: bir karşı değerleme bulmak.
Geçerlilik teoreminin kontrapozisyonu, sonsuz sayıdaki olası türetimin
hiçbirinin $A$'yı vermeyeceğini garanti eder.}
\label{fig:karsi-model}
\end{figure}

\begin{teorem}{Tutarlılık}{tutarlilik-klasik}
Klasik doğal çıkarım tutarlıdır: $\nvdash \bot$.
\end{teorem}

\noindent Bölüm~5'te yapıcı mantığın tutarlılığını normalleştirme ve alt
formül özelliğiyle, yani \emph{sözdizimsel} yollarla göstermiştik. Burada
aynı sonuca, klasik sistem için, \emph{anlambilimsel} yoldan ulaşıyoruz:
$\bot$ hiçbir değerlemede doğru olmadığı için türetilemez. İki yöntem
birbirini tamamlar; ispat kuramında her ikisi de kullanılır.

% ---------------------------------------------------------------------
\section{Yapısal özellikler}\label{sec:yapisal}

Doğal çıkarımın kuralları belirli bir bağlaca bağlıdır. Bunların yanında
bir de bağlamın kendisiyle ilgili \emph{yapısal} özellikler vardır:

\begin{center}
\renewcommand{\arraystretch}{1.5}
\rowcolors{2}{tablosatir}{white}
\begin{tabular}{>{\raggedright\arraybackslash}p{2.8cm}
                >{\centering\arraybackslash}p{4.6cm}
                >{\raggedright\arraybackslash}p{4.6cm}}
\rowcolor{tabloust}
\tablobaslik{Özellik} & \tablobaslik{Kural} & \tablobaslik{Bizim sistemimizde} \\
zayıflatma  & $\Gamma \vdash A \;\Rightarrow\; B, \Gamma \vdash A$ & bir teorem (\lc{zayiflatma}) \\
değiştirme  & $A, B, \Gamma \vdash C \;\Rightarrow\; B, A, \Gamma \vdash C$ & zayıflatmanın özel hâli \\
daraltma    & $A, A, \Gamma \vdash C \;\Rightarrow\; A, \Gamma \vdash C$ & zayıflatmanın özel hâli \\
\end{tabular}
\end{center}

\noindent Bağlamı bir liste olarak tutsak da Var kuralı yalnızca
\emph{üyeliğe} baktığı için sıra ve tekrar önemli değildir. Bu yüzden
üç özelliğin üçü de tek bir teoremden çıkar: Bağlam genişletilirse
(daha doğrusu, bağlamın her elemanı yeni bağlamda da bulunuyorsa)
türetilebilirlik korunur.

\leandosya[Zayıflatma]{Bolum11}{zayiflatma}

\noindent \leantac{induction h generalizing Δ} satırına dikkat edin. Bağlamı
genişleten kurallarda (\lc{veyaE}, \lc{iseG}, \lc{raa}) tümevarım
hipotezini \emph{farklı} bir \lc{Δ} ile, yani \lc{A :: Δ} ile kullanmamız
gerekir. Bunun için tümevarımdan önce \lc{Δ}'yı genelleştiririz; böylece
tümevarım hipotezi her \lc{Δ} için geçerli olur. Bölüm~7'deki ``önermeyi
güçlendirme'' tekniğinin bir başka örneği budur.

% ---------------------------------------------------------------------
\section{Geçerlilik neyi söylemez?}\label{sec:gecerlilik-sinir}

Geçerlilik teoremi bir yöndedir: Türetilebilen her şey doğrudur. Ters
yön, yani doğru olan her şeyin türetilebilir olması, \emph{tamlık}
teoremidir ve ispatı çok daha zordur (\cref{bol:tamlik}).

Bu noktada RAA kuralının rolü önemlidir. RAA'yı sistemden çıkarırsak,
geri kalan on kural \emph{yapıcı} doğal çıkarımı oluşturur. Bu sistem de
iki değerli anlambilime göre geçerlidir: Ne de olsa RAA'sız her türetim,
RAA'lı sistemde de bir türetimdir. Ama tam değildir: Peirce yasası bir
totolojidir, yine de RAA olmadan türetilemez. Bunu ispatlamak için iki
değerli anlambilim yetmez; Kripke modellerine ihtiyaç vardır
(\ciltiki{2}).

\begin{onemlifikir}[Önemli Fikir: İspat sistemi ile anlambilimin uyumu]
Bir ispat sistemi ile bir anlambilim arasında iki ilişki olabilir:
Sistem anlambilime göre \emph{geçerli} olabilir (çok fazla şey
türetmez) ve \emph{tam} olabilir (yeterince şey türetir). Klasik doğal
çıkarım, iki değerli anlambilime göre hem geçerli hem tamdır. Yapıcı
doğal çıkarım ise iki değerli anlambilime göre geçerli ama tam
değildir; Kripke anlambilimine göre hem geçerli hem tamdır. Her mantığın
kendine uygun anlambilimi vardır.
\end{onemlifikir}

% ---------------------------------------------------------------------
\begin{arastirmanotu}
\textbf{Yapısal kuralların ötesi.} Zayıflatma ve daraltma kurallarının
\emph{kısıtlandığı} mantıklara \emph{alt yapısal mantıklar} denir.
Jean-Yves Girard'ın \emph{doğrusal mantığı} \citep{girard1987}, her
varsayımın tam olarak bir kez kullanılmasını ister ve kaynak kullanımını
modellemek için bilgisayar biliminde önemli bir rol oynar; programlama
dillerindeki kaynak ve sahiplik tiplerinin kuramsal arka planında yer alır. Bu tür
mantıkları \Lean{}'de formalize etmek için bağlamı liste yerine çoklu
küme (\emph{multiset}) olarak tutmak gerekir.

\textbf{Metamantığın formalizasyonu.} Önermeler mantığı için geçerlilik
ve tamlık teoremlerinin formalizasyonu, ispat asistanları için klasik
bir alıştırmadır; birinci mertebe mantık ve Gödel teoremleri ise çok daha
zordur. John Harrison'ın HOL Light'taki çalışmaları \citep{harrison2009}
ve Lean~4'teki Foundation projesi \citep{foundation, saitou2026} bu
alanın önemli örnekleridir.
\end{arastirmanotu}

\begin{bolumozeti}
\begin{itemize}
  \item Klasik doğal çıkarım, \Lean{}'de on bir kurucusu olan
        tümevarımsal bir yüklemdir: \lc{Turetilir Γ A} ($\Gamma \vdash A$).
        Bir türetim, bu yüklemin bir kanıtıdır.
  \item Nesne dilinde türetimler \leantac{apply} ile kurucuları
        uygulayarak kurulur; bağlam üyelikleri \lc{by decide} ile
        denetlenir.
  \item Geçerlilik teoremi ($\Gamma \vdash A \Rightarrow \Gamma \models A$)
        kural tümevarımıyla ispatlanır: Her kural doğruluğu korur.
  \item Geçerlilik, türetilemezliği karşı değerlemelerle göstermeyi ve
        sistemin tutarlılığını ($\nvdash \bot$) sağlar.
  \item Zayıflatma, değiştirme ve daraltma tek bir teoremden çıkar;
        ispatında bağlam genelleştirilir.
  \item RAA, iki değerli anlambilime göre geçerlidir; RAA'sız yapıcı
        sistem geçerlidir ama bu anlambilime göre tam değildir.
\end{itemize}
\end{bolumozeti}

\begin{alistirmalar}
\alistirma{1} $p \vdash p$. (\lc{a11\_1})

\alistirma{1} $K$ aksiyomu: her $A, B$ için $\vdash A \oise B \oise A$.
Bağlamın hangi elemanını kullandığınıza dikkat edin. (\lc{a11\_2})

\alistirma{2} $p \ove q \vdash q \ove p$. (\lc{a11\_3})

\alistirma{2} Hipotetik kıyas, nesne dilinde:
$p \oise q,\; q \oise r \vdash p \oise r$. Türetim ağacını da çizin ve
Bölüm~2'deki \lc{hipotetik\_kiyas} ile karşılaştırın. (\lc{a11\_4})

\alistirma{2} $\nvdash p \oise q$. Hangi değerlemeyi seçtiniz?
(\lc{a11\_5})

\alistirma{2} $p \oveya q \nvdash p$. (\lc{a11\_6})

\alistirma{3} Üçüncü hâlin olanaksızlığı: $\vdash p \oveya {\sim}p$.
Bölüm~3'teki \lc{em\_yanlis\_olamaz} ispatıyla karşılaştırın: RAA hangi
adımın yerini tutuyor? (\lc{a11\_7})

\alistirma{3} Çift değilleme eleme: ${\sim}{\sim}A \vdash A$. (\lc{a11\_8})

\alistirma{2} (Kâğıt-kalem) \Cref{tab:gecerlilik-durumlar}'deki $\vee$E
satırını ayrıntılı bir matematiksel ispata dönüştürün.

\alistirma{3} (Proje) RAA kuralı olmayan bir \lc{TuretilirY} yüklemi
tanımlayın. \lc{TuretilirY Γ A → Turetilir Γ A} önermesini ispatlayın ve
buradan yapıcı sistemin geçerliliğini tek satırda çıkarın.

\alistirma{3} (Tartışma) Bir arkadaşınız, geçerlilik teoreminin
\lc{simp} kullanmadan ispatlanmasını istiyor. Hangi durumlar zorlaşır?
\lc{Bool} yerine \lc{Prop} değerli bir anlambilim (\lc{deger v A : Prop})
kullansaydık RAA durumunda ne değişirdi?
\end{alistirmalar}
   % Geçerlilik teoremi
% =====================================================================
%  Bölüm 12 — Tamlık Teoremi
% =====================================================================
\chapter{Tamlık Teoremi}\label{bol:tamlik}

\begin{bolumhedefleri}
\begin{itemize}
  \item Tamlık teoreminin ne söylediğini, zayıf ve güçlü biçimlerini
        ayırt etmeyi,
  \item Uzun bir metamantık ispatını küçük, bağımsız önsavlara bölmeyi,
  \item Kalmár'ın yöntemiyle doğruluk tablosunun her satırını bir
        türetime dönüştürmeyi,
  \item Atomları teker teker eleyerek önermeler mantığının tamlık
        teoremini \Lean{}'de eksiksiz ispatlamayı,
  \item Lindenbaum önsavı ve Henkin yöntemiyle birinci mertebe mantığın
        tamlık teoremine giden yolu görmeyi.
\end{itemize}
\end{bolumhedefleri}

Bölüm~11'de köprünün bir ayağını kurduk: Türetilebilen her şey
doğrudur. Bu bölümde öteki ayağı kuruyoruz:
\[
  \Gamma \models A \quad\Longrightarrow\quad \Gamma \vdash A.
\]
İki teorem birlikte, ispat sistemimizin iki değerli anlambilimi
\emph{tam olarak} yakaladığını söyler. Ne fazla şey türetir, ne eksik.
Geçerlilik ispatı, türetim ağacı üzerinde düz bir tümevarımdı. Tamlık
ise bambaşka bir iş ister: Elimizde yalnızca ``her değerlemede doğru''
bilgisi vardır ve bu bilgiden, hiç görmediğimiz bir türetim ağacını
\emph{inşa etmemiz} gerekir.

% ---------------------------------------------------------------------
\section{Tamlığın ifadesi}\label{sec:tamlik-ifade}

\begin{tanim}{Zayıf ve güçlü tamlık}{tamlik-turleri}
Bir ispat sistemi, bir anlambilime göre
\begin{itemize}
  \item \textbf{zayıf tamdır}, eğer her totoloji türetilebiliyorsa:
        $\models A \Rightarrow\; \vdash A$;
  \item \textbf{güçlü tamdır}, eğer her mantıksal sonuç
        türetilebiliyorsa: $\Gamma \models A \Rightarrow \Gamma \vdash
        A$. Burada $\Gamma$ sonsuz bir küme de olabilir.
\end{itemize}
\end{tanim}

\noindent Bizim \lc{Sonuc} ve \lc{Turetilir} tanımlarımızda $\Gamma$ bir
\emph{liste}, dolayısıyla sonludur. Sonlu $\Gamma$ için güçlü tamlık
zayıf tamlıktan kolayca çıkar: Varsayımları birer öncül olarak formülün
içine taşırız (\cref{sec:tamlik-teoremi}). Sonsuz $\Gamma$ için ise ek
bir fikir, \emph{tıkızlık} gerekir (\cref{sec:lindenbaum}).

\begin{tarihce}[Tarihçe: Tamlığın kısa tarihi]
Önermeler mantığının tamlığı ilk olarak Paul Bernays'in 1918 tarihli
(yayımlanmamış) habilitasyon tezinde ve bağımsız olarak Emil Post'un
1921'deki makalesinde ispatlandı \citep{post1921}. Post'un ispatı,
doğruluk tablolarının bugünkü biçimini de ilk kez sistemli olarak
kullanır. Bu bölümde izleyeceğimiz yöntem, Macar matematikçi László
Kalmár'ın 1935'te verdiği zarif ispattır \citep{kalmar1935}. Mendelson'ın
klasik ders kitabı da bu ispatı kullanır \citep{mendelson2015}. Birinci
mertebe mantığın tamlığını ise Kurt Gödel 1929'daki doktora teziyle
ispatladı \citep{godel1930}. Bugün ders kitaplarında okutulan ispat, Leon
Henkin'in 1949'daki ispatıdır \citep{henkin1949}.
\end{tarihce}

% ---------------------------------------------------------------------
\section{İspatın haritası}\label{sec:tamlik-harita}

Tamlık ispatı, kitabın şimdiye kadarki en uzun \Lean{} ispatı olacak.
Uzun bir ispatı yönetmenin yolu, onu her biri tek başına anlaşılabilen ve
tek başına derlenen \emph{önsavlara} bölmektir.
\Cref{fig:tamlik-harita}, ispatın beş aşamasını ve birbirine nasıl
bağlandığını gösteriyor.

\begin{figure}[htbp]
\centering
\begin{tikzpicture}[
    asama/.style={draw=#1, fill=#1!8, very thick, rounded corners=5pt,
                  text width=3.0cm, align=center, inner sep=5pt,
                  font=\small, minimum height=1.35cm},
    ok/.style={-{Stealth[length=2.6mm]}, very thick, black!45},
    no/.style={font=\baslikfont\scriptsize\bfseries, text=white,
               fill=#1, circle, inner sep=1.6pt},
  ]
  \node[asama=anarenk]     (a1) at (0,0)    {\textbf{Türetilmiş kurallar}\\{\scriptsize \lc{zay}, \lc{ucuncu\_hal}, \dots}};
  \node[asama=ornekframe]  (a2) at (4.2,0)  {\textbf{Kalmár önsavı}\\{\scriptsize her satır bir türetim}};
  \node[asama=ikincil]     (a3) at (8.4,0)  {\textbf{Atom eleme}\\{\scriptsize $p \oveya {\sim}p$ ile birleştir}};
  \node[asama=dikkatframe] (a4) at (8.4,-2.4) {\textbf{Zayıf tamlık}\\{\scriptsize $\models A \Rightarrow\; \vdash A$}};
  \node[asama=teoremframe] (a5) at (4.2,-2.4) {\textbf{Tamlık}\\{\scriptsize $\Gamma \models A \Rightarrow \Gamma \vdash A$}};
  \node[asama=arastirmaframe] (a6) at (0,-2.4) {\textbf{$\vdash \;=\; \models$}\\{\scriptsize geçerlilik ile birlikte}};
  \draw[ok] (a1) -- (a2); \draw[ok] (a2) -- (a3); \draw[ok] (a3) -- (a4);
  \draw[ok] (a4) -- node[above, font=\baslikfont\scriptsize]{\lc{kapat}} (a5);
  \draw[ok] (a5) -- (a6);
  \foreach \n/\c/\k in {a1/anarenk/1, a2/ornekframe/2, a3/ikincil/3, a4/dikkatframe/4, a5/teoremframe/5}
    \node[no=\c] at (\n.north west) {\k};
\end{tikzpicture}
\caption{Tamlık ispatının haritası. Her kutu, \lc{Bolum12.lean}
dosyasında bir ya da birkaç önsava karşılık gelir.}
\label{fig:tamlik-harita}
\end{figure}

\noindent Ana fikir şudur. Bir $A$ totolojisinin doğruluk tablosunun her
satırı, atomlara birer değer verir. Kalmár, \emph{her satırın kendisinin
bir türetime dönüştürülebildiğini} gözlemledi. Satırda doğru olan
atomları, yanlış olanların ise değillerini varsayarsak, bu varsayımlardan
$A$ türetilebilir. Sonra satırları ikişer ikişer birleştirerek
varsayımlardan kurtuluruz.

\subsection{Türetilmiş kurallar}

\Lean{}'deki \lc{Turetilir} yüklemi yalnızca on bir temel kurala sahip.
Uzun türetimler yazarken aynı küçük kalıpları tekrar tekrar kullanırız.
Matematikte bir önsavı bir kez ispatlayıp sonra kullanırız; aynı şeyi
ispat sistemimizde de yapabiliriz. Bir \emph{türetilmiş kural},
temel kurallardan ispatlanan ve sonra temel kural gibi kullanılan bir
teoremdir.

En sık kullanacağımız iki araç, bağlamın başındaki varsayımı kullanmak ve
bağlamın başına yeni bir varsayım eklemektir:

\leandosya[Temel araçlar]{Bolum12}{temel_araclar}

\noindent Klasik mantığın kalbi olan üçüncü hâlin olanaksızlığını
Bölüm~11'in alıştırmalarında boş bağlamda türetmiştik (\lc{a11\_7}).
Şimdi onu \emph{her} bağlamda ve \emph{her} formül için türetiyor, ondan
da \emph{durum ayrımı} kuralını çıkarıyoruz:

\leandosya[Üçüncü hâl ve durum ayrımı]{Bolum12}{ucuncu_hal}

\begin{center}
\begin{prooftree}
  \hypo{B, \Gamma \vdash A}
  \hypo{{\sim}B, \Gamma \vdash A}
  \infer2[Durum]{\Gamma \vdash A}
\end{prooftree}
\end{center}

\noindent Bu kural, ispatın son aşamasında atomları elemek için
kullanılacak. Kalmár önsavı içinse her bağlacın doğruluk tablosunu
taklit eden kurallara ihtiyacımız var (\cref{tab:turetilmis-kurallar}).

\leandosya[Bağlaçlar için türetilmiş kurallar]{Bolum12}{turetilmis_kurallar}

\begin{table}[htbp]
\centering
\small
\renewcommand{\arraystretch}{1.5}
\rowcolors{2}{tablosatir}{white}
\begin{tabular}{>{\raggedright\arraybackslash}p{3.0cm}
                >{\centering\arraybackslash}p{4.4cm}
                >{\raggedright\arraybackslash}p{5.0cm}}
\rowcolor{tabloust}
\tablobaslik{Lean adı} & \tablobaslik{Kural} & \tablobaslik{Doğruluk tablosundaki karşılığı} \\
\lc{ve\_yanlis1}    & ${\sim}A \;\vdash\; {\sim}(A \ove B)$          & $A$ yanlışsa $A \ove B$ yanlış \\
\lc{ve\_yanlis2}    & ${\sim}B \;\vdash\; {\sim}(A \ove B)$          & $B$ yanlışsa $A \ove B$ yanlış \\
\lc{veya\_yanlis}   & ${\sim}A,\, {\sim}B \;\vdash\; {\sim}(A \oveya B)$ & ikisi de yanlışsa $A \oveya B$ yanlış \\
\lc{ise\_on\_yanlis} & ${\sim}A \;\vdash\; A \oise B$                 & öncül yanlışsa $A \oise B$ doğru \\
\lc{ise\_son\_dogru} & $B \;\vdash\; A \oise B$                      & sonuç doğruysa $A \oise B$ doğru \\
\lc{ise\_yanlis}    & $A,\, {\sim}B \;\vdash\; {\sim}(A \oise B)$    & tek yanlış satır \\
\end{tabular}
\caption{Kalmár önsavı için gereken türetilmiş kurallar. Doğru durumlar
($A, B \vdash A \ove B$ gibi) zaten temel kurallardır.}
\label{tab:turetilmis-kurallar}
\end{table}

\begin{onemlifikir}[Önemli Fikir: İspat sistemi içinde ispat sistemi]
Türetilmiş kurallar, \Lean{}'in kendi ispatlarında yaptığımız şeyin nesne
dilindeki yansımasıdır. \lc{ve\_yanlis1} bir kez ispatlandıktan sonra
\lc{Turetilir}'in on ikinci kuralıymış gibi kullanılabilir. Sistemi
değiştirmeyiz; yalnızca onun hakkında bir teorem ispatlarız. Bu, ispat
sistemini küçük tutup (geçerlilik ispatı kısa kalsın diye) kullanımını
zenginleştirmenin standart yoludur.
\end{onemlifikir}

% ---------------------------------------------------------------------
\section{Kalmár yöntemi}\label{sec:kalmar}

\subsection{İşaretli atomlar}

Bir $v$ değerlemesi ve bir $p_n$ atomu için \emph{işaretli atomu}
\[
  p_n^v \;=\;
  \begin{cases}
    p_n        & v(n) = \mathsf{true} \text{ ise},\\
    {\sim}p_n  & v(n) = \mathsf{false} \text{ ise}
  \end{cases}
\]
diye tanımlarız. $p_n^v$, $v$ altında her zaman doğrudur. İlk $k$ atomun
işaretli hâllerinden oluşan bağlama $\Delta_v^k = [p_{k-1}^v, \dots,
p_1^v, p_0^v]$ diyelim. Bu bağlam, doğruluk tablosunun bir satırının
nesne dilindeki ifadesidir.

\leandosya[İşaretli atomlar]{Bolum12}{isaretli}

\noindent \lc{atomBaglami}'nın sayı üzerinde özyinelemeyle tanımlanması,
ileride atomları teker teker elemeyi kolaylaştıracak: \lc{atomBaglami v
(k + 1)} bağlamının başında her zaman $p_k^v$ durur. İki basit özellik
daha ispatlıyoruz. İşaretli atomun açık biçimi ve bağlamda bulunması:

\leandosya[İşaretli atomların özellikleri]{Bolum12}{isaretli_ozellik}

\subsection{Kalmár önsavı}

\begin{teorem}{Kalmár önsavı}{kalmar}
$A$'nın bütün atomları $p_0, \dots, p_{k-1}$ arasındaysa, her $v$ için
\[
  \Delta_v^k \vdash A^v,
  \qquad\text{burada}\qquad
  A^v =
  \begin{cases}
    A        & \lBrack A \rBrack_v = \mathsf{true} \text{ ise},\\
    {\sim}A  & \lBrack A \rBrack_v = \mathsf{false} \text{ ise}.
  \end{cases}
\]
\end{teorem}

\noindent Önsav, \emph{doğru olanı da yanlış olanı da} türetebildiğimizi
söylüyor. Bu güçlendirme, tümevarımın işlemesi için şarttır. $A \oise B$
yanlış olduğunda ${\sim}(A \oise B)$'yi türetmek için $A$'yı ve
${\sim}B$'yi türetebiliyor olmamız gerekir. Yani tümevarım hipotezinin
hem doğru hem yanlış yüzüne ihtiyaç vardır (Bölüm~7'deki ``önermeyi
güçlendirme'' tekniği). \Cref{fig:kalmar-tablo}, önsavın $p \oise q$ için
ne söylediğini satır satır gösteriyor.

\begin{figure}[htbp]
\centering
\renewcommand{\arraystretch}{1.6}
\setlength{\tabcolsep}{9pt}
\begin{tabular}{cc c >{\centering\arraybackslash}p{3.0cm} c >{\raggedright\arraybackslash}p{3.3cm}}
\rowcolor{tabloust}
\tablobaslik{$p$} & \tablobaslik{$q$} & \tablobaslik{$p \oise q$} &
\tablobaslik{Bağlam $\Delta_v$} & & \tablobaslik{Türetilen ($A^v$)} \\
\cellcolor{ikincil!15}\textsf{D} & \cellcolor{ikincil!15}\textsf{D} & \textsf{D} &
  $q,\; p$ & $\vdash$ & $p \oise q$ \quad{\scriptsize\textcolor{leantac}{\lc{ise\_son\_dogru}}} \\
\cellcolor{ikincil!15}\textsf{D} & \cellcolor{vurgu!15}\textsf{Y} & \cellcolor{vurgu!15}\textsf{Y} &
  ${\sim}q,\; p$ & $\vdash$ & \textcolor{teoremframe}{${\sim}(p \oise q)$} \quad{\scriptsize\textcolor{leantac}{\lc{ise\_yanlis}}} \\
\cellcolor{vurgu!15}\textsf{Y} & \cellcolor{ikincil!15}\textsf{D} & \textsf{D} &
  $q,\; {\sim}p$ & $\vdash$ & $p \oise q$ \quad{\scriptsize\textcolor{leantac}{\lc{ise\_on\_yanlis}}} \\
\cellcolor{vurgu!15}\textsf{Y} & \cellcolor{vurgu!15}\textsf{Y} & \textsf{D} &
  ${\sim}q,\; {\sim}p$ & $\vdash$ & $p \oise q$ \quad{\scriptsize\textcolor{leantac}{\lc{ise\_on\_yanlis}}} \\
\end{tabular}
\caption{Kalmár önsavı $p \oise q$ için. Doğruluk tablosunun her satırı
bir türetime dönüşür: Bağlam satırdaki atom değerlerini, sağ taraf ise
formülün o satırdaki değerini kodlar. Son sütun, satırın hangi türetilmiş
kuralla elde edildiğini gösterir.}
\label{fig:kalmar-tablo}
\end{figure}

\noindent Önsavı \Lean{}'de ifade etmek için ``$b$ değerine uygun
biçimde türetir'' anlamına gelen küçük bir yüklem tanımlıyoruz. $A^v$'yi
bir \leankw{if} ifadesiyle tanımlamak yerine iki yönü bir \lc{∧} ile
yazmak, ispatlarda \leankw{if} ifadelerini açma zahmetinden kurtarır:

\leandosya[Kalmár yüklemi]{Bolum12}{kalmar_tanim}

\noindent \leankw{abbrev} kullanmamızın sebebi, \lc{Kalmar}'ın \lc{∧}
olarak açılabilmesidir. Böylece \lc{⟨\_, \_⟩} ile kurulabilir ve
\lc{hA.1}, \lc{hA.2} ile parçalanabilir. Şimdi her bağlaç için ayrı bir
önsav ispatlıyoruz. Her önsav, o bağlacın doğruluk tablosunu satır
satır ele alır:

\leandosya[Bağlaç adımları]{Bolum12}{kalmar_baglaclar}

\noindent Bu önsavlarda \lc{a} ve \lc{b}, alt formüllerin doğruluk
değerleridir. \leantac{cases a with} her satırı ayrı bir hedefe
dönüştürür. Olanaksız yönler (örneğin \lc{(false \&\& b) = true}) tek bir
\leantac{simp at h} ile kapanır. Geri kalan her hedef, ya bir temel
kuralla ya da \cref{tab:turetilmis-kurallar}'deki bir türetilmiş
kuralla tek satırda biter.

Önsavın kendisi artık yalnızca bir yapısal tümevarımdır. Her durum, ilgili
bağlaç önsavına yönlendirilir:

\leandosya[Kalmár önsavı]{Bolum12}{kalmar}

\begin{dikkat}[Dikkat: Tanımsal eşitlik iş başında]
\lc{kalmar\_ve}'nin sonucu \lc{Kalmar Δ (a \&\& b) (A {\sembolkodfont ⋀} B)} biçimindedir.
Hedef ise \lc{Kalmar Δ (deger v (A {\sembolkodfont ⋀} B)) (A {\sembolkodfont ⋀} B)}'dir. \Lean{} bu ikisini
\emph{tanımsal olarak} eşit görür, çünkü \lc{deger}'in tanımı açıldığında
ikisi aynı terim olur. Bu sayede \leantac{exact} hiçbir yeniden yazma
olmadan çalışır. \lc{deger}'i bir tablo yerine özyinelemeli bir fonksiyon
olarak tanımlamanın bir kazancı da budur.
\end{dikkat}

% ---------------------------------------------------------------------
\section{Atomların elenmesi}\label{sec:atom-eleme}

$A$ bir totolojiyse, Kalmár önsavı \emph{her} $v$ için $\Delta_v^k
\vdash A$ verir; çünkü $A^v$ her zaman $A$'nın kendisidir. Elimizde
$2^k$ tane türetim var, her satır için bir tane. Bunları birleştirip
bağlamı boşaltmak istiyoruz.

Fikir şudur: Yalnızca $p_{k-1}$'in değerinde farklılaşan iki satırı
alalım. Bu satırların bağlamları
\[
  p_{k-1},\, \Delta \vdash A
  \qquad\text{ve}\qquad
  {\sim}p_{k-1},\, \Delta \vdash A
\]
biçimindedir; durum ayrımı kuralı $\Delta \vdash A$ verir. Böylece
$p_{k-1}$'i bağlamdan çıkarmış oluruz. Bunu $k$ kez tekrarlarız
(\cref{fig:atom-eleme}).

\begin{figure}[htbp]
\centering
\begin{tikzpicture}[
    yaprak/.style={draw=ornekframe, fill=ornekbg, thick, rounded corners=3pt,
                   font=\small, inner sep=4pt, minimum width=2.6cm},
    ara/.style={draw=ikincil, fill=fikirbg, thick, rounded corners=3pt,
                font=\small, inner sep=4pt, minimum width=2.4cm},
    kok/.style={draw=teoremframe, fill=teorembg, very thick, rounded corners=4pt,
                font=\small\bfseries, inner sep=5pt},
    ok/.style={-{Stealth[length=2.2mm]}, thick, black!45},
    etiket/.style={font=\baslikfont\scriptsize\bfseries, text=ikincil},
  ]
  \node[yaprak] (l1) at (-5.1,0) {$q,\; p \vdash A$};
  \node[yaprak] (l2) at (-1.7,0) {${\sim}q,\; p \vdash A$};
  \node[yaprak] (l3) at (1.7,0)  {$q,\; {\sim}p \vdash A$};
  \node[yaprak] (l4) at (5.1,0)  {${\sim}q,\; {\sim}p \vdash A$};
  \node[ara] (m1) at (-3.4,-1.7) {$p \vdash A$};
  \node[ara] (m2) at (3.4,-1.7)  {${\sim}p \vdash A$};
  \node[kok] (r) at (0,-3.3) {$\vdash A$};
  \draw[ok] (l1) -- (m1); \draw[ok] (l2) -- (m1);
  \draw[ok] (l3) -- (m2); \draw[ok] (l4) -- (m2);
  \draw[ok] (m1) -- (r);  \draw[ok] (m2) -- (r);
  \node[etiket] at (-3.4,-0.85) {$q \oveya {\sim}q$};
  \node[etiket] at (3.4,-0.85)  {$q \oveya {\sim}q$};
  \node[etiket] at (0,-2.35)    {$p \oveya {\sim}p$};
  \node[font=\baslikfont\scriptsize, text=ornekframe, anchor=west] at (6.6,0)
       {\begin{tabular}{@{}l@{}}Kalmár:\\ $2^k$ satır\end{tabular}};
  \node[font=\baslikfont\scriptsize, text=ikincil, anchor=west] at (6.6,-1.7)
       {\begin{tabular}{@{}l@{}}\lc{atom\_ele}\\(1. adım)\end{tabular}};
  \node[font=\baslikfont\scriptsize, text=teoremframe, anchor=west] at (6.6,-3.3)
       {\begin{tabular}{@{}l@{}}\lc{atom\_ele}\\(2. adım)\end{tabular}};
\end{tikzpicture}
\caption{İki atomlu bir $A$ totolojisi için atomların elenmesi.
Yapraklar Kalmár önsavından gelir. Her birleştirme, bir atom için durum
ayrımıdır. $k$ atom için ağaç $k$ katlıdır ve $2^k$ yaprağı vardır.}
\label{fig:atom-eleme}
\end{figure}

\noindent Bunu \Lean{}'de yapmak için, bir $v$ değerlemesinden yalnızca
$j$'inci atomu değiştirilmiş yeni değerlemeler üretebilmeliyiz. Ayrıca
$j$'den küçük atomlarda anlaşan değerlemelerin aynı atom bağlamını
verdiğini bilmeliyiz:

\leandosya[Değerlemeyi güncellemek]{Bolum12}{ata}

\noindent Artık eleme adımını yazabiliriz. Hipotez, $j+1$ atomlu
\emph{her} satırın $A$'yı türettiğini söylüyor. $v$'yi verilmiş kabul
edip $j$'inci atomu bir kez \lc{true}, bir kez \lc{false} yapıyoruz.
İki türetimi durum ayrımıyla birleştirince $j$ atomlu satır elde ediliyor:

\leandosya[Atom eleme]{Bolum12}{atom_ele}

\noindent \lc{atomlari\_ele}, \lc{atom\_ele}'yi $k$ kez uygular. Sayı
üzerinde yapısal özyinelemeyle tanımlanmış bir \emph{teoremdir}:
\lc{k + 1} durumu, \lc{k} durumunu çağırır. $k = 0$ olduğunda atom
bağlamı boştur, dolayısıyla herhangi bir değerleme (burada hep yanlış
olanı) işimizi görür.

% ---------------------------------------------------------------------
\section{Tamlık teoremi}\label{sec:tamlik-teoremi}

Parçaları birleştirelim. $A$'nın atomları $p_0, \dots, p_m$
arasındadır; burada $m$, \lc{enBuyukAtom A} değeridir. Dolayısıyla
$k = m + 1$ yeterlidir.

\begin{teorem}{Zayıf tamlık}{zayif-tamlik}
Her $A$ totolojisi türetilebilir: $\models A \Rightarrow\; \vdash A$.
\end{teorem}

\leandosya[Zayıf tamlık]{Bolum12}{zayif_tamlik}

\noindent Varsayımlı hâl için, varsayımları öncül olarak formülün içine
taşıyan bir \lc{kapat} fonksiyonu tanımlıyoruz. Sonra iki şey
gösteriyoruz: Anlamsal sonuç ilişkisi bu taşımayla korunur, türetim ise
geri alınabilir. İkincisi, Bölüm~5'teki \emph{çıkarım teoreminin} ($\to$G
kuralının) tersidir:

\leandosya[Varsayımları içeri taşımak]{Bolum12}{kapat}

\begin{teorem}{Tamlık (önermeler mantığı)}{tamlik}
Her sonlu $\Gamma$ ve her $A$ için, $\Gamma \models A$ ise $\Gamma
\vdash A$. Geçerlilikle birlikte:
\[
  \Gamma \vdash A \quad\Longleftrightarrow\quad \Gamma \models A .
\]
\end{teorem}

\leandosya[Tamlık teoremi]{Bolum12}{tamlik}

\begin{figure}[htbp]
\centering
\begin{tikzpicture}[
    kutu/.style={draw=#1, fill=#1!8, very thick, rounded corners=6pt,
                 text width=3.6cm, align=center, inner sep=8pt},
  ]
  \node[kutu=anarenk] (s) at (0,0)
       {\baslikfont\bfseries Sözdizim\\[2pt]\normalfont $\Gamma \vdash A$\\[2pt]
        {\scriptsize türetim ağaçları}\\{\scriptsize \lc{Turetilir}}};
  \node[kutu=ikincil] (a) at (8,0)
       {\baslikfont\bfseries Anlambilim\\[2pt]\normalfont $\Gamma \models A$\\[2pt]
        {\scriptsize değerlemeler}\\{\scriptsize \lc{Sonuc}}};
  \draw[-{Stealth[length=3.2mm]}, very thick, teoremframe]
       ([yshift=6pt]s.east) -- node[above, font=\baslikfont\small\bfseries]{geçerlilik}
       node[below, font=\scriptsize]{Bölüm 11} ([yshift=6pt]a.west);
  \draw[-{Stealth[length=3.2mm]}, very thick, dikkatframe]
       ([yshift=-18pt]a.west) -- node[below, font=\baslikfont\small\bfseries]{tamlık}
       node[above, font=\scriptsize]{Bölüm 12} ([yshift=-18pt]s.east);
\end{tikzpicture}
\caption{Önermeler mantığının iki yüzü. Geçerlilik ve tamlık birlikte,
türetilebilirlik ile mantıksal sonucun \emph{aynı} ilişki olduğunu
söyler: \lc{turetilir\_iff\_sonuc}.}
\label{fig:iki-yuz}
\end{figure}

\begin{onemlifikir}[Önemli Fikir: Tamlık bir ispat \emph{üretir}]
Tamlık ispatımız yapıcıdır: \lc{tamlik\_bos}, verilen bir totoloji
için gerçek bir türetim ağacı inşa eder. Bu ağaç büyüktür. $k$ atomlu bir
formül için $2^k$ Kalmár türetimini birleştirir. Ama vardır ve
mekaniktir. Bu yüzden önermeler mantığında ``türetilebilir mi?''
sorusu, ``doğruluk tablosu her satırda doğru mu?'' sorusuna indirgenir.
Bölüm~10'daki \lc{totolojiMi} fonksiyonu, aslında bir \emph{ispat arama}
algoritmasıdır.
\end{onemlifikir}

% ---------------------------------------------------------------------
\section{Sonuçlar}\label{sec:tamlik-sonuclar}

Tamlığın ilk getirisi pratiktir: Artık bir türetimi elle kurmak zorunda
değiliz. Bölüm~11'de Peirce yasasını on adımlık bir türetimle elde
etmiştik. Şimdi doğruluk tablosunu kontrol etmek yeterli. İkinci getiri
kuramsaldır: Tutarlılık ile sağlanabilirlik aynı şeydir.

\leandosya[Tamlığın sonuçları]{Bolum12}{sonuclar}

\begin{teorem}{Model varlığı}{model-varligi}
$\Gamma \nvdash \bot$ ise $\Gamma$'nın bütün formüllerini doğrulayan bir
$v$ değerlemesi vardır.
\end{teorem}

\noindent Bu teoremin ispatında \lc{Classical.byContradiction}
kullandık, çünkü değerlemenin \emph{varlığını} çelişki yoluyla
gösteriyoruz. Aslında burada da yapıcı bir ispat mümkündür:
Formülün atomları sonlu olduğu için doğruluk tablosunun satırlarını tek
tek deneyebiliriz (\lc{tumDegerlemeler}). Bu ispat \lc{a12\_8}'in
ötesinde bir alıştırma olarak bırakılmıştır.

\begin{dikkat}[Dikkat: Hesaplama karmaşıklığı]
\lc{tamlik\_bos}'un ürettiği türetim, atom sayısında üsteldir. Daha
kısa türetimler var mı? Bu, \emph{ispat karmaşıklığı} alanının temel
sorusudur. Cook ve Reckhow, her totolojinin polinom boyutlu bir türetimi
olan bir ispat sistemi bulunmasının $\mathrm{NP} = \mathrm{coNP}$'ye denk
olduğunu gösterdi \citep{cookreckhow1979}. Bu, açık bir problemdir. Belirli
sistemler için alt sınırlar bilinir: Örneğin güvercin yuvası ilkesinin
\emph{çözünürlük} (resolution) türetimleri üstel uzunluktadır
\citep{haken1985}. Yani tamlık ``bir türetim vardır'' der; ``kısa bir
türetim vardır'' demez.
\end{dikkat}

% ---------------------------------------------------------------------
\section{Maksimal tutarlı kümeler ve Lindenbaum önsavı}\label{sec:lindenbaum}

Kalmár yöntemi, sonlu sayıda atom ve sonlu bir $\Gamma$ için çalışır.
$\Gamma$ sonsuz olduğunda, ya da birinci mertebe mantığa geçtiğimizde,
doğruluk tablosu kullanılamaz. Bu durumlarda kullanılan genel yöntem,
modeli \emph{sözdizimin kendisinden} inşa etmektir. Bu kesimin
ispatlarını kâğıt üzerinde veriyoruz. Birinci mertebe mantık için
\Lean{} formalizasyonu Cilt~II'nin ilk bölümünün konusudur.

\begin{tanim}{Maksimal tutarlı küme}{maksimal-tutarli}
Bir $\Delta$ formül kümesi
\begin{itemize}
  \item \textbf{tutarlıdır}, eğer $\Delta \nvdash \bot$ ise;
  \item \textbf{maksimal tutarlıdır}, eğer tutarlıysa ve kendisine
        eklenen her yeni formül onu tutarsız yapıyorsa.
\end{itemize}
\end{tanim}

\noindent Maksimal tutarlı bir $\Delta$, ``her soruya kesin bir cevap
veren'' bir kümedir. Her $A$ için ya $A \in \Delta$ ya da ${\sim}A \in
\Delta$'dır, ikisi birden asla. Üstelik $\Delta$ türetim altında
kapalıdır ve bağlaçlarla uyumludur:
\[
  A \oise B \in \Delta \iff (A \notin \Delta \text{ ya da } B \in \Delta),
  \qquad
  A \ove B \in \Delta \iff (A \in \Delta \text{ ve } B \in \Delta).
\]
Bu yüzden $v_\Delta(n) = \mathsf{true} \iff p_n \in \Delta$ diye
tanımlanan \emph{kanonik değerleme}, yapısal tümevarımla şunu sağlar:
$\lBrack A \rBrack_{v_\Delta} = \mathsf{true} \iff A \in \Delta$
(\emph{doğruluk önsavı}).

\begin{teorem}{Lindenbaum önsavı}{lindenbaum}
Her tutarlı $\Gamma$ kümesi, maksimal tutarlı bir $\Delta \supseteq
\Gamma$ kümesine genişletilebilir.
\end{teorem}

\noindent \emph{İspat taslağı.} Formüller sayılabilir olduğu için onları
$\varphi_0, \varphi_1, \varphi_2, \dots$ diye sıralayabiliriz
(\lc{Formul} için böyle bir sıralama, Bölüm~8'deki sayılabilirlik
tekniğiyle kurulur). Sonra adım adım genişletiriz
(\cref{fig:lindenbaum}):
\[
  \Delta_0 = \Gamma, \qquad
  \Delta_{n+1} =
  \begin{cases}
    \Delta_n \cup \{\varphi_n\} & \text{bu küme tutarlıysa},\\
    \Delta_n \cup \{{\sim}\varphi_n\} & \text{değilse},
  \end{cases}
  \qquad
  \Delta = \bigcup_n \Delta_n .
\]
Her adım tutarlılığı korur. Çünkü hem $\Delta_n, \varphi_n \vdash \bot$
hem $\Delta_n, {\sim}\varphi_n \vdash \bot$ olsaydı, durum ayrımı
$\Delta_n \vdash \bot$ verirdi. Birleşim de tutarlıdır, çünkü bir
türetim yalnızca sonlu sayıda varsayım kullanır ve bunlar bir
$\Delta_n$'nin içindedir. \hfill\textcolor{teoremframe}{\rule{0.55em}{0.55em}}

\begin{figure}[htbp]
\centering
\begin{tikzpicture}[
    kume/.style={draw=#1, fill=#1!7, thick, rounded corners=6pt},
  ]
  \draw[kume=teoremframe] (-0.3,-1.25) rectangle (11.0,1.25);
  \draw[kume=ikincil]     (-0.15,-1.0) rectangle (7.6,1.0);
  \draw[kume=ornekframe]  (0.0,-0.8) rectangle (5.1,0.8);
  \draw[kume=anarenk]     (0.15,-0.6) rectangle (2.6,0.6);
  \node[font=\small] at (1.35,0) {$\Delta_0 = \Gamma$};
  \node[font=\small] at (3.85,0.2) {$\Delta_1$};
  \node[font=\small] at (6.35,0.2) {$\Delta_2$\;\;$\cdots$};
  \node[font=\small\bfseries, text=teoremframe] at (9.3,0) {$\Delta = \bigcup_n \Delta_n$};
  \node[font=\baslikfont\scriptsize, text=ornekframe] at (3.85,-0.45) {$+\,\varphi_0$ ya da ${\sim}\varphi_0$};
  \node[font=\baslikfont\scriptsize, text=ikincil] at (6.35,-0.55) {$+\,\varphi_1$ ya da ${\sim}\varphi_1$};
  \node[font=\baslikfont\scriptsize, text=teoremframe, anchor=north] at (9.3,-0.35) {maksimal tutarlı};
\end{tikzpicture}
\caption{Lindenbaum yapısı. Her adımda sıradaki formül ya da değili
eklenir; tutarlılık her adımda korunur. Sonsuz birleşim, her formül
hakkında karar vermiş maksimal tutarlı bir kümedir.}
\label{fig:lindenbaum}
\end{figure}

\noindent Lindenbaum önsavı ile kanonik değerleme birleşince, model
varlığı teoremi \emph{sonsuz} $\Gamma$ için de elde edilir: Tutarlı
$\Gamma$'yı maksimal tutarlı $\Delta$'ya genişletiriz ve $v_\Delta$,
$\Gamma$'nın her formülünü doğrular. Buradan güçlü tamlık çıkar:
$\Gamma \nvdash A$ ise $\Gamma \cup \{{\sim}A\}$ tutarlıdır, bir modeli
vardır ve bu model $\Gamma \nvDash A$'ya tanıklık eder.

\begin{teorem}{Tıkızlık}{tikizlik}
Bir $\Gamma$ kümesinin her sonlu alt kümesi sağlanabilirse, $\Gamma$'nın
kendisi de sağlanabilirdir.
\end{teorem}

\noindent \emph{İspat.} $\Gamma$ sağlanamasaydı, güçlü tamlıkla $\Gamma
\vdash \bot$ olurdu. Bu türetim sonlu sayıda varsayım kullanır; o sonlu
alt küme geçerlilik teoremiyle sağlanamaz olurdu. \hfill\textcolor{teoremframe}{\rule{0.55em}{0.55em}}

\begin{dikkat}[Dikkat: Lindenbaum ve seçim]
Formüller sayılabilir olduğu için yukarıdaki yapı seçim aksiyomu
gerektirmez. Ama ``$\Delta_n \cup \{\varphi_n\}$ tutarlı mı?'' sorusuna
her adımda cevap verebilmemiz gerekir. Önermeler mantığında bu karar
verilebilirdir (tamlık ile doğruluk tablosuna indirgenir). Birinci
mertebe mantıkta değildir (\ciltiki{4}). Bu yüzden \Lean{}'de Lindenbaum
yapısı genellikle \leankw{Classical} ile, ya da sayılamaz diller için
Zorn önsavıyla yapılır. Mathlib'deki \lc{zorn\_subset} buna hazır bir
araçtır.
\end{dikkat}

% ---------------------------------------------------------------------
\section{Birinci mertebe mantık: Henkin yöntemi}\label{sec:henkin}

Birinci mertebe mantıkta bir model, bir \emph{tanım evreni} ve bu evren
üzerinde yüklemlerin yorumlarından oluşur. Model varlığını göstermek
için Henkin'in fikri, evreni \emph{terimlerin kendisinden} kurmaktır.
Ancak bir sorun vardır. $\exists x\, \varphi(x) \in \Delta$ olabilir,
ama hiçbir $t$ terimi için $\varphi(t) \in \Delta$ olmayabilir. Terim
modelinde varoluş iddiasının bir tanığı yoktur.

Henkin'in çözümü, dile her varoluş formülü için yeni bir \emph{tanık
sabiti} eklemektir: $c_{\exists x\varphi}$. Bu sabitle birlikte
\[
  \exists x\, \varphi(x) \;\oise\; \varphi(c_{\exists x\varphi})
\]
\emph{Henkin aksiyomu} da kümeye katılır. Bu aksiyomların eklenmesi
tutarlılığı korur. Sonra Lindenbaum önsavı uygulanır ve terim modeli
tanımlanır (\cref{fig:henkin}).

\begin{figure}[htbp]
\centering
\begin{tikzpicture}[
    adim/.style={draw=#1, fill=#1!8, very thick, rounded corners=5pt,
                 text width=2.55cm, align=center, inner sep=5pt,
                 font=\small, minimum height=1.6cm},
    ok/.style={-{Stealth[length=2.6mm]}, very thick, black!45},
  ]
  \node[adim=anarenk]     (a) at (0,0)     {tutarlı\\$\Gamma$};
  \node[adim=ornekframe]  (b) at (3.3,0)   {Henkin\\genişlemesi\\{\scriptsize tanık sabitleri}};
  \node[adim=ikincil]     (c) at (6.6,0)   {maksimal\\tutarlı $\Delta$\\{\scriptsize Lindenbaum}};
  \node[adim=teoremframe] (d) at (9.9,0)   {terim\\modeli $\mathcal{M}_\Delta$\\{\scriptsize $\mathcal{M}_\Delta \models \Gamma$}};
  \draw[ok] (a) -- (b); \draw[ok] (b) -- (c); \draw[ok] (c) -- (d);
\end{tikzpicture}
\caption{Henkin'in tamlık ispatının dört adımı. Evren, terimlerden oluşur
(eşitlik varsa eşit terimlerin denklik sınıflarından). $\Delta$'nın
kapalılık özellikleri, terim modelinin $\Delta$'yı, dolayısıyla
$\Gamma$'yı doğrulamasını sağlar.}
\label{fig:henkin}
\end{figure}

\begin{teorem}{Gödel'in tamlık teoremi}{godel-tamlik}
Birinci mertebe mantıkta, her $\Gamma$ cümle kümesi ve her $\varphi$
cümlesi için $\Gamma \models \varphi$ ise $\Gamma \vdash \varphi$.
\end{teorem}

\noindent Bu teoremin şaşırtıcı sonuçları vardır. Tıkızlık birinci mertebe
mantık için de geçerlidir. Sayılabilir bir dildeki her tutarlı kuramın
sayılabilir bir modeli vardır (aşağı Löwenheim--Skolem teoremi): Terim
modeli sayılabilirdir. Gerçel sayılar kuramının bile sayılabilir bir
modeli vardır! Bu ``paradoks'' \ciltiki{5}'te küme kuramıyla birlikte ele
alınacak.

\begin{dikkat}[Dikkat: Tamlık ve eksiklik]
Gödel'in \emph{tamlık} teoremi (1930) ile \emph{eksiklik} teoremleri
(1931) çelişmez, çünkü farklı şeylerden söz ederler. Tamlık, \emph{ispat
sistemi} hakkındadır: Bir kuramın bütün modellerinde doğru olan her cümle
o kuramdan türetilebilir. Eksiklik ise \emph{kuramlar} hakkındadır:
Aritmetiğin, örneğin Peano aritmetiğinin, hem $\varphi$'yi hem
${\sim}\varphi$'yi doğrulayan farklı modelleri olduğu için ikisi de
türetilemez. ``Standart modelde doğru'' ile ``her modelde doğru''
arasındaki fark, \ciltiki{4}'ün konusudur.
\end{dikkat}

% ---------------------------------------------------------------------
\section{Formalizasyondaki tasarım kararları}\label{sec:tamlik-tasarim}

Bu bölümün \Lean{} ispatı birkaç tasarım kararının ürünüdür. Bunlar,
kendi metamantık formalizasyonunuzu yaparken karşılaşacağınız
seçimlerdir.

\begin{table}[htbp]
\centering
\small
\renewcommand{\arraystretch}{1.45}
\rowcolors{2}{tablosatir}{white}
\begin{tabularx}{\linewidth}{>{\raggedright\arraybackslash}p{3.0cm} >{\raggedright\arraybackslash}X >{\raggedright\arraybackslash}X}
\rowcolor{tabloust}
\tablobaslik{Karar} & \tablobaslik{Bizim seçimimiz} & \tablobaslik{Alternatif ve bedeli} \\
Anlambilim değerleri & \lc{Bool}: \leantac{cases} ile satır satır ilerlemek kolay & \lc{Prop}: daha genel, ama RAA ve Kalmár için klasik mantık gerekir \\
Bağlam & \lc{List Formul}: sonlu, \leantac{decide} ile denetlenebilir & \lc{Set Formul}: sonsuz $\Gamma$ mümkün, ama tıkızlık ayrıca ispatlanmalı \\
Tamlık yöntemi & Kalmár: yapıcı, sonlu, ek kütüphane gerektirmez & Lindenbaum: birinci mertebeye genellenir, ama sayma ve seçim gerektirir \\
Atomlar & \lc{Nat} indisli: \lc{enBuyukAtom} ile sınırlanır & Herhangi bir tip: atomların sonlu kümesini ayrıca hesaplamak gerekir \\
Önsav biçimi & \lc{Kalmar Δ b A}: iki yönlü \lc{∧} & \leankw{if}--\leankw{then}: daha kısa ifade, ama her adımda \leankw{if} açılmalı \\
\end{tabularx}
\caption{Tamlık formalizasyonunda yapılan seçimler.}
\label{tab:tamlik-tasarim}
\end{table}

\noindent Bu kararların hiçbiri ``doğru'' ya da ``yanlış'' değildir. Her
biri, ispatın bir bölümünü kolaylaştırırken başka bir bölümünü
zorlaştırır. Formalizasyonda deneyim, büyük ölçüde bu dengeleri önceden
görebilmektir. Michaelis ve Nipkow, Isabelle'de önermeler mantığının dört
farklı ispat sistemini (doğal çıkarım, ardışık hesap, Hilbert sistemi,
çözünürlük) aynı çatı altında formalize etmiş ve bu tür seçimleri
ayrıntılı biçimde tartışmıştır \citep{michaelis2018}.

% ---------------------------------------------------------------------
\begin{arastirmanotu}
\textbf{Birinci mertebe tamlığın formalizasyonları.} Gödel'in tamlık
teoremi birçok ispat asistanında formalize edilmiştir. Lean~3'teki
\emph{Flypitch} projesi, süreklilik hipotezinin bağımsızlığını
ispatlarken birinci mertebe mantığın tamlığını da formalize etti
\citep{han2020flypitch}. Lean~4'teki Foundation kütüphanesi
\citep{foundation} tamlık, tıkızlık ve Gödel'in eksiklik teoremlerini
birlikte içerir \citep{saitou2026}. Bu projelerde en çok emek isteyen
kısım, çoğu zaman tamlık ispatının kendisi değil, \emph{bağlı
değişkenlerin} ve \emph{yerine koymanın} doğru yönetimidir. Bölüm~14'te
bu sorunu de~Bruijn indisleriyle ele alacağız \citep{debruijn1970}.

\textbf{İspat karmaşıklığı.} Cook--Reckhow programı \citep{cookreckhow1979},
$\mathrm{NP} \neq \mathrm{coNP}$'yi ispat sistemleri için alt sınırlar
yoluyla göstermeyi hedefler. SAT çözücülerinin ürettiği çözünürlük
türetimleri bu alanın pratik yüzüdür. \Lean{}'deki \leantac{bv\_decide}
taktiği, bir SAT çözücünün sonucunu \lc{LRAT} sertifikası olarak
denetler ve böylece tamlık teoreminin ``bir türetim vardır'' vaadini
verimli bir biçimde gerçekleştirir (\ciltiki{6}).
\end{arastirmanotu}

\begin{bolumozeti}
\begin{itemize}
  \item Tamlık teoremi, geçerliliğin tersidir: $\Gamma \models A
        \Rightarrow \Gamma \vdash A$. Geçerlilikle birlikte $\vdash$ ve
        $\models$ aynı ilişkidir (\lc{turetilir\_iff\_sonuc}).
  \item Uzun bir ispat, türetilmiş kurallar, Kalmár önsavı, atom eleme
        ve varsayımları içeri taşıma gibi bağımsız önsavlara bölünerek
        yönetilir.
  \item Kalmár önsavı: Doğruluk tablosunun her satırı, işaretli
        atomlardan $A^v$'nin bir türetimidir. Önsav, hem doğru hem yanlış
        durumu birlikte taşıyacak biçimde güçlendirilir.
  \item Atom eleme: Bir atomda farklılaşan iki satır, $p \oveya {\sim}p$
        üzerinden durum ayrımıyla birleşir. $k$ adımda bağlam boşalır.
  \item Tamlık ispatı yapıcıdır ve gerçek bir türetim üretir; ancak
        türetimin boyutu atom sayısında üsteldir.
  \item Sonsuz $\Gamma$ ve birinci mertebe mantık için Lindenbaum önsavı,
        maksimal tutarlı kümeler ve Henkin'in terim modeli kullanılır.
        Tıkızlık ve Löwenheim--Skolem teoremleri bu yöntemin
        sonuçlarıdır.
\end{itemize}
\end{bolumozeti}

\begin{alistirmalar}
\alistirma{1} $\vdash p \oise p \oveya q$. Türetim yazmayın: Doğruluk
tablosunu kontrol edip \lc{tamlik\_bos} kullanın. (\lc{a12\_1})

\alistirma{1} Çift değilleme giriş: $\Delta \vdash A$ ise $\Delta \vdash
{\sim}{\sim}A$. Bu kural RAA gerektiriyor mu? (\lc{a12\_2})

\alistirma{2} $p \ove q \vdash q \ove p$. Bölüm~11'deki \lc{a11\_3}'ü
bu kez \lc{tamlik} ile ispatlayın ve iki ispatı karşılaştırın.
(\lc{a12\_3})

\alistirma{2} Kontrapozisyon, türetilmiş kural olarak: $\Delta \vdash A
\oise B$ ise $\Delta \vdash {\sim}B \oise {\sim}A$. (\lc{a12\_4})

\alistirma{2} $p \oise q \nvdash q \oise p$. (\lc{a12\_5})

\alistirma{1} Kesme kuralı: $\Gamma \vdash A$ ve $A, \Gamma \vdash B$ ise
$\Gamma \vdash B$. Bu kuralın sistemde neden \emph{temel} kural olarak
bulunmadığını tartışın. (\lc{a12\_6})

\alistirma{2} Denklik bağlacı $A \leftrightarrow B$ (\lc{esd A B}) için
Kalmár adımını ispatlayın. \lc{esd}'nin tanımı sayesinde tek satır
yeterli olmalı. (\lc{a12\_7})

\alistirma{3} Tutarlılık ile sağlanabilirlik: $[A] \nvdash \bot \iff A$
sağlanabilirdir. (\lc{a12\_8})

\alistirma{2} (Kâğıt-kalem) $(p \ove q) \oise p$ için
\cref{fig:kalmar-tablo} benzeri dört satırlık tabloyu yazın. Her satırda
hangi türetilmiş kuralların kullanıldığını belirtin.

\alistirma{3} (Proje) Model varlığı teoremini \lc{Classical} kullanmadan
ispatlayın. İpucu: \lc{tumDegerlemeler}'in her değerlemeyi
\lc{enBuyukAtom}'a kadar ``kapsadığını'' gösterin, sonra listeyi
\lc{List.find?} ile tarayın.

\alistirma{3} (Proje) \lc{Turetilir}'in RAA'sız yapıcı versiyonu
\lc{TuretilirY} için Kalmár önsavının hangi adımı başarısız olur? İpucu:
\lc{ucuncu\_hal}.

\alistirma{3} (Araştırma) \lc{tamlik\_bos}'un ürettiği türetimin
boyutunu (kural sayısını) hesaplayan bir fonksiyon yazmak için
\lc{Turetilir}'i \lc{Prop} yerine \lc{Type} değerli yapmak gerekir.
Neden? Kanıt-ilgisizliği (Bölüm~9) burada nasıl bir engel oluşturuyor?
\end{alistirmalar}
   % Tamlık teoremi
% =====================================================================
%  Bölüm 13 — Ardışık Hesap ve Kesme Eleme
% =====================================================================
\chapter{Ardışık Hesap ve Kesme Eleme}\label{bol:ardisik}

\begin{bolumhedefleri}
\begin{itemize}
  \item Gentzen'in ardışık hesabını, LK ve LJ sistemlerini ve aralarındaki
        farkı,
  \item Çok sonuçlu ardışıkların klasik mantığı nasıl ``kendiliğinden''
        verdiğini,
  \item Kuralların tersinirliğini ve bunun bir ispat arama algoritmasına
        nasıl dönüştüğünü,
  \item Kesme kuralını, Gentzen'in \emph{Hauptsatz}'ını ve kesmenin
        \Lean{}'de anlambilimsel yoldan elenmesini,
  \item Alt formül özelliğini ve bunun tutarlılık, karar verilebilirlik ve
        otomatik ispat için sonuçlarını.
\end{itemize}
\end{bolumhedefleri}

Kısım~III'ün ilk üç bölümünde önermeler mantığının \emph{bir} ispat
sistemini, doğal çıkarımı, inceledik ve onun iki değerli anlambilime göre
hem geçerli hem tam olduğunu gösterdik. Bu bölümde aynı mantık için
ikinci bir ispat sistemi tanıtıyoruz: Gerhard Gentzen'in \emph{ardışık
hesabı}. İki sistem aynı formülleri türetir. Ama ardışık hesabın
kurallarının biçimi ona doğal çıkarımda olmayan bir özellik kazandırır:
Türetimler \emph{aşağıdan yukarıya}, mekanik olarak aranabilir. Bu bölümün
sonunda, \Lean{}'de çalışan küçük bir otomatik ispatlayıcımız olacak.

% ---------------------------------------------------------------------
\section{Ardışıklar}\label{sec:ardisiklar}

Bölüm~5'te doğal çıkarım kurallarını $\Gamma \vdash A$ biçimindeki
\emph{ardışıklarla} yazmıştık. Solda bir varsayım listesi, sağda tek bir
sonuç vardı. Gentzen'in ardışık hesabında sağ taraf da bir listedir:

\begin{tanim}{Ardışık}{ardisik}
Bir \textbf{ardışık}, iki formül listesinden oluşan bir çifttir ve
\[
  A_1, \dots, A_m \;\Rightarrow\; B_1, \dots, B_n
\]
diye yazılır. Anlamı şudur: \emph{Soldaki formüllerin hepsi doğruysa,
sağdaki formüllerden en az biri doğrudur.} Yani sol taraf bir ``ve'', sağ
taraf bir ``veya'' gibi okunur:
\[
  \Gamma \Rightarrow \Delta
  \quad\text{geçerlidir}\quad\iff\quad
  \models \textstyle\bigwedge\Gamma \oise \bigvee\Delta .
\]
\end{tanim}

\noindent İki uç durum öğreticidir. $\Rightarrow A$ (sol boş), $A$'nın bir
totoloji olduğunu söyler. $A \Rightarrow$ (sağ boş) ise $A$'nın
\emph{sağlanamaz} olduğunu söyler, çünkü boş bir ``veya'' hiçbir zaman
doğru değildir. Boş ardışık $\Rightarrow$ ise hiçbir zaman geçerli
olamaz; onu türetebilen bir sistem tutarsızdır.

\begin{tarihce}[Tarihçe: Gentzen ve Hauptsatz]
Gerhard Gentzen, 1934--35'te yayımlanan ``Mantıksal çıkarım üzerine
incelemeler'' adlı tezinde \citep{gentzen1935} iki sistem birden tanıttı:
\emph{doğal çıkarım} (NK, NJ) ve \emph{ardışık hesap} (LK, LJ). Doğal
çıkarımın matematikçilerin gerçek akıl yürütmesine yakın olduğunu, ama
onun hakkında teoremler ispatlamanın zor olduğunu gördü. Ardışık hesabı bu
yüzden geliştirdi: Tezin ana sonucu olan \emph{Hauptsatz} (ana teorem),
ardışık hesapta \emph{kesme} kuralının elenebileceğini söyler. Gentzen bir
yıl sonra aynı teknikleri kullanarak Peano aritmetiğinin tutarlılığını
ispatladı \citep{gentzen1936}. Ardışık hesabın bugünkü standart sunumu
için \citet{troelstra2000} ve \citet{negri2001} başvuru kaynaklarıdır.
\end{tarihce}

% ---------------------------------------------------------------------
\section{Kesmesiz bir ardışık hesap}\label{sec:g3}

Ardışık hesabın her bağlaç için iki kuralı vardır: bağlacı \emph{solda}
tanıtan bir kural ve \emph{sağda} tanıtan bir kural. Eleme kuralı yoktur.
Bunun yerine doğal çıkarımdaki eleme kurallarının işini sol kurallar
görür. \Cref{fig:g3-kurallar}, \Lean{}'de tanımlayacağımız sistemin
kurallarını gösteriyor.

\begin{figure}[htbp]
\centering
\setlength{\tabcolsep}{14pt}
\renewcommand{\arraystretch}{3.0}
\begin{tabular}{cc}
\begin{prooftree}
  \hypo{p_n \in \Gamma}
  \hypo{p_n \in \Delta}
  \infer2[Aks]{\Gamma \Rightarrow \Delta}
\end{prooftree}
&
\begin{prooftree}
  \hypo{\bot \in \Gamma}
  \infer1[$\bot$L]{\Gamma \Rightarrow \Delta}
\end{prooftree}
\\
\begin{prooftree}
  \hypo{A, B, \Gamma \Rightarrow \Delta}
  \infer1[$\wedge$L]{A \ove B, \Gamma \Rightarrow \Delta}
\end{prooftree}
&
\begin{prooftree}
  \hypo{\Gamma \Rightarrow A, \Delta}
  \hypo{\Gamma \Rightarrow B, \Delta}
  \infer2[$\wedge$R]{\Gamma \Rightarrow A \ove B, \Delta}
\end{prooftree}
\\
\begin{prooftree}
  \hypo{A, \Gamma \Rightarrow \Delta}
  \hypo{B, \Gamma \Rightarrow \Delta}
  \infer2[$\vee$L]{A \oveya B, \Gamma \Rightarrow \Delta}
\end{prooftree}
&
\begin{prooftree}
  \hypo{\Gamma \Rightarrow A, B, \Delta}
  \infer1[$\vee$R]{\Gamma \Rightarrow A \oveya B, \Delta}
\end{prooftree}
\\
\begin{prooftree}
  \hypo{\Gamma \Rightarrow A, \Delta}
  \hypo{B, \Gamma \Rightarrow \Delta}
  \infer2[$\to$L]{A \oise B, \Gamma \Rightarrow \Delta}
\end{prooftree}
&
\begin{prooftree}
  \hypo{A, \Gamma \Rightarrow B, \Delta}
  \infer1[$\to$R]{\Gamma \Rightarrow A \oise B, \Delta}
\end{prooftree}
\\
\multicolumn{2}{c}{%
\colorbox{fikirbg}{\begin{prooftree}
  \hypo{\Gamma \Rightarrow \Delta}
  \hypo{\Gamma \subseteq \Gamma'}
  \hypo{\Delta \subseteq \Delta'}
  \infer3[Yapı]{\Gamma' \Rightarrow \Delta'}
\end{prooftree}\quad\raisebox{0.4em}{\textcolor{fikirframe}{\baslikfont\scriptsize\bfseries yapısal kural}}}}
\\
\end{tabular}
\caption{Kesmesiz klasik ardışık hesap. Her bağlaç için bir sol (L) ve
bir sağ (R) kuralı vardır. Son kural, Gentzen'in zayıflatma, değiştirme
ve daraltma kurallarını tek bir kuralda toplar. Kesme kuralı \emph{yoktur}.}
\label{fig:g3-kurallar}
\end{figure}

\noindent Kurallara yakından bakalım. $\to$L kuralı, doğal çıkarımdaki
modus ponensin ardışık hesaptaki karşılığıdır. Solda $A \oise B$ varsa ve
$A$'yı ``sağlayabiliyorsak'' (sol öncül), $B$'yi varsayabiliriz (sağ öncül).
$\vee$R kuralı ise çok sonuçlu olmanın gücünü gösterir: $A \oveya B$'yi
kanıtlamak için hangisinin doğru olduğuna \emph{karar vermek} zorunda
değiliz; ikisini birden sağa yazarız. \Lean{}'deki tanım şöyle:

\leandosya[Kesmesiz ardışık hesap]{Bolum13}{ardisik}

\noindent \lc{Ardisik Γ Δ} tanımı, Bölüm~11'deki \lc{Turetilir}
tanımıyla aynı kalıptadır: Her kural bir kurucudur. Tek fark, ana
formülün listenin \emph{başında} durmasıdır. Ana formül başka bir yerdeyse,
\lc{degis} yapısal kuralıyla öne alınır.

\begin{dikkat}[Dikkat: \lc{::} ve bağlaçların önceliği]
\Lean{}'de \lc{::} işlecinin önceliği (67), bağlaçlarımızınkinden
(25--35) yüksektir. Bu yüzden \lc{A {\sembolkodfont ⋀} B :: Γ}, \lc{A {\sembolkodfont ⋀} (B :: Γ)} diye
okunur ve bir tip hatası verir. Ana formülü her zaman paranteze alıyoruz:
\lc{(A {\sembolkodfont ⋀} B) :: Γ}.
\end{dikkat}

\subsection{LK ve LJ}

Gentzen'in klasik sistemi LK'dır. Sezgici sistem LJ ise LK'dan tek bir
kısıtlamayla elde edilir: \emph{Sağ tarafta en fazla bir formül
bulunabilir.} Bu küçük sözdizimsel kısıtlama, Bölüm~4'te uzun uzun
tartıştığımız yapıcı/klasik ayrımını tam olarak yakalar. Bunu görmek için
Peirce yasasının türetimine bakalım (\cref{fig:peirce-ardisik}).

\leandosya[İki örnek türetim]{Bolum13}{ornekler}

\begin{figure}[htbp]
\centering
\begin{prooftree}
  \infer0[Aks]{p \Rightarrow \colorbox{teorembg}{$q,\; p$}}
  \infer1[$\to$R]{\Rightarrow \colorbox{teorembg}{$p \oise q,\; p$}}
  \infer0[Aks]{p \Rightarrow p}
  \infer2[$\to$L]{(p \oise q) \oise p \Rightarrow p}
  \infer1[$\to$R]{\Rightarrow ((p \oise q) \oise p) \oise p}
\end{prooftree}
\caption{Peirce yasasının ardışık hesaptaki türetimi, \Lean{}'deki
\lc{peirce\_ardisik} ispatının adım adım karşılığı. Renkli ardışıklarda
sağ tarafta \emph{iki} formül vardır; LJ'de bu adımlara izin verilmez. Bu
yüzden Peirce yasası LJ'de türetilemez.}
\label{fig:peirce-ardisik}
\end{figure}

\begin{onemlifikir}[Önemli Fikir: Klasik mantık = çok sonuçlu ardışıklar]
Doğal çıkarımda klasik mantığa geçmek için yeni bir \emph{kural}
eklememiz gerekmişti (RAA). Ardışık hesapta kurallar aynı kalır; değişen
yalnızca sağ tarafın kaç formül taşıyabildiğidir. Klasik mantığın ``ya
öyle ya böyle, ama hangisi olduğunu söylemeden'' akıl yürütmesi, sağ
taraftaki virgülde saklıdır. Girard'ın deyişiyle, klasik ve sezgici mantık
arasındaki fark bağlaçlarda değil, \emph{yapıdadır} \citep{girard1989}.
\end{onemlifikir}

% ---------------------------------------------------------------------
\section{Anlambilim ve geçerlilik}\label{sec:ardisik-gecerlilik}

Bir ardışığın anlamını \Lean{}'de iki küçük yüklemle ifade ediyoruz. Sol
tarafın hepsi doğru (\lc{HepsiDogru}), sağ taraftan biri doğru
(\lc{BiriDogru}):

\leandosya[Ardışıkların anlamı]{Bolum13}{gecerli_tanim}

\noindent Listelerin başına bir eleman eklendiğinde bu iki yüklemin nasıl
davrandığını söyleyen iki önsav, bölümün geri kalanında sürekli
kullanılacak. Üçüncü önsav ise geçerliliğin listeler büyüdükçe
korunduğunu söyler (bu, yapısal kuralın geçerliliğidir):

\leandosya[Liste önsavları]{Bolum13}{liste_onsavlari}

\noindent Geçerlilik teoremi Bölüm~11'deki gibi kural tümevarımıyla
ispatlanır. Her kural, öncüllerindeki geçerliliği sonucuna taşır:

\leandosya[Ardışık hesabın geçerliliği]{Bolum13}{ardisik_gecerli}

\noindent Son satırdaki \lc{bos\_ardisik\_yok}, boş ardışığın
türetilemediğini, yani sistemin tutarlı olduğunu söyler. Geçerlilik
ispatı burada da kısa bir sonuç verir: Boş ardışık hiçbir değerlemede
doğru değildir.

% ---------------------------------------------------------------------
\section{Tersinirlik ve ispat arama}\label{sec:tersinirlik}

Doğal çıkarımda bir formülü türetmeye çalışırken çoğu zaman ne yapacağımızı
\emph{tahmin etmemiz} gerekir. Modus ponens hangi $A$ ile uygulanmalı?
Hangi ara önsav işe yarar? Ardışık hesabın kurallarında böyle bir tahmin
yoktur. Her kuralın öncüllerinde yalnızca sonucunda geçen formüllerin
\emph{parçaları} bulunur. Üstelik bizim sistemimizin her kuralı
\emph{tersinirdir}:

\begin{tanim}{Tersinir kural}{tersinir}
Bir kural, sonucu geçerli olduğunda öncüllerinin hepsi de geçerliyse
\textbf{tersinirdir}.
\end{tanim}

\noindent Tersinirlik, ispat aramada hiçbir zaman yanlış bir seçim
yapamayacağımız anlamına gelir. Geçerli bir ardışığa \emph{hangi} kuralı
uygularsak uygulayalım, öncüller de geçerli kalır. Bu yüzden arama hiçbir
zaman geri dönmek (\emph{backtracking}) zorunda kalmaz.

\leandosya[Tersinirlik önsavları]{Bolum13}{tersinirlik}

\begin{figure}[htbp]
\centering
\begin{tikzpicture}[
    ard/.style={draw=#1, fill=#1!8, thick, rounded corners=3pt,
                font=\small, inner sep=4pt},
    ok/.style={-{Stealth[length=2.2mm]}, thick, black!45},
    kural/.style={font=\baslikfont\scriptsize\bfseries, text=leantac},
  ]
  \node[ard=teoremframe] (k) at (0,0) {$\Rightarrow (p \oveya q) \oise (q \oveya p)$};
  \node[ard=anarenk] (a1) at (0,1.4) {$p \oveya q \Rightarrow q \oveya p$};
  \node[ard=anarenk] (a2) at (0,2.8) {$p \oveya q \Rightarrow q,\; p$};
  \node[ard=ikincil] (l) at (-2.6,4.3) {$p \Rightarrow q,\; p$};
  \node[ard=ikincil] (r) at (2.6,4.3) {$q \Rightarrow q,\; p$};
  \node[font=\baslikfont\scriptsize\bfseries, text=ikincil] at (-2.6,5.0) {Aks ($p$)};
  \node[font=\baslikfont\scriptsize\bfseries, text=ikincil] at (2.6,5.0) {Aks ($q$)};
  \draw[ok] (k) -- node[right, kural]{$\to$R} (a1);
  \draw[ok] (a1) -- node[right, kural]{$\vee$R} (a2);
  \draw[ok] (a2) -- node[left, kural, pos=0.45]{$\vee$L} (l);
  \draw[ok] (a2) -- (r);
  \node[font=\baslikfont\scriptsize, text=black!55, align=left, anchor=west] at (4.6,1.4)
       {her adımda:\\en soldaki bağlacı\\parçala};
  \node[font=\baslikfont\scriptsize, text=black!55, align=left, anchor=west] at (4.6,4.3)
       {yapraklar:\\yalnızca atomlar};
  \draw[-{Stealth[length=3mm]}, very thick, black!25] (-5.2,0) -- (-5.2,4.6)
       node[midway, left, font=\baslikfont\scriptsize, text=black!55, align=right]{arama\\yönü};
\end{tikzpicture}
\caption{$(p \oveya q) \oise (q \oveya p)$ için ispat arama. Kök, ispatlamak
istediğimiz ardışıktır. Her adımda bir bağlaç parçalanır. Yapraklarda
yalnızca atomlar kaldığında, solda ve sağda ortak bir atom olup olmadığına
bakılır.}
\label{fig:ispat-arama}
\end{figure}

\subsection{Kesmesiz tamlık}

Tersinirlik, bir tamlık ispatına doğrudan dönüşür. Geçerli bir ardışık
verilsin. Bağlaçları birer birer parçalayalım; tersinirlik sayesinde her
yeni ardışık da geçerlidir. Bağlaç kalmadığında elimizde yalnızca
atomlardan oluşan geçerli ardışıklar kalır. Böyle bir ardışık ya solda ve
sağda ortak bir atom içerir (o zaman Aks kuralı uygulanır), ya da
geçersizdir. Çünkü soldaki atomları doğru, diğer her şeyi yanlış yapan
değerleme bir karşı model olurdu.

Bunu \Lean{}'de yazmak için önce aramanın sona ereceğini garanti eden bir
\emph{ölçü} tanımlıyoruz. Her kural uygulaması bu ölçüyü en az bir
azaltır:

\leandosya[Ardışığın ölçüsü]{Bolum13}{olcu}

\noindent Ölçüdeki $2 \cdot \mathrm{boyut}$ çarpanı önemlidir. Örneğin
$\wedge$L kuralı $A \ove B$'yi $A$ ve $B$ ile değiştirir. Bağlaç sayısı
bir azalır ama formül sayısı bir artar. Bağlaçlara iki kat ağırlık
verilmezse toplam değişmezdi.

Atomik formülleri işlerken onları listenin sonundaki bir ``atom
bölgesine'' taşırız. \lc{tasi1} ve \lc{tasi2}, bu taşımanın üyeliği
değiştirmediğini söyler. Böylece \lc{degis} kuralı taşımayı türetime
dönüştürebilir:

\leandosya[Atomları kenara almak]{Bolum13}{tasi}

\noindent Aramanın yaprakları, yalnızca atomlardan oluşan ardışıklardır.
Burada karşı model yöntemi devreye girer:

\leandosya[Atomik ardışıklar]{Bolum13}{atomik}

\noindent Artık arama algoritmasını bir tamlık ispatı olarak yazabiliriz.
\lc{tam\_aux}, ölçü üzerinde (\lc{n} ile sınırlanmış) özyinelemeyle
tanımlanmış bir teoremdir. Önce soldaki formülleri, sol boşalınca da
sağdakileri işler. Her durumda tersinirlik önsavı yeni ardışığın geçerli
olduğunu, özyinelemeli çağrı ise türetilebildiğini söyler:

\leandosya[İspat arama, bir tamlık ispatı olarak]{Bolum13}{tam_aux}

\noindent \lc{tam\_aux}'un örüntüleri, \cref{fig:ispat-arama}'daki
aramanın birebir tarifidir. \lc{P} ve \lc{Q} kenara alınmış atomları
tutar. Ardından son adım gelir: Başlangıçta kenarda hiç atom yoktur.

\leandosya[Kesmesiz tamlık]{Bolum13}{ardisik_tam}

\begin{teorem}{Kesmesiz tamlık}{kesmesiz-tamlik}
Geçerli her $\Gamma \Rightarrow \Delta$ ardışığı, kesme kuralı
kullanılmadan türetilebilir.
\end{teorem}

% ---------------------------------------------------------------------
\section{Kesme kuralı}\label{sec:kesme}

Matematikte bir teoremi ispatlarken önce bir \emph{önsav} ispatlar, sonra
onu kullanırız. Ardışık hesapta bu adımın adı \emph{kesme}dir:
\begin{center}
\begin{prooftree}
  \hypo{\Gamma \Rightarrow A, \Delta}
  \hypo{A, \Gamma \Rightarrow \Delta}
  \infer2[\textcolor{teoremframe}{Kesme}]{\Gamma \Rightarrow \Delta}
\end{prooftree}
\end{center}
Sol öncül önsavı ispatlar. Sağ öncül önsavı varsayarak asıl sonuca
ulaşır. $A$ formülü sonuçta görünmez; ``kesilip atılmıştır''. Kesme,
tüm diğer kurallardan farklıdır: Sonuçta \emph{bulunmayan} bir formül
(önsav) öncüllerde ortaya çıkar. İspat arama açısından bu bir felakettir.
Kesmeyi aşağıdan yukarıya uygulayacak olsak, sonsuz sayıdaki olası
$A$'lardan birini tahmin etmemiz gerekirdi.

\begin{teorem}{Hauptsatz (Gentzen, 1935)}{hauptsatz}
Kesme kuralı elenebilir: Kesme ile türetilebilen her ardışık, kesme
kullanılmadan da türetilebilir.
\end{teorem}

\noindent Elimizdeki araçlarla bu teoremin önermeler mantığı için
ispatı birkaç satırdır. Kesmenin iki öncülü geçerliyse sonucu da
geçerlidir. Kesmesiz tamlık ise geçerli her ardışığın kesmesiz bir
türetimi olduğunu söyler:

\leandosya[Kesme elenebilir]{Bolum13}{kesme}

\noindent Bu ispat, kesmeyi sisteme bir kural olarak eklememize gerek
olmadığını gösterir. \lc{kesme} bir \emph{teorem} olarak elimizdedir ve
kesmesiz sistemin türetebildiklerini değiştirmez. Böyle kurallara
\emph{kabul edilebilir} (\emph{admissible}) kurallar denir.

\begin{figure}[htbp]
\centering
\begin{tikzpicture}[
    kutu/.style={draw=#1, fill=#1!8, very thick, rounded corners=5pt,
                 text width=2.3cm, align=center, inner sep=6pt, font=\small},
    ok/.style={-{Stealth[length=3mm]}, very thick, #1},
  ]
  \node[kutu=anarenk]     (k) at (0,0)    {Kesmeli türetim\\$\Gamma \Rightarrow \Delta$};
  \node[kutu=ikincil]     (g) at (5.3,0)  {$\Gamma \Rightarrow \Delta$\\geçerli};
  \node[kutu=teoremframe] (s) at (10.6,0) {Kesmesiz türetim\\$\Gamma \Rightarrow \Delta$};
  \draw[ok=ikincil] (k) -- node[above, font=\baslikfont\scriptsize\bfseries]{geçerlilik} (g);
  \draw[ok=teoremframe] (g) -- node[above, font=\baslikfont\scriptsize\bfseries]{kesmesiz tamlık} (s);
  \draw[ok=black!30, dashed] (k.south) to[bend right=18]
       node[below, font=\baslikfont\scriptsize, text=black!60]{Gentzen: türetimi doğrudan dönüştür (sözdizimsel)} (s.south);
\end{tikzpicture}
\caption{Kesme elemenin iki yolu. Üstteki anlambilimsel yol bu bölümde
\Lean{}'de yapıldı. Alttaki sözdizimsel yol, Gentzen'in yöntemidir.}
\label{fig:kesme-yollari}
\end{figure}

\subsection{Sözdizimsel kesme eleme}

Gentzen'in ispatı anlambilim kullanmaz. Kesmeli bir türetimi adım adım
kesmesiz bir türetime \emph{dönüştürür}. Ana fikir, bir kesmeyi
öncüllerine doğru ``yukarı itmektir''. Kesme formülü her iki öncülde de
yeni tanıtılmışsa, kesme daha küçük formüller üzerinde kesmelere bölünür.
Örneğin $A \ove B$ için:
\begin{center}
\scriptsize
\begin{prooftree}
  \hypo{\Gamma \Rightarrow A, \Delta}
  \hypo{\Gamma \Rightarrow B, \Delta}
  \infer2[$\wedge$R]{\Gamma \Rightarrow A \ove B, \Delta}
  \hypo{A, B, \Gamma \Rightarrow \Delta}
  \infer1[$\wedge$L]{A \ove B, \Gamma \Rightarrow \Delta}
  \infer2[Kesme]{\Gamma \Rightarrow \Delta}
\end{prooftree}
\quad\raisebox{1.2em}{\normalsize$\rightsquigarrow$}\quad
\begin{prooftree}
  \hypo{\Gamma \Rightarrow A, \Delta}
  \hypo{\Gamma \Rightarrow B, \Delta}
  \hypo{A, B, \Gamma \Rightarrow \Delta}
  \infer2[Kesme$_B$]{A, \Gamma \Rightarrow \Delta}
  \infer2[Kesme$_A$]{\Gamma \Rightarrow \Delta}
\end{prooftree}
\end{center}
(Sağdaki türetimde öncüllere gereken zayıflatmalar gösterilmemiştir.)
Kesme formülü bir öncülde yeni tanıtılmamışsa, kesme o öncülün kuralının
üstüne taşınır. İspat, kesme formülünün boyutu ve türetimin yüksekliği
üzerinde \emph{iç içe} bir tümevarımdır. Bu tümevarımın \Lean{}'deki
formalizasyonu, \cref{bol:proje}'deki projelerden biridir. Pfenning'in
``yapısal kesme eleme'' yöntemi, bu ispatın ispat asistanlarına uygun bir
biçimidir \citep{pfenning2000}.

\begin{dikkat}[Dikkat: Kesme elemenin bedeli]
Kesme eleme ``bedava'' değildir. Kesmesiz türetimler, kesmeli olanlardan
\emph{akıl almaz} derecede uzun olabilir. Bazı formül aileleri için
kesmeli türetimin boyutu $n$ iken, kesmesiz türetimin boyutu $2^{2^{\cdot^{\cdot^{2}}}}$
(yüksekliği $n$ olan bir üs kulesi) mertebesindedir \citep{statman1979}.
Matematikte önsavların neden bu kadar önemli olduğunun biçimsel bir
açıklaması budur. Önsavlar türetimleri küçültür. Bu bölümün
\lc{ardisik\_tam} ispatı da, kesmesiz türetimleri doğruluk tablosu
büyüklüğünde üretir.
\end{dikkat}

% ---------------------------------------------------------------------
\section{Alt formül özelliği ve sonuçları}\label{sec:alt-formul}

\Cref{fig:g3-kurallar}'deki kurallara bir kez daha bakın. Her kuralın
öncüllerinde geçen her formül, sonucundaki bir formülün \emph{alt
formülüdür}. Kesme bu özelliği taşımayan tek kuraldı.

\begin{teorem}{Alt formül özelliği}{alt-formul}
Kesmesiz bir türetimde geçen her formül, kök ardışıktaki bir formülün alt
formülüdür.
\end{teorem}

\noindent Bu sözdizimsel özelliğin şaşırtıcı derecede güçlü sonuçları
vardır:

\begin{itemize}
  \item \textbf{Tutarlılık.} Boş ardışığın kesmesiz bir türetimi
        olsaydı, alt formül özelliğine göre bu türetimde hiçbir formül
        geçemezdi. Ama Aks ve $\bot$L kurallarının ikisi de bir formül
        gerektirir. Gentzen'in Peano aritmetiği için tutarlılık ispatının
        iskeleti budur. Bölüm~5'te doğal çıkarım için aynı sonuca
        normalleştirmeyle ulaşmıştık.
  \item \textbf{Karar verilebilirlik.} Bir ardışığın olası kesmesiz
        türetimlerinde yalnızca sonlu sayıda formül kullanılabilir. Bu
        yüzden arama sonludur. Önermeler mantığı için bunu
        \lc{tam\_aux}'ta gördük. LJ için de geçerlidir, ama daraltma
        kuralı yüzünden çok daha incedir (aşağıdaki araştırma notuna
        bakın).
  \item \textbf{Ara değer (interpolasyon).} $A \Rightarrow B$ geçerliyse,
        yalnızca $A$ ile $B$'nin \emph{ortak} atomlarını içeren bir $C$
        formülü vardır ki $A \Rightarrow C$ ve $C \Rightarrow B$ geçerlidir
        \citep{craig1957}. Kesmesiz türetim üzerinde tümevarımla $C$
        doğrudan inşa edilebilir.
\end{itemize}

\subsection{Çalıştırılabilir bir ispatlayıcı}

\lc{tam\_aux} bir \emph{teoremdir}; bir türetim olduğunu söyler. Aynı
özyinelemeyi \lc{Bool} döndüren bir \emph{fonksiyon} olarak yazarsak,
çalıştırabileceğimiz bir karar yöntemi elde ederiz:

\leandosya[Çalıştırılabilir ispat arama]{Bolum13}{ara}

\noindent \lc{ara}'nın doğruluğunu (yani \lc{ara} \lc{true} döndürdüğünde
gerçekten bir türetim olduğunu) ispatlamak, \lc{tam\_aux}'un yapısını
izleyen güzel bir alıştırmadır. \ciltiki{6}'da bu fikri \emph{yansıma}
(\emph{reflection}) adıyla genelleştireceğiz. Orada bir karar yönteminin
doğruluğunu bir kez ispatlayıp, sonra onu \leantac{decide} ile çalışan bir
taktik olarak kullanacağız.

% ---------------------------------------------------------------------
\section{Doğal çıkarım ile ardışık hesap}\label{sec:dogal-ardisik}

İki sistemimiz de iki değerli anlambilime göre geçerli ve tamdır. Öyleyse
aynı şeyleri türetirler. \Lean{}'de bu, iki tamlık teoreminin
birleştirilmesiyle tek satırda ispatlanır:

\leandosya[Doğal çıkarım ile denklik]{Bolum13}{denklik}

\noindent Ama iki sistemin ilişkisi ``aynı şeyleri türetmek''ten daha
derindir. \Cref{fig:nd-lk}'de görüldüğü gibi, kurallar arasında
yapısal bir eşleşme vardır.

\begin{figure}[htbp]
\centering
\begin{tikzpicture}[
    kutu/.style={draw=#1, fill=#1!8, thick, rounded corners=4pt,
                 text width=3.6cm, align=center, inner sep=5pt, font=\small,
                 minimum height=1.0cm},
    ok/.style={{Stealth[length=2.4mm]}-{Stealth[length=2.4mm]}, thick, black!40},
  ]
  \node[font=\baslikfont\bfseries, text=anarenk] at (0,1.2) {Doğal çıkarım};
  \node[font=\baslikfont\bfseries, text=ikincil] at (7,1.2) {Ardışık hesap};
  \node[kutu=anarenk] (g)  at (0,0)    {giriş kuralları};
  \node[kutu=anarenk] (e)  at (0,-1.4) {eleme kuralları};
  \node[kutu=anarenk] (d)  at (0,-2.8) {dolambaç\\{\scriptsize (giriş + eleme)}};
  \node[kutu=anarenk] (n)  at (0,-4.2) {normalleştirme};
  \node[kutu=ikincil] (r)  at (7,0)    {sağ kurallar (R)};
  \node[kutu=ikincil] (l)  at (7,-1.4) {sol kurallar (L)};
  \node[kutu=ikincil] (k)  at (7,-2.8) {kesme};
  \node[kutu=ikincil] (ke) at (7,-4.2) {kesme eleme};
  \draw[ok] (g) -- (r); \draw[ok] (e) -- (l);
  \draw[ok] (d) -- (k); \draw[ok] (n) -- (ke);
\end{tikzpicture}
\caption{Doğal çıkarım ile ardışık hesap arasındaki eşleşme. Doğal
çıkarımda bir girişin hemen ardından gelen eleme (dolambaç, Bölüm~5),
ardışık hesapta bir kesmeye karşılık gelir. Prawitz'in normalleştirme
teoremi ile Gentzen'in Hauptsatz'ı aynı madalyonun iki yüzüdür.}
\label{fig:nd-lk}
\end{figure}

\noindent Curry--Howard açısından bakınca bu eşleşme daha da
anlamlıdır. Doğal çıkarım türetimleri $\lambda$-terimleridir;
normalleştirme ise bir programın \emph{çalıştırılmasıdır}. Ardışık hesaptaki
kesme eleme de bir hesaplamadır. Bu hesaplamanın programlama dillerindeki
karşılıkları (örneğin ``\emph{let}'' ifadelerinin açılması) Bölüm~14'te
$\lambda$-hesabı ile birlikte ele alınacak.

% ---------------------------------------------------------------------
\begin{arastirmanotu}
\textbf{Sezgici mantıkta ispat arama.} LJ'de $\to$L kuralı ana formülü
öncüllerde de tutmak zorundadır. Aksi hâlde sistem tam olmaz. Bu yüzden
saf geriye arama sonsuza kadar sürebilir. Roy Dyckhoff'un \emph{LJT}
sistemi \citep{dyckhoff1992}, $\to$L kuralını öncülün biçimine göre dört
kurala bölerek bu sorunu çözer ve her türetimde ölçünün azaldığı,
daraltmasız bir sezgici ardışık hesap verir. Mathlib'deki \leantac{itauto}
gibi sezgici karar yöntemleri bu çizgidedir.

\textbf{Formalize edilmiş kesme eleme.} Kesme eleme, ispat asistanlarında
metamantık formalizasyonunun ölçütlerinden biridir. Chaudhuri, Lima ve
Reis doğrusal mantıklar için kesme elemeyi formalize ettiler
\citep{chaudhuri2019}. Shillito, van~der~Giessen, Goré ve Iemhoff ise
sezgici bir kanıtlanabilirlik mantığı için yeni bir ardışık hesap
önerip kesme elemesini ve ispat aramasının sonlanmasını Coq'ta
formalize ettiler \citep{shillito2023}. Lean~4 tarafında Foundation
kütüphanesi \citep{foundation}, kip mantıkları ve kanıtlanabilirlik
mantığı için çeşitli ispat sistemlerini formalize eder.

\textbf{İspat aramadan otomasyona.} Mathlib'deki \leantac{tauto} ve
genel amaçlı \leantac{aesop} taktiği \citep{limperg2023aesop}, bu
bölümdeki fikirlerin \Lean{}'in kendi mantığına uygulanmış hâlleridir:
Tersinir kurallar ``güvenli'' kabul edilip hemen uygulanır, tersinir
olmayanlar için ise en iyi öncelikli bir arama yapılır.
\end{arastirmanotu}

\begin{bolumozeti}
\begin{itemize}
  \item Bir ardışık $\Gamma \Rightarrow \Delta$, ``$\Gamma$'nın hepsi
        doğruysa $\Delta$'dan biri doğrudur'' demektir. Ardışık hesapta
        her bağlacın bir sol ve bir sağ kuralı vardır.
  \item LK (klasik) ile LJ (sezgici) arasındaki fark, sağ tarafta kaç
        formül bulunabileceğidir. Peirce yasası ve üçüncü hâl, sağda iki
        formül gerektirir.
  \item Sistemimizin her kuralı tersinirdir. Bu yüzden geriye doğru ispat
        arama hiç yanlış seçim yapmaz ve bir ölçü sayesinde sona erer.
  \item Arama, bir tamlık ispatına dönüşür: Geçerli her ardışığın
        kesmesiz bir türetimi vardır (\lc{ardisik\_tam}).
  \item Kesme kuralı elenebilir (\emph{Hauptsatz}). Anlambilimsel ispat,
        geçerlilik ile kesmesiz tamlığın birleşimidir. Gentzen'in
        sözdizimsel ispatı ise türetimleri doğrudan dönüştürür; bunun
        bedeli türetim boyutunda büyük bir artıştır.
  \item Kesmesiz türetimlerin alt formül özelliği; tutarlılık, karar
        verilebilirlik ve interpolasyon sonuçlarının kaynağıdır.
  \item Doğal çıkarım ve ardışık hesap aynı formülleri türetir. Dolambaç
        ile kesme ve normalleştirme ile kesme eleme birbirine karşılık
        gelir.
\end{itemize}
\end{bolumozeti}

\begin{alistirmalar}
\alistirma{1} $p \ove q \Rightarrow q \ove p$. Türetim ağacını da çizin
ve \lc{a11\_3} ile karşılaştırın: Hangisi daha kısa? (\lc{a13\_1})

\alistirma{1} ${\sim}p, p \Rightarrow$\,. Sağ tarafın boş olmasının
anlamı nedir? (\lc{a13\_2})

\alistirma{1} Zayıflatma, yapısal kuralın özel bir hâlidir.
(\lc{a13\_3})

\alistirma{2} Dummett formülü: $\Rightarrow (p \oise q) \oveya (q \oise
p)$. Ana formülü öne almak için \lc{degis}'i kullanmanız gerekecek. Bu
formül sezgici mantıkta türetilemez. Türetiminizin hangi adımı LJ'de
yasaktır? (\lc{a13\_4})

\alistirma{2} $p \oveya q \not\Rightarrow p$: Bu ardışık türetilemez.
(\lc{a13\_5})

\alistirma{2} Sözdizimsel tersinirlik: $A \ove B, \Gamma \Rightarrow
\Delta$ türetilebiliyorsa $A, B, \Gamma \Rightarrow \Delta$ da
türetilebilir. Tamlığı kullanın. (\lc{a13\_6})

\alistirma{2} Doğal çıkarımda ispatlanan bir önsavı, ardışık hesapta
kesme ile kullanın. (\lc{a13\_7})

\alistirma{3} $\Rightarrow A$ türetilebilir $\iff$ $A$ bir totolojidir.
(\lc{a13\_8})

\alistirma{2} (Kâğıt-kalem) $(p \oise q) \oise ({\sim}q \oise {\sim}p)$
için \cref{fig:ispat-arama} benzeri bir arama ağacı çizin. Kaç yaprak
var?

\alistirma{3} (Proje) \lc{ara n Γ Δ P Q = true} ise \lc{Ardisik (Γ ++
P.map atom) (Δ ++ Q.map atom)} olduğunu ispatlayın. İpucu: \lc{tam\_aux}
ile aynı örüntüleri kullanın; atomik durumda \lc{List.any} için bir
önsav gerekecek.

\alistirma{3} (Proje) LJ'yi \lc{ArdisikJ : List Formul → Option Formul
→ Prop} olarak tanımlayın (sağda en fazla bir formül). Geçerliliğini
ispatlayın ve \lc{ArdisikJ [] (some peirce)}'in türetilemediğini
tartışın. Bu türetilemezliği ispatlamak için hangi anlambilim gerekir?
(İpucu: \ciltiki{2}.)

\alistirma{3} (Araştırma) Gentzen'in sözdizimsel kesme eleme ispatını
\lc{Ardisik} sistemine kesme kuralı eklenmiş bir \lc{ArdisikK} için
formalize etmeye çalışın. Tümevarımın ölçüsü ne olmalı? \lc{degis}
kuralı ispatı nasıl etkiliyor?
\end{alistirmalar}
   % Ardışık hesap ve kesme eleme

\part{Tip Kuramı}
% =====================================================================
%  Bölüm 14 — Basit Tipli λ-Hesabı ve Curry–Howard
% =====================================================================
\chapter{Basit Tipli $λ$-Hesabı ve Curry--Howard}\label{bol:lambda}

\begin{bolumhedefleri}
\begin{itemize}
  \item Tipsiz $\lambda$-terimlerini, bağlı değişkenleri ve de~Bruijn
        indislerini,
  \item Yerine koymayı ve ``yakalama'' sorununu,
  \item $\beta$-indirgemeyi, normal biçimleri, indirgeme stratejilerini ve
        normal biçimi olmayan $\Omega$ terimini,
  \item Church sayılarıyla $\lambda$-hesabında hesaplamayı,
  \item Basit tipleri, S--K--I kombinatörlerini ve bir tip denetleyiciyi,
  \item Curry--Howard eşlemesinin kesin ifadesini ve \Lean{}'deki
        ispatını.
\end{itemize}
\end{bolumhedefleri}

Kitabın ilk bölümünden beri bir iddiayı tekrarlıyoruz: \Lean{}'de bir
ispat bir programdır, bir önerme ise bir tiptir. Bölüm~2'de bunu
gözlemledik, Bölüm~5'te doğal çıkarım kurallarıyla somutlaştırdık.
Kısım~IV'ün bu ilk bölümünde iddiayı bir \emph{teorem} hâline
getiriyoruz. Bunun için önce \Lean{}'in arkasındaki hesaplama modelini,
Alonzo Church'ün $\lambda$-hesabını, sıfırdan kuracağız. Sonra ona tipler
ekleyecek ve tiplenebilen her terimin Bölüm~11'deki doğal çıkarım
sisteminde bir türetim olduğunu ispatlayacağız.

% ---------------------------------------------------------------------
\section{$λ$-terimleri}\label{sec:lambda-terimler}

$\lambda$-hesabında yalnızca üç tür ifade vardır:
\[
  t \;::=\; x \;\mid\; \lambda x.\, t \;\mid\; t\; u
\]
Bir \emph{değişken}, bir \emph{soyutlama} ($x$'i alıp $t$'yi döndüren
fonksiyon) ve bir \emph{uygulama} ($t$ fonksiyonunu $u$'ya uygulamak).
Sayılar, mantıksal değerler, listeler, özyineleme yoktur. Yine de bu
dilde hesaplanabilen her şey ifade edilebilir. \Lean{}'deki \leankw{fun}
\lc{x => t}, $\lambda x.\, t$'nin ta kendisidir.

\begin{tarihce}[Tarihçe: Church, Kleene ve Rosser]
Alonzo Church $\lambda$-hesabını 1930'ların başında matematiğin temelleri
için bir mantık sistemi olarak tasarladı. Öğrencileri Stephen Kleene ve
J.~Barkley Rosser 1935'te bu ilk sistemin tutarsız olduğunu gösterdi.
Church bunun üzerine sistemin mantık kısmını bırakıp yalnızca
\emph{hesaplama} kısmını, bugünkü tipsiz $\lambda$-hesabını korudu. 1936'da
bu hesabın Turing makineleriyle aynı gücü taşıdığı anlaşıldı
\citep{turing1936}. Tutarlı bir mantık elde etmek için Church 1940'ta
\emph{basit tip kuramını} önerdi \citep{church1940}. \Lean{}'in tip kuramı,
bu sistemin çok geliştirilmiş bir torunudur. Kapsamlı başvuru kaynağı
\citet{barendregt1984}'tür.
\end{tarihce}

\subsection{Bağlı değişkenler ve de Bruijn indisleri}

$\lambda x.\, x$ ile $\lambda y.\, y$ aynı fonksiyondur. Bağlı
değişkenin adı önemsizdir (\emph{$\alpha$-denklik}). Kâğıt üzerinde bunu
``adları gerektiğinde değiştiririz'' diyerek geçiştiririz. Bir
bilgisayarda ise adlarla çalışmak sürekli hataya açıktır. Nicolaas
de~Bruijn'in çözümü adları tamamen atmaktır \citep{debruijn1972}: Her
değişkenin yerine, onu bağlayan $\lambda$'nın \emph{kaç $\lambda$ dışarıda}
olduğunu yazarız (\cref{fig:de-bruijn}).

\begin{figure}[htbp]
\centering
\begin{tikzpicture}[
    d/.style={font=\large},
    ok/.style={-{Stealth[length=2mm]}, thick, #1},
  ]
  \node[d, anchor=base, inner sep=1pt] (l1) at (0.00,0) {$\lambda$};
  \node[d, anchor=base, inner sep=1pt] (bx) at (0.48,0) {$x$};
  \node[d, anchor=base, inner sep=1pt] (d1) at (0.87,0) {$.$};
  \node[d, anchor=base, inner sep=1pt] (l2) at (1.42,0) {$\lambda$};
  \node[d, anchor=base, inner sep=1pt] (by) at (1.91,0) {$y$};
  \node[d, anchor=base, inner sep=1pt] (d2) at (2.29,0) {$.$};
  \node[d, anchor=base, inner sep=1pt] (x1) at (3.00,0) {$x$};
  \node[d, anchor=base, inner sep=1pt] (p1) at (3.60,0) {$($};
  \node[d, anchor=base, inner sep=1pt] (y1) at (3.97,0) {$y$};
  \node[d, anchor=base, inner sep=1pt] (x2) at (4.50,0) {$x$};
  \node[d, anchor=base, inner sep=1pt] (p2) at (4.91,0) {$)$};
  \node[d, anchor=west] (b) at (7.0,0.12)
    {$\lambda\, .\; \lambda\, .\;\; \textcolor{anarenk}{\mathbf{1}}\;\; (\textcolor{ikincil}{\mathbf{0}}\;\; \textcolor{anarenk}{\mathbf{1}})$};
  \draw[ok=anarenk] (x1.north) to[out=110,in=70,looseness=1.1] (bx.north);
  \draw[ok=ikincil] (y1.north) to[out=105,in=75,looseness=1.0] (by.north);
  \draw[ok=anarenk] (x2.north) to[out=100,in=60,looseness=1.3] (bx.north);
  \node[font=\baslikfont\scriptsize, text=black!55] at (2.4,-0.65) {adlı gösterim};
  \node[font=\baslikfont\scriptsize, text=black!55] at (8.9,-0.65) {de~Bruijn gösterimi};
  \draw[-{Stealth[length=3mm]}, very thick, black!25] (5.4,0.12) -- (6.7,0.12);
  \node[font=\baslikfont\scriptsize, text=black!60, align=center] at (8.9,0.95)
       {\lc{lam (lam (app (var 1)}\\\lc{(app (var 0) (var 1))))}};
\end{tikzpicture}
\caption{De Bruijn indisleri. Her değişken, kendisini bağlayan
$\lambda$'ya olan ``uzaklığıyla'' yazılır: $y$ en yakın $\lambda$'dır (0),
$x$ bir dıştaki (1). $\alpha$-denk terimlerin de~Bruijn gösterimleri
\emph{aynıdır}.}
\label{fig:de-bruijn}
\end{figure}

\leandosya[$\lambda$-terimleri]{Bolum14}{terim}

\noindent \lc{Terim} tipinin \leankw{deriving} \lc{DecidableEq} ile
karar verilebilir eşitliğe sahip olması de~Bruijn gösteriminin bir
armağanıdır. İki terimin $\alpha$-denk olup olmadığını sormak, artık
yalnızca sözdizimsel eşitliği sormaktır.

% ---------------------------------------------------------------------
\section{Yerine koyma}\label{sec:yerine-koyma}

Hesaplamanın tek kuralı, bir fonksiyonu argümanına uygulamaktır:
\[
  (\lambda x.\, t)\; s \;\longrightarrow_\beta\; t[x := s].
\]
$t[x := s]$, $t$'deki serbest $x$'lerin yerine $s$ koymaktır. Bölüm~10'da
formüllerde yerine koymayı tanımlamıştık (\lc{yerineKoy}). Orada bağlayıcı
yoktu ve iş kolaydı. Burada ise bir tuzak var:
\[
  (\lambda x.\, \lambda y.\, x)\; y
  \;\longrightarrow\; \lambda y.\, y \quad\text{\textcolor{hataframe}{(yanlış!)}}
\]
Dışarıdaki serbest $y$, içerideki $\lambda y$ tarafından \emph{yakalandı}.
Sonuç, sabit fonksiyon olması gerekirken özdeşlik fonksiyonu oldu. Doğrusu
$\lambda y'.\, y$'dir. De Bruijn gösteriminde bu sorun, $s$ bir
bağlayıcının altına girdiğinde serbest indislerini \emph{kaydırarak}
çözülür:

\leandosya[Kaydırma ve yerine koyma]{Bolum14}{yerine}

\noindent \lc{yerine k s t}'nin üç durumu vardır: \lc{k}'dan küçük
indisler $t$'nin içinde bağlıdır, dokunulmaz. \lc{k} indisi yerine
konacak değişkendir; oraya, aradaki \lc{k} bağlayıcı kadar kaydırılmış
\lc{s} gelir. \lc{k}'dan büyük indisler ise dışarıdaki serbest
değişkenlerdir; aradaki bir $\lambda$ ortadan kalktığı için bir azalırlar.
Bu küçük indis aritmetiği, adlı gösterimdeki bütün ``değişken adı
değiştirme'' işinin yerini tutar.

\leandosya[Sıfır kaydırma]{Bolum14}{kaydir_sifir}

\begin{dikkat}[Dikkat: Bağlayıcılar formalizasyonun en zor kısmıdır]
Bir mantık ya da programlama dilini formalize ederken en çok zaman alan iş,
çoğu zaman asıl teorem değil, \lc{kaydir} ve \lc{yerine} hakkındaki onlarca
küçük önsavdır. Örneğin ``kaydırmak ile yerine koymak sıra değiştirir''
türünden önsavlar. 2005'teki \emph{POPLmark} yarışması, ispat
asistanlarını tam olarak bu sorun üzerinden karşılaştırmayı amaçladı
\citep{aydemir2005}. Bu bölümde bu önsavlardan yalnızca gerektiği kadarını
ispatlıyoruz.
\end{dikkat}

% ---------------------------------------------------------------------
\section{$β$-indirgeme ve normal biçimler}\label{sec:beta}

Tek adımlık $\beta$-indirgeme, \Lean{}'de dört kurallı bir tümevarımsal
bağıntıdır. Bir $\beta$-adımı terimin \emph{herhangi bir yerinde}
atılabilir. Buna bir terimin içindeki bir alt terim de, bir $\lambda$'nın
gövdesi de dahildir:

\leandosya[$\beta$-indirgeme]{Bolum14}{adim}

\begin{tanim}{Normal biçim}{normal-bicim}
Hiçbir $\beta$-adımı atılamayan bir terim \textbf{normal biçimdedir}.
$t \to^* u$ ve $u$ normal biçimdeyse, $u$'ya $t$'nin \textbf{normal
biçimi} denir.
\end{tanim}

\leandosya[İlk örnekler]{Bolum14}{ornek_adim}

\noindent Normal biçim, bir hesaplamanın \emph{sonucudur}. Ama her terimin
bir sonucu yoktur. En ünlü örnek, kendisini kendisine uygulayan
$\Omega = (\lambda x.\, x\,x)(\lambda x.\, x\,x)$'dır. Bu terim tek bir
adımda kendisine indirgenir (\cref{fig:omega}). Sonsuz bir döngüdür.

\leandosya[$\Omega$: normal biçimi olmayan terim]{Bolum14}{omega}

\begin{figure}[htbp]
\centering
\begin{tikzpicture}[
    ter/.style={draw=#1, fill=#1!8, thick, rounded corners=4pt,
                font=\small, inner sep=5pt},
  ]
  \node[ter=hataframe] (o) at (0,0) {$(\lambda x.\, x\,x)\;(\lambda x.\, x\,x)$};
  \draw[-{Stealth[length=3mm]}, very thick, hataframe]
       (o.north east) .. controls +(1.6,1.4) and +(-1.6,1.4) .. (o.north west)
       node[pos=0.5, above, font=\baslikfont\scriptsize\bfseries]{$\beta$};
  \node[ter=ikincil] (k) at (7.0,0.9) {$K\; I\; \Omega$};
  \node[ter=ikincil] (i) at (7.0,-1.0) {$I$};
  \draw[-{Stealth[length=3mm]}, thick, ikincil] (k) -- node[right, font=\baslikfont\scriptsize]{normal sıra: 2 adım} (i);
  \draw[-{Stealth[length=3mm]}, thick, hataframe] (k.north east) .. controls +(1.2,1.0) and +(1.6,0.2) .. (k.east)
       node[pos=0.55, right, font=\baslikfont\scriptsize, text=hataframe, align=left]{uygulayıcı sıra:\\sonsuz döngü};
\end{tikzpicture}
\caption{Solda: $\Omega$ kendisine indirgenir ve başka hiçbir şeye
indirgenmez (\lc{omega\_tek}). Sağda: $K\,I\,\Omega$'nın normal biçimi
$I$'dır. Ama önce argümanı ($\Omega$'yı) hesaplamaya çalışan bir strateji
bu normal biçimi asla bulamaz.}
\label{fig:omega}
\end{figure}

\subsection{İndirgeme stratejileri}

Bir terimde birden fazla indirgenebilir ifade (\emph{redeks}) olabilir.
Hangisini önce indirgeyeceğimiz bir \emph{strateji} meselesidir.
\Cref{fig:omega}'nin sağ tarafı, seçimin önemli olabileceğini gösteriyor.
Programlama dilleri açısından bu seçim, \emph{değerle çağırma}
(\emph{call-by-value}) ile \emph{adla çağırma} (\emph{call-by-name})
arasındaki farktır.

\begin{center}
\renewcommand{\arraystretch}{1.45}
\rowcolors{2}{tablosatir}{white}
\begin{tabular}{>{\raggedright\arraybackslash}p{3.2cm}
                >{\raggedright\arraybackslash}p{4.6cm}
                >{\raggedright\arraybackslash}p{4.4cm}}
\rowcolor{tabloust}
\tablobaslik{Strateji} & \tablobaslik{Kural} & \tablobaslik{Nerede?} \\
normal sıra & en soldaki, en dıştaki redeks & Haskell (tembel), \lc{adim?} \\
uygulayıcı sıra & önce argümanlar normalleştirilir & çoğu dil, \Lean{}'in derlenmiş kodu \\
\end{tabular}
\end{center}

\begin{teorem}{Normalleştirme (Curry--Feys)}{normal-sira}
Bir terimin normal biçimi varsa, normal sıra stratejisi onu bulur.
\end{teorem}

\noindent Bu teorem, aşağıdaki normalleştiricinin neden normal sırayı
kullandığını açıklıyor. \lc{adim?} bir adım atar ya da terim normal
biçimdeyse \lc{none} döndürür. \lc{normallestir} ise en fazla \lc{n} adım
atar. Sonlanmayan terimler yüzünden yakıt (\emph{fuel}) parametresi
gereklidir: \Lean{}'de her fonksiyonun sonlanması gerekir, $\lambda$-hesabında
ise gerekmez.

\leandosya[Bir normalleştirici]{Bolum14}{normallestir}

\noindent İkinci \leankw{\#eval}, $S\,K\,K$'nın normal biçiminin $I$
olduğunu söylüyor. Bu, kombinatör mantığının ünlü bir gözlemidir:
$I$ aslında gereksizdir, $S$ ve $K$ yeterlidir.

\begin{teorem}{Church--Rosser}{church-rosser}
$t \to^* u_1$ ve $t \to^* u_2$ ise, $u_1 \to^* w$ ve $u_2 \to^* w$
olacak biçimde bir $w$ vardır. Dolayısıyla bir terimin \emph{en fazla bir}
normal biçimi vardır.
\end{teorem}

\noindent Church--Rosser teoremi (\emph{birleşebilirlik}), hesaplamanın
sonucunun stratejiye bağlı olmadığını garanti eder. Strateji yalnızca
sonuca \emph{ulaşıp ulaşmadığımızı} etkiler. İspatı Tait ve Martin-Löf'ün
``paralel indirgeme'' yöntemiyle yapılır ve bu bölümün sonundaki
projelerden biridir.

% ---------------------------------------------------------------------
\section{$λ$-hesabında hesaplama: Church sayıları}\label{sec:church}

Sayı olmayan bir dilde sayıları nasıl ifade ederiz? Church'ün fikri, $n$
sayısını ``bir fonksiyonu $n$ kez uygulamak'' eylemiyle özdeşleştirmektir:
\[
  \overline{n} \;=\; \lambda f.\,\lambda x.\; \underbrace{f\,(f\,(\cdots(f}_{n}\; x)\cdots)).
\]
Toplama ``önce $n$ kez, sonra $m$ kez daha uygulamak'', çarpma ise ``$n$
kez uygulamayı $m$ kez tekrarlamak''tır:

\leandosya[Church sayıları]{Bolum14}{church}

\noindent Son satırdaki \leankw{example}, \lc{decide} ile kanıtlanıyor:
\Lean{} $1 + 1$'i $\lambda$-hesabında gerçekten \emph{hesaplıyor} ve
sonucun \lc{church 2} olduğunu denetliyor. Aynı fikri \Lean{}'in kendi
$\lambda$-hesabıyla, yani sığ gömmeyle de ifade edebiliriz:

\leandosya[Lean'in kendi $\lambda$'sıyla Church sayıları]{Bolum14}{sig_church}

\noindent Buradaki \lc{ChurchSayi} tipinin \lc{Type 1}'de yaşaması
dikkat çekicidir. Bir Church sayısı \emph{her} \lc{α} tipi için çalışır,
bu yüzden \lc{Type}'tan bir üst evrene çıkar. Evrenler Bölüm~15'in
konusudur.

% ---------------------------------------------------------------------
\section{Basit tipler}\label{sec:basit-tipler}

Tipsiz $\lambda$-hesabı çok güçlüdür. Her şey her şeye uygulanabilir, $\Omega$
gibi anlamsız döngüler kurulabilir. Church'ün 1940'taki çözümü, her
terime bir \emph{tip} vermek ve yalnızca tiplerin uyduğu uygulamalara
izin vermekti. Basit tipler, temel tiplerden ok ile kurulur:
\[
  A, B \;::=\; \alpha_n \;\mid\; A \Rightarrow B .
\]
Tipleme bağıntısı $\Gamma \vdash t : A$, ``$\Gamma$ bağlamında $t$'nin tipi
$A$'dır'' diye okunur. De Bruijn gösteriminde bağlam bir tip
listesidir. \lc{var n}'nin tipi listenin \lc{n}'inci elemanıdır.

\leandosya[Basit tipler ve tipleme]{Bolum14}{tip}

\noindent \Cref{fig:tipleme-kurallari}'de üç tipleme kuralı, Bölüm~11'deki
doğal çıkarım kurallarıyla yan yana duruyor. Benzerlik bir tesadüf
değildir. Bu bölümün geri kalanı, bu benzerliği bir teoreme
dönüştürmektir.

\begin{figure}[htbp]
\centering
\setlength{\tabcolsep}{10pt}
\renewcommand{\arraystretch}{2.8}
\begin{tabular}{r c c}
& \textcolor{anarenk}{\baslikfont\bfseries Tipleme} & \textcolor{ikincil}{\baslikfont\bfseries Doğal çıkarım} \\
\baslikfont\scriptsize değişken &
\begin{prooftree}
  \hypo{\Gamma(n) = A}
  \infer1{\Gamma \vdash \mathtt{var}\;n : A}
\end{prooftree}
&
\begin{prooftree}
  \hypo{A \in \Gamma}
  \infer1[Var]{\Gamma \vdash A}
\end{prooftree}
\\
\baslikfont\scriptsize soyutlama &
\begin{prooftree}
  \hypo{A, \Gamma \vdash t : B}
  \infer1{\Gamma \vdash \lambda.\,t : A \Rightarrow B}
\end{prooftree}
&
\begin{prooftree}
  \hypo{A, \Gamma \vdash B}
  \infer1[$\to$G]{\Gamma \vdash A \oise B}
\end{prooftree}
\\
\baslikfont\scriptsize uygulama &
\begin{prooftree}
  \hypo{\Gamma \vdash t : A \Rightarrow B}
  \hypo{\Gamma \vdash u : A}
  \infer2{\Gamma \vdash t\;u : B}
\end{prooftree}
&
\begin{prooftree}
  \hypo{\Gamma \vdash A \oise B}
  \hypo{\Gamma \vdash A}
  \infer2[$\to$E]{\Gamma \vdash B}
\end{prooftree}
\\
\end{tabular}
\caption{Basit tipli $\lambda$-hesabının kuralları (solda) ve doğal
çıkarımın $\to$ parçası (sağda). Terimleri silerseniz, soldaki kurallar
sağdakilere dönüşür.}
\label{fig:tipleme-kurallari}
\end{figure}

\subsection{Kombinatörler ve Hilbert sistemleri}

Üç kombinatörün tipleri tanıdık gelmeli:

\leandosya[S, K ve I'nın tipleri]{Bolum14}{ski_tip}

\noindent $K$'nın tipi $A \Rightarrow B \Rightarrow A$, $S$'ninki ise
$(A \Rightarrow B \Rightarrow C) \Rightarrow (A \Rightarrow B)
\Rightarrow A \Rightarrow C$'dir. Bunlar, önermeler mantığının
\emph{Hilbert tarzı} aksiyom sisteminin iki aksiyomudur. Hilbert
sisteminin tek kuralı modus ponenstir ve uygulamaya karşılık gelir.
Yukarıda gördüğümüz $S\,K\,K = I$ eşitliği, Hilbert sistemindeki ünlü
$\vdash A \oise A$ türetiminin ta kendisidir. Curry ve Feys bu gözlemi
1958'de yaptılar \citep{curry1958}. Curry--Howard eşlemesinin ilk yarısı
budur.

% ---------------------------------------------------------------------
\section{$Ω$ tiplenemez}\label{sec:omega-tip}

$\lambda x.\, x\,x$ terimine bir tip vermeye çalışalım. $x$'in tipi $A$
olsun. $x\,x$ uygulamasının anlamlı olması için $x$'in bir ok tipi
olması gerekir: $A = A \Rightarrow B$. Ama hiçbir sonlu tip kendisini
içeren bir ok tipine eşit olamaz:

\leandosya[$\lambda x.\, x\,x$ tiplenemez]{Bolum14}{omega_tip}

\noindent Bu küçük gözlemin büyük bir sonucu vardır:

\begin{teorem}{Güçlü normalleşme (Tait, 1967)}{guclu-normallesme}
Basit tipli $\lambda$-hesabında tiplenebilen her terim \emph{güçlü
normalleşir}: Hangi stratejiyle indirgenirse indirgensin, her indirgeme
dizisi sonludur.
\end{teorem}

\noindent İspatın ana fikri, Tait'in \emph{indirgenebilirlik} (ya da
\emph{mantıksal bağıntılar}) yöntemidir \citep{tait1967}. Her tip için
``iyi davranan terimler'' kümesi tip üzerinde özyinelemeyle tanımlanır,
sonra her tiplenebilir terimin kendi tipinin kümesinde olduğu gösterilir.
Bu yöntemin \Lean{}'deki formalizasyonu \cref{bol:proje}'deki projelerden
biridir. Bir uyarı: Basit tipli $\lambda$-hesabı her terimin sonlandığı
bir dil olduğu için artık Turing-tam \emph{değildir}. \Lean{}'in her
fonksiyonun sonlanmasını istemesinin, ama karşılığında mantık olarak
tutarlı kalmasının nedeni budur.

\begin{onemlifikir}[Önemli Fikir: Tipler tutarlılık getirir]
Tipsiz $\lambda$-hesabının mantık olarak tutarsız olması ile $\Omega$ gibi
sonlanmayan terimlerin varlığı aynı olgunun iki yüzüdür. Bir mantıkta her
önermenin bir ``kanıtı'' olabiliyorsa mantık işe yaramaz. Sonsuz bir döngü
ise herhangi bir tipe ``sahipmiş gibi'' davranabilir. Tipler döngüleri
yasaklayarak mantığı kurtarır. Kitabın başından beri \Lean{}'in bazı
tanımlarımızı reddetmesinin (sonlanma denetimi, Bölüm~7; pozitiflik
koşulu, Bölüm~9) altında bu ilke yatar.
\end{onemlifikir}

% ---------------------------------------------------------------------
\section{Bir tip denetleyici}\label{sec:tip-denetleyici}

Şimdiye kadar kullandığımız tipleme \emph{Curry tarzıdır}: Terimlerde
tip yazılmaz, bir terimin birden çok tipi olabilir ($I$, her $A$ için
$A \Rightarrow A$ tipindedir). \emph{Church tarzında} ise her $\lambda$,
bağladığı değişkenin tipini taşır: $\lambda x{:}A.\, t$. \Lean{}'in kendisi
Church tarzındadır. \leankw{fun} \lc{(x : Nat) => x + 1} yazarız. Church
tarzında tip bulmak bir \emph{algoritmadır}:

\leandosya[Church tarzı terimler ve tip çıkarımı]{Bolum14}{denetleyici}

\noindent Bir tip denetleyicinin en önemli özelliği \emph{doğru}
olmasıdır. Denetleyici bir tip bulduğunda, o tip tipleme kurallarına göre
gerçekten geçerli olmalıdır:

\leandosya[Denetleyici doğrudur]{Bolum14}{denetleyici_dogru}

\noindent Bu küçük teorem, ispat asistanlarının güvenilirliği hakkındaki
büyük bir sorunun minyatürüdür. \Lean{}'e bir ispatı kabul ettiren şey,
çekirdeğindeki tip denetleyicidir. Peki o denetleyicinin kendisi doğru
mu? Lean4Lean projesi, \Lean{}'in tip denetleyicisini \Lean{}'de yeniden
yazıp doğruluğunu \Lean{}'de ispatlamayı hedefler \citep{carneiro2026}.

% ---------------------------------------------------------------------
\section{Curry--Howard eşlemesi}\label{sec:curry-howard}

Artık bölümün ana teoremine hazırız. Her tipi bir formüle çeviriyoruz:
Temel tipler atomlara, ok tipi ise $\oise$'ya gider. Sonra tiplenebilen
her terimin, Bölüm~11'deki \lc{Turetilir} sisteminde bir türetime
dönüştüğünü gösteriyoruz:

\leandosya[Curry--Howard teoremi]{Bolum14}{curry_howard}

\begin{teorem}{Curry--Howard (basit tipler)}{curry-howard}
$\Gamma \vdash t : A$ ise $\Gamma^\circ \vdash A^\circ$. Burada $^\circ$,
tiplerin formüllere çevirisidir. Tersine, $\to$'lu bir formülün her
sezgici doğal çıkarım türetimi, o formüle karşılık gelen tipte bir
terimdir.
\end{teorem}

\noindent İspatın üç satırı, \cref{fig:tipleme-kurallari}'deki üç
satırdır. Teoremin ters yönü (türetimden terim üretmek) de aynı kolaylıkla
ispatlanabilir. Ama hedefteki sistemin RAA'sız olması gerekir. Çünkü
Peirce yasasının tipinde hiçbir terim yoktur. \Cref{tab:curry-howard},
eşlemenin kitapta karşılaştığımız bütün katmanlarını özetliyor.

\begin{table}[htbp]
\centering
\small
\renewcommand{\arraystretch}{1.45}
\rowcolors{2}{tablosatir}{white}
\begin{tabular}{>{\raggedright\arraybackslash}p{5.2cm}
                >{\raggedright\arraybackslash}p{5.2cm}
                >{\centering\arraybackslash}p{1.4cm}}
\rowcolor{tabloust}
\tablobaslik{Mantık} & \tablobaslik{Programlama} & \tablobaslik{Bölüm} \\
önerme $A$                     & tip $A$ & 2 \\
$A$'nın ispatı                 & $A$ tipinde bir terim & 2 \\
$A \to B$                      & fonksiyon tipi & 2, 14 \\
$\to$G (varsayımı boşaltma)    & $\lambda$-soyutlama & 5, 14 \\
$\to$E (modus ponens)          & uygulama & 5, 14 \\
$A \wedge B$, $A \vee B$       & çarpım (yapı), toplam (seçenek) tipi & 3, 9 \\
$\forall x,\, P\,x$, $\exists x,\, P\,x$ & bağımlı fonksiyon, bağımlı çift & 6, 15 \\
dolambaç (giriş + eleme)       & redeks $(\lambda x.\,t)\,s$ & 5, 14 \\
türetim normalleştirme         & program çalıştırma ($\beta$) & 5, 14 \\
tutarlılık                     & güçlü normalleşme & 14 \\
Hilbert aksiyomları K, S       & kombinatörler K, S & 14 \\
klasik mantık (Peirce)         & devamlar, \lc{call/cc} \citep{griffin1990} & 4 \\
\end{tabular}
\caption{Curry--Howard eşlemesinin katmanları \citep{howard1980,wadler2015}.}
\label{tab:curry-howard}
\end{table}

\noindent Teoremin hemen bir sonucu vardır: Bir tipte kapalı bir terim
varsa, o tipin formülü bir totolojidir. Bu, bazı tiplerin \emph{boş}
olduğunu göstermenin bir yoludur:

\leandosya[Boş tipler]{Bolum14}{sonuclar}

\begin{figure}[htbp]
\centering
\renewcommand{\arraystretch}{1.2}
\setlength{\tabcolsep}{12pt}
\begin{tabular}{c c c}
\textcolor{anarenk}{\baslikfont\small\bfseries Program} & & \textcolor{ikincil}{\baslikfont\small\bfseries İspat} \\[6pt]
\fcolorbox{anarenk}{tanimbg}{$(\lambda x{:}A.\; t)\; s$} &
\raisebox{-0.2em}{$\longleftrightarrow$} &
\begin{prooftree}
  \hypo{[A]}
  \ellipsis{$t$}{B}
  \infer1[$\to$G]{A \oise B}
  \hypo{}
  \ellipsis{$s$}{A}
  \infer2[$\to$E]{B}
\end{prooftree}
\\[10pt]
$\Big\downarrow\;\textcolor{anarenk}{\beta}$ & &
$\Big\downarrow\;\textcolor{ikincil}{\text{\baslikfont\scriptsize normalleştir}}$
\\[10pt]
\fcolorbox{anarenk}{tanimbg}{$t[x := s]$} &
\raisebox{-0.2em}{$\longleftrightarrow$} &
\begin{prooftree}
  \hypo{}
  \ellipsis{$s$}{A}
  \ellipsis{$t$}{B}
\end{prooftree}
\\
\end{tabular}
\caption{Hesaplama = ispat normalleştirme. Solda bir $\beta$-adımı, sağda
ona karşılık gelen dolambaç eleme: $\to$G'nin hemen ardından gelen
$\to$E silinir ve $A$'nın varsayıldığı yerlere $s$'nin türetimi
yerleştirilir.}
\label{fig:beta-dolambac}
\end{figure}

% ---------------------------------------------------------------------
\begin{arastirmanotu}
\textbf{Bağlayıcılarla çalışmak.} De~Bruijn indisleri bu bölümde
kullandığımız ``yerel'' çözümdür. Araştırma literatüründe başka
yaklaşımlar da vardır: \emph{yerel adsız} (\emph{locally nameless})
gösterim, \emph{nominal} teknikler, \emph{yüksek mertebe soyut sözdizim}
(HOAS) ve yerine koyma önsavlarını otomatik üreten araçlar (Coq için
Autosubst gibi). POPLmark yarışması \citep{aydemir2005} bu yaklaşımları
karşılaştırmak için hâlâ bir ölçüttür.

\textbf{Tip denetleyicilerin doğrulanması.} Coq'un çekirdeğinin
doğruluğu MetaCoq projesinde Coq'un kendisi içinde ele alındı. \Lean{}
için benzer çalışma Lean4Lean'dir \citep{carneiro2026}. Bu çalışma
\Lean{}'in tip kuramının (tutarlılık hariç) önemli özelliklerini \Lean{}'de
formalize eder ve \Lean{}'in çekirdeğinden bağımsız bir ikinci tip
denetleyici sunar.

\textbf{Curry--Howard'ın ötesi.} Eşleme yalnızca $\to$ ile sınırlı
değildir. Bağımlı tipler yüklemler mantığına (Bölüm~15), doğrusal tipler
doğrusal mantığa \citep{girard1987}, devamlar klasik mantığa
\citep{griffin1990}, kip tipleri ise kip mantığına karşılık gelir. Wadler'in
``Propositions as Types'' makalesi \citep{wadler2015} bu öykünün okunabilir
bir özetidir.
\end{arastirmanotu}

\begin{bolumozeti}
\begin{itemize}
  \item $\lambda$-hesabının üç yapısı vardır: değişken, soyutlama,
        uygulama. De~Bruijn indisleri, bağlı değişken adlarını ortadan
        kaldırarak $\alpha$-denkliği sözdizimsel eşitliğe indirger.
  \item Yerine koyma, yakalamayı önlemek için serbest indislerin
        kaydırılmasını gerektirir. Bağlayıcılar formalizasyonun en zor
        kısmıdır.
  \item $\beta$-indirgeme hesaplamanın tek kuralıdır. $\Omega$'nın normal
        biçimi yoktur. Normal sıra, varsa normal biçimi bulur.
        Church--Rosser teoremi normal biçimin tekliğini garanti eder.
  \item Church sayıları, sayıların bile yalnızca fonksiyonlarla ifade
        edilebildiğini gösterir.
  \item Basit tipler $\Omega$ gibi terimleri yasaklar. Tiplenebilen her
        terim güçlü normalleşir.
  \item Church tarzı terimler için tip bulmak bir algoritmadır, ve bu
        algoritmanın doğruluğu \Lean{}'de ispatlanabilir.
  \item Curry--Howard: Tiplenebilen her terim, tipinin bir doğal çıkarım
        türetimidir. $\beta$-indirgeme, türetimlerdeki dolambaçların
        elenmesidir.
\end{itemize}
\end{bolumozeti}

\begin{alistirmalar}
\alistirma{1} $K$ normal biçimdedir. (\lc{a14\_1})

\alistirma{1} $I\,(I\,s) \to^* s$. (\lc{a14\_2})

\alistirma{2} Bileşke kombinatörü $B = \lambda f.\,\lambda g.\,\lambda x.\,
f\,(g\,x)$'in tipi $(B \Rightarrow C) \Rightarrow (A \Rightarrow B)
\Rightarrow A \Rightarrow C$'dir. Bu tip hangi mantık yasasına karşılık
geliyor? (\lc{a14\_3})

\alistirma{2} Her Church sayısı $(A \Rightarrow A) \Rightarrow A \Rightarrow
A$ tipindedir. Önce \lc{tekrarla} için bir önsav ispatlayın. (\lc{a14\_4})

\alistirma{2} Bir temel tipte kapalı terim yoktur. (\lc{a14\_5})

\alistirma{2} Curry tarzında bir terimin tipi tek değildir. Church tarzında
neden tektir? (\lc{a14\_6})

\alistirma{2} $(\alpha_0 \Rightarrow \alpha_1) \Rightarrow \alpha_0$
tipinde kapalı terim yoktur. (\lc{a14\_7})

\alistirma{3} $\Omega$ hiçbir bağlamda tiplenemez. (\lc{a14\_8})

\alistirma{1} (Kâğıt-kalem) $\lambda x.\,\lambda y.\,\lambda z.\, x\,z\,(y\,z)$
terimini de~Bruijn gösterimine çevirin ve \lc{S} tanımıyla
karşılaştırın.

\alistirma{2} (Kâğıt-kalem) $K\,I\,\Omega$ terimini normal sırayla adım
adım indirgeyin. \lc{\#eval normallestir 10 (app (app K I) Omega)} ile
denetleyin.

\alistirma{3} (Proje) \lc{adim? t = some u} ise \lc{Adim t u} olduğunu
ispatlayın. İpucu: \lc{t} üzerinde tümevarım; \lc{app} durumunda
\lc{t}'nin bir \lc{lam} olup olmadığına göre ayırın.

\alistirma{3} (Proje) \emph{Tip korunumu}: \lc{Tiplenir Γ t A} ve
\lc{Adim t u} ise \lc{Tiplenir Γ u A}. Önce \lc{kaydir} için bir
zayıflatma önsavı, sonra \lc{yerine} için bir yerine koyma önsavı
gerekecek.

\alistirma{3} (Araştırma) Curry--Howard'ın ters yönünü ispatlayın: Bölüm~11'deki
sistemin yalnızca \lc{varsayim}, \lc{iseG}, \lc{iseE} kurallarını kullanan
bir \lc{TuretilirOk} parçasını tanımlayın ve her türetimden bir terim
üretin.
\end{alistirmalar}
   % Basit tipli λ-hesabı
% =====================================================================
%  Bölüm 15 — Bağımlı Tipler ve Lean'in Çekirdeği
% =====================================================================
\chapter{Bağımlı Tipler ve Lean'in Çekirdeği}\label{bol:bagimli}

\begin{bolumhedefleri}
\begin{itemize}
  \item Bağımlı fonksiyon ($Π$) ve bağımlı çift ($Σ$) tiplerini ve bunların
        $\forall$ ve $\exists$ ile ilişkisini,
  \item Tip aileleriyle ``doğruluğu tipine yazılmış'' veri yapıları kurmayı
        (uzunluk indisli vektörler, \lc{Fin}),
  \item Martin-Löf tip kuramını, İnşalar Hesabını ve $\lambda$-küpünü,
  \item \Lean{}'in evren hiyerarşisini ve \lc{Prop}'un impredikatifliğini,
  \item Kanıt-ilgisizliğini, tanımsal ve önermesel eşitlik ayrımını,
  \item \Lean{}'in güvenilir çekirdeğini ve \leankw{\#print} \lc{axioms}
        komutunu.
\end{itemize}
\end{bolumhedefleri}

Bölüm~14'teki basit tiplerde terimler ve tipler iki ayrı dünyada
yaşıyordu: Tipler terimlere bağlı olamıyordu. $A \Rightarrow B$ yazabiliyor,
ama ``$n$ uzunluğunda listeler'' gibi bir tip yazamıyorduk. Bu bölümde iki
dünyayı birleştiriyoruz. Tiplerin terimlere bağlı olabildiği bir dilde,
yani \emph{bağımlı tip kuramında}, Curry--Howard eşlemesi önermeler
mantığından yüklemler mantığına genişler: $\forall$ bir bağımlı fonksiyon,
$\exists$ ise bir bağımlı çift olur. \Lean{}'in kendisi böyle bir dildir.
Bu bölümün sonunda \Lean{}'in çekirdeğinin neye inandığını ve neden ona
güvenebileceğimizi anlatacağız.

% ---------------------------------------------------------------------
\section{$Π$-tipleri: bağımlı fonksiyonlar}\label{sec:pi}

Sıradan bir fonksiyon $f : A \to B$, her $a : A$ için \emph{aynı} $B$
tipinde bir değer döndürür. Bağımlı bir fonksiyonun dönüş tipi ise
argümana göre değişebilir:
\[
  f : \prod_{x : A} B(x)
  \qquad\text{\Lean{}'de}\qquad
  f : (x : A) \to B\;x .
\]
Burada $B : A \to \mathsf{Type}$ bir \emph{tip ailesidir}: Her $x$'e bir
tip atar. Basit bir örnekle başlayalım:

\leandosya[Π-tipleri]{Bolum15}{pi}

\noindent \lc{ornekDeger true} bir doğal sayı, \lc{ornekDeger false} ise
bir metindir. \Cref{fig:pi-sigma-lif}'nın sol tarafı bunu geometrik olarak
gösteriyor. Bir tip ailesini, taban tipin her noktasının üzerinde duran
bir ``lif'' olarak düşünebiliriz. Bağımlı bir fonksiyon, her lif üzerinden
bir nokta seçer.

\begin{figure}[htbp]
\centering
\begin{tikzpicture}[
    lif/.style={draw=#1, fill=#1!8, thick, rounded corners=4pt,
                minimum width=1.6cm, minimum height=2.6cm},
    nokta/.style={circle, fill=#1, inner sep=2.2pt},
  ]
  % --- Π ---
  \node[font=\baslikfont\bfseries, text=anarenk] at (1.6,3.6) {$Π$: her liften bir nokta};
  \node[lif=anarenk] (l1) at (0.6,1.6) {};
  \node[lif=ikincil] (l2) at (2.6,1.6) {};
  \node[font=\baslikfont\scriptsize, text=anarenk] at (0.6,3.05) {\lc{Nat}};
  \node[font=\baslikfont\scriptsize, text=ikincil] at (2.6,3.05) {\lc{String}};
  \node[nokta=teoremframe] (p1) at (0.6,1.9) {};
  \node[nokta=teoremframe] (p2) at (2.6,1.2) {};
  \node[font=\scriptsize, right=1pt of p1] {42};
  \node[font=\scriptsize, right=1pt of p2] {``kırk iki''};
  \draw[very thick, black!60] (-0.2,-0.2) -- (3.4,-0.2);
  \node[nokta=black!70] (b1) at (0.6,-0.2) {};
  \node[nokta=black!70] (b2) at (2.6,-0.2) {};
  \node[font=\scriptsize, below=2pt of b1] {\lc{true}};
  \node[font=\scriptsize, below=2pt of b2] {\lc{false}};
  \node[font=\baslikfont\scriptsize, text=black!55] at (4.0,-0.2) {\lc{Bool}};
  \draw[teoremframe, thick, dashed] (p1) -- (p2);
  % --- Σ ---
  \begin{scope}[xshift=7.2cm]
  \node[font=\baslikfont\bfseries, text=ikincil] at (1.6,3.6) {$Σ$: liflerin ayrık birleşimi};
  \node[lif=anarenk] at (0.6,1.6) {};
  \node[lif=ikincil] at (2.6,1.6) {};
  \node[font=\baslikfont\scriptsize, text=anarenk] at (0.6,3.05) {\lc{Vek Nat 0}};
  \node[font=\baslikfont\scriptsize, text=ikincil] at (2.6,3.05) {\lc{Vek Nat 2}};
  \node[nokta=teoremframe] (q) at (2.6,1.6) {};
  \node[font=\scriptsize, below=2pt of q, align=center] {\lc{⟨2, [7,8]⟩}};
  \draw[very thick, black!60] (-0.2,-0.2) -- (3.4,-0.2);
  \node[nokta=black!70] (c1) at (0.6,-0.2) {};
  \node[nokta=black!70] (c2) at (2.6,-0.2) {};
  \node[font=\scriptsize, below=2pt of c1] {0};
  \node[font=\scriptsize, below=2pt of c2] {2};
  \node[font=\scriptsize, text=black!55] at (1.6,1.6) {$\cdots$};
  \node[font=\baslikfont\scriptsize, text=black!55] at (4.0,-0.2) {\lc{Nat}};
  \end{scope}
\end{tikzpicture}
\caption{Bir tip ailesini, taban tipin her noktası üzerinde duran bir
lif olarak düşünmek. Solda: bir $Π$-tipi elemanı, \emph{her} liften bir
nokta seçer (bir ``kesit''). Sağda: bir $Σ$-tipi elemanı, \emph{bir} lif
ve o liften bir noktadır.}
\label{fig:pi-sigma-lif}
\end{figure}

\noindent Bu bakış açısı Bölüm~6'da $\forall$ için verdiğimiz okumayla
örtüşür: $\forall x,\, P\,x$'in bir kanıtı, her $x$ için $P\,x$'in bir
kanıtını veren bir fonksiyondur. Yani \emph{$\forall$, dönüş tipi
\lc{Prop} olan bir $Π$-tipidir}. \lc{pi} kesitindeki son satır bunu
doğruluyor. \Lean{}'de \lc{∀ n : Nat, P n} ile \lc{(n : Nat) → P n} aynı
terimdir. $A \to B$ ise $B$'nin $x$'e bağlı olmadığı özel durumdur.

% ---------------------------------------------------------------------
\section{Doğruluğu tipine yazılmış veri: vektörler}\label{sec:vektor}

Bağımlı tiplerin programlamadaki en bilinen kullanımı, bir veri yapısının
bir özelliğini tipine yazmaktır. Bölüm~9'daki \lc{Liste} tipinin boş bir
listenin başını almaya çalışınca ne yapacağı belli değildi; \lc{Option}
döndürmek zorundaydık. Uzunluğu tipinde taşıyan bir vektör tipi bu
sorunu ortadan kaldırır:

\leandosya[Uzunluk indisli vektörler]{Bolum15}{vek}

\noindent \lc{Vek α n}, Bölüm~9'da gördüğümüz \emph{tümevarımsal
ailelerin} bir örneğidir: \lc{α} bir \emph{parametre}, \lc{n} ise bir
\emph{indistir}. Kurucular indisi belirler: \lc{bos}'un uzunluğu 0,
\lc{ekle a v}'ninki \lc{v}'ninkinden bir fazladır. Fonksiyonlara
bakalım:

\begin{itemize}
  \item \lc{bas}, yalnızca \lc{Vek α (n + 1)} kabul eder. Boş bir vektörün
        tipi \lc{Vek α 0}'dır ve \lc{0}, \lc{n + 1} biçiminde değildir.
        Bu yüzden \lc{bas}'ın tanımında \lc{bos} durumu yoktur: \Lean{} bu
        durumun \emph{olanaksız} olduğunu kendisi görür.
  \item \lc{fermuar}, iki vektörün aynı uzunlukta olmasını \emph{tipinde}
        ister. Uzunlukları farklı iki vektörü fermuarlamaya çalışmak bir
        çalışma zamanı hatası değil, bir tip hatasıdır.
  \item \lc{birlestir}'in sonuç tipi \lc{Vek α (n + m)}'dir. Tanım,
        \lc{Nat.add}'ın ikinci argümanı üzerinde özyinelemeli olmasından
        yararlanır. \lc{bos} durumunda \lc{n + 0} ile \lc{n}, \lc{ekle}
        durumunda \lc{n + (k + 1)} ile \lc{(n + k) + 1} \emph{tanımsal
        olarak} eşittir.
\end{itemize}

\leandosya[Vektörlerle çalışmak]{Bolum15}{vek_ornek}

\begin{dikkat}[Dikkat: İndis aritmetiği]
\lc{birlestir}'in sonuç tipini \lc{Vek α (m + n)} diye yazsaydık
tanım reddedilirdi. \lc{bos} durumunda \lc{0 + n} ile \lc{n}
\emph{tanımsal} olarak eşit değildir; yalnızca \emph{önermesel} olarak
eşittir (\lc{Nat.zero\_add}). Bağımlı tiplerle programlamanın en sık
karşılaşılan güçlüğü budur. Tip denetleyici yalnızca hesaplayarak görülen
eşitlikleri kabul eder; geri kalanlar için eşitliği açıkça kanıtlayıp
\lc{h ▸ v} ile ``taşımak'' gerekir (\cref{sec:tanimsal}).
\end{dikkat}

\subsection{\lc{Fin}: güvenli indisleme}

\lc{Fin n}, $n$'den küçük doğal sayıların tipidir. Bir sayı ile onun $n$'den
küçük olduğunun kanıtından oluşan bir yapıdır. \lc{Vek α n}'yi \lc{Fin n}
ile indislemek ``dizi sınırı aşıldı'' hatasını tip düzeyinde olanaksız
kılar:

\leandosya[Fin ile indisleme]{Bolum15}{al}

\noindent \lc{al}'ın tanımında \lc{bos} durumunda \lc{Fin 0}'ın elemanı
olmadığı için \lc{Fin.elim0} kullanıyoruz. Özyinelemeli durumda, $i + 1 <
k + 1$ kanıtından $i < k$ kanıtını \lc{Nat.lt\_of\_succ\_lt\_succ} ile
elde ediyoruz. Veri ile kanıt bu fonksiyonda iç içedir.

% ---------------------------------------------------------------------
\section{Bağımlı tipler üzerinde ispat}\label{sec:bagimli-ispat}

Vektörleri ters çeviren bir fonksiyon yazalım ve iki kez ters
çevirmenin vektörü değiştirmediğini ispatlayalım. Bölüm~9'da aynı teoremi
listeler için ispatlamıştık. Burada tipler, ters çevirmenin uzunluğu
koruduğunu zaten garanti ediyor:

\leandosya[Ters çevirme]{Bolum15}{ters}

\noindent İspatlar, listelerdekilerle neredeyse aynıdır. Fark, tümevarımın
artık bir \emph{tümevarımsal aile} üzerinde yapılmasıdır. \leantac{induction}
\lc{v}, \lc{v}'nin uzunluğunu da otomatik olarak genelleştirir. \lc{ekle}
durumunda uzunluk \lc{k + 1}, tümevarım hipotezinde ise \lc{k}'dır.

\leandosya[Eşitlik boyunca taşıma]{Bolum15}{donustur}

\noindent \lc{h ▸ v} ifadesi, \lc{h : n = m} eşitliğini kullanarak \lc{v}'yi
\lc{Vek α n} tipinden \lc{Vek α m} tipine ``taşır''. Bölüm~9'daki \lc{Esit}
tanımının eleme ilkesi tam olarak budur: Eşit şeyler, birbirinin yerine
konabilir.

% ---------------------------------------------------------------------
\section{$Σ$-tipleri: bağımlı çiftler}\label{sec:sigma}

$Σ$-tipleri, $Π$-tiplerinin ikizidir. $\sum_{x : A} B(x)$ tipinin bir
elemanı, bir $a : A$ ile bir $b : B(a)$'dan oluşan bir çifttir. İkinci
bileşenin tipi birinci bileşene bağlıdır. \Lean{}'de üç çeşidi vardır:

\leandosya[Σ-tipleri, alt tipler ve $\exists$]{Bolum15}{sigma}

\begin{table}[htbp]
\centering
\small
\renewcommand{\arraystretch}{1.45}
\rowcolors{2}{tablosatir}{white}
\begin{tabular}{>{\raggedright\arraybackslash}p{3.1cm}
                >{\raggedright\arraybackslash}p{2.3cm}
                >{\raggedright\arraybackslash}p{2.5cm}
                >{\raggedright\arraybackslash}p{3.8cm}}
\rowcolor{tabloust}
\tablobaslik{Yazım} & \tablobaslik{İkinci bileşen} & \tablobaslik{Sonuç} & \tablobaslik{Tanık çıkarılabilir mi?} \\
\lc{(x : α) × β x}     & veri (\lc{Type})  & veri (\lc{Type}) & evet: \lc{.1} \\
\lc{\{x : α // P x\}}  & kanıt (\lc{Prop}) & veri (\lc{Type}) & evet: \lc{.val} \\
\lc{∃ x : α, P x}      & kanıt (\lc{Prop}) & önerme (\lc{Prop}) & hayır (yalnızca bir kanıt içinde) \\
\end{tabular}
\caption{\Lean{}'deki üç $Σ$-tipi. Fark, sonucun hangi evrende yaşadığıdır.}
\label{tab:sigma-cesitleri}
\end{table}

\noindent Tablonun son satırı önemlidir. \lc{∃ n, n > 3}'ün bir kanıtından,
\emph{bir program içinde} tanığı (4'ü) çıkaramayız:

\begin{leanhata}[Lean hata mesajı (özet)]
def tanik (h : ∃ n : Nat, n > 3) : Nat :=
  match h with | ⟨n, _⟩ => n
-- hata: Exists.casesOn yalnızca Prop'a eleme yapabilir
\end{leanhata}

\noindent Bunun nedeni bir sonraki kesimde göreceğimiz
\emph{kanıt-ilgisizliğidir}: Bir önermenin iki kanıtı \Lean{} için ayırt
edilemez. \lc{∃ n, n > 3}'ün biri 4'ü, diğeri 5'i tanık gösteren iki kanıtı
\emph{aynı} sayılır. Kanıttan tanığı çıkaran bir fonksiyon, aynı girdiye
iki farklı cevap vermek zorunda kalırdı. Bu yüzden \lc{Prop}'tan
\lc{Type}'a ``büyük eleme'' genel olarak yasaktır. Bir \emph{ispat içinde}
ise tanığı kullanmak serbesttir (\leantac{obtain} \lc{⟨n, hn⟩ := h}), çünkü
sonuç yine bir önermedir. Bu kısıtlama, Bölüm~4'teki yapıcı/klasik
tartışmasının bir başka yüzüdür: \Lean{}'de varlık kanıtları
\emph{hesaplama içeriği taşımaz}. Hesaplama içeriği isteniyorsa $Σ$ ya da
alt tip kullanılır.

% ---------------------------------------------------------------------
\section{Martin-Löf tip kuramı}\label{sec:mltt}

Bağımlı tiplerin mantık için sistemli bir kuramını İsveçli mantıkçı Per
Martin-Löf 1970'lerde geliştirdi \citep{martinlof1984}. Martin-Löf tip
kuramında (MLTT) her tip oluşturucu dört tür kuralla tanıtılır. Bu kalıp,
Bölüm~5'teki giriş/eleme kurallarının tiplere genişletilmiş hâlidir:

\begin{table}[htbp]
\centering
\small
\renewcommand{\arraystretch}{1.6}
\rowcolors{2}{tablosatir}{white}
\begin{tabularx}{\linewidth}{>{\raggedright\arraybackslash}p{2.4cm} X X}
\rowcolor{tabloust}
\tablobaslik{Kural} & \tablobaslik{$Π$} & \tablobaslik{$Σ$} \\
oluşum   & $A$ tip, $x{:}A \vdash B$ tip $\Rightarrow$ $Π_{x:A} B$ tip & $A$ tip, $x{:}A \vdash B$ tip $\Rightarrow$ $Σ_{x:A} B$ tip \\
giriş    & $x{:}A \vdash b : B$ $\Rightarrow$ $\lambda x.\, b : Π_{x:A} B$ & $a : A$, $b : B[a/x]$ $\Rightarrow$ $\langle a, b\rangle : Σ_{x:A}B$ \\
eleme    & $f : Π_{x:A}B$, $a : A$ $\Rightarrow$ $f\,a : B[a/x]$ & $p : Σ_{x:A}B$ $\Rightarrow$ $p.1 : A$, $p.2 : B[p.1/x]$ \\
hesaplama & $(\lambda x.\, b)\,a \equiv b[a/x]$ \hfill($\beta$) & $\langle a, b\rangle.1 \equiv a$, $\langle a, b\rangle.2 \equiv b$ \\
\end{tabularx}
\caption{Martin-Löf tip kuramında $Π$ ve $Σ$ için kurallar. $\equiv$,
tanımsal eşitliktir.}
\label{tab:mltt}
\end{table}

\noindent MLTT'nin üçüncü temel tipi \emph{özdeşlik tipidir}: $a =_A b$.
Bölüm~9'da \lc{Esit} olarak tanımladığımız tümevarımsal aile budur. Tek
kurucusu $\mathsf{refl} : a = a$'dır, eleme ilkesi ise ``eşitleri
birbirinin yerine koyma''dır (\lc{▸}). Özdeşlik tiplerinin beklenmedik
derecede zengin bir yapısı vardır. Bir $p : a = b$ kanıtını $a$'dan $b$'ye
bir \emph{yol} olarak düşünmek, \emph{homotopi tip kuramına} götürür
\citep{hott2013}.

\begin{tarihce}[Tarihçe: Martin-Löf'ten Lean'e]
Martin-Löf'ün ilk sistemi (1971) $\mathsf{Type} : \mathsf{Type}$ kuralını
içeriyordu. Jean-Yves Girard aynı yıl bu sistemin tutarsız olduğunu
gösterdi \citep{girard1972}. Martin-Löf bunun üzerine evren hiyerarşisini
getirdi. Thierry Coquand ve Gérard Huet 1988'de \emph{İnşalar Hesabını}
(Calculus of Constructions, CoC) tanıttılar \citep{coquand1988}. Christine
Paulin-Mohring ve Coquand bu sisteme tümevarımsal tipleri ekleyerek
\emph{Tümevarımsal İnşalar Hesabını} (CIC) kurdular \citep{paulin1993}.
CIC, Coq'un (bugün Rocq) temelidir. \Lean{}'in tip kuramı da CIC'nin bir
çeşitlemesidir \citep{demoura2021lean4,carneiro2019}.
\end{tarihce}

% ---------------------------------------------------------------------
\section{İnşalar Hesabı ve evrenler}\label{sec:evrenler}

\subsection{$\lambda$-küpü}

Basit tipli $\lambda$-hesabından \Lean{}'e giden yolu, Barendregt'in
$\lambda$-küpü ile görselleştirebiliriz \citep{barendregt1991}. Basit
tiplerde yalnızca \emph{terimlere bağlı terimler} vardır (sıradan
fonksiyonlar). Küpün üç ekseni üç yeni bağımlılık ekler
(\cref{fig:lambda-kup}):

\begin{itemize}
  \item \textbf{Tiplere bağlı terimler} (çok biçimlilik): \lc{id : (α :
        Type) → α → α}. Sistem F \citep{girard1989}.
  \item \textbf{Tiplere bağlı tipler} (tip oluşturucular): \lc{List :
        Type → Type}.
  \item \textbf{Terimlere bağlı tipler} (bağımlı tipler): \lc{Vek α :
        Nat → Type}. $\lambda P$, mantıkta yüklemler.
\end{itemize}

\begin{figure}[htbp]
\centering
\begin{tikzpicture}[
    kose/.style={draw=#1, fill=#1!10, thick, rounded corners=3pt,
                 font=\small, inner sep=4pt, minimum width=1.3cm},
    kenar/.style={-{Stealth[length=2mm]}, thick, black!45},
    x={(3.6cm,0cm)}, y={(0cm,3.0cm)}, z={(1.7cm,1.2cm)},
  ]
  \node[kose=anarenk]     (st)  at (0,0,0) {$\lambda{\to}$};
  \node[kose=anarenk]     (f)   at (0,1,0) {$\lambda 2$ (F)};
  \node[kose=anarenk]     (w)   at (0,0,1) {$\lambda\underline{\omega}$};
  \node[kose=anarenk]     (fw)  at (0,1,1) {$\lambda\omega$};
  \node[kose=ikincil]     (p)   at (1,0,0) {$\lambda P$};
  \node[kose=ikincil]     (p2)  at (1,1,0) {$\lambda P2$};
  \node[kose=ikincil]     (pw)  at (1,0,1) {$\lambda P\underline{\omega}$};
  \node[kose=teoremframe] (coc) at (1,1,1) {\textbf{CoC}};
  \draw[kenar] (st) -- (f);  \draw[kenar] (st) -- (w);  \draw[kenar] (st) -- (p);
  \draw[kenar] (f) -- (fw);  \draw[kenar] (f) -- (p2);
  \draw[kenar] (w) -- (fw);  \draw[kenar] (w) -- (pw);
  \draw[kenar] (p) -- (p2);  \draw[kenar] (p) -- (pw);
  \draw[kenar] (fw) -- (coc); \draw[kenar] (p2) -- (coc); \draw[kenar] (pw) -- (coc);
  \node[font=\baslikfont\scriptsize, text=black!60, anchor=north] at (0.5,-0.12,0) {terimlere bağlı tipler $\to$};
  \node[font=\baslikfont\scriptsize, text=black!60, anchor=south, rotate=90] at (-0.2,0.5,0) {çok biçimlilik $\rightarrow$};
  \node[font=\baslikfont\scriptsize, text=black!60, anchor=north west] at (0.02,-0.02,0.5) {tip oluşturucular};
  \node[font=\baslikfont\scriptsize, text=teoremframe, anchor=west, align=left] at (1.15,1.1,1)
       {+ tümevarımsal tipler\\+ evrenler\\$\Rightarrow$ CIC, \Lean{}};
\end{tikzpicture}
\caption{Barendregt'in $\lambda$-küpü. Sol alt köşe Bölüm~14'ün basit tipli
$\lambda$-hesabıdır. Her eksen bir bağımlılık türü ekler. Sağ üst köşe,
Coquand ve Huet'nin İnşalar Hesabıdır. \Lean{}'in tip kuramı buna
tümevarımsal tipleri ve bir evren hiyerarşisini ekler.}
\label{fig:lambda-kup}
\end{figure}

\subsection{Evrenler}

\lc{Nat}'ın tipi \lc{Type}'tır. Peki \lc{Type}'ın tipi nedir?
\lc{Type : Type} olamaz; Girard paradoksu bunu tutarsız kılar. \Lean{}'in
çözümü sonsuz bir evren hiyerarşisidir (\cref{fig:evrenler}):

\leandosya[Evrenler]{Bolum15}{evrenler}

\begin{figure}[htbp]
\centering
\begin{tikzpicture}[
    evren/.style={draw=#1, fill=#1!8, thick, rounded corners=4pt,
                  minimum width=2.6cm, minimum height=0.8cm, font=\small},
    ok/.style={-{Stealth[length=2mm]}, thick, black!40},
    ornek/.style={font=\scriptsize, text=black!60, anchor=west},
  ]
  \node[evren=teoremframe] (p)  at (0,0)   {\lc{Prop} = \lc{Sort 0}};
  \node[evren=anarenk]     (t0) at (0,1.4) {\lc{Type} = \lc{Sort 1}};
  \node[evren=anarenk]     (t1) at (0,2.8) {\lc{Type 1} = \lc{Sort 2}};
  \node[font=\large, text=black!40] (dd) at (0,3.9) {$\vdots$};
  \draw[ok] (p) -- node[right, font=\scriptsize]{:} (t0);
  \draw[ok] (t0) -- node[right, font=\scriptsize]{:} (t1);
  \draw[ok] (t1) -- (dd);
  \node[ornek] at (1.7,0)   {\lc{2 + 2 = 4}, \lc{∀ p : Prop, p → p}, \dots \quad (önermeler)};
  \node[ornek] at (1.7,1.4) {\lc{Nat}, \lc{List Nat}, \lc{Vek Nat 3}, \lc{Prop}, \dots};
  \node[ornek] at (1.7,2.8) {\lc{Type}, \lc{ChurchSayi}, \lc{(α : Type) → α → α}, \dots};
\end{tikzpicture}
\caption{\Lean{}'in evren hiyerarşisi. Her evren bir üsttekinin elemanıdır.
\lc{Prop} en alttadır ve \emph{impredikatiftir}: \lc{Prop}'lar üzerinde
niceleme yapan bir önerme yine \lc{Prop}'tadır.}
\label{fig:evrenler}
\end{figure}

\noindent İki niceleme arasındaki fark incedir ama önemlidir.
\lc{(α : Type) → α → α} tipi \emph{bütün} tipler üzerinde nicelediği için
\lc{Type}'ın içinde olamaz; bir üst evrene, \lc{Type 1}'e çıkar. Buna
\emph{predikatif} niceleme denir. \lc{∀ p : Prop, p → p} ise \emph{bütün}
önermeler üzerinde niceler, kendisi de bunlardan biri olmasına rağmen
\lc{Prop}'ta kalır. Bu \emph{impredikatif} niceleme, \ciltiki{5}'te göreceğimiz
gibi klasik matematiğin büyük bir kısmını \Lean{}'de ifade etmeyi mümkün
kılar. Tutarlılığı ise \lc{Prop}'tan büyük elemenin kısıtlanmasına bağlıdır
(\cref{sec:sigma}).

\noindent Evren çok biçimliliği (\leankw{universe} \lc{u}) sayesinde
\lc{List} ya da \lc{Vek} gibi tanımlar her evrende bir kez yazılır.
\lc{\#check @List.map}'in çıktısındaki \lc{u\_1}, \lc{u\_2} bu evren
değişkenleridir.

% ---------------------------------------------------------------------
\section{Kanıt-ilgisizliği ve tanımsal eşitlik}\label{sec:tanimsal}

\Lean{}'de \lc{Prop}'un elemanları için \emph{tanımsal kanıt-ilgisizliği}
geçerlidir: Aynı önermenin herhangi iki kanıtı tanımsal olarak eşittir.

\leandosya[Kanıt-ilgisizliği]{Bolum15}{kanit_ilgisiz}

\noindent Bu, Bölüm~9'da \lc{Prop} ile \lc{Bool}'u karşılaştırırken
vurguladığımız ayrımın ta kendisidir. Bir önermede önemli olan
\emph{doğru olup olmadığıdır}, nasıl kanıtlandığı değil. Pratikte bunun
büyük bir yararı vardır: \lc{\{n : Nat // n > 0\}} ya da \lc{Fin n} gibi
tiplerin elemanları yalnızca değerleriyle belirlenir. Kanıt bileşenleri
eşitlik ispatlarında hiç dert olmaz.

\subsection{İki tür eşitlik}

Kitap boyunca iki farklı eşitlik kullandık. Şimdi aralarındaki farkı
kesin olarak söyleyebiliriz:

\begin{table}[htbp]
\centering
\small
\renewcommand{\arraystretch}{1.5}
\rowcolors{2}{tablosatir}{white}
\begin{tabularx}{\linewidth}{>{\raggedright\arraybackslash}p{2.6cm} X X}
\rowcolor{tabloust}
\tablobaslik{} & \tablobaslik{Tanımsal eşitlik ($\equiv$)} & \tablobaslik{Önermesel eşitlik ($=$)} \\
nedir?          & bir \emph{yargı}: iki terim hesaplanınca aynı olur & bir \emph{önerme}: \lc{a = b} bir tiptir \\
kim karar verir? & tip denetleyici, otomatik olarak & biz, bir ispatla \\
örnek           & \lc{n + 0 ≡ n}, \lc{(fun x => x + 1) 2 ≡ 3} & \lc{0 + n = n} (tümevarım gerekir) \\
\Lean{}'de      & \leantac{rfl}, \leantac{show}, \leantac{change}, \leantac{decide} & \leantac{rw}, \leantac{simp}, \lc{▸}, \leantac{calc} \\
karar verilebilir mi? & ilke olarak evet (tip denetimi buna dayanır) & genel olarak hayır \\
\end{tabularx}
\caption{Tanımsal ve önermesel eşitlik.}
\label{tab:esitlikler}
\end{table}

\leandosya[Tanımsal ve önermesel eşitlik]{Bolum15}{tanimsal}

\noindent Tanımsal eşitlik ``bedavadır'': \Lean{} onu görmek için hiçbir
kanıt istemez. Ama onu kullanan bir tanım, tanımın \emph{nasıl} yazıldığına
bağımlı hâle gelir. \lc{n + 0 = n}'nin \leantac{rfl} ile, \lc{0 + n = n}'nin
ise yalnızca tümevarımla ispatlanabilmesinin nedeni, \lc{Nat.add}'ın
ikinci argümanı üzerinde özyinelemeli tanımlanmış olmasıdır.

\begin{dikkat}[Dikkat: Tanımsal eşitliğin sınırları]
Kuramsal olarak \Lean{}'in tanımsal eşitliği karar verilebilir
\emph{değildir}: Kanıt-ilgisizliği ile alt tekil eleme
(\emph{subsingleton elimination}) birleşince, tanımsal eşitliğin geçişli
olmadığı yapay örnekler kurulabilir \citep{carneiro2019}. Uygulamada bu
örneklere neredeyse hiç rastlanmaz. \Lean{}'in çekirdeği tanımsal eşitliği
bir algoritmayla denetler; algoritma bazı durumlarda ``eşit değil''
diyebilir, ama asla yanlış yere ``eşit'' demez. Tutarlılık için önemli
olan da budur.
\end{dikkat}

% ---------------------------------------------------------------------
\section{Lean'in çekirdeği}\label{sec:cekirdek}

Bir ispat asistanına neden güvenelim? Asistanın kendisi binlerce satırlık
bir programdır ve hatalar içerebilir. Nicolaas de~Bruijn'in Automath
projesinden \citep{debruijn1970} kalan cevap, \emph{de~Bruijn ölçütü}
olarak bilinir:

\begin{onemlifikir}[Önemli Fikir: de Bruijn ölçütü]
Bir ispat asistanı, ürettiği her ispatı küçük ve bağımsız olarak
denetlenebilir bir \emph{çekirdeğe} doğrulatmalıdır. Taktikler, otomasyon,
notasyon ve tip sınıfı çözümlemesi ne kadar karmaşık olursa olsun, sonuçta
çekirdeğin anlayacağı basit bir terim üretilir. Güvenmemiz gereken tek
şey, bu terimi denetleyen çekirdektir.
\end{onemlifikir}

\begin{figure}[htbp]
\centering
\begin{tikzpicture}[
    kat/.style={draw=#1, fill=#1!8, thick, rounded corners=5pt,
                text width=3.0cm, align=center, inner sep=6pt, font=\small,
                minimum height=1.6cm},
    ok/.style={-{Stealth[length=2.6mm]}, very thick, black!40},
  ]
  \node[kat=arastirmaframe] (k) at (0,0)   {\textbf{Kaynak kod}\\{\scriptsize \leantac{simp}, \leantac{omega}, notasyon, \lc{x + 1}}};
  \node[kat=anarenk]        (e) at (4.4,0) {\textbf{Elaboratör}\\{\scriptsize taktikler, tip sınıfları, örtük argümanlar}};
  \node[kat=teoremframe]    (c) at (8.8,0) {\textbf{Çekirdek}\\{\scriptsize CIC terimleri, tip denetimi, 3 aksiyom}};
  \draw[ok] (k) -- (e);
  \draw[ok] (e) -- node[above, font=\baslikfont\scriptsize]{terim} (c);
  \node[font=\baslikfont\scriptsize, text=hataframe, anchor=north] at (4.4,-1.05) {güvenilmez (hata yapabilir)};
  \node[font=\baslikfont\scriptsize, text=bittiframe, anchor=north] at (8.8,-1.05) {güvenilir (küçük, denetlenebilir)};
  \node[draw=bittiframe, dashed, thick, rounded corners=4pt, fit=(c), inner sep=8pt] {};
  \node[font=\baslikfont\scriptsize, text=black!55, align=center] at (8.8,1.45) {bağımsız denetleyiciler:\\Lean4Lean, nanoda, \dots};
\end{tikzpicture}
\caption{\Lean{}'in mimarisi. Elaboratörün hataları en kötü ihtimalle
geçerli bir ispatın reddedilmesine yol açar. Geçersiz bir ispatın kabul
edilmesi için çekirdekte bir hata olması gerekir.}
\label{fig:lean-mimari}
\end{figure}

\noindent Elaboratörün ne kadar iş yaptığını görmek için
\lc{pp.all} seçeneğiyle bir terimin çekirdekteki hâline bakabiliriz:

\leandosya[Elaborasyon]{Bolum15}{elaborasyon}

\noindent Masum görünen \lc{x + 1}, aslında üç tip sınıfı örneği ve bir
sayı sabiti kodlaması içerir. Bir taktik ispatı ise çok daha büyük bir
terime dönüşür. Çekirdeğin gördüğü her tümevarımsal tip, kendi
\emph{özyineleyicisi} ile birlikte gelir:

\leandosya[Özyineleyiciler]{Bolum15}{ozyineleyici}

\noindent \lc{Vek.rec}'in tipi Bölüm~9'daki tümevarım ilkelerinin bağımlı
hâlidir. \emph{Motive}, hem indise (\lc{n}) hem de vektörün kendisine
bağlıdır. \leankw{match} ile yazdığımız her tanım ve \leantac{induction}
ile yaptığımız her ispat, çekirdekte bu özyineleyiciye çevrilir.

\subsection{Aksiyomlar}

\Lean{}'in çekirdeği CIC'ye yalnızca üç aksiyom ekler:

\begin{center}
\renewcommand{\arraystretch}{1.45}
\rowcolors{2}{tablosatir}{white}
\begin{tabular}{>{\raggedright\arraybackslash}p{3.4cm}
                >{\raggedright\arraybackslash}p{5.2cm}
                >{\raggedright\arraybackslash}p{3.6cm}}
\rowcolor{tabloust}
\tablobaslik{Aksiyom} & \tablobaslik{Söylediği} & \tablobaslik{Nerede?} \\
\lc{propext}          & denk önermeler eşittir: \lc{(p ↔ q) → p = q} & \leantac{simp}, \leantac{rw} ile ↔ \\
\lc{Quot.sound}       & bölüm tiplerinde ilişkili elemanlar eşittir & fonksiyon dışsallığı (\lc{funext}) \\
\lc{Classical.choice} & boş olmayan her tipten bir eleman seçilebilir & üçüncü hâl, \lc{Classical} \\
\end{tabular}
\end{center}

\noindent Bir teoremin hangi aksiyomlara dayandığını \leankw{\#print}
\lc{axioms} ile öğrenebiliriz:

\leandosya[Aksiyomları denetlemek]{Bolum15}{aksiyomlar}

\noindent Vektör teoremimiz hiçbir aksiyoma dayanmıyor: Tamamen yapıcıdır.
Üçüncü hâlin olanaksızlığı ise üç aksiyomun üçüne de dayanır. Bölüm~4'te
yapıcı ile klasik arasında çizdiğimiz çizgi, \Lean{}'de tam olarak bu
komutla denetlenebilir. Diaconescu teoremi \citep{diaconescu1975},
seçim aksiyomunun \lc{propext} ve \lc{funext} ile birlikte üçüncü hâli
gerektirdiğini söyler. Bu yüzden \lc{Classical.em} bir aksiyom değil, bir
\emph{teoremdir}. Aksiyomları ve onların kuramsal sonuçlarını Bölüm~16'da
ayrıntılı olarak ele alacağız.

% ---------------------------------------------------------------------
\begin{arastirmanotu}
\textbf{Tip kuramının çeşitleri.} \Lean{}, kanıt-ilgisizliği ve
\lc{propext} sayesinde ``klasik matematik dostu'' bir tasarım seçmiştir.
Başka seçimler de mümkündür. \emph{Homotopi tip kuramı} \citep{hott2013},
kanıt-ilgisizliği yerine eşitlik kanıtlarına yollar olarak zengin bir yapı
verir ve \emph{tekdeğerlilik} (\emph{univalence}) aksiyomunu ekler:
Denk tipler eşittir. Kübik tip kuramı bu aksiyomu hesaplanabilir kılar ve
Agda'nın kübik kipinde kullanılır. \emph{Gözlemsel tip kuramı}
(\emph{observational type theory}) ise fonksiyon dışsallığını aksiyom
olmadan tanımsal olarak sağlamayı hedefler.

\textbf{Çekirdeğin doğrulanması.} Mario Carneiro'nun yüksek lisans tezi
\citep{carneiro2019}, \Lean{}'in tip kuramının ilk kesin tanımını verdi ve
tutarlılığını küme kuramında bir modelle gösterdi. Lean4Lean
\citep{carneiro2026}, bu çalışmayı \Lean{}'in içine taşır: \Lean{}'in
çekirdeğinin bir kopyası \Lean{}'de yazılmıştır ve özelliklerinin
ispatları da \Lean{}'dedir. Gödel'in ikinci eksiklik teoremi (\ciltiki{4})
nedeniyle bir sistemin kendi tutarlılığını ispatlayamayacağını
hatırlatalım. Bu yüzden bu tür çalışmalar tutarlılığı ancak daha güçlü bir
sistemde (örneğin ek bir evrenle) ifade edebilir.
\end{arastirmanotu}

\begin{bolumozeti}
\begin{itemize}
  \item $Π$-tipleri bağımlı fonksiyonlardır. $\forall$, dönüş tipi
        \lc{Prop} olan bir $Π$-tipidir. $A \to B$ bağımsız özel
        durumdur.
  \item $Σ$-tipleri bağımlı çiftlerdir. \Lean{}'de üç çeşidi vardır:
        \lc{Sigma} (veri), alt tip (veri + kanıt) ve \lc{Exists} (önerme).
        \lc{Exists}'ten tanık bir program içinde çıkarılamaz.
  \item Tümevarımsal aileler (\lc{Vek α n}, \lc{Fin n}) bir özelliği tipe
        yazarak çalışma zamanı hatalarını tip hatalarına dönüştürür. Bunun
        bedeli indis aritmetiği ve eşitlik boyunca taşımadır.
  \item Martin-Löf tip kuramı her tipi oluşum, giriş, eleme ve hesaplama
        kurallarıyla tanıtır. CIC ve \Lean{}, bu çizginin
        devamıdır. $\lambda$-küpü bağımlılık türlerini düzenler.
  \item Evrenler bir hiyerarşi oluşturur. \lc{Prop} impredikatiftir,
        \lc{Type u}'lar predikatiftir.
  \item \lc{Prop}'ta kanıt-ilgisizliği tanımsaldır. Tanımsal eşitlik tip
        denetleyicinin yargısıdır, önermesel eşitlik ise ispat ister.
  \item de~Bruijn ölçütü: Güven yalnızca küçük çekirdeğe dayanır.
        \Lean{}'in üç aksiyomu vardır ve \leankw{\#print} \lc{axioms} bir
        teoremin hangilerine dayandığını gösterir.
\end{itemize}
\end{bolumozeti}

\begin{alistirmalar}
\alistirma{1} \lc{tekrarlaV a n}'yi \lc{f} ile eşlemek, \lc{tekrarlaV (f a)
n}'yi verir. (\lc{a15\_1})

\alistirma{1} \lc{son (v.sonaEkle a) = a}. \lc{son}'un tanımında neden
uzunluk üzerinde örüntü eşleme yapılıyor? (\lc{a15\_2})

\alistirma{2} Fermuarın birinci bileşenleri ilk vektörü verir. (\lc{a15\_3})

\alistirma{2} \lc{birlestir} ile \lc{toList} uyumludur. (\lc{a15\_4})

\alistirma{2} Eşitlik boyunca taşımak, vektörün listesini değiştirmez.
\leantac{subst}'ün burada ne yaptığını açıklayın. (\lc{a15\_5})

\alistirma{1} \lc{Fin}'de kanıt-ilgisizliği. (\lc{a15\_6})

\alistirma{2} Her liste, uzunluğu kadar bir vektöre dönüşür. Sonuç neden bir
$Σ$-tipi? (\lc{a15\_7})

\alistirma{3} Sabit bir vektörün her elemanı aynıdır. (\lc{a15\_8})

\alistirma{2} (Kâğıt-kalem) \lc{(α : Type) → List α → Nat} ve
\lc{∀ α : Type, Nonempty α → True} tiplerinin hangi evrende yaşadığını
belirleyin ve \leankw{\#check} ile doğrulayın.

\alistirma{2} (Kâğıt-kalem) \lc{birlestir}'in sonuç tipini \lc{Vek α (m +
n)} olarak yazıp \Lean{}'in hata mesajını okuyun. Tanımı
\lc{Nat.add\_comm} ve \lc{▸} ile düzeltin.

\alistirma{3} (Proje) \lc{Vek} için bir \lc{ters} fonksiyonunu
\lc{sonaEkle} kullanmadan, bir biriktirici ile tanımlayın:
\lc{tersAux : Vek α m → Vek α n → Vek α (m + n)}. Hangi indis eşitliğine
ihtiyaç duyuyorsunuz?

\alistirma{3} (Araştırma) Bölüm~14'teki basit tipli $\lambda$-hesabına
bağımlı fonksiyon tipleri ekleyerek $\lambda P$'yi tanımlayın. Tipler artık
terimler içerebildiği için \lc{Tip} ile \lc{Terim} neden aynı tümevarımsal
tip içinde tanımlanmak zorundadır?
\end{alistirmalar}
   % Bağımlı tipler ve Lean'in çekirdeği
% =====================================================================
%  Bölüm 16 — Lean'in Aksiyomları
% =====================================================================
\chapter{Lean'in Aksiyomları}\label{bol:aksiyomlar}

\begin{bolumhedefleri}
\begin{itemize}
  \item \Lean{}'in üç aksiyomunu (\lc{propext}, \lc{Quot.sound},
        \lc{Classical.choice}) ve her birinin ne işe yaradığını,
  \item Bölüm tipleriyle yeni matematiksel nesneler kurmayı (tamsayılar),
  \item Fonksiyon dışsallığının \lc{Quot.sound}'dan nasıl türetildiğini,
  \item Seçim aksiyomunu, hesaplanamayan tanımları ve \leankw{noncomputable}
        anahtar sözcüğünü,
  \item Diaconescu teoremini: seçimin üçüncü hâli gerektirdiğini,
  \item \Lean{}'in aksiyomlarını homotopi tip kuramının tek-değerlilik
        aksiyomuyla karşılaştırmayı.
\end{itemize}
\end{bolumhedefleri}

Bölüm~15'in sonunda \Lean{}'in çekirdeğinin, İnşalar Hesabına yalnızca
üç aksiyom eklediğini söyledik. Bu bölümde o üç aksiyomu tek tek ele
alıyoruz. Her biri bir ihtiyaçtan doğmuştur: \lc{propext} önermeleri
``doğruluk değerleri'' gibi kullanabilmek için, \lc{Quot.sound} bölüm
yapıları kurabilmek için, \lc{Classical.choice} ise klasik matematiği
yapabilmek için. Bölümün ana teoremi, bu üç aksiyomun birlikte üçüncü
hâlin olanaksızlığını \emph{gerektirdiğini} söyleyen Diaconescu
teoremidir. Bölüm~4'te ``\Lean{} klasik mantığa izin verir'' diye
başlattığımız öykünün sonu budur.

\begin{figure}[htbp]
\centering
\begin{tikzpicture}[
    aks/.style={draw=teoremframe, fill=teorembg, very thick, rounded corners=4pt,
                font=\small\ttfamily, inner sep=5pt, minimum width=3.0cm},
    sonuc/.style={draw=#1, fill=#1!8, thick, rounded corners=4pt,
                  font=\small, inner sep=4pt, align=center},
    ok/.style={-{Stealth[length=2.2mm]}, thick, black!45},
  ]
  \node[aks] (pe) at (0,0)   {propext};
  \node[aks] (qs) at (4.6,0) {Quot.sound};
  \node[aks] (ch) at (9.2,0) {Classical.choice};
  \node[sonuc=anarenk]  (iff) at (0,-1.8)   {\lc{↔} ile \leantac{rw},\\\leantac{simp}};
  \node[sonuc=anarenk]  (bol) at (3.2,-1.8) {bölüm tipleri\\(\lc{Int}, \lc{ℚ}, \lc{ℝ}, \dots)};
  \node[sonuc=anarenk]  (fe)  at (6.0,-1.8) {\lc{funext}};
  \node[sonuc=anarenk]  (nc)  at (9.2,-1.8) {\leankw{noncomputable}\\tanımlar};
  \node[sonuc=dikkatframe] (em) at (4.6,-3.8) {\textbf{\lc{Classical.em}}\\{\scriptsize Diaconescu teoremi}};
  \node[sonuc=ikincil] (kl) at (4.6,-5.4) {klasik matematik: \leantac{by\_contra}, \leantac{by\_cases}, \dots};
  \draw[ok] (pe) -- (iff); \draw[ok] (qs) -- (bol); \draw[ok] (qs) -- (fe);
  \draw[ok] (ch) -- (nc);
  \draw[ok, dikkatframe] (pe) -- (em); \draw[ok, dikkatframe] (fe) -- (em);
  \draw[ok, dikkatframe] (ch) -- (em);
  \draw[ok] (em) -- (kl);
\end{tikzpicture}
\caption{\Lean{}'in üç aksiyomu ve en önemli sonuçları. Üçüncü hâlin
olanaksızlığı bir aksiyom değil, üç aksiyomun birlikte verdiği bir
teoremdir.}
\label{fig:aksiyom-haritasi}
\end{figure}

% ---------------------------------------------------------------------
\section{Önerme dışsallığı}\label{sec:propext}

Matematikte ``$A \cap B = B \cap A$'' derken iki kümenin \emph{aynı
elemanlara sahip olduğu için} eşit olduğunu söyleriz. Buna
\emph{dışsallık} denir: Bir nesne, ``dışarıdan'' gözlemlenebilen
davranışıyla belirlenir. \emph{Önerme dışsallığı} aynı ilkeyi önermelere
uygular: Birbirine denk iki önerme eşittir.

\leandosya[propext]{Bolum16}{propext}

\noindent İkinci örnek önemli bir gözlem içeriyor: \lc{h : p ↔ q} ile
\leantac{rw} yapmak \lc{propext} gerektirir. \leantac{rw} yalnızca
\emph{eşitliklerle} yeniden yazabilir. Bir denkliği eşitliğe çeviren şey
\lc{propext}'tir. \leantac{simp} de aynı nedenle \lc{propext}'i sık
kullanır. Bu yüzden Mathlib'deki teoremlerin büyük çoğunluğu bu aksiyoma
dayanır.

\begin{onemlifikir}[Önemli Fikir: Prop, doğruluk değerlerinin tipi olur]
\lc{propext} olmadan \lc{Prop}'un elemanları, kanıtları birbirinden ayırt
edilemeyen ama \emph{kendileri} ayırt edilebilen ``tipler''dir.
\lc{propext} ile birlikte, denk önermeler özdeşleşir. Üçüncü hâli de
eklersek \lc{Prop} tam olarak iki elemanlı bir tip olur: Her önerme ya
\lc{True}'ya ya \lc{False}'a eşittir. Bölüm~10'daki iki değerli anlambilim,
bu durumda \lc{Prop}'un \emph{kendisidir}.
\end{onemlifikir}

\leandosya[Prop'un iki değeri]{Bolum16}{iki_deger}

% ---------------------------------------------------------------------
\section{Bölüm tipleri}\label{sec:bolum-tipleri}

Matematikte yeni nesneleri sık sık eski nesnelerin \emph{denklik
sınıfları} olarak kurarız. Rasyonel sayılar kesirlerin, tamsayılar ise
doğal sayı çiftlerinin denklik sınıflarıdır. Bölüm~8'de denklik bağıntılarını
ve denklik sınıflarını küme kuramı içinde görmüştük. Şimdi aynı yapıyı
\Lean{}'in çekirdek düzeyinde, bir \emph{tip} olarak kuruyoruz.

\begin{tanim}{Bölüm tipi}{bolum-tipi}
$\alpha$ bir tip, $r$ de $\alpha$ üzerinde bir bağıntı olsun. \textbf{Bölüm
tipi} \lc{Quot r}, $r$ ile ilişkili elemanların özdeşleştirildiği tiptir.
Dört ilkel işlemi vardır:
\begin{itemize}
  \item \lc{Quot.mk r a : Quot r} \hfill (bir elemanın sınıfı)
  \item \lc{Quot.lift f h : Quot r → β}, eğer \lc{h : ∀ a b, r a b → f a = f b} \hfill (bölümden dışarı)
  \item \lc{Quot.ind} \hfill (her eleman bir \lc{Quot.mk}'dir)
  \item \lc{Quot.sound : r a b → Quot.mk r a = Quot.mk r b} \hfill \textcolor{teoremframe}{\textbf{(aksiyom)}}
\end{itemize}
Ayrıca \lc{Quot.lift f h (Quot.mk r a) ≡ f a} bir \emph{hesaplama
kuralıdır}: Tanımsal olarak geçerlidir.
\end{tanim}

\noindent İlk üç işlem ve hesaplama kuralı çekirdeğe yerleşik olarak
gelir ve bir aksiyom sayılmaz. Bunlar yalnızca yeni bir tip ve onun
eleme kuralıdır, tümevarımsal tipler gibi. \lc{Quot.sound} ise yeni bir
\emph{eşitlik} ilan eder. Bu yüzden bir aksiyomdur.

\subsection{Tamsayıları sıfırdan kurmak}

Bir $(a, b)$ doğal sayı çiftini ``$a - b$'' sayısı olarak düşünelim.
$(3, 1)$ ile $(5, 3)$ aynı sayıyı, yani 2'yi temsil etmelidir. Genel
olarak $(a, b) \sim (c, d) \iff a + d = c + b$. Bu bağıntı bir denklik
bağıntısıdır ve bölümü tamsayılardır (\cref{fig:tamsayi-izgara}).

\begin{figure}[htbp]
\centering
\begin{tikzpicture}[scale=0.62]
  \foreach \k/\c in {-3/vurgu, -2/dikkatframe, -1/tarihframe, 0/black!60, 1/ikincil, 2/anarenk, 3/ornekframe} {
    \pgfmathsetmacro{\xs}{max(0,\k)}
    \pgfmathsetmacro{\ys}{max(0,-\k)}
    \pgfmathsetmacro{\xe}{min(6,6+\k)}
    \pgfmathsetmacro{\ye}{min(6,6-\k)}
    \draw[\c, very thick, opacity=0.55] (\xs,\ys) -- (\xe,\ye);
  }
  \foreach \x in {0,...,6} \foreach \y in {0,...,6} \fill[black!70] (\x,\y) circle (2.2pt);
  \draw[-{Stealth[length=2mm]}, thick] (-0.5,0) -- (6.8,0) node[right, font=\small] {$a$};
  \draw[-{Stealth[length=2mm]}, thick] (0,-0.5) -- (0,6.8) node[above, font=\small] {$b$};
  \foreach \k/\c/\l in {3/ornekframe/{+3}, 2/anarenk/{+2}, 1/ikincil/{+1}, 0/black!60/0} {
    \node[font=\scriptsize\bfseries, text=\c, anchor=west] at (6.25, 6-\k) {\l};
  }
  \foreach \k/\c/\l in {-1/tarihframe/{$-1$}, -2/dikkatframe/{$-2$}, -3/vurgu/{$-3$}} {
    \node[font=\scriptsize\bfseries, text=\c, anchor=south] at (6+\k, 6.25) {\l};
  }
  \fill[teoremframe] (3,1) circle (4pt); \fill[teoremframe] (5,3) circle (4pt);
  \node[font=\scriptsize, text=teoremframe, anchor=north west] at (3,1) {$(3,1)$};
  \node[font=\scriptsize, text=teoremframe, anchor=north west] at (5,3) {$(5,3)$};
\end{tikzpicture}
\caption{$\mathbb{N} \times \mathbb{N}$ ızgarasında tamsayılar. Her köşegen
bir denklik sınıfıdır: Aynı köşegendeki çiftler aynı tamsayıyı temsil eder.
$(3,1)$ ve $(5,3)$, $+2$ köşegenindedir.}
\label{fig:tamsayi-izgara}
\end{figure}

\leandosya[Tamsayılar bir bölüm tipi olarak]{Bolum16}{tamsayi}

\noindent \lc{Setoid}, bir bağıntı ile onun bir denklik bağıntısı olduğunun
kanıtını bir araya getiren bir yapıdır. \lc{Quotient}, \lc{Quot}'un
\lc{Setoid}'li hâlidir. Son satırdaki \lc{example} \lc{Quot.sound}'u
kullanıyor: $3 + 3 = 5 + 1$ olduğu için \lc{Tamsayi.mk 3 1} ile
\lc{Tamsayi.mk 5 3} \emph{eşittir}. Bu eşitlik bir hesaplamayla elde
edilemezdi. İki terim farklıdır ve aynı oldukları ancak aksiyomla
söylenebilir.

\subsection{Bölümden dışarı fonksiyonlar}

Bir tamsayı üzerinde fonksiyon tanımlamak için bir \emph{temsilci} üzerinde
tanımlayıp sonucun temsilciye bağlı olmadığını kanıtlamamız gerekir.
Bu kanıt \lc{Quotient.lift}'in ikinci argümanıdır. Matematikte ``iyi
tanımlılık'' denilen şeyin biçimsel hâlidir:

\leandosya[Tamsayılarda işlemler]{Bolum16}{tamsayi_islem}

\noindent Toplamada iki argümanın da temsilcisi değişebilir. Bu yüzden
\lc{Quotient.lift₂}'nin iyi tanımlılık kanıtı iki varsayım alır. Her
kanıt, \leantac{show} ile hedefi açık bir doğal sayı eşitliğine çevirip
\leantac{omega} ile biter. Bölümler üzerinde teorem ispatlamak için
\lc{Quotient.inductionOn} kullanırız. Bir tamsayı hakkındaki her önermeyi,
temsilcileri hakkındaki bir önermeye indirger:

\leandosya[Tamsayılarda ispat]{Bolum16}{tamsayi_teorem}

\begin{dikkat}[Dikkat: İyi tanımlılığı unutmak]
Kâğıt üzerinde ``$[(a,b)] + [(c,d)] = [(a+c, b+d)]$ olarak tanımlayalım''
deyip iyi tanımlılığı atlamak kolaydır. \Lean{} bunu atlatmaz:
\lc{Quotient.lift}'in ikinci argümanı verilmeden tanım kabul edilmez.
Ünlü bir örnek: Kesirler üzerinde $\frac{a}{b} \mapsto a$ ``fonksiyonu''
iyi tanımlı değildir ($\frac12 = \frac24$, ama $1 \neq 2$). Bu fonksiyonu
\lc{Quotient.lift} ile tanımlamaya çalışan herkes, ispatlanamaz bir hedefle
karşılaşır.
\end{dikkat}

% ---------------------------------------------------------------------
\section{Fonksiyon dışsallığı}\label{sec:funext}

İki fonksiyon her girdide aynı çıktıyı veriyorsa eşit midir? Matematikte
bu açıktır, çünkü küme kuramında bir fonksiyon, girdi--çıktı çiftlerinin
kümesidir. Tip kuramında ise fonksiyonlar \emph{programlardır}:
\lc{fun n => 0 + n} ile \lc{fun n => n} farklı programlardır ve
tanımsal olarak eşit değildir. Martin-Löf'ün yoğunlamalı (\emph{intensional})
tip kuramında fonksiyon dışsallığı ispatlanamaz.

\Lean{}'de ise \lc{funext} bir teoremdir. Şaşırtıcı biçimde, ispatı
\lc{Quot.sound}'u kullanır. Fikir şudur: Fonksiyonları ``noktasal eşitlik''
bağıntısına göre bölelim. Bu bölümün elemanlarını bir noktada
değerlendirmek iyi tanımlıdır. Noktasal eşit iki fonksiyonun bölümdeki
sınıfları \lc{Quot.sound} ile eşittir. Bu eşitliğe \lc{uygula}'yı
uyguladığımızda, iki fonksiyonun kendisinin eşitliğini elde ederiz:

\leandosya[Quot.sound'dan funext]{Bolum16}{funext}

\noindent İspatın kilit adımı \leantac{show} satırıdır.
\lc{uygula (Quot.mk NoktaEsit f)}, \lc{fun x => f x}'e, o da $\eta$-kuralı
ile \lc{f}'ye tanımsal olarak eşittir. Yani hedefi ``iki bölüm elemanına
aynı fonksiyonu uygulamak'' biçiminde yeniden yazabiliriz. Gerisi
\lc{congrArg} ve \lc{Quot.sound}'dur. \leankw{\#print} \lc{axioms}, bu
ispatın yalnızca \lc{Quot.sound}'a dayandığını doğrular.

\leandosya[funext kullanmak]{Bolum16}{funext_ornek}

% ---------------------------------------------------------------------
\section{Seçim aksiyomu}\label{sec:secim}

\lc{Classical.choice}, \Lean{}'in en güçlü ve en ``klasik'' aksiyomudur:
\[
  \texttt{Classical.choice} : \texttt{Nonempty}\;\alpha \to \alpha .
\]
\lc{Nonempty α}, \lc{α}'nın boş olmadığını söyleyen bir \emph{önermedir}.
Aksiyom, bu önermeden \lc{α}'nın bir \emph{elemanını} üretir. Bölüm~15'te
\lc{Prop}'tan \lc{Type}'a ``büyük elemenin'' yasak olduğunu, bu yüzden bir
\lc{∃} kanıtından tanık çıkaramayacağımızı görmüştük. Seçim aksiyomu bu
yasağı \emph{tek bir noktada}, bir aksiyomla deler:

\leandosya[Seçim aksiyomu]{Bolum16}{secim}

\noindent \lc{secilen} bir fonksiyondur ama \leankw{noncomputable} olarak
işaretlenmiştir. \Lean{}'in derleyicisi \lc{Classical.choice} için kod
üretemez, çünkü aksiyom hangi elemanın seçildiğini söylemez. Seçilen eleman
hakkında bildiğimiz tek şey, istenen özelliği sağladığıdır
(\lc{secilen\_ozellik}). \lc{ucten\_buyuk} 4 müdür, 5 midir, 1000 midir?
Bilemeyiz ve bilmemize gerek yoktur.

\begin{dikkat}[Dikkat: Kanıt ile hesaplama]
\leankw{noncomputable} bir tanım, mantıksal olarak tamamen geçerlidir:
Teoremlerde serbestçe kullanılabilir. Yalnızca \leankw{\#eval} ile
çalıştırılamaz. Mathlib'de gerçel sayılar, ölçüler ve birçok soyut yapı
bu türdendir. Hesaplama gerektiğinde (örneğin bir karar yöntemi
yazarken) seçimden kaçınmak gerekir. \leankw{\#print} \lc{axioms} burada
da yol göstericidir.
\end{dikkat}

% ---------------------------------------------------------------------
\section{Diaconescu teoremi}\label{sec:diaconescu}

\begin{teorem}{Diaconescu (1975)}{diaconescu}
Seçim aksiyomu, önerme dışsallığı ve fonksiyon dışsallığı birlikte
üçüncü hâlin olanaksızlığını gerektirir: Her $p$ için $p \lor \lnot p$.
\end{teorem}

\noindent Radu Diaconescu bu sonucu topos kuramında ispatladı
\citep{diaconescu1975}. Goodman ve Myhill aynı fikri yapıcı küme kuramına
taşıdılar \citep{goodman1978}. İspatın fikri zariftir
(\cref{fig:diaconescu}). Bir $p$ önermesi verilsin. İki yüklem tanımlayalım:
\[
  U(x) \;:=\; (x = \mathsf{True}) \lor p,
  \qquad
  V(x) \;:=\; (x = \mathsf{False}) \lor p .
\]
İkisi de boş değildir: $U(\mathsf{True})$ ve $V(\mathsf{False})$
doğrudur. Seçim aksiyomuyla birer tanık seçelim: $u$ ve $v$.

\begin{itemize}
  \item $u$'nun seçildiği yerde ya $u = \mathsf{True}$ ya da $p$ doğrudur;
        $v$ için de ya $v = \mathsf{False}$ ya da $p$. Dolayısıyla ya
        $p$ doğrudur ya da $u = \mathsf{True}$ ve $v = \mathsf{False}$,
        yani $u \neq v$.
  \item Öte yandan $p$ doğruysa, $U$ ile $V$ her $x$ için doğrudur. propext
        ile her noktada eşittirler, funext ile $U = V$'dir. Aynı yüklemden
        yapılan seçim aynı sonucu verir: $u = v$.
\end{itemize}

\noindent Demek ki $u \neq v$ ise $p$ yanlıştır. Böylece her durumda ya $p$
ya da $\lnot p$ elde edilir. Seçim, burada bir ``karar verici'' gibi
çalışır: $u$ ile $v$'nin eşit olup olmadığına bakarak $p$ hakkında karar
verir.

\begin{figure}[htbp]
\centering
\begin{tikzpicture}[
    kutu/.style={draw=#1, fill=#1!8, thick, rounded corners=5pt,
                 text width=4.6cm, align=center, inner sep=6pt, font=\small},
  ]
  \node[kutu=ikincil] (a) at (0,0)
     {\textbf{$p$ doğru}\\[3pt] $U(x)$ ve $V(x)$ her $x$ için doğru\\
      $\Downarrow$ propext + funext\\ $U = V$\\ $\Downarrow$\\ seçimler aynı: $u = v$};
  \node[kutu=hataframe] (b) at (7,0)
     {\textbf{$p$ yanlış}\\[3pt] $U(x) \iff x = \mathsf{True}$\\ $V(x) \iff x = \mathsf{False}$\\
      $\Downarrow$\\ $u = \mathsf{True}$, $v = \mathsf{False}$\\ seçimler farklı: $u \neq v$};
  \node[draw=dikkatframe, fill=dikkatbg, very thick, rounded corners=5pt,
        font=\small, align=center, inner sep=6pt] (c) at (3.5,-3.3)
     {$u = v$ mi?\quad bu soru $p$'yi belirler\\
      {\scriptsize (eşitlik sorusu, seçilmiş iki \emph{önerme} hakkındadır)}};
  \draw[-{Stealth[length=2.6mm]}, very thick, black!40] (a.south) -- (c);
  \draw[-{Stealth[length=2.6mm]}, very thick, black!40] (b.south) -- (c);
\end{tikzpicture}
\caption{Diaconescu teoreminin fikri. Seçim aksiyomu, $p$'nin doğru olup
olmadığını $u$ ile $v$'nin eşitliği sorusuna dönüştürür.}
\label{fig:diaconescu}
\end{figure}

\leandosya[Diaconescu teoremi]{Bolum16}{diaconescu}

\noindent \Lean{} ispatı, kâğıt üzerindeki ispatın satır satır karşılığıdır.
\lc{h0} adımı dikkat çekicidir: \lc{U = V} eşitliğiyle yeniden yazdıktan
sonra iki seçim, \emph{aynı} yüklemin iki varlık kanıtı üzerinden yapılır.
Kanıt-ilgisizliği sayesinde bu iki kanıt da aynıdır, dolayısıyla seçimler
\leantac{rfl} ile eşittir. Kanıt-ilgisizliğinin bu ispattaki gizli rolü,
yapıcı tip kuramlarında Diaconescu teoreminin neden ``kaçınılabildiğini''
de açıklar.

\begin{onemlifikir}[Önemli Fikir: Lean'de klasik mantık bir teoremdir]
Bölüm~4'te \lc{Classical.em}'i bir ``kapı'' olarak tanıtmıştık. Şimdi
kapının arkasını görüyoruz: \lc{Classical.em}, \Lean{}'in standart
kütüphanesinde \emph{tam olarak bu ispatla} kanıtlanmış bir teoremdir. Ayrı
bir aksiyom değildir. Bu yüzden ``\Lean{} klasik midir?'' sorusunun cevabı
``evet, çünkü seçim aksiyomunu kabul eder'' olur. Seçimi kullanmayan bir
teorem, \leankw{\#print} \lc{axioms} çıktısında \lc{Classical.choice}
görünmüyorsa, yapıcıdır.
\end{onemlifikir}

% ---------------------------------------------------------------------
\section{Aksiyomlar ve tutarlılık}\label{sec:aksiyom-tutarlilik}

Bir aksiyom eklemek tehlikelidir: Yanlış bir aksiyom her şeyi
kanıtlanabilir kılar. \Lean{}'de herkes \leankw{axiom} anahtar
sözcüğüyle yeni bir aksiyom ekleyebilir:

\begin{leanhata}[Kötü bir fikir (derlemeyin)]
axiom kotu : False

theorem bir_iki : 1 = 2 := kotu.elim     -- kabul edilir!

#print axioms bir_iki
-- 'bir_iki' depends on axioms: [kotu]
\end{leanhata}

\noindent Bu örnek, \leankw{\#print} \lc{axioms}'un neden önemli olduğunu
gösteriyor: Bir teoremin güvenilirliği, dayandığı aksiyomların
güvenilirliği kadardır. Mathlib'de yeni bir \leankw{axiom} bildirimi
kabul edilmez. \ciltiki{8}'de göreceğimiz otomatik ispat yarışmalarında da
ilk denetim, çözümün yalnızca bu üç standart aksiyoma dayanmasıdır.

\Lean{}'in üç aksiyomunun birlikte tutarlı olduğu, \Lean{}'in tip kuramının
küme kuramında bir modeli kurularak gösterilmiştir \citep{carneiro2019}.
Bu model ZFC'ye ve sayılabilir sayıda erişilemez kardinale dayanır. Gödel'in
ikinci eksiklik teoremi (\ciltiki{4}), böyle bir göreli tutarlılık sonucundan
daha iyisini bekleyemeyeceğimizi söyler.

\begin{table}[htbp]
\centering
\small
\renewcommand{\arraystretch}{1.45}
\rowcolors{2}{tablosatir}{white}
\begin{tabularx}{\linewidth}{>{\raggedright\arraybackslash}p{3.2cm} X X}
\rowcolor{tabloust}
\tablobaslik{Aksiyom} & \tablobaslik{Olmadan ne olur?} & \tablobaslik{Kim reddeder?} \\
\lc{propext}          & denk önermeler ayrı kalır; \leantac{rw} ↔ ile çalışmaz & bazı yapıcı sistemler (ama çoğu kabul eder) \\
\lc{Quot.sound}       & bölüm tipleri ve \lc{funext} yoktur & yoğunlamalı MLTT \\
\lc{Classical.choice} & üçüncü hâl yoktur; her tanım hesaplanabilir & Bishop tarzı yapıcı matematik \citep{bishop1967} \\
\end{tabularx}
\caption{Her aksiyomun rolü. Yapıcı matematik genellikle yalnızca
\lc{Classical.choice}'u reddeder.}
\label{tab:aksiyom-rolleri}
\end{table}

% ---------------------------------------------------------------------
\section{Bir karşılaştırma: tek-değerlilik}\label{sec:univalence}

Homotopi tip kuramı (HoTT) \citep{hott2013}, \Lean{}'in aksiyomlarına
farklı bir cevap verir. HoTT'de eşitlik kanıtları birer ``yol''dur ve
birbirinden farklı olabilir; kanıt-ilgisizliği \emph{yoktur}. Vladimir
Voevodsky'nin \emph{tek-değerlilik} (\emph{univalence}) aksiyomu, iki tip
arasındaki her denkliğin bir eşitliğe karşılık geldiğini söyler:
\[
  (A \simeq B) \;\simeq\; (A = B).
\]
Bu aksiyom, \lc{propext}'in tiplere genellenmesidir. \lc{propext} denk
\emph{önermelerin} eşit olduğunu söyler, tek-değerlilik ise denk
\emph{tiplerin}. Üstelik tek-değerlilikten \lc{funext} de türetilir.

\begin{table}[htbp]
\centering
\small
\renewcommand{\arraystretch}{1.45}
\rowcolors{2}{tablosatir}{white}
\begin{tabularx}{\linewidth}{>{\raggedright\arraybackslash}p{3.2cm} X X}
\rowcolor{tabloust}
\tablobaslik{} & \tablobaslik{\Lean{} 4 / Mathlib} & \tablobaslik{HoTT} \\
eşitlik kanıtları & hepsi aynı (kanıt-ilgisizliği) & yollar, birbirinden farklı olabilir \\
denk nesneler     & \lc{propext}: denk önermeler eşit & tek-değerlilik: denk tipler eşit \\
\lc{funext}       & \lc{Quot.sound}'dan teorem & tek-değerlilikten teorem \\
bölümler          & ilkel (\lc{Quot}) & yüksek tümevarımsal tipler \\
üçüncü hâl        & teorem (Diaconescu) & önermeler için tutarlı; tüm tipler için tek-değerlilikle çelişir \\
amaç              & klasik matematiğin büyük ölçekli formalizasyonu & homotopi kuramı, yapısal matematik \\
\end{tabularx}
\caption{\Lean{}'in aksiyomları ile homotopi tip kuramının karşılaştırılması.}
\label{tab:lean-hott}
\end{table}

\noindent İki yaklaşım arasındaki seçim bir zevk meselesi değil, bir mühendislik
kararıdır. \Lean{}'in tasarımı, klasik matematikçilerin alıştığı akıl
yürütmeyi mümkün olduğunca az sürtünmeyle formalize etmeyi hedefler. HoTT
ise matematiğin temellerine yeni bir bakış önerir.

% ---------------------------------------------------------------------
\begin{arastirmanotu}
\textbf{Ters matematik.} Hangi teoremin hangi aksiyomlara gerçekten ihtiyaç
duyduğu sorusu, \emph{ters matematik} (\emph{reverse mathematics})
programının konusudur \citep{simpson2009}. İkinci mertebe aritmetiğin beş
alt sistemi, sıradan matematiğin teoremlerinin büyük çoğunluğunu
sınıflandırmaya yeter. \Lean{}'de \leankw{\#print} \lc{axioms}, bu sorunun
çok kaba bir versiyonunu otomatik olarak cevaplar.

\textbf{Yapıcı formalizasyon.} Mathlib klasiktir, ama yapıcı matematiği
\Lean{}'de yapmak mümkündür: \lc{Classical} kullanılmaz ve
\leankw{\#print} \lc{axioms} ile denetlenir. Agda ve Rocq toplulukları
yapıcı ve homotopi tip kuramı tabanlı kütüphaneler geliştirir. Kübik tip
kuramı, tek-değerliliği bir aksiyom olmaktan çıkarıp \emph{hesaplanabilir}
bir teorem hâline getirir. Bu, Bölüm~15'teki ``aksiyomlar hesaplamayı
durdurur'' sorununa bir çözümdür.
\end{arastirmanotu}

\begin{bolumozeti}
\begin{itemize}
  \item \Lean{}'in üç aksiyomu vardır: \lc{propext}, \lc{Quot.sound},
        \lc{Classical.choice}. \leankw{\#print} \lc{axioms} bir teoremin
        hangilerine dayandığını gösterir.
  \item \lc{propext} denk önermeleri eşitler. ↔ ile yeniden yazmak buna
        dayanır. Üçüncü hâlle birlikte \lc{Prop}'u iki değerli yapar.
  \item Bölüm tipleri denklik sınıflarını kurar. Bölümden dışarı her
        fonksiyon iyi tanımlılık kanıtı ister. Tamsayılar doğal sayı
        çiftlerinin bir bölümüdür.
  \item \lc{funext}, \lc{Quot.sound}'dan türetilen bir teoremdir.
  \item \lc{Classical.choice}, \lc{Prop}'tan veri çıkarmanın tek yoludur.
        Kullanan tanımlar \leankw{noncomputable}'dır.
  \item Diaconescu teoremi: seçim + propext + funext $\Rightarrow$ üçüncü
        hâl. \Lean{}'deki \lc{Classical.em} tam olarak bu ispattır.
  \item HoTT'nin tek-değerlilik aksiyomu, \lc{propext}'in tiplere
        genellenmesidir.
\end{itemize}
\end{bolumozeti}

\begin{alistirmalar}
\alistirma{1} $(\mathsf{True} \to p) = p$. (\lc{a16\_1})

\alistirma{1} \lc{(fun n => n * 1) = (fun n => n)}. Hangi aksiyom
kullanıldı? (\lc{a16\_2})

\alistirma{2} İki yüklemin eşitliği için propext ile funext'i birlikte
kullanın. (\lc{a16\_3})

\alistirma{2} \lc{Tamsayi.mk (a + 1) (b + 1) = Tamsayi.mk a b}.
(\lc{a16\_4})

\alistirma{2} Tamsayılarda toplama birleşmelidir. \lc{Quotient.inductionOn₃}
kullanın. (\lc{a16\_5})

\alistirma{2} Seçilen tanık istenen özelliği sağlar. (\lc{a16\_6})

\alistirma{2} \lc{p ≠ True} ise \lc{p = False}. Bu teorem yapıcı mıdır?
\leankw{\#print} \lc{axioms} ile denetleyin. (\lc{a16\_7})

\alistirma{3} \lc{Tamsayi.isaret} birebirdir. Buradan tamsayıların
\lc{Int}'e eşyapılı olduğunu çıkarmak için başka ne gerekir? (\lc{a16\_8})

\alistirma{2} (Kâğıt-kalem) Kesirleri \lc{Int × Nat} çiftlerinin bir bölümü
olarak tanımlayın. Denklik bağıntısı nedir? Paydası sıfır olan çiftlerle ne
yapılmalı?

\alistirma{3} (Proje) \lc{Tamsayi} için çarpmayı tanımlayın ve dağılma
özelliğini ispatlayın. İyi tanımlılık kanıtında \leantac{omega} neden yetmez?
(İpucu: doğrusal olmayan terimler.)

\alistirma{3} (Araştırma) \lc{Classical.choice} kullanmadan \lc{propext} ve
\lc{Quot.sound} ile üçüncü hâli ispatlayamayacağımızı gösteren bir model
düşünün. Bunu \Lean{}'de ispatlamak neden mümkün değildir?
\end{alistirmalar}
   % Lean'in aksiyomları

\part{Projeler}
% =====================================================================
%  Bölüm 17 — Projeler (Cilt I)
% =====================================================================
\chapter{Projeler}\label{bol:proje}

\begin{bolumhedefleri}
\begin{itemize}
  \item Bir formalizasyon projesini ifadeden temiz bir ispata kadar
        yürütmenin yöntemini,
  \item Cilt~I'in konularını derinleştiren on iki projeyi ve her biri için
        önerilen bir yol haritasını,
  \item Projeler arasındaki bağımlılıkları ve Cilt~II'ye uzanan köprüleri.
\end{itemize}
\end{bolumhedefleri}

Bir alıştırma birkaç satırlık bir ispatla biter. Bir \emph{proje} ise
günler, bazen haftalar sürer. Önce doğru tanımları bulmak, sonra ispatı
önsavlara bölmek ve en sonunda tanımlarınızın ispatı ne kadar
kolaylaştırdığını ya da zorlaştırdığını görmek gerekir. Kitap boyunca bazı
sonuçları ``bir projedir'' diyerek kâğıt üzerinde bıraktık. Church--Rosser
teoremi, güçlü normalleşme ve sözdizimsel kesme eleme bunlardandır. Bu
bölüm, bunları on iki proje olarak bir araya getiriyor. Her projenin hedef
teoremi eşlik eden \Lean{} projesinde \lc{sorry} ile hazır bekliyor:
\lc{ModernMantik/}\allowbreak\lc{Alistirmalar/}\allowbreak\lc{Projeler.lean}.

% ---------------------------------------------------------------------
\section{Bir formalizasyon projesi nasıl yürütülür?}\label{sec:proje-yontem}

Deneyimli formalizasyoncuların neredeyse hepsi aynı döngüyü izler
(\cref{fig:proje-dongusu}). Bölüm~12'nin tamlık ispatını bu döngüyle
yazmıştık: Önce \lc{tamlik} ifadesi, sonra onu besleyen önsavlar, en son
da her önsavın kendi ispatı geldi.

\begin{figure}[htbp]
\centering
\begin{tikzpicture}[
    adim/.style={draw=#1, fill=#1!8, very thick, rounded corners=5pt,
                 text width=2.5cm, align=center, inner sep=5pt,
                 font=\small, minimum height=1.5cm},
    ok/.style={-{Stealth[length=2.6mm]}, very thick, black!40},
    no/.style={font=\baslikfont\scriptsize\bfseries, text=white,
               fill=#1, circle, inner sep=1.6pt},
  ]
  \node[adim=anarenk]        (a) at (0,0)    {\textbf{İfade}\\{\scriptsize tanımlar ve hedef teorem, \lc{sorry} ile}};
  \node[adim=ornekframe]     (b) at (3.4,0)  {\textbf{İskelet}\\{\scriptsize önsavlar, her biri \lc{sorry}}};
  \node[adim=ikincil]        (c) at (6.8,0)  {\textbf{İspat}\\{\scriptsize önsavları birer birer kapat}};
  \node[adim=teoremframe]    (d) at (10.2,0) {\textbf{Temizlik}\\{\scriptsize \leankw{\#print} \lc{axioms}, belgeler}};
  \draw[ok] (a) -- (b); \draw[ok] (b) -- (c); \draw[ok] (c) -- (d);
  \draw[ok, dikkatframe, dashed] (c.south) to[bend left=35]
       node[below, font=\baslikfont\scriptsize, text=dikkatframe]{ispat zorlaşıyor mu? tanımları değiştir} (a.south);
  \foreach \n/\c/\k in {a/anarenk/1, b/ornekframe/2, c/ikincil/3, d/teoremframe/4}
    \node[no=\c] at (\n.north west) {\k};
\end{tikzpicture}
\caption{Bir formalizasyon projesinin döngüsü. Kesikli ok en önemli
adımdır: Zor bir ispat çoğu zaman yanlış seçilmiş bir tanımın
işaretidir.}
\label{fig:proje-dongusu}
\end{figure}

\begin{onemlifikir}[Önemli Fikir: \lc{sorry} ile geliştirme]
\lc{sorry}, \Lean{}'e ``bu kısmı sonra ispatlayacağım'' demenin yoludur.
Bir projenin iskeletini önce \lc{sorry}'lerle kurup sonra her \lc{sorry}'yi
ayrı ayrı doldurmak, ispatın genel yapısının doğru olup olmadığını
erkenden görmenizi sağlar. İskeletteki bir önsav yanlış ifade edilmişse
(örneğin bir varsayım eksikse), bunu ana teoremi ispatlamaya çalışırken
fark edersiniz. Önsavı ispatlamak için saatler harcamadan önce.
\end{onemlifikir}

\noindent Kitaptaki tasarım tartışmalarından çıkan birkaç pratik kural:

\begin{itemize}
  \item \textbf{Hesaplayan tanımlar seçin.} \lc{deger}'i \lc{Bool} değerli
        tanımlamak (Bölüm~10), Kalmár önsavını \leantac{cases} ile satır satır
        ilerletmeyi mümkün kıldı.
  \item \textbf{Önermeyi güçlendirin.} Kalmár önsavı hem doğruyu hem yanlışı
        taşıdı. Zayıflatma ispatı \lc{Δ}'yı genelleştirdi. Tümevarım
        takıldığında ilk soru budur.
  \item \textbf{İndis aritmetiğine dikkat.} Bağımlı tiplerde \lc{n + 0} ile
        \lc{0 + n} farklıdır (Bölüm~15). Tanımı özyinelemenin yönüne göre
        yazın.
  \item \textbf{Aksiyomları denetleyin.} Yapıcı olması gereken bir sonuç
        için \leankw{\#print} \lc{axioms} son adımdır (Bölüm~16).
\end{itemize}

% ---------------------------------------------------------------------
\section{Projelerin haritası}\label{sec:proje-haritasi}

\Cref{fig:proje-haritasi}, on iki projeyi dayandıkları bölümlere ve
zorluklarına göre yerleştiriyor. Oklar projeler arasındaki
bağımlılıkları gösteriyor: Örneğin P6'yı yapmak için önce P4'ün ilk
yarısı ve P5'in fikirleri gerekir.

\begin{figure}[htbp]
\centering
\begin{tikzpicture}[
    proje/.style={draw=#1, fill=#1!12, thick, rounded corners=3pt,
                  font=\baslikfont\small\bfseries, inner sep=4pt, minimum width=1.0cm},
    ok/.style={-{Stealth[length=2mm]}, thick, black!40},
    x=1.75cm, y=1.25cm,
  ]
  % eksenler
  \foreach \b/\x in {10/0, 11/1, 12/2, 13/3, 14/4, 15/5, 16/6}
    \node[font=\baslikfont\scriptsize, text=black!55] at (\x,-0.75) {Bölüm \b};
  \draw[black!20] (-0.5,-0.45) -- (6.5,-0.45);
  \node[font=\baslikfont\scriptsize, text=ikincil, anchor=east]     at (-0.55,0) {\zorluk{2}};
  \node[font=\baslikfont\scriptsize, text=dikkatframe, anchor=east] at (-0.55,1.5) {\zorluk{3}};
  \node[font=\baslikfont\scriptsize, text=vurgu, anchor=east]       at (-0.55,3) {\zorluk{3}+};
  \foreach \y in {0,1.5,3} \draw[black!10, dashed] (-0.5,\y) -- (6.5,\y);
  % projeler
  \node[proje=ikincil]     (p1)  at (0,0)   {P1};
  \node[proje=dikkatframe] (p4)  at (1,1.5) {P4};
  \node[proje=dikkatframe] (p12) at (2,1.5) {P12};
  \node[proje=dikkatframe] (p2)  at (3,1.5) {P2};
  \node[proje=vurgu]       (p5)  at (3,3)   {P5};
  \node[proje=vurgu]       (p6)  at (2.3,3) {P6};
  \node[proje=ikincil]     (p3)  at (4,0)   {P3};
  \node[proje=dikkatframe] (p7)  at (4,1.5) {P7};
  \node[proje=vurgu]       (p8)  at (4.6,3) {P8};
  \node[proje=vurgu]       (p9)  at (3.8,3) {P9};
  \node[proje=ikincil]     (p10) at (5,0)   {P10};
  \node[proje=ikincil]     (p11) at (6,0)   {P11};
  % bağımlılıklar
  \draw[ok] (p1) -- (p12);
  \draw[ok] (p4) -- (p6);
  \draw[ok] (p5) -- (p6);
  \draw[ok] (p7) -- (p9);
  \draw[ok] (p3) -- (p8);
  \draw[ok, dashed, teoremframe] (p4.north) -- ++(0,0.65)
       node[above, font=\baslikfont\scriptsize, text=teoremframe]{Cilt~II: Kripke};
\end{tikzpicture}
\caption{On iki projenin haritası. Yatay eksen projenin dayandığı ana
bölümü, dikey eksen zorluğu gösterir. Oklar ``önce bunu yapın''
anlamındadır.}
\label{fig:proje-haritasi}
\end{figure}

\begin{table}[htbp]
\centering
\small
\renewcommand{\arraystretch}{1.35}
\rowcolors{2}{tablosatir}{white}
\begin{tabularx}{\linewidth}{>{\bfseries}p{0.8cm} X p{1.7cm} p{1.4cm} p{2.1cm}}
\rowcolor{tabloust}
\tablobaslik{} & \tablobaslik{Proje} & \tablobaslik{Bölüm} & \tablobaslik{Zorluk} & \tablobaslik{Emek} \\
P1  & Doğruluk tablosu yönteminin doğruluğu        & 10     & \zorluk{2} & 1--2 gün \\
P2  & Ardışık hesapta ispat aramanın doğruluğu     & 13     & \zorluk{3} & 3--5 gün \\
P3  & Normalleştiricinin doğruluğu                 & 14     & \zorluk{2} & 1--2 gün \\
P4  & Yapıcı doğal çıkarım ve Peirce yasası         & 4, 11  & \zorluk{3} & 1 hafta \\
P5  & Sözdizimsel kesme eleme                      & 13     & \zorluk{3} & 2--4 hafta \\
P6  & LJ ile yapıcı doğal çıkarımın denkliği       & 5, 13  & \zorluk{3} & 1--2 hafta \\
P7  & Tip korunumu                                 & 14     & \zorluk{3} & 1 hafta \\
P8  & Church--Rosser teoremi                       & 14     & \zorluk{3} & 1--2 hafta \\
P9  & Güçlü normalleşme                            & 14     & \zorluk{3} & 2--3 hafta \\
P10 & Biriktiricili ters çevirme                   & 15     & \zorluk{2} & 1 gün \\
P11 & Tamsayıların halka yapısı                    & 16     & \zorluk{2} & 2--3 gün \\
P12 & Yapıcı model varlığı                         & 10, 12 & \zorluk{3} & 1 hafta \\
\end{tabularx}
\caption{Projelerin özeti. Emek tahminleri, kitabı buraya kadar
alıştırmalarıyla birlikte izlemiş bir okur içindir.}
\label{tab:projeler}
\end{table}

% ---------------------------------------------------------------------
\section{A. Karar yöntemleri ve yansıma}\label{sec:projeler-a}

Bu gruptaki üç proje, kitapta yazdığımız \emph{çalışan} fonksiyonların
(\lc{totolojiMi}, \lc{ara}, \lc{adim?}) gerçekten söylediklerini yaptığını
ispatlamayı ister. Bir fonksiyonun doğruluğunu ispatlamak, onu bir
\emph{karar yöntemine} dönüştürmenin ilk adımıdır. Cilt~II'nin 6. bölümünde
bu tekniği, \emph{yansıma} adıyla, \Lean{}'in kendi önermelerine
uygulayacağız.

\begin{projekarti}{P1}{Doğruluk tablosu yöntemi doğrudur}
\projebilgi{Bölüm~10}{\zorluk{2}}{1--2 gün}
Bölüm~10'daki \lc{totolojiMi}, formülü \lc{tumDegerlemeler}'deki her
değerlemede hesaplar. Bunun gerçekten totolojiliği yakaladığını ispatlayın.

\leandosya[P1]{Alistirmalar/Projeler}{p1}

\textbf{Önerilen yol.}
\begin{enumerate}[nosep]
  \item \emph{Kapsama önsavı}: Her \lc{w} için, \lc{tumDegerlemeler n}
        içinde ilk \lc{n} atomda \lc{w} ile anlaşan bir \lc{v} vardır.
        \lc{n} üzerinde tümevarımla ve \lc{genislet}'in tanımıyla.
  \item Kapsama önsavını \lc{yerellik} ile birleştirin: \lc{deger v A =
        deger w A}.
  \item \lc{List.all} için \lc{List.all\_eq\_true} önsavını kullanın.
\end{enumerate}
\end{projekarti}

\begin{projekarti}{P2}{İspat arama doğrudur}
\projebilgi{Bölüm~13}{\zorluk{3}}{3--5 gün}
Bölüm~13'teki \lc{ara} fonksiyonu \lc{true} döndürdüğünde, ardışığın gerçekten
kesmesiz bir türetimi vardır.

\leandosya[P2]{Alistirmalar/Projeler}{p2}

\textbf{Önerilen yol.}
\begin{enumerate}[nosep]
  \item \lc{tam\_aux}'un örüntülerini aynen izleyin; her durumda \lc{ara}'nın
        tanımı tek bir özyinelemeli çağrıya ya da bir \lc{\&\&}'e açılır.
  \item \lc{(a \&\& b) = true}'yu \lc{Bool.and\_eq\_true} ile ikiye bölün.
  \item Atomik durum için \lc{List.any\_eq\_true} ve \lc{List.contains}
        önsavlarından ortak bir atom elde edin.
  \item \emph{Uzantı:} Ters yönü (\lc{Gecerli} ise \lc{ara = true}) de
        ispatlayıp \lc{ara}'yı bir karar yöntemine dönüştürün.
\end{enumerate}
\end{projekarti}

\begin{projekarti}{P3}{Normalleştiricinin her adımı bir β-adımıdır}
\projebilgi{Bölüm~14}{\zorluk{2}}{1--2 gün}
\leandosya[P3]{Alistirmalar/Projeler}{p3}

\textbf{Önerilen yol.}
\begin{enumerate}[nosep]
  \item \lc{t} üzerinde tümevarım yapın. \lc{app t₁ t₂} durumunda
        \lc{t₁}'in bir \lc{lam} olup olmadığına göre ayırın; çünkü
        \lc{adim?}'ın örüntüleri örtüşüyor.
  \item \lc{Option.map f x = some u} hipotezini \lc{x = some u'} ve \lc{u =
        f u'} biçimine çeviren küçük bir önsav yazın.
  \item \emph{Uzantı:} \lc{adim? t = none} ise \lc{Normal t}.
\end{enumerate}
\end{projekarti}

% ---------------------------------------------------------------------
\section{B. İspat kuramı}\label{sec:projeler-b}

\begin{projekarti}{P4}{Yapıcı doğal çıkarım ve Peirce yasası}
\projebilgi{Bölüm~4, 11}{\zorluk{3}}{1 hafta}
\lc{Turetilir}'den RAA kuralını çıkarınca yapıcı doğal çıkarım elde edilir
(Bölüm~11, \cref{sec:gecerlilik-sinir}). Peirce yasası bu sistemde
türetilemez. Ama bunu göstermek için iki değerli anlambilim yetmez.

\leandosya[P4]{Alistirmalar/Projeler}{p4}

\textbf{Önerilen yol.}
\begin{enumerate}[nosep]
  \item (a) kural tümevarımıyla birkaç satırdır.
  \item (b) için bir \emph{Kripke modeli} tanımlayın: kısmi sıralı bir
        dünyalar kümesi ve tekdüze bir atom yorumu. Zorlama bağıntısını
        ($w \Vdash A$) formül üzerinde özyinelemeyle tanımlayın (Bölüm~4).
  \item \lc{TuretilirY}'nin Kripke anlambilimine göre geçerliliğini ispatlayın.
        \lc{iseG} durumunda tekdüzelik önsavı gerekecek.
  \item Peirce yasası için iki dünyalı bir karşı model kurun:
        $w_0 \le w_1$, $p$ yalnızca $w_1$'de doğru.
\end{enumerate}
\emph{Kaynak:} \citet{kripke1965,troelstra1988}. Kripke anlambilimi
Cilt~II'nin 2. bölümünde ayrıntılı olarak işlenecek.
\end{projekarti}

\begin{projekarti}{P5}{Sözdizimsel kesme eleme}
\projebilgi{Bölüm~13}{\zorluk{3}}{2--4 hafta}
Bölüm~13'te kesmenin elenebilirliğini anlambilim yoluyla, birkaç satırda
ispatladık. Gentzen'in yolu ise kesmeli bir türetimi doğrudan kesmesiz bir
türetime \emph{dönüştürmektir}. Bu proje, bu dönüşümü yazmayı ister.

\leandosya[P5]{Alistirmalar/Projeler}{p5}

\textbf{Önerilen yol.}
\begin{enumerate}[nosep]
  \item Kesmesiz \lc{Ardisik} için \emph{sözdizimsel} tersinirlik önsavları
        ispatlayın (Bölüm~13'tekiler anlambilimseldi). Örneğin \lc{Ardisik
        ((A {\sembolkodfont ⋀} B) :: Γ) Δ → Ardisik (A :: B :: Γ) Δ}, türetim üzerinde
        tümevarımla.
  \item \emph{Kesme kabul edilebilirliği}: \lc{Ardisik Γ (A :: Δ)} ve
        \lc{Ardisik (A :: Γ) Δ} ise \lc{Ardisik Γ Δ}. \lc{A}'nın boyutu
        üzerinde dış tümevarım, birinci türetim üzerinde iç tümevarım.
        Ana formül her iki tarafta da yeni tanıtılmışsa tersinirlik
        önsavlarını kullanın.
  \item Son olarak \lc{ArdisikK} üzerinde tümevarımla, her kesmeyi bir
        önceki adımla eleyin.
\end{enumerate}
\emph{Kaynak:} \citet{gentzen1935,pfenning2000,negri2001}.
\end{projekarti}

\begin{projekarti}{P6}{LJ ile yapıcı doğal çıkarım aynıdır}
\projebilgi{Bölüm~5, 13}{\zorluk{3}}{1--2 hafta}
Gentzen'in sezgici ardışık hesabı LJ, sağ tarafta en fazla bir formüle izin
verir. LJ'nin, P4'teki \lc{TuretilirY} ile aynı formülleri türettiğini
gösterin.

\leandosya[P6]{Alistirmalar/Projeler}{p6}

\textbf{Önerilen yol.}
\begin{enumerate}[nosep]
  \item LJ $\Rightarrow$ doğal çıkarım: Sol kurallar eleme kurallarıyla
        taklit edilir. Kolay yöndür.
  \item Doğal çıkarım $\Rightarrow$ LJ: Eleme kuralları için LJ'de
        \emph{kesme} gerekir. Önce LJ için kesme kabul edilebilirliğini
        (P5'in tekniğiyle) ispatlayın.
  \item \lc{iseL} kuralının ana formülü öncülde de tuttuğuna dikkat edin.
        Bu, LJ'de daraltmanın kabul edilebilirliği için gereklidir.
\end{enumerate}
\emph{Kaynak:} \citet{negri2001,prawitz1965,dyckhoff1992}.
\end{projekarti}

% ---------------------------------------------------------------------
\section{C. $λ$-hesabı}\label{sec:projeler-c}

\begin{projekarti}{P7}{Tip korunumu}
\projebilgi{Bölüm~14}{\zorluk{3}}{1 hafta}
Bir terim indirgendiğinde tipi değişmez. Bu teorem, ``iyi tiplenmiş
programlar yanlış gitmez'' sloganının yarısıdır.

\leandosya[P7]{Alistirmalar/Projeler}{p7}

\textbf{Önerilen yol.}
\begin{enumerate}[nosep]
  \item \emph{Zayıflatma}: \lc{Tiplenir (Γ₁ ++ Γ₂) t A} ise
        \lc{Tiplenir (Γ₁ ++ B :: Γ₂) (kaydir 1 Γ₁.length t) A}.
        \lc{t} üzerinde tümevarım, \lc{Γ₁}'i genelleştirerek.
  \item \emph{Yerine koyma}: \lc{Tiplenir (Γ₁ ++ B :: Γ₂) t A} ve
        \lc{Tiplenir Γ₂ s B} ise \lc{Tiplenir (Γ₁ ++ Γ₂) (yerine
        Γ₁.length s t) A}. \lc{bul} için liste birleştirme önsavları
        gerekecek.
  \item Son olarak \lc{Adim} üzerinde tümevarım. Yalnızca \lc{beta}
        durumu yerine koyma önsavını kullanır.
\end{enumerate}
\emph{Kaynak:} \citet{barendregt1984,sorensen2006}.
\end{projekarti}

\begin{projekarti}{P8}{Church--Rosser teoremi}
\projebilgi{Bölüm~14}{\zorluk{3}}{1--2 hafta}
\leandosya[P8]{Alistirmalar/Projeler}{p8}

\textbf{Önerilen yol.}
\begin{enumerate}[nosep]
  \item \emph{Paralel indirgeme} \lc{Paralel t u}'yu tanımlayın: Bir
        terimdeki redekslerin herhangi bir alt kümesi aynı anda
        indirgenir.
  \item \lc{Adim ⊆ Paralel ⊆ Adimlar} olduğunu gösterin.
  \item Takahashi'nin yöntemi \citep{takahashi1995}: Her terimin bütün
        redekslerini indirgeyen \lc{tamGelisme t}'yi tanımlayın ve
        \lc{Paralel t u → Paralel u (tamGelisme t)} önsavını ispatlayın.
        Bu, \lc{Paralel}'in elmas özelliğini hemen verir.
  \item Elmas özelliğini \lc{Adimlar}'a taşıyın. \lc{kaydir} ve \lc{yerine}
        için bir ``yerine koyma ile paralel indirgeme'' önsavı gerekecek.
\end{enumerate}
\end{projekarti}

\begin{projekarti}{P9}{Güçlü normalleşme}
\projebilgi{Bölüm~14}{\zorluk{3}}{2--3 hafta}
Tiplenebilen her terimin her indirgeme dizisi sonludur. \lc{GucluNormal},
bu özelliğin \emph{tümevarımsal} bir tanımıdır: Bir terim, bütün tek adımlık
indirgenmişleri güçlü normalleşiyorsa güçlü normalleşir.

\leandosya[P9]{Alistirmalar/Projeler}{p9}

\textbf{Önerilen yol.}
\begin{enumerate}[nosep]
  \item Isınma: \lc{omega\_adim} ile \lc{p9\_omega}.
  \item Tait'in yöntemi: Tip üzerinde özyinelemeyle bir
        \lc{Indirgenebilir : Tip → Terim → Prop} tanımlayın. Ok tipinde
        ``her indirgenebilir argümana uygulandığında indirgenebilir''.
  \item Üç özellik (CR1--CR3): İndirgenebilir terimler güçlü normalleşir,
        indirgenme altında kapalıdır ve ``nötr'' terimler için genişleme
        altında kapalıdır.
  \item \emph{Temel önsav}: Tiplenebilen bir terime indirgenebilir terimler
        yerine konursa sonuç indirgenebilirdir. P7'nin yerine koyma
        önsavları burada yeniden kullanılır.
\end{enumerate}
\emph{Kaynak:} \citet{tait1967}; \citet[Bölüm~6]{girard1989}.
\end{projekarti}

% ---------------------------------------------------------------------
\section{D. Bağımlı tipler ve bölümler}\label{sec:projeler-d}

\begin{projekarti}{P10}{Biriktiricili ters çevirme}
\projebilgi{Bölüm~15}{\zorluk{2}}{1 gün}
\leandosya[P10]{Alistirmalar/Projeler}{p10}

\textbf{Önerilen yol.}
\begin{enumerate}[nosep]
  \item \lc{p10\_tersAux bos w}'nin tipi \lc{Vek α (0 + n)}'dir, \lc{w}'nin
        ise \lc{Vek α n}. Ya tipi \lc{Vek α (n + m)} olarak değiştirin
        (hangi durumda sorun çıkıyor?) ya da \lc{Nat.zero\_add n ▸ w}
        ile taşıyın.
  \item \lc{toList}'in taşımayla değişmediğini \lc{a15\_5}'teki gibi
        \leantac{subst} ile gösterin.
  \item Genelleştirilmiş önsav: \lc{(tersAux v w).toList = v.toList.reverse
        ++ w.toList}.
\end{enumerate}
\end{projekarti}

\begin{projekarti}{P11}{Tamsayıların halka yapısı}
\projebilgi{Bölüm~16}{\zorluk{2}}{2--3 gün}
\leandosya[P11]{Alistirmalar/Projeler}{p11}

\textbf{Önerilen yol.}
\begin{enumerate}[nosep]
  \item $(a - b)(c - d) = (ac + bd) - (ad + bc)$ olduğundan çarpmayı
        temsilciler üzerinde \lc{(a*c + b*d, a*d + b*c)} olarak tanımlayın.
  \item İyi tanımlılık doğrusal olmayan bir eşitlik ister; \leantac{omega}
        yetmez. \lc{Nat.mul\_add}, \lc{Nat.add\_mul} ile elle ya da Mathlib'in
        \leantac{ring}/\leantac{nlinarith} taktikleriyle ispatlayın.
        Önce bir argümanı, sonra diğerini değiştirmek işi kolaylaştırır.
  \item Halka aksiyomlarını \lc{Quotient.inductionOn₃} ile ispatlayın.
  \item \emph{Uzantı:} \lc{Tamsayi.isaret}'in bir halka eşyapı dönüşümü
        olduğunu ve \lc{Int}'e bir eşleme verdiğini gösterin.
\end{enumerate}
\end{projekarti}

\begin{projekarti}{P12}{Yapıcı model varlığı}
\projebilgi{Bölüm~10, 12}{\zorluk{3}}{1 hafta}
Bölüm~12'deki \lc{model\_varligi}, \lc{Classical.byContradiction} kullanıyordu.
Doğruluk tablosunu gerçekten tarayan yapıcı bir sürüm yazın.

\leandosya[P12]{Alistirmalar/Projeler}{p12}

\textbf{Önerilen yol.}
\begin{enumerate}[nosep]
  \item \lc{Γ}'nın formüllerindeki en büyük atom indisini bulun ve
        \lc{tumDegerlemeler} listesini \lc{List.find?} ile tarayın.
  \item \lc{p12\_dogru}, \lc{List.find?}'ın bulduğu elemanın koşulu
        sağladığını söyleyen önsavdan çıkar.
  \item \lc{p12\_tam} için P1'deki kapsama önsavı ve Bölüm~12'deki
        \lc{tamlik} gerekir.
  \item Son olarak \leankw{\#print} \lc{axioms} ile \lc{Classical.choice}'un
        kullanılmadığını doğrulayın.
\end{enumerate}
\end{projekarti}

% ---------------------------------------------------------------------
\section{Projeyi bitirdiğinizde}\label{sec:proje-bitis}

\begin{itemize}
  \item \textbf{\lc{sorry} kalmasın.} Projenizi \lc{Alistirmalar/}'dan
        kendi dosyanıza taşıyın ve sürekli tümleştirmenin \lc{sorry}
        denetimine dahil edin.
  \item \textbf{Aksiyomları listeleyin.} Sonucunuzun dayandığı aksiyomları
        belgeleyin. Yapıcı olması gereken bir sonuç \lc{Classical.choice}'a
        dayanıyorsa nedenini bulun.
  \item \textbf{Belgeleyin.} Her tanımın ve ana teoremin başına,
        neyi neden öyle tanımladığınızı açıklayan bir yorum yazın. Altı ay
        sonra kodu okuyacak kişi muhtemelen siz olacaksınız.
  \item \textbf{Paylaşın.} \Lean{} topluluğunun Zulip sohbet sunucusu ve
        Reservoir paket dizini, bitmiş projeler için doğal yerlerdir. Bir
        sonucun Mathlib'e katkı olarak uygun olup olmadığını orada
        sorabilirsiniz.
\end{itemize}

\begin{arastirmanotu}
Bu projelerin çoğu, ispat asistanlarının kıyaslama ölçütü olarak
kullanılmıştır. Church--Rosser ve güçlü normalleşme, POPLmark yarışmasının
\citep{aydemir2005} ruhuna uygun ``bağlayıcı içeren metateori''
örnekleridir. Kesme eleme ise formalize edilmiş ispat kuramının merkezindedir
\citep{chaudhuri2019,shillito2023}. Lean~4'te bu sonuçların bir kısmı
Foundation kütüphanesinde \citep{foundation} genel biçimleriyle bulunur.
Kendi çözümünüzü yazdıktan sonra oradaki çözümlerle karşılaştırmak,
tanım seçimlerinin etkisini görmenin en iyi yoludur.
\end{arastirmanotu}

% ---------------------------------------------------------------------
\chapter*{Sonsöz: Cilt II'ye Bakış}
\addcontentsline{toc}{chapter}{Sonsöz: Cilt II'ye Bakış}
\markboth{Sonsöz}{Sonsöz}

Cilt~I'e ``Mantık nedir?'' sorusuyla başladık ve \Lean{}'in çekirdeğinin
neye inandığını sorarak bitirdik. Bu yolda önermeler mantığını hem bir
araç hem de bir inceleme nesnesi olarak gördük. İspat sistemlerinin
geçerli ve tam olduğunu, kesmenin elenebildiğini ve tiplenebilen
programların ispat olduğunu \Lean{}'de ispatladık. Bölüm~16'da üçüncü hâlin
olanaksızlığının \Lean{}'de bir teorem olduğunu gördüğümüzde, Bölüm~4'te
açtığımız yapıcı/klasik sorusu da kapanmış oldu.

Cilt~I'in büyük ölçüde \emph{önermeler} mantığında kaldığını fark etmiş
olabilirsiniz. Niceleyicileri \Lean{}'in kendi mantığında bolca kullandık
(sığ gömme). Ama niceleyicili bir nesne dilini ve onun metateorisini derin
gömmeyle kurmadık. Bölüm~12'de Henkin yöntemini kâğıt üzerinde bırakmamızın
nedeni budur. Bağlı değişkenlerin bütün zorluğu, Bölüm~14'te
$\lambda$-hesabında gördüğümüz gibi, orada bizi bekliyor.

\textbf{Cilt~II} tam bu noktadan başlıyor. Planlanan içerik şöyle:

\begin{center}
\renewcommand{\arraystretch}{1.4}
\rowcolors{2}{tablosatir}{white}
\begin{tabularx}{0.92\linewidth}{>{\bfseries}p{0.7cm} X}
\rowcolor{tabloust}
\tablobaslik{} & \tablobaslik{Cilt II: İleri Konular} \\
1 & Birinci mertebe mantık: sözdizim, anlambilim ve Gödel--Henkin tamlığı, \Lean{}'de \\
2 & Kip mantığı ve Kripke anlambilimi; sezgici mantığın S4'e gömülmesi \\
3 & Kanıtlanabilirlik mantığı (GL) \\
4 & Gödel'in eksiklik teoremleri \\
5 & Küme kuramı, modeller ve bağımsızlık \\
6 & Karar verilebilirlik ve otomasyon: SAT/SMT, \leantac{omega}, \leantac{bv\_decide}, yansıma \\
7 & Metaprogramlama: kendi taktiğinizi yazmak \\
8 & Yapay zekâ ve otomatik teorem ispatı \\
9 & Büyük formalizasyon projeleri ve araştırma projeleri \\
\end{tabularx}
\end{center}

Cilt~I boyunca ``Cilt~II, Bölüm~$n$'' diye işaret ettiğimiz her konu bu
tablodadır. Cilt~I'deki P4 projesi (Kripke modelleri ile Peirce yasası) ve
P2 projesi (ispat aramanın doğruluğu), sırasıyla 2. ve 6. bölümlerin doğal
başlangıç noktalarıdır.

\vspace{1em}
\begin{flushright}
\yazar\\
\textit{Ekim 2026}
\end{flushright}
   % Projeler ve Sonsöz: Cilt II'ye bakış

\appendix
\chapter{Lean 4 Hızlı Başvuru}\label{ek:basvuru}

Bu ek, kitap boyunca kullanılan taktikleri ve komutları tek bir yerde toplar.
Her yeni bölümle birlikte genişletilecektir.

\section{Taktikler ve doğal çıkarım kuralları}

\begin{center}
\renewcommand{\arraystretch}{1.4}
\rowcolors{2}{tablosatir}{white}
\begin{tabularx}{\linewidth}{>{\raggedright\arraybackslash}p{2.6cm}
                             >{\raggedright\arraybackslash}p{3.2cm}
                             >{\raggedright\arraybackslash}X
                             >{\centering\arraybackslash}p{1.4cm}}
\rowcolor{tabloust}
\tablobaslik{Taktik} & \tablobaslik{Hedef / durum} & \tablobaslik{Etkisi} & \tablobaslik{Bölüm} \\
\leantac{intro h}       & \lc{⊢ A → B}           & $A$'yı \lc{h} adıyla varsayar; hedef $B$ olur. ($\to$G) & \ref{bol:onermeler1} \\
\leantac{apply h}       & \lc{h : A → B}, \lc{⊢ B} & Hedefi $A$'ya çevirir. ($\to$E, geriye doğru) & \ref{bol:onermeler1} \\
\leantac{exact t}       & \lc{t : B}, \lc{⊢ B}   & İspatı bitirir. & \ref{bol:giris} \\
\leantac{constructor}   & \lc{⊢ A ∧ B}           & İki hedef açar: $A$ ve $B$. ($\wedge$G) & \ref{bol:onermeler1} \\
\leantac{obtain ⟨a, b⟩ := h} & \lc{h : A ∧ B}    & \lc{a : A} ve \lc{b : B} hipotezlerini verir. ($\wedge$E) & \ref{bol:onermeler1} \\
\leantac{left} / \leantac{right} & \lc{⊢ A ∨ B}  & Hedefi $A$ (ya da $B$) yapar. ($\vee$G) & \ref{bol:onermeler1} \\
\leantac{cases h}       & \lc{h : A ∨ B}         & İki duruma ayırır. ($\vee$E) & \ref{bol:onermeler1} \\
\leantac{rcases h with a | b} & \lc{h : A ∨ B}   & Durumlara ayırır ve adlandırır. & \ref{bol:onermeler1} \\
\leantac{assumption}    & hedef bir hipotezde    & Uygun hipotezi bulup kullanır. & \ref{bol:onermeler1} \\
\leantac{sorry}         & her durumda            & İspatı \emph{kapatmadan} geçer; uyarı üretir. & \ref{bol:giris} \\
\leantac{trivial}       & \lc{⊢ True}            & Hedefi kapatır. ($\top$G) & \ref{bol:degilleme} \\
\leantac{exfalso}       & her hedef              & Hedefi \lc{False} yapar. ($\bot$E) & \ref{bol:degilleme} \\
\leantac{contradiction} & bağlamda çelişki       & İspatı bitirir. & \ref{bol:degilleme} \\
\leantac{rw [h]}        & \lc{h : P ↔ Q} ya da \lc{a = b} & Hedefi yeniden yazar. & \ref{bol:degilleme} \\
\lc{Classical.em P}     & her durumda            & \lc{P ∨ ¬P} verir (klasik). & \ref{bol:yapici-klasik} \\
\lc{Classical.byContradiction} & her hedef       & Çelişki ile ispat (klasik). & \ref{bol:yapici-klasik} \\
\leantac{decide}        & karar verilebilir hedef & Hesaplayarak ispatlar. & \ref{bol:yapici-klasik} \\
\leantac{have h : A := t} & her durumda          & Bağlama ara sonuç ekler (ileriye doğru). & \ref{bol:dogal-cikarim} \\
\leantac{suffices h : A by t} & her hedef        & ``\lc{A}'yı göstermek yeter''; hedef \lc{A} olur. & \ref{bol:dogal-cikarim} \\
\leantac{show A}        & hedef \lc{A}           & Hedefi açıkça yazar; değiştirmez. & \ref{bol:dogal-cikarim} \\
\leantac{refine ⟨?\_, ?\_⟩} & her hedef          & Terim iskeleti yazar; \lc{?\_} yeni hedef olur. & \ref{bol:dogal-cikarim} \\
\lc{t₁ <;> t₂}          & her hedef              & \lc{t₁}'den çıkan her hedefe \lc{t₂}'yi uygular. & \ref{bol:dogal-cikarim} \\
\leantac{all\_goals t}  & birden çok hedef       & Bütün hedeflere \lc{t}'yi uygular. & \ref{bol:dogal-cikarim} \\
\leantac{first | t₁ | t₂} & her hedef            & İlk başarılı taktiği uygular. & \ref{bol:dogal-cikarim} \\
\leantac{intro x}       & \lc{⊢ ∀ x, P x}        & Keyfi bir \lc{x} alır. ($\forall$G) & \ref{bol:yuklemler} \\
\leantac{specialize h t} & \lc{h : ∀ x, P x}     & \lc{h}'yi \lc{h : P t} yapar. ($\forall$E) & \ref{bol:yuklemler} \\
\lc{⟨t, ht⟩}            & \lc{⊢ ∃ x, P x}        & Tanık ve kanıt verir. ($\exists$G) & \ref{bol:yuklemler} \\
\leantac{obtain ⟨a, ha⟩ := h} & \lc{h : ∃ x, P x} & Tanığa ad verir. ($\exists$E) & \ref{bol:yuklemler} \\
\leantac{rw [← h]}      & \lc{h : a = b}         & Hedefte \lc{b}'yi \lc{a} yapar. & \ref{bol:yuklemler} \\
\leankw{calc}           & eşitlik zinciri        & Adım adım eşitlik ispatı. & \ref{bol:yuklemler} \\
\leantac{induction n with} & \lc{n : Nat}, \lc{n : Dogal} & Tümevarım; her kurucu için bir durum. & \ref{bol:tumevarim} \\
\leantac{injection h with h'} & \lc{h : c a = c b} & \lc{h' : a = b} verir. & \ref{bol:tumevarim} \\
\leantac{change t at h} & her hipotez           & Hipotezi tanımsal eşdeğerine çevirir. & \ref{bol:tumevarim} \\
\leantac{omega}         & doğrusal aritmetik     & \lc{Nat}/\lc{Int} üzerinde karar verir. & \ref{bol:tumevarim} \\
\leantac{simp}          & her hedef              & Bilinen eşitliklerle sadeleştirir. & \ref{bol:tumevarim} \\
\leantac{ext x}         & \lc{⊢ s = t}           & Elemanlara indirger. (\Mathlib{}) & \ref{bol:kumeler} \\
\leantac{ring}          & polinom eşitliği       & Halka eşitliklerini kanıtlar. (\Mathlib{}) & \ref{bol:kumeler} \\
\leantac{linarith}      & doğrusal eşitsizlik    & Hipotezlerden türetir. (\Mathlib{}) & \ref{bol:kumeler} \\
\leantac{norm\_num}     & sayısal hedef          & Hesaplar, asallık denetler. (\Mathlib{}) & \ref{bol:kumeler} \\
\leantac{positivity}    & \lc{⊢ 0 ≤ e}, \lc{0 < e} & Pozitifliği gösterir. (\Mathlib{}) & \ref{bol:kumeler} \\
\leantac{tauto}         & totoloji               & Klasik önermeler mantığı. (\Mathlib{}) & \ref{bol:kumeler} \\
\lc{exact?}             & her hedef              & Kütüphanede teorem arar. & \ref{bol:kumeler} \\
\leantac{induction h with} & \lc{h} bir yüklem kanıtı & Kural tümevarımı. & \ref{bol:tumevarimsal-tipler} \\
\lc{cases r <;> rfl}    & sonlu tip              & Her durumu hesaplayarak kapatır. & \ref{bol:tumevarimsal-tipler} \\
\leantac{show t}        & her hedef              & Hedefi tanımsal eşdeğeriyle açıkça yazar. & \ref{bol:sozdizim} \\
\leantac{cases h : e}   & \lc{e : Bool}          & \lc{e}'nin değerine göre ayırır, \lc{h : e = true} vb. verir. & \ref{bol:sozdizim} \\
\leantac{induction h generalizing Δ} & kural tümevarımı & Hipotezi her \lc{Δ} için güçlendirir. & \ref{bol:gecerlilik} \\
\leantac{cases a with | false => … | true => …} & \lc{a : Bool} & Adlandırılmış durumlar; sıradan bağımsız. & \ref{bol:tamlik} \\
\leantac{simp at h}     & \lc{h : false = true} & Çelişkili hipotezle hedefi kapatır. & \ref{bol:tamlik} \\
\leantac{unfold f}      & tanımlı \lc{f}         & \lc{f}'nin tanımını açar. & \ref{bol:tamlik} \\
\leankw{abbrev}         & tanım                  & Açılabilir (indirgenebilir) tanım; \lc{⟨\_, \_⟩} ile kurulabilir. & \ref{bol:tamlik} \\
\leantac{refine t ?\_}  & her hedef              & Terimi verir; \lc{?\_} delikleri yeni hedef olur. & \ref{bol:ardisik} \\
\leantac{apply f (x := t)} & örtük argümanlı \lc{f} & Örtük argümanı adıyla belirleyerek uygular. & \ref{bol:ardisik} \\
\leantac{rcases h with h1 | h2} & \lc{h : a ∈ s ++ t} & \lc{List.mem\_append} ile durumlara ayırır. & \ref{bol:ardisik} \\
\leantac{injection h with h1 h2} & \lc{h : c a b = c a' b'} & Kurucu argümanlarının eşitliklerini verir. & \ref{bol:lambda} \\
\leantac{by\_cases h : p} & her hedef         & \lc{p} ve \lc{¬p} durumlarına ayırır. & \ref{bol:lambda} \\
\leantac{subst h}       & \lc{h : x = e}         & Değişkeni her yerde terimle değiştirir. & \ref{bol:lambda} \\
\lc{h ▸ e}              & \lc{h : a = b}         & \lc{e}'yi eşitlik boyunca taşır. & \ref{bol:bagimli} \\
\leankw{\#print axioms t} & bir teorem           & \lc{t}'nin dayandığı aksiyomları listeler. & \ref{bol:bagimli} \\
\lc{Fin.elim0 i}        & \lc{i : Fin 0}         & Olanaksız durumu kapatır. & \ref{bol:bagimli} \\
\leantac{funext x}      & \lc{⊢ f = g}           & Noktasal eşitliğe indirger. & \ref{bol:aksiyomlar} \\
\lc{propext ⟨h1, h2⟩}   & \lc{⊢ p = q}           & Denklikten önerme eşitliği. & \ref{bol:aksiyomlar} \\
\lc{Quotient.inductionOn x} & \lc{x : Quotient s} & Temsilci üzerinde ispata indirger. & \ref{bol:aksiyomlar} \\
\lc{Quotient.sound h}   & \lc{⊢ ⟦a⟧ = ⟦b⟧}       & İlişkili temsilcilerin sınıfları eşittir. & \ref{bol:aksiyomlar} \\
\end{tabularx}
\end{center}

\section{Komutlar}

\begin{center}
\renewcommand{\arraystretch}{1.4}
\rowcolors{2}{tablosatir}{white}
\begin{tabularx}{\linewidth}{>{\raggedright\arraybackslash}p{3.4cm} X}
\rowcolor{tabloust}
\tablobaslik{Komut} & \tablobaslik{Ne yapar?} \\
\lc{\textcolor{leancmd}{\bfseries \#check} e}   & \lc{e} ifadesinin türünü gösterir. \\
\lc{\textcolor{leancmd}{\bfseries \#eval} e}    & \lc{e} ifadesini hesaplar ve sonucu yazar. \\
\lc{\textcolor{leancmd}{\bfseries \#print} ad}  & Bir tanımın içeriğini gösterir. \\
\lc{\textcolor{leancmd}{\bfseries \#print axioms} ad} & Bir teoremin hangi aksiyomlara dayandığını listeler. \\
\end{tabularx}
\end{center}

\section{Unicode sembollerin yazılışı}

VS~Code'daki Lean eklentisinde semboller ters bölü ile yazılır ve boşluk
tuşuna basıldığında dönüşür.

\begin{center}
\renewcommand{\arraystretch}{1.3}
\begin{tabular}{cl@{\hspace{2.5em}}cl@{\hspace{2.5em}}cl}
\toprule
\lc{→} & \lc{\textbackslash to}  & \lc{∧} & \lc{\textbackslash and} & \lc{∀} & \lc{\textbackslash forall} \\
\lc{←} & \lc{\textbackslash l}   & \lc{∨} & \lc{\textbackslash or}  & \lc{∃} & \lc{\textbackslash ex} \\
\lc{↔} & \lc{\textbackslash iff} & \lc{¬} & \lc{\textbackslash not} & \lc{λ} & \lc{\textbackslash fun} \\
\lc{⟨ ⟩} & \lc{\textbackslash< \textbackslash>} & \lc{⊢} & \lc{\textbackslash |-} & \lc{ℕ} & \lc{\textbackslash N} \\
\bottomrule
\end{tabular}
\end{center}

\chapter{Türkçe--İngilizce Terim Sözlüğü}\label{ek:sozluk}

Mantık ve tip kuramı terimlerinin Türkçe karşılıkları henüz tam olarak
yerleşmiş değildir. Bu kitapta aşağıdaki karşılıklar tutarlı biçimde
kullanılmaktadır. Parantez içindeki karşılıklar da literatürde rastlanan
seçeneklerdir.

\begin{center}
\renewcommand{\arraystretch}{1.25}
\rowcolors{2}{tablosatir}{white}
\begin{tabularx}{\linewidth}{>{\raggedright\arraybackslash}X >{\raggedright\arraybackslash\itshape}X}
\rowcolor{tabloust}
\tablobaslik{Türkçe} & \tablobaslik{İngilizce} \\
adla / değerle çağırma & call-by-name / call-by-value \\
$\alpha$-denklik & alpha-equivalence \\
alt formül özelliği & subformula property \\
alt küme & subset \\
alt tip & subtype \\
alt yapısal mantık & substructural logic \\
ana formül & principal formula \\
anlambilim & semantics \\
ara değer (interpolasyon) & interpolation \\
ardıl & successor \\
ardışık hesap & sequent calculus \\
ardışık (sekant) & sequent \\
bağımlı çift & dependent pair \\
bağımlı fonksiyon & dependent function \\
bağımlı tip & dependent type \\
bağımlı tip kuramı & dependent type theory \\
bağlı değişken & bound variable \\
$\beta$-indirgeme & beta-reduction \\
bileşke & composition \\
birebir & injective \\
birleşebilirlik & confluence, Church--Rosser property \\
birleşim & union \\
bölüm tipi & quotient type \\
büyük eleme & large elimination \\
çelişki ile ispat & proof by contradiction, reductio ad absurdum \\
çelişmezlik ilkesi & law of non-contradiction \\
çıkarım kuralı & inference rule \\
çift değilleme çevirisi & double-negation translation \\
çift değilleme eleme & double negation elimination \\
çok sonuçlu ardışık & multiple-conclusion sequent \\
çözünürlük & resolution \\
daraltma & contraction \\
de Bruijn indisi & de Bruijn index \\
değiştirme & exchange \\
denklik & equivalence, biconditional \\
denklik ilişkisi & equivalence relation \\
denklik sınıfı & equivalence class \\
derin gömme & deep embedding \\
dışsallık & extensionality \\
doğal çıkarım & natural deduction \\
doğruluk değerlemesi & valuation \\
doğruluk tablosu & truth table \\
doğrusal mantık & linear logic \\
dolambaç & detour, maximal formula \\
durum ayrımı & proof by cases \\
elaboratör & elaborator \\
eleme kuralı & elimination rule \\
eşleme & bijection \\
evren & universe \\
evren hiyerarşisi & universe hierarchy \\
evrensel niceleyici & universal quantifier \\
ex falso (çelişkiden her şey) & ex falso quodlibet, principle of explosion \\
fonksiyon dışsallığı & function extensionality \\
geçerlilik (sağlamlık) & soundness \\
geri dönme & backtracking \\
geriye / ileriye akıl yürütme & backward / forward reasoning \\
giriş kuralı & introduction rule \\
güçlü normalleşme & strong normalization \\
güçlü tümevarım & strong induction \\
güvenilir çekirdek & trusted kernel \\
hedef & goal \\
hesaplanamaz & noncomputable \\
hipotetik kıyas & hypothetical syllogism \\
impredikatif & impredicative \\
indeks & index \\
İnşalar Hesabı & Calculus of Constructions \\
ispat arama & proof search \\
ispat asistanı & proof assistant \\
ispat karmaşıklığı & proof complexity \\
işaretli atom & signed atom, literal \\
iyi sıralama ilkesi & well-ordering principle \\
iyi tanımlılık & well-definedness \\
kabul edilebilir kural & admissible rule \\
kanıt-ilgisizliği & proof irrelevance \\
kanonik değerleme & canonical valuation \\
karar verilebilir & decidable \\
karşı model & countermodel \\
kaydırma & shifting \\
kesişim & intersection \\
kesme eleme & cut elimination \\
kesme kuralı & cut rule \\
kısmi sıralama & partial order \\
kip mantığı & modal logic \\
kombinatör & combinator \\
kontrapozisyon & contraposition \\
köşegen yöntemi & diagonal argument \\
kural tümevarımı & rule induction \\
kurucu & constructor \\
lif & fibre \\
Lindenbaum önsavı & Lindenbaum's lemma \\
maksimal tutarlı küme & maximal consistent set \\
mantıksal sonuç & logical consequence \\
metamantık & metalogic \\
model varlığı & model existence \\
motive & motive \\
nesne dili & object language \\
niceleyici & quantifier \\
normal türetim & normal derivation \\
normalleştirme & normalization \\
önerme & proposition \\
önerme dışsallığı & propositional extensionality \\
önermesel eşitlik & propositional equality \\
örten & surjective \\
öz değişken koşulu & eigenvariable condition \\
özdeşlik tipi & identity type \\
özyineleme & recursion \\
özyineleyici & recursor \\
parametre & parameter \\
pozitiflik koşulu & positivity condition \\
predikatif & predicative \\
redeks & redex \\
sağ / sol kural & right / left rule \\
sağlanabilir & satisfiable \\
seçim aksiyomu & axiom of choice \\
serbest değişken & free variable \\
sezgicilik & intuitionism \\
sığ gömme & shallow embedding \\
soyutlama & abstraction \\
sözdizim ağacı & syntax tree \\
sözdizimi & syntax \\
taktik & tactic \\
tamlık & completeness \\
tanık & witness \\
tanık sabiti & witness constant, Henkin constant \\
tanım evreni & domain of discourse \\
tanımsal eşitlik & definitional equality \\
taşıma & transport \\
tekdeğerlilik & univalence \\
tekdüzelik & monotonicity \\
temsilci & representative \\
terim modeli & term model \\
ters matematik & reverse mathematics \\
tersinir kural & invertible rule \\
tıkızlık & compactness \\
tip ailesi & type family \\
tip denetleyici & type checker \\
tip korunumu & subject reduction, type preservation \\
totoloji & tautology \\
tutarlılık & consistency \\
tümevarım & induction \\
tümevarım hipotezi & induction hypothesis \\
tümevarımsal aile & inductive family \\
tümevarımsal tip & inductive type \\
tümevarımsal yüklem & inductive predicate \\
tümleyen & complement \\
türetilebilirlik & derivability \\
türetilmiş kural & derived rule \\
türetim & derivation \\
uygulama & application \\
üçüncü hâlin olanaksızlığı & law of excluded middle \\
üst dil & metalanguage \\
varoluşsal niceleyici & existential quantifier \\
yakalama & variable capture \\
yansımalı-geçişli kapanış & reflexive-transitive closure \\
yapıcı (sezgici) mantık & constructive (intuitionistic) logic \\
yapısal kural & structural rule \\
yapısal özyineleme & structural recursion \\
yapısal tümevarım & structural induction \\
yerellik & locality, coincidence lemma \\
yerine koyma & substitution \\
yoğunlamalı & intensional \\
yüklem & predicate \\
zayıf / güçlü tamlık & weak / strong completeness \\
zayıflatma & weakening \\
\end{tabularx}
\end{center}


\backmatter
\bibliographystyle{plainnat}
\renewcommand{\bibname}{Kaynakça}
\bibliography{kaynakca}

\end{document}
